\section{Diagonalization --- NTime Diagonalizer}
%%% generated by the blueprint pipeline from TCSlib.Complexity.Diagonalization.NTimeDiagonalizer
This module builds the lazy diagonalizer of the nondeterministic time hierarchy theorem
[AB09, \S3.2, Figure 3.1]. It fixes the pure stage structure (the ladder $\ell_0 = 2$,
$\ell_{i+1} = 2^{(T^*_i)^2}$ with budgets $T^*_i$, the stage of an input length, the code each
stage diagonalizes against, and the stage query), and develops the machine-level toolkit used to
compute the stage query in time linear in $g(n) + 1$: normal-form configurations, hosted phases,
copy, clear, rewind, comparison and emission rows, released and clocked calls, and their
embeddings into a larger host machine.

\begin{definition}[Budget of a stage]
\lean{Turing.NTimeDiag.budget}
\label{Turing.NTimeDiag.budget}
\leanok
For a function $f : \mathbb{N} \to \mathbb{N}$ and a natural number $\ell$ (the bottom ladder value
of a stage), let $a = \ell + 1$. The budget of the stage is $T^* = (f(a) + a + 1)^2$.
\end{definition}

\begin{theorem}[Budget formula]
\lean{Turing.NTimeDiag.budget_eq}
\label{Turing.NTimeDiag.budget_eq}
\leanok
\uses{Turing.NTimeDiag.budget}
\difficulty{1}
For every function $f : \mathbb{N} \to \mathbb{N}$ and every $\ell \in \mathbb{N}$, the budget of
the stage with bottom ladder value $\ell$ equals $(f(\ell + 1) + (\ell + 1) + 1)^2$.
\end{theorem}

\begin{definition}[The ladder]
\lean{Turing.NTimeDiag.ladder}
\label{Turing.NTimeDiag.ladder}
\leanok
\uses{Turing.NTimeDiag.budget}
For a function $f : \mathbb{N} \to \mathbb{N}$, write $B_f(\ell) = (f(\ell+1) + \ell + 2)^2$ for
the stage budget. The ladder of $f$ is the sequence $(\ell_i)_{i \ge 0}$ defined by $\ell_0 = 2$
and $\ell_{i+1} = 2^{B_f(\ell_i)^2}$.
\end{definition}

\begin{theorem}[The ladder starts at two]
\lean{Turing.NTimeDiag.ladder_zero}
\label{Turing.NTimeDiag.ladder_zero}
\leanok
\uses{Turing.NTimeDiag.ladder}
\difficulty{1}
For every function $f : \mathbb{N} \to \mathbb{N}$, the first rung of the ladder of $f$ (the
sequence $\ell_0 = 2$, $\ell_{i+1} = 2^{(f(\ell_i+1) + \ell_i + 2)^4}$) is $\ell_0 = 2$.
\end{theorem}

\begin{theorem}[Next rung of the ladder]
\lean{Turing.NTimeDiag.ladder_succ}
\label{Turing.NTimeDiag.ladder_succ}
\leanok
\uses{Turing.NTimeDiag.budget, Turing.NTimeDiag.ladder}
\difficulty{1}
Let $f : \mathbb{N} \to \mathbb{N}$, let $B_f(\ell) = (f(\ell+1) + \ell + 2)^2$ be the stage
budget, and let $(\ell_i)$ be the ladder of $f$, with $\ell_0 = 2$. Then for every $i$,
$\ell_{i+1} = 2^{B_f(\ell_i)^2}$.
\end{theorem}

\begin{theorem}[The ladder more than doubles]
\lean{Turing.NTimeDiag.two_mul_lt_ladder_succ}
\label{Turing.NTimeDiag.two_mul_lt_ladder_succ}
\leanok
\uses{Turing.NTimeDiag.budget, Turing.NTimeDiag.ladder, Turing.NTimeDiag.ladder_succ}
\difficulty{4}
Let $f : \mathbb{N} \to \mathbb{N}$ and let $(\ell_i)$ be the ladder of $f$, given by
$\ell_0 = 2$ and $\ell_{i+1} = 2^{B^2}$ with $B = (f(\ell_i+1) + \ell_i + 2)^2$. Then
$2\ell_i < \ell_{i+1}$ for every $i \in \mathbb{N}$.
\end{theorem}

\begin{theorem}[The ladder dominates its index]
\lean{Turing.NTimeDiag.le_ladder}
\label{Turing.NTimeDiag.le_ladder}
\leanok
\uses{Turing.NTimeDiag.ladder, Turing.NTimeDiag.ladder_zero, Turing.NTimeDiag.two_mul_lt_ladder_succ}
\difficulty{3}
Let $f : \mathbb{N} \to \mathbb{N}$ and let $(\ell_i)$ be the ladder of $f$, given by
$\ell_0 = 2$ and $\ell_{i+1} = 2^{(f(\ell_i+1) + \ell_i + 2)^4}$. Then $i + 2 \le \ell_i$ for
every $i \in \mathbb{N}$.
\end{theorem}

\begin{theorem}[The ladder is strictly increasing]
\lean{Turing.NTimeDiag.ladder_strictMono}
\label{Turing.NTimeDiag.ladder_strictMono}
\leanok
\uses{Turing.NTimeDiag.ladder, Turing.NTimeDiag.two_mul_lt_ladder_succ}
\difficulty{2}
For every $f : \mathbb{N} \to \mathbb{N}$, the ladder of $f$, given by $\ell_0 = 2$ and
$\ell_{i+1} = 2^{(f(\ell_i+1) + \ell_i + 2)^4}$, is strictly increasing: $i < j$ implies
$\ell_i < \ell_j$.
\end{theorem}

\begin{theorem}[Every length lies below some rung]
\lean{Turing.NTimeDiag.exists_le_ladder_succ}
\label{Turing.NTimeDiag.exists_le_ladder_succ}
\leanok
\uses{Turing.NTimeDiag.ladder, Turing.NTimeDiag.le_ladder}
\difficulty{2}
Let $f : \mathbb{N} \to \mathbb{N}$ and let $(\ell_i)$ be the ladder of $f$, given by
$\ell_0 = 2$ and $\ell_{i+1} = 2^{(f(\ell_i+1) + \ell_i + 2)^4}$. For every $n \in \mathbb{N}$
there is an index $i$ with $n \le \ell_{i+1}$.
\end{theorem}

\begin{definition}[Stage of a length]
\lean{Turing.NTimeDiag.stage}
\label{Turing.NTimeDiag.stage}
\leanok
\uses{Turing.NTimeDiag.exists_le_ladder_succ}
Let $f : \mathbb{N} \to \mathbb{N}$ and let $(\ell_i)$ be the ladder of $f$, given by
$\ell_0 = 2$ and $\ell_{i+1} = 2^{(f(\ell_i+1) + \ell_i + 2)^4}$. The stage $s_f(n)$ of a length
$n \in \mathbb{N}$ is the least index $i$ with $n \le \ell_{i+1}$.
\end{definition}

\begin{theorem}[The stage brackets its length]
\lean{Turing.NTimeDiag.stage_spec}
\label{Turing.NTimeDiag.stage_spec}
\leanok
\uses{Turing.NTimeDiag.exists_le_ladder_succ, Turing.NTimeDiag.ladder, Turing.NTimeDiag.ladder_zero, Turing.NTimeDiag.stage}
\difficulty{4}
Let $f : \mathbb{N} \to \mathbb{N}$, let $(\ell_i)$ be the ladder of $f$ (with $\ell_0 = 2$ and
$\ell_{i+1} = 2^{(f(\ell_i+1) + \ell_i + 2)^4}$), and let $s_f(n)$ be the least $i$ with
$n \le \ell_{i+1}$. If $n \ge 3$, then $\ell_{s_f(n)} < n \le \ell_{s_f(n)+1}$.
\end{theorem}

\begin{theorem}[The stage is determined by its bracket]
\lean{Turing.NTimeDiag.stage_eq}
\label{Turing.NTimeDiag.stage_eq}
\leanok
\uses{Turing.NTimeDiag.exists_le_ladder_succ, Turing.NTimeDiag.ladder, Turing.NTimeDiag.ladder_strictMono, Turing.NTimeDiag.stage}
\difficulty{4}
Let $f : \mathbb{N} \to \mathbb{N}$, let $(\ell_i)$ be the ladder of $f$ (with $\ell_0 = 2$ and
$\ell_{i+1} = 2^{(f(\ell_i+1) + \ell_i + 2)^4}$), and let $s_f(n)$ be the least $i$ with
$n \le \ell_{i+1}$. For all $n, i \in \mathbb{N}$, if $\ell_i < n \le \ell_{i+1}$ then
$s_f(n) = i$.
\end{theorem}

\begin{theorem}[Short lengths lie in stage zero]
\lean{Turing.NTimeDiag.stage_of_le_two}
\label{Turing.NTimeDiag.stage_of_le_two}
\leanok
\uses{Turing.NTimeDiag.exists_le_ladder_succ, Turing.NTimeDiag.le_ladder, Turing.NTimeDiag.stage}
\difficulty{2}
Let $f : \mathbb{N} \to \mathbb{N}$, let $(\ell_i)$ be the ladder of $f$ (with $\ell_0 = 2$ and
$\ell_{i+1} = 2^{(f(\ell_i+1) + \ell_i + 2)^4}$), and let $s_f(n)$ be the least $i$ with
$n \le \ell_{i+1}$. If $n \le 2$, then $s_f(n) = 0$.
\end{theorem}

\begin{definition}[Code diagonalized at a stage]
\lean{Turing.NTimeDiag.stageCode}
\label{Turing.NTimeDiag.stageCode}
\leanok
For a stage index $i \in \mathbb{N}$, let $j$ be the first component of the image of $i$ under the
standard unpairing bijection $\mathbb{N} \to \mathbb{N} \times \mathbb{N}$. The code of stage $i$
is the binary word $\alpha_i$ obtained from the binary expansion of $j + 1$ (least significant
bit first) by deleting its last, most significant, bit.
\end{definition}

\begin{theorem}[Stage code formula]
\lean{Turing.NTimeDiag.stageCode_eq}
\label{Turing.NTimeDiag.stageCode_eq}
\leanok
\uses{Turing.NTimeDiag.stageCode}
\difficulty{1}
For every $i \in \mathbb{N}$, the code of stage $i$ equals the binary expansion (least
significant bit first) of $j + 1$ with its last bit removed, where $j$ is the first component
of the standard unpairing of $i$.
\end{theorem}

\begin{theorem}[Every code recurs at unboundedly many stages]
\lean{Turing.NTimeDiag.stageCode_frequently}
\label{Turing.NTimeDiag.stageCode_frequently}
\leanok
\uses{Turing.NTimeDiag.stageCode}
\difficulty{5}
For a stage index $i$, let $\alpha_i$ denote the binary expansion (least significant bit first)
of $j + 1$ with its most significant bit removed, where $j$ is the first component of the
standard unpairing of $i$. Then for every binary word $\alpha$ and every $N \in \mathbb{N}$
there is a stage $i \ge N$ with $\alpha_i = \alpha$.
\end{theorem}

\begin{definition}[The stage query]
\lean{Turing.NTimeDiag.query}
\label{Turing.NTimeDiag.query}
\leanok
\uses{Turing.NTimeDiag.budget, Turing.NTimeDiag.ladder, Turing.NTimeDiag.stage, Turing.NTimeDiag.stageCode, Turing.pairEncode}
Let $f, g : \mathbb{N} \to \mathbb{N}$, let $(\ell_i)$ be the ladder of $f$, let $s_f(n)$ be the
stage of a length $n$, let $\alpha_i$ be the code of stage $i$, and let
$B_f(\ell) = (f(\ell+1) + \ell + 2)^2$ be the stage budget. Write $\langle u, v \rangle$ for the
pairing encoding of binary words and $\mathrm{bin}(m)$ for the binary expansion of $m$. The stage
query $q_{f,g}(x)$ of a binary word $x$ is defined as follows. If $x = 1^n$ with $n \ge 3$, put
$i = s_f(n)$; then $q_{f,g}(x)$ is the bit $1$ followed by
$\langle\langle \mathrm{bin}(B_f(\ell_i)), \alpha_i\rangle, 1^{\ell_i + 1}\rangle$ when
$n = \ell_{i+1}$, and the bit $0$ followed by
$\langle\langle \mathrm{bin}(g(n)), \alpha_i\rangle, 1^{n + 1}\rangle$ otherwise. On every other
input, $q_{f,g}(x)$ is the empty word.
\end{definition}

\begin{theorem}[Stage query on unary inputs]
\lean{Turing.NTimeDiag.query_eq}
\label{Turing.NTimeDiag.query_eq}
\leanok
\uses{Turing.NTimeDiag.budget, Turing.NTimeDiag.ladder, Turing.NTimeDiag.query, Turing.NTimeDiag.stage, Turing.NTimeDiag.stageCode, Turing.pairEncode}
\difficulty{1}
Let $f, g : \mathbb{N} \to \mathbb{N}$, with ladder $(\ell_i)$ of $f$, stage function $s_f$,
stage codes $\alpha_i$, stage budget $B_f(\ell) = (f(\ell+1) + \ell + 2)^2$, pairing encoding
$\langle \cdot, \cdot \rangle$ and binary expansion $\mathrm{bin}$. For every $n \ge 3$, with
$i = s_f(n)$, the stage query of $1^n$ is the bit $1$ followed by
$\langle\langle \mathrm{bin}(B_f(\ell_i)), \alpha_i\rangle, 1^{\ell_i+1}\rangle$ if
$n = \ell_{i+1}$, and is the bit $0$ followed by
$\langle\langle \mathrm{bin}(g(n)), \alpha_i\rangle, 1^{n+1}\rangle$ otherwise.
\end{theorem}

\begin{theorem}[Stage query off its guard is empty]
\lean{Turing.NTimeDiag.query_of_not_guard}
\label{Turing.NTimeDiag.query_of_not_guard}
\leanok
\uses{Turing.NTimeDiag.query}
\difficulty{1}
Let $f, g : \mathbb{N} \to \mathbb{N}$ and let $x$ be a binary word that is not of the form $1^n$
with $n \ge 3$. Then the stage query $q_{f,g}(x)$ of the lazy diagonalizer is the empty word.
\end{theorem}

\begin{theorem}[Bounded language of a halting NDTM is in NTIME]
\lean{Turing.NTimeDiag.acceptLang_mem_NTIME}
\label{Turing.NTimeDiag.acceptLang_mem_NTIME}
\leanok
\uses{Complexity.NTIME, Turing.FinNDTM, Turing.FinNDTM.AcceptsWithin, Turing.NDTM.HaltsWithin}
\difficulty{1}
Let $D$ be a finite nondeterministic Turing machine over the binary alphabet, let
$K \in \mathbb{N}$ and $g : \mathbb{N} \to \mathbb{N}$. If on every binary input $x$ every
computation branch of $D$ halts within $K(g(|x|) + 1)$ steps, then the language of all $x$ that
$D$ accepts within $K(g(|x|) + 1)$ steps belongs to $\mathbf{NTIME}(g(n) + 1)$.
\end{theorem}

\begin{theorem}[Stage query reduces to unary words]
\lean{Turing.NTimeDiag.qy_guard}
\label{Turing.NTimeDiag.qy_guard}
\leanok
\uses{Turing.NTimeDiag.query, Turing.NTimeDiag.query_of_not_guard}
\difficulty{3}
Let $f, g : \mathbb{N} \to \mathbb{N}$ and let $q_{f,g}$ be the stage query of the lazy
diagonalizer. For every binary word $x$, $q_{f,g}(x) = q_{f,g}(1^{|x|})$ if $x = 1^{|x|}$ and
$|x| \ge 3$, and $q_{f,g}(x)$ is the empty word otherwise.
\end{theorem}

\begin{theorem}[Linear-time unary guard decider]
\lean{Turing.NTimeDiag.qy_guard_decider}
\label{Turing.NTimeDiag.qy_guard_decider}
\leanok
\uses{Turing.FinTM, Turing.FinTM.ComputesFunInTime, Turing.FinTM.ComputesInTime.mono, Turing.FinTM.computesFunInTime_allTrue, Turing.FinTM.computesFunInTime_cond, Turing.FinTM.computesFunInTime_const, Turing.FinTM.computesFunInTime_decideEq, Turing.FinTM.computesFunInTime_iff}
\difficulty{5}
There exist a finite deterministic Turing machine $D$ over the binary alphabet and a constant
$c \in \mathbb{N}$ such that, on every binary input $x$, $D$ halts within $c(|x| + 1)$ steps
with the one-bit output $1$ if $x$ consists only of ones and $|x| \ge 3$, and $0$ otherwise.
\end{theorem}

\begin{theorem}[Full stage query from its unary part]
\lean{Turing.NTimeDiag.qy_query_of_main}
\label{Turing.NTimeDiag.qy_query_of_main}
\leanok
\uses{Complexity.TimeConstructible, Complexity.timeConstructible_iff, Turing.FinTM, Turing.FinTM.ComputesFunInTime, Turing.FinTM.ComputesInTime.mono, Turing.FinTM.computesFunInTime_cond, Turing.FinTM.computesFunInTime_const, Turing.FinTM.computesFunInTime_iff, Turing.NTimeDiag.query, Turing.NTimeDiag.qy_guard, Turing.NTimeDiag.qy_guard_decider}
\difficulty{4}
Let $f, g : \mathbb{N} \to \mathbb{N}$ with $g$ time-constructible, and let $q_{f,g}$ be the stage
query of the lazy diagonalizer. Suppose a finite deterministic Turing machine $M$ computes the
map $x \mapsto q_{f,g}(1^{|x|})$ within $A(g(n) + 1)$ steps on inputs of length $n$, for some
constant $A$. Then some finite deterministic Turing machine $G$ and constant $C$ compute
$q_{f,g}$ itself within $C(g(n) + 1)$ steps on inputs of length $n$.
\end{theorem}

\begin{theorem}[Next rung as a fourth power]
\lean{Turing.NTimeDiag.qy_ladder_succ_eq}
\label{Turing.NTimeDiag.qy_ladder_succ_eq}
\leanok
\uses{Turing.NTimeDiag.budget_eq, Turing.NTimeDiag.ladder, Turing.NTimeDiag.ladder_succ}
\difficulty{1}
Let $f : \mathbb{N} \to \mathbb{N}$ and let $(\ell_i)$ be the ladder of $f$, given by
$\ell_0 = 2$ and $\ell_{i+1} = 2^{B^2}$ with budget $B = (f(\ell_i + 1) + \ell_i + 2)^2$. Then for
every $i$, $\ell_{i+1} = 2^{X^4}$ where $X = f(\ell_i + 1) + (\ell_i + 1) + 1$.
\end{theorem}

\begin{theorem}[Sum of earlier rungs]
\lean{Turing.NTimeDiag.qy_ledger}
\label{Turing.NTimeDiag.qy_ledger}
\leanok
\uses{Turing.NTimeDiag.ladder, Turing.NTimeDiag.two_mul_lt_ladder_succ}
\difficulty{3}
Let $f : \mathbb{N} \to \mathbb{N}$ and let $(\ell_i)$ be the ladder of $f$, given by
$\ell_0 = 2$ and $\ell_{i+1} = 2^{(f(\ell_i+1) + \ell_i + 2)^4}$. For every $s \in \mathbb{N}$,
\[ \sum_{h < s} \ell_{h+1} \le 2\,\ell_s. \]
\end{theorem}

\begin{theorem}[Value of the all-ones counter]
\lean{Turing.NTimeDiag.qy_ones_val}
\label{Turing.NTimeDiag.qy_ones_val}
\leanok
\difficulty{3}
For every $m \in \mathbb{N}$, the word $1^m$, read as a binary numeral, has value $2^m - 1$.
\end{theorem}

\begin{theorem}[Value of zeros followed by a one]
\lean{Turing.NTimeDiag.qy_rung_val}
\label{Turing.NTimeDiag.qy_rung_val}
\leanok
\difficulty{3}
For every $m \in \mathbb{N}$, the word $0^m 1$, read as a binary numeral with the least
significant digit first, has value $2^m$.
\end{theorem}

\begin{theorem}[The all-ones counter is large enough]
\lean{Turing.NTimeDiag.qy_ones_ge}
\label{Turing.NTimeDiag.qy_ones_ge}
\leanok
\uses{Complexity.LogProg.lt_two_pow_bits}
\difficulty{4}
For all $n, D \in \mathbb{N}$, writing $|\mathrm{bin}(m)|$ for the length of the binary expansion
of $m$,
\[ D(n + 1) \le 2^{|\mathrm{bin}(n)| + |\mathrm{bin}(D)|} - 1. \]
\end{theorem}

\begin{theorem}[Install-call bound for small arguments]
\lean{Turing.NTimeDiag.qy_call_bound}
\label{Turing.NTimeDiag.qy_call_bound}
\leanok
\uses{Turing.length_bits_le_self}
\difficulty{4}
Let $f : \mathbb{N} \to \mathbb{N}$ and $c, c', a, n \in \mathbb{N}$ with $a \le n$ and
$f(a) \le n$. Writing $|\mathrm{bin}(m)|$ for the length of the binary expansion of $m$,
\[ c'\bigl(c(f(a) + 1) + a + |\mathrm{bin}(f(a))| + 1\bigr) \le c'(c + 2)(n + 1). \]
\end{theorem}

\begin{theorem}[Exponential lower bound on the ladder]
\lean{Turing.NTimeDiag.qy_two_pow_le_ladder}
\label{Turing.NTimeDiag.qy_two_pow_le_ladder}
\leanok
\uses{Turing.NTimeDiag.ladder, Turing.NTimeDiag.ladder_zero, Turing.NTimeDiag.two_mul_lt_ladder_succ}
\difficulty{3}
Let $f : \mathbb{N} \to \mathbb{N}$ and let $(\ell_i)$ be the ladder of $f$, given by
$\ell_0 = 2$ and $\ell_{i+1} = 2^{(f(\ell_i+1) + \ell_i + 2)^4}$. Then $2^{i+1} \le \ell_i$ for
every $i \in \mathbb{N}$.
\end{theorem}

\begin{theorem}[Stage index is below the logarithm]
\lean{Turing.NTimeDiag.qy_stage_lt_log}
\label{Turing.NTimeDiag.qy_stage_lt_log}
\leanok
\uses{Turing.NTimeDiag.qy_two_pow_le_ladder, Turing.NTimeDiag.stage, Turing.NTimeDiag.stage_spec}
\difficulty{3}
Let $f : \mathbb{N} \to \mathbb{N}$, let $(\ell_i)$ be the ladder of $f$ (with $\ell_0 = 2$ and
$\ell_{i+1} = 2^{(f(\ell_i+1) + \ell_i + 2)^4}$), and let $s_f(n)$ be the least $i$ with
$n \le \ell_{i+1}$. For every $n \ge 3$, $s_f(n) < \lfloor \log_2 n \rfloor$.
\end{theorem}

\begin{theorem}[Timeout certificate]
\lean{Turing.NTimeDiag.qy_timeout_cert}
\label{Turing.NTimeDiag.qy_timeout_cert}
\leanok
\uses{Turing.NTimeDiag.qy_call_bound, Turing.NTimeDiag.qy_ones_ge}
\difficulty{3}
Let $f : \mathbb{N} \to \mathbb{N}$ and $c, c', a, n, t \in \mathbb{N}$, and write
$|\mathrm{bin}(m)|$ for the length of the binary expansion of $m$. Suppose $a \le n$,
$t \le c'(c(f(a) + 1) + a + |\mathrm{bin}(f(a))| + 1)$, and
$2^{|\mathrm{bin}(n)| + |\mathrm{bin}(c'(c+2))|} - 1 < t$. Then $n < f(a)$.
\end{theorem}

\begin{definition}[Normal-form configuration]
\lean{Turing.NTimeDiag.qyC}
\label{Turing.NTimeDiag.qyC}
\leanok
\uses{Turing.Cfg}
For a multi-tape machine with $K$ work tapes over the binary alphabet, an input $x$, a state $s$,
an input-head position $p \in \{0, \dots, |x| + 1\}$, work-tape contents $T$ (a partial map from
tape index and integer cell to bits), work-head positions $P$ and an output word $o$, the
normal-form configuration $\langle s, p, T, P, o \rangle$ is the live configuration with state
$s$ and exactly these input head, tapes, heads and output.
\end{definition}

\begin{theorem}[Renaming states of normal-form configurations]
\lean{Turing.NTimeDiag.qyC_mapState}
\label{Turing.NTimeDiag.qyC_mapState}
\leanok
\uses{Turing.Cfg, Turing.Cfg.mapState, Turing.NTimeDiag.qyC}
\difficulty{2}
Let $\varphi : S \to S'$ be a map of state sets. For a live configuration
$\langle s, p, T, P, o \rangle$ of a $K$-tape machine on input $x$, with state $s$, input head
$p$, work tapes $T$, work heads $P$ and output $o$, renaming its state along $\varphi$ yields
$\langle \varphi(s), p, T, P, o \rangle$.
\end{theorem}

\begin{theorem}[A hosted phase is a segment]
\lean{Turing.NTimeDiag.qy_hosted'}
\label{Turing.NTimeDiag.qy_hosted'}
\leanok
\uses{Turing.MultiTapeTM, Turing.MultiTapeTM.runFrom, Turing.MultiTapeTM.runFrom_state_ne_none_of_le, Turing.NTimeDiag.qyC, Turing.NTimeDiag.qyC_mapState, Turing.Segment, Turing.phaseHostTM, Turing.phaseHostTM_run, Turing.segment_iff}
\difficulty{4}
Let $(M_i)_{i \in I}$ be a family of $K$-tape machines with state sets $\Phi_i$, and let $H$ be
the phase host with start state $q_0$ and redirect table $r$: $H$ runs on tagged states
$(i, s)$, takes one stationary silent step to $r_i(s)$ when this is defined, and otherwise
follows $M_i$. Fix a phase $k$ and an input $x$, and write $\langle s, p, T, P, o\rangle$ for the
live configuration with state $s$, input head $p$, work tapes $T$, heads $P$ and output $o$.
Suppose $M_k$ runs from $\langle e, p, T, P, o\rangle$ to $\langle s', p', T', P', o'\rangle$ in
$t$ steps, $r_k(s') = q'$, and $r_k$ is undefined on every state of this run before time $t$.
Let $F$ be a set of host states containing no tagged state $(k, s)$ for any $s$ visited by
$M_k$'s run up to time $t$. Then $H$ runs from $\langle (k, e), p, T, P, o\rangle$ to
$\langle q', p', T', P', o'\rangle$ in exactly $t + 1$ steps, with no state in $F$ before time
$t + 1$.
\end{theorem}

\begin{theorem}[A hosted phase avoids foreign states]
\lean{Turing.NTimeDiag.qy_hosted}
\label{Turing.NTimeDiag.qy_hosted}
\leanok
\uses{Turing.MultiTapeTM, Turing.MultiTapeTM.runFrom, Turing.NTimeDiag.qyC, Turing.NTimeDiag.qy_hosted', Turing.Segment, Turing.phaseHostTM}
\difficulty{1}
Let $(M_i)_{i \in I}$ be a family of $K$-tape machines with state sets $\Phi_i$, and let $H$ be
the phase host with start state $q_0$ and redirect table $r$ (a redirected state takes one
stationary silent step to its target; any other tagged state $(i, s)$ follows $M_i$). Write
$\langle s, p, T, P, o\rangle$ for the live configuration on input $x$ with state $s$, input head
$p$, work tapes $T$, heads $P$ and output $o$. Suppose $M_k$ runs from
$\langle e, p, T, P, o\rangle$ to $\langle s', p', T', P', o'\rangle$ in $t$ steps,
$r_k(s') = q'$, and $r_k$ is undefined on every state of the run before time $t$. Then for every
set $F$ of host states that contains no state $(k, s)$, $H$ runs from
$\langle (k, e), p, T, P, o\rangle$ to $\langle q', p', T', P', o'\rangle$ in exactly $t + 1$ steps
without visiting $F$ before time $t + 1$.
\end{theorem}

\begin{theorem}[A hosted last phase is a segment]
\lean{Turing.NTimeDiag.qy_hosted_last}
\label{Turing.NTimeDiag.qy_hosted_last}
\leanok
\uses{Turing.MultiTapeTM, Turing.MultiTapeTM.runFrom, Turing.MultiTapeTM.runFrom_state_ne_none_of_le, Turing.NTimeDiag.qyC, Turing.NTimeDiag.qyC_mapState, Turing.Segment, Turing.phaseHostTM, Turing.phaseHostTM_run_le, Turing.segment_iff}
\difficulty{4}
Let $(M_i)_{i \in I}$ be a family of $K$-tape machines with state sets $\Phi_i$, and let $H$ be
the phase host with start state $q_0$ and redirect table $r$ (a redirected state takes one
stationary silent step to its target; any other tagged state $(i, s)$ follows $M_i$). Write
$\langle s, p, T, P, o\rangle$ for the live configuration on input $x$ with state $s$, input head
$p$, work tapes $T$, heads $P$ and output $o$. Suppose $M_k$ runs from
$\langle e, p, T, P, o\rangle$ to $\langle s', p', T', P', o'\rangle$ in $t$ steps and $r_k$ is
undefined on every state of the run before time $t$. Let $F$ be a set of host states containing
no tagged state $(k, s)$ with $s$ visited by the run before time $t$. Then $H$ runs from
$\langle (k, e), p, T, P, o\rangle$ to $\langle (k, s'), p', T', P', o'\rangle$ in exactly $t$
steps, with no state in $F$ before time $t$.
\end{theorem}

\begin{theorem}[Run of a hosted phase]
\lean{Turing.NTimeDiag.qy_hosted_run}
\label{Turing.NTimeDiag.qy_hosted_run}
\leanok
\uses{Turing.MultiTapeTM, Turing.MultiTapeTM.runFrom, Turing.NTimeDiag.qyC, Turing.NTimeDiag.qy_hosted, Turing.phaseHostTM, Turing.segment_iff}
\difficulty{1}
Let $(M_i)_{i \in I}$ be a family of $K$-tape machines with state sets $\Phi_i$, and let $H$ be
the phase host with start state $q_0$ and redirect table $r$ (a redirected state takes one
stationary silent step to its target; any other tagged state $(i, s)$ follows $M_i$). Write
$\langle s, p, T, P, o\rangle$ for the live configuration on input $x$ with state $s$, input head
$p$, work tapes $T$, heads $P$ and output $o$. If $M_k$ runs from
$\langle e, p, T, P, o\rangle$ to $\langle s', p', T', P', o'\rangle$ in $t$ steps,
$r_k(s') = q'$, and $r_k$ is undefined on every state of the run before time $t$, then $H$ run
for $t + 1$ steps from $\langle (k, e), p, T, P, o\rangle$ reaches
$\langle q', p', T', P', o'\rangle$.
\end{theorem}

\begin{theorem}[Writing at the end of a buffer appends]
\lean{Turing.NTimeDiag.qy_patch_end}
\label{Turing.NTimeDiag.qy_patch_end}
\leanok
\uses{Turing.FinTM.bufferTape, Turing.FinTM.bufferTape_append_shift, Turing.FinTM.bufferTape_eq, Turing.FinTM.bufferTape_eq_none_of_not_mem}
\difficulty{4}
For a binary word $w$, let $\mathrm{buf}(w)$ be the tape holding $w$ in cells
$0, \dots, |w| - 1$ and blank elsewhere. Let $u, w$ be binary words and $p = |u|$. The tape whose
cell $q$ holds $\mathrm{buf}(w)(q - p)$ for $p \le q < p + |w|$ and $\mathrm{buf}(u)(q)$ otherwise
equals $\mathrm{buf}(uw)$.
\end{theorem}

\begin{theorem}[Buffer of a mapped word]
\lean{Turing.NTimeDiag.qy_bufferTape_map}
\label{Turing.NTimeDiag.qy_bufferTape_map}
\leanok
\uses{Turing.FinTM.bufferTape, Turing.FinTM.bufferTape_eq}
\difficulty{2}
For a binary word $w$, let $\mathrm{buf}(w)$ be the tape holding $w$ in cells
$0, \dots, |w| - 1$ and blank elsewhere. For every map $g$ on bits and every cell
$z \in \mathbb{Z}$, the cell $z$ of $\mathrm{buf}(g(w))$, where $g$ is applied letterwise, is the
image under $g$ of cell $z$ of $\mathrm{buf}(w)$ (blank if that cell is blank).
\end{theorem}

\begin{theorem}[Redirect guard from a first-visit cut]
\lean{Turing.NTimeDiag.qy_guard_of_cut}
\label{Turing.NTimeDiag.qy_guard_of_cut}
\leanok
\uses{Turing.Cfg, Turing.MultiTapeTM, Turing.MultiTapeTM.runFrom}
\difficulty{1}
Let $M$ be a multi-tape machine with state set $S$, let $c$ be a configuration, $t \in \mathbb{N}$,
$d \in S$, and let $r$ be a partial map on $S$ that is undefined on every state other than $d$.
If the run of $M$ from $c$ is not in state $d$ at any time $u < t$, then $r$ is undefined on
every state that this run occupies at any time $u < t$.
\end{theorem}

\begin{theorem}[Mapped copy appending to a word]
\lean{Turing.NTimeDiag.qy_copy}
\label{Turing.NTimeDiag.qy_copy}
\leanok
\uses{Turing.FinTM.bufferTape, Turing.MultiTapeTM.runFrom, Turing.NTimeDiag.qyC, Turing.NTimeDiag.qy_bufferTape_map, Turing.NTimeDiag.qy_patch_end, Turing.copyMapTM, Turing.copyMapTM_run_ofCfg}
\difficulty{4}
Let $K \in \mathbb{N}$, let $\mathrm{src} \ne \mathrm{dst}$ be two of $K$ work tapes, and let $g$
be a map on bits. Consider the mapped-copy machine that copies the word under the head of
$\mathrm{src}$, letterwise through $g$, onto $\mathrm{dst}$ starting at its head, and returns both
heads. Write $\mathrm{buf}(w)$ for the tape holding $w$ in cells $0, \dots, |w| - 1$ and blank
elsewhere, and $\langle s, p, T, P, o\rangle$ for the live configuration on input $x$ with state
$s$, input head $p$, work tapes $T$, heads $P$ and output $o$. If tape $\mathrm{src}$ is
$\mathrm{buf}(w)$ with its head at $0$ and tape $\mathrm{dst}$ is $\mathrm{buf}(u)$ with its head
at $|u|$, then from the start state the machine reaches, after exactly $2|w| + 2$ steps and not
before, its exit state with tape $\mathrm{dst}$ replaced by $\mathrm{buf}(u \, g(w))$ and the
input head, all other tapes, all heads and the output unchanged.
\end{theorem}

\begin{theorem}[Mapped copy overwriting a word]
\lean{Turing.NTimeDiag.qy_copy_over}
\label{Turing.NTimeDiag.qy_copy_over}
\leanok
\uses{Turing.FinTM.bufferTape, Turing.FinTM.bufferTape_eq_none_of_not_mem, Turing.MultiTapeTM.runFrom, Turing.NTimeDiag.qyC, Turing.NTimeDiag.qy_bufferTape_map, Turing.copyMapTM, Turing.copyMapTM_run_ofCfg}
\difficulty{4}
Let $K \in \mathbb{N}$, let $\mathrm{src} \ne \mathrm{dst}$ be two of $K$ work tapes, and let $g$
be a map on bits. Consider the mapped-copy machine that copies the word under the head of
$\mathrm{src}$, letterwise through $g$, onto $\mathrm{dst}$ starting at its head, and returns both
heads. Write $\mathrm{buf}(w)$ for the tape holding $w$ in cells $0, \dots, |w| - 1$ and blank
elsewhere, and $\langle s, p, T, P, o\rangle$ for the live configuration on input $x$ with state
$s$, input head $p$, work tapes $T$, heads $P$ and output $o$. Let $u, w$ be binary words with
$|u| \le |w|$. If tape $\mathrm{src}$ is $\mathrm{buf}(w)$ and tape $\mathrm{dst}$ is
$\mathrm{buf}(u)$, both heads at $0$, then from the start state the machine reaches, after exactly
$2|w| + 2$ steps and not before, its exit state with tape $\mathrm{dst}$ replaced by
$\mathrm{buf}(g(w))$ and everything else unchanged.
\end{theorem}

\begin{theorem}[Mapped move appending to a word]
\lean{Turing.NTimeDiag.qy_transfer}
\label{Turing.NTimeDiag.qy_transfer}
\leanok
\uses{Turing.FinTM.bufferTape, Turing.FinTM.bufferTape_eq_none_of_not_mem, Turing.FinTM.bufferTape_nil, Turing.MultiTapeTM.runFrom, Turing.NTimeDiag.qyC, Turing.NTimeDiag.qy_bufferTape_map, Turing.NTimeDiag.qy_patch_end, Turing.transferMapTM, Turing.transferMapTM_run_ofCfg}
\difficulty{4}
Let $K \in \mathbb{N}$, let $\mathrm{src} \ne \mathrm{dst}$ be two of $K$ work tapes, and let $g$
be a map on bits. Consider the mapped-transfer machine that copies the word under the head of
$\mathrm{src}$, letterwise through $g$, onto $\mathrm{dst}$ starting at its head, blanking each
source cell on the return pass. Write $\mathrm{buf}(w)$ for the tape holding $w$ in cells
$0, \dots, |w| - 1$ and blank elsewhere, and $\langle s, p, T, P, o\rangle$ for the live
configuration on input $x$ with state $s$, input head $p$, work tapes $T$, heads $P$ and output
$o$. If tape $\mathrm{src}$ is $\mathrm{buf}(w)$ with its head at $0$ and tape $\mathrm{dst}$ is
$\mathrm{buf}(u)$ with its head at $|u|$, then from the start state the machine reaches, after
exactly $2|w| + 2$ steps and not before, its exit state with tape $\mathrm{src}$ blank, tape
$\mathrm{dst}$ equal to $\mathrm{buf}(u \, g(w))$, and everything else unchanged.
\end{theorem}

\begin{theorem}[Clearing a word]
\lean{Turing.NTimeDiag.qy_clear}
\label{Turing.NTimeDiag.qy_clear}
\leanok
\uses{Turing.FinTM.bufferTape, Turing.FinTM.bufferTape_eq_none_of_not_mem, Turing.FinTM.bufferTape_nil, Turing.MultiTapeTM.runFrom, Turing.NTimeDiag.qyC, Turing.clearTM, Turing.clearTM_run_ofCfg}
\difficulty{4}
Let $K \in \mathbb{N}$, let $i$ be one of $K$ work tapes, and consider the clearing machine that
blanks the word under the head of tape $i$ and returns its head. Write $\mathrm{buf}(w)$ for the
tape holding $w$ in cells $0, \dots, |w| - 1$ and blank elsewhere, and
$\langle s, p, T, P, o\rangle$ for the live configuration on input $x$ with state $s$, input head
$p$, work tapes $T$, heads $P$ and output $o$. If tape $i$ is $\mathrm{buf}(w)$ with its head at
$0$, then from the start state the machine reaches, after exactly $2|w| + 2$ steps and not
before, its exit state with tape $i$ entirely blank and everything else unchanged.
\end{theorem}

\begin{theorem}[Appending a constant word]
\lean{Turing.NTimeDiag.qy_writeConst}
\label{Turing.NTimeDiag.qy_writeConst}
\leanok
\uses{Turing.FinTM.bufferTape, Turing.MultiTapeTM.runFrom, Turing.NTimeDiag.qyC, Turing.NTimeDiag.qy_patch_end, Turing.WritePhase, Turing.writeConstTM, Turing.writeConstTM_run_ofCfg}
\difficulty{4}
Let $K \in \mathbb{N}$, let $i$ be one of $K$ work tapes, let $w_0$ be a fixed binary word, and
consider the constant-writer machine that writes $w_0$ on tape $i$ from its head and returns the
head. Write $\mathrm{buf}(w)$ for the tape holding $w$ in cells $0, \dots, |w| - 1$ and blank
elsewhere, and $\langle s, p, T, P, o\rangle$ for the live configuration on input $x$ with state
$s$, input head $p$, work tapes $T$, heads $P$ and output $o$. If tape $i$ is $\mathrm{buf}(u)$
with its head at $|u|$, then from the first writing state the machine reaches, after exactly
$2|w_0| + 2$ steps and not before, its exit state with tape $i$ equal to $\mathrm{buf}(u w_0)$ and
everything else, including the head of tape $i$, unchanged.
\end{theorem}

\begin{theorem}[Rewinding a head]
\lean{Turing.NTimeDiag.qy_rewind}
\label{Turing.NTimeDiag.qy_rewind}
\leanok
\uses{Turing.FinTM.bufferTape, Turing.MultiTapeTM.runFrom, Turing.NTimeDiag.qyC, Turing.rewindTM, Turing.rewindTM_run_ofCfg}
\difficulty{3}
Let $K \in \mathbb{N}$, let $i$ be one of $K$ work tapes, and consider the rewinding machine that
moves the head of tape $i$ left to the start of its word. Write $\mathrm{buf}(u)$ for the tape
holding $u$ in cells $0, \dots, |u| - 1$ and blank elsewhere, and $\langle s, p, T, P, o\rangle$
for the live configuration on input $x$ with state $s$, input head $p$, work tapes $T$, heads $P$
and output $o$. If tape $i$ is $\mathrm{buf}(u)$ with its head at $n \le |u|$, then from the start
state the machine reaches, after exactly $n + 2$ steps and not before, its exit state with the
head of tape $i$ at $0$ and everything else unchanged.
\end{theorem}

\begin{theorem}[Embedded run as a transported run]
\lean{Turing.NTimeDiag.qy_embed_gen}
\label{Turing.NTimeDiag.qy_embed_gen}
\leanok
\uses{Turing.Cfg, Turing.MultiTapeTM, Turing.MultiTapeTM.runFrom, Turing.NTimeDiag.qyC, Turing.embedEmitCfg, Turing.embedEmitCfg_eq, Turing.embedEmitTM, Turing.embedEmitTM_runFrom}
\difficulty{2}
Let $M$ be an $m$-tape machine with state set $S$, and let $\sigma$ be an injection from its $m$
tapes into $K$ host tapes; the embedded machine runs $M$ on the host tapes $\sigma(0), \dots,
\sigma(m - 1)$ and forwards every emission to the host output. Write
$\langle s, p, T, P, o\rangle$ for the live configuration on input $x$ with state $s$, input head
$p$, work tapes $T$, heads $P$ and output $o$. For every $t$, the embedded run for $t$ steps from
$\langle s, p, T, P, o\rangle$ is the configuration whose state and input head are those of
$M$'s run for $t$ steps from $\langle s, p, T \circ \sigma, P \circ \sigma, \varepsilon\rangle$,
whose tape and head $\sigma(i)$ are that run's tape and head $i$, whose other tapes and heads are
those of $T$ and $P$, and whose output is $o$ followed by that run's output.
\end{theorem}

\begin{theorem}[Embedded run in normal form]
\lean{Turing.NTimeDiag.qy_embed}
\label{Turing.NTimeDiag.qy_embed}
\leanok
\uses{Turing.MultiTapeTM, Turing.MultiTapeTM.runFrom, Turing.NTimeDiag.qyC, Turing.NTimeDiag.qy_embed_gen, Turing.embedEmitCfg_eq, Turing.embedEmitTM}
\difficulty{2}
Let $M$ be an $m$-tape machine and $\sigma$ an injection from its tapes into $K$ host tapes; the
embedded machine runs $M$ on the host tapes selected by $\sigma$ and forwards its emissions to
the host output. Write $\langle s, p, T, P, o\rangle$ for the live configuration on input $x$ with
state $s$, input head $p$, work tapes $T$, heads $P$ and output $o$. Suppose $M$ runs in $t$ steps
from $\langle s, p, T \circ \sigma, P \circ \sigma, \varepsilon\rangle$ to
$\langle s', p', T_0, P_0, o'\rangle$, and let $T', P'$ agree with $T_0, P_0$ on the selected
tapes ($T'(\sigma(i)) = T_0(i)$, $P'(\sigma(i)) = P_0(i)$) and with $T, P$ on every unselected
tape. Then the embedded machine runs in $t$ steps from $\langle s, p, T, P, o\rangle$ to
$\langle s', p', T', P', o\, o'\rangle$.
\end{theorem}

\begin{theorem}[States of an embedded run]
\lean{Turing.NTimeDiag.qy_embed_state}
\label{Turing.NTimeDiag.qy_embed_state}
\leanok
\uses{Turing.MultiTapeTM, Turing.MultiTapeTM.runFrom, Turing.NTimeDiag.qyC, Turing.NTimeDiag.qy_embed_gen, Turing.embedEmitTM}
\difficulty{1}
Let $M$ be an $m$-tape machine and $\sigma$ an injection from its tapes into $K$ host tapes; the
embedded machine runs $M$ on the host tapes selected by $\sigma$ and forwards its emissions to
the host output. Write $\langle s, p, T, P, o\rangle$ for the live configuration on input $x$ with
state $s$, input head $p$, work tapes $T$, heads $P$ and output $o$. For every time $u$, the state
of the embedded run from $\langle s, p, T, P, o\rangle$ after $u$ steps equals the state of $M$'s
run from $\langle s, p, T \circ \sigma, P \circ \sigma, \varepsilon\rangle$ after $u$ steps.
\end{theorem}

\begin{theorem}[A released call in normal form]
\lean{Turing.NTimeDiag.qy_call}
\label{Turing.NTimeDiag.qy_call}
\leanok
\uses{Turing.Cfg, Turing.Cfg.mapState, Turing.Cfg.ofWords, Turing.Cfg.ofWords_inputPos, Turing.Cfg.ofWords_output, Turing.Cfg.ofWords_state, Turing.Cfg.ofWords_workTapePos, Turing.Cfg.ofWords_workTapes, Turing.FinTM, Turing.FinTM.bufferTape, Turing.MultiTapeTM.runFrom, Turing.MultiTapeTM.runFrom_zero, Turing.NTimeDiag.qyC, Turing.seamReleaseTM, Turing.seamReleaseTM_runFrom_pos, Turing.stateWord}
\difficulty{4}
Let $C$ be a finite deterministic $k$-tape machine over bits with states $\mathit{entry}$ and
$\mathit{exit}$, and let $x, a, b$ be binary words. For a state $q$ and a word $w$, the seam
configuration $\mathrm{seam}(q, w)$ on input $x$ has state $q$, input head at $1$, tape $0$
holding $w$ from cell $0$, all other tapes blank, all heads at $0$ and empty output. Suppose that
for some $t > 0$ the run of $C$ from $\mathrm{seam}(\mathit{entry}, a)$ is not in $\mathit{exit}$
at any time $t'$ with $0 < t' < t$ and equals $\mathrm{seam}(\mathit{exit}, b)$ at time $t$.
Consider the release adapter of $C$ at $\mathit{entry}$: a machine with a fresh start state that
performs $\mathit{entry}$'s action and then runs a tagged copy of $C$. Started in its fresh state
with the tapes, heads and output of $\mathrm{seam}(\mathit{entry}, a)$, it is at time $t$ in the
tagged copy of $\mathit{exit}$ with the tapes, heads and output of $\mathrm{seam}(\mathit{exit}, b)$,
and it is not in that state at any earlier time.
\end{theorem}

\begin{theorem}[A hosted released call]
\lean{Turing.NTimeDiag.qy_callHost}
\label{Turing.NTimeDiag.qy_callHost}
\leanok
\uses{Turing.Cfg.ofWords, Turing.FinTM, Turing.FinTM.bufferTape, Turing.MultiTapeTM.runFrom, Turing.NTimeDiag.qyC, Turing.NTimeDiag.qy_call, Turing.NTimeDiag.qy_embed, Turing.NTimeDiag.qy_embed_state, Turing.embedEmitTM, Turing.seamReleaseTM, Turing.stateWord, Turing.stateWord_eq}
\difficulty{4}
Let $C$ be a finite deterministic machine over bits with $k > 0$ work tapes and states
$\mathit{entry}$, $\mathit{exit}$, and let $x, a, b$ be binary words. Write
$\mathrm{seam}(q, w)$ for the configuration on input $x$ with state $q$, input head at $1$, tape
$0$ holding $w$, other tapes blank, heads at $0$ and empty output. Suppose that for some $t > 0$
the run of $C$ from $\mathrm{seam}(\mathit{entry}, a)$ avoids $\mathit{exit}$ at all times
strictly between $0$ and $t$ and equals $\mathrm{seam}(\mathit{exit}, b)$ at time $t$. Let
$\sigma$ inject the $k$ tapes into $K$ host tapes, and run the release adapter of $C$ at
$\mathit{entry}$ (a fresh start state performing $\mathit{entry}$'s action, then a tagged copy of
$C$) embedded along $\sigma$. Start from the host configuration in the fresh state with input
head $1$, arbitrary output $o$, tapes $T$ and heads $P$ such that tape $\sigma(i)$ holds tape $i$
of $\mathrm{seam}(\mathit{entry}, a)$ and head $\sigma(i)$ is at $0$. Then at time $t$ the host is
in the tagged copy of $\mathit{exit}$ with input head $1$, tape $\sigma(0)$ holding exactly $b$,
and all other tapes, all heads and the output unchanged; and it is not in that state at any
earlier time.
\end{theorem}

\begin{theorem}[One silent step in normal form]
\lean{Turing.NTimeDiag.qy_step}
\label{Turing.NTimeDiag.qy_step}
\leanok
\uses{Turing.Action.apply_workTapes, Turing.Cfg, Turing.Cfg.inputSymbol, Turing.Cfg.workTapeSymbols, Turing.Cfg.workTapeSymbols_eq, Turing.MultiTapeTM, Turing.MultiTapeTM.step, Turing.MultiTapeTM.step_of_some, Turing.NTimeDiag.qyC, Turing.moveInputPos_zero}
\difficulty{3}
Let $M$ be a $K$-tape machine with state set $S$, and write $\langle s, p, T, P, o\rangle$ for the
live configuration on input $x$ with state $s$, input head $p$, work tapes $T$, heads $P$ and
output $o$. Suppose that at state $s$, reading the input symbol under $p$ and the work symbols
$T_j(P_j)$, the transition of $M$ writes nothing, leaves the input head in place, moves each work
head $j$ by $\mu_j \in \{-1, 0, 1\}$, emits the optional bit $e$ and enters state $s'$. Then one
step of $M$ from $\langle s, p, T, P, o\rangle$ yields
$\langle s', p, T, (P_j + \mu_j)_j, o\,e\rangle$, where $o\,e$ is $o$ followed by the emitted bit
if any.
\end{theorem}

\begin{definition}[States of the length comparison]
\lean{Turing.NTimeDiag.QyCmp}
\label{Turing.NTimeDiag.QyCmp}
\leanok
The control states of the length-comparison machine: a forward scanning state, a return state
for each verdict in $\{<, =, >\}$, and a live exit state for each verdict.
\end{definition}

\begin{definition}[Length-comparison machine]
\lean{Turing.NTimeDiag.qyCmpTM}
\label{Turing.NTimeDiag.qyCmpTM}
\leanok
\uses{Turing.MultiTapeTM, Turing.NTimeDiag.QyCmp}
For $K \in \mathbb{N}$ and two work tapes $a, b$ among $K$, the length-comparison machine starts
with the heads of $a$ and $b$ on the first cells of two words. In the scanning state it moves
both heads right in lockstep while both read a symbol; at the first blank under either head it
records the verdict ($<$ if only $a$ reads a blank, $=$ if both do, $>$ if only $b$ does) and
moves both heads left in lockstep until tape $a$ reads a blank, then moves them right once and
enters the exit state of the verdict, where it idles. It never writes, never moves the input head
and never emits.
\end{definition}

\begin{definition}[Lockstep head vector]
\lean{Turing.NTimeDiag.qyAt2}
\label{Turing.NTimeDiag.qyAt2}
\leanok
For a head-position vector $P$ over $K$ work tapes, two tapes $a, b$ and an integer $i$, the
lockstep head vector is $P$ with the heads of $a$ and $b$ both set to $i$.
\end{definition}

\begin{theorem}[Moving both compared heads]
\lean{Turing.NTimeDiag.qyAt2_move}
\label{Turing.NTimeDiag.qyAt2_move}
\leanok
\uses{Turing.NTimeDiag.qyAt2}
\difficulty{3}
Let $P$ be a head-position vector over $K$ work tapes, $a, b$ two tapes, $i \in \mathbb{Z}$, and
$d \in \{-1, 0, 1\}$ with integer value $d_{\mathbb{Z}}$. Let $P^{a,b}_i$ be $P$ with the heads
of $a$ and $b$ set to $i$. Moving the heads of $a$ and $b$ of $P^{a,b}_i$ by $d$ and leaving all
other heads in place gives $P^{a,b}_{i + d_{\mathbb{Z}}}$.
\end{theorem}

\begin{theorem}[Run of the length comparison]
\lean{Turing.NTimeDiag.qy_cmp_run}
\label{Turing.NTimeDiag.qy_cmp_run}
\leanok
\uses{Turing.FinTM.bufferTape, Turing.FinTM.bufferTape_left, Turing.FinTM.bufferTape_nat, Turing.MultiTapeTM.runFrom, Turing.MultiTapeTM.runFrom_succ_eq_step', Turing.MultiTapeTM.runFrom_zero, Turing.NTimeDiag.qyAt2, Turing.NTimeDiag.qyAt2_move, Turing.NTimeDiag.qyC, Turing.NTimeDiag.qyCmpTM, Turing.NTimeDiag.qy_step}
\difficulty{6}
Let $a \ne b$ be two of $K$ work tapes, and consider the length-comparison machine on $a, b$
(lockstep scan right while both heads read symbols, then lockstep return, then the exit state of
the verdict). Write $\mathrm{buf}(w)$ for the tape holding $w$ in cells $0, \dots, |w| - 1$ and
blank elsewhere, and $\langle s, p, T, P, o\rangle$ for the live configuration on input $x$ with
state $s$, input head $p$, work tapes $T$, heads $P$ and output $o$. If tape $a$ is
$\mathrm{buf}(u)$ and tape $b$ is $\mathrm{buf}(v)$, both heads at $0$, then from the scanning
state the machine reaches, after exactly $2\min(|u|, |v|) + 2$ steps, the exit state for the
comparison of $|u|$ with $|v|$, with input head, tapes, heads and output unchanged; and at no
earlier time is it in any exit state.
\end{theorem}

\begin{definition}[States of the emitter]
\lean{Turing.NTimeDiag.QyEm}
\label{Turing.NTimeDiag.QyEm}
\leanok
For natural numbers $n_p, n_s, r$, the control states of the emitter are: a prefix state for each
position $0, \dots, n_p$; a reading state; a holding state for each bit; a copying state for each
bit and each count $0, \dots, r$; a suffix state for each position $0, \dots, n_s$; and an exit
state.
\end{definition}

\begin{definition}[The emitter machine]
\lean{Turing.NTimeDiag.qyEmitTM}
\label{Turing.NTimeDiag.qyEmitTM}
\leanok
\uses{Turing.MultiTapeTM, Turing.NTimeDiag.QyEm}
For $K \in \mathbb{N}$, a work tape $i$, fixed binary words $\mathit{pre}$ and $\mathit{suf}$,
$r \in \mathbb{N}$ and a flag $\delta$, the emitter is the $K$-tape machine that first emits
$\mathit{pre}$; then scans the word under the head of tape $i$, emitting each symbol $r + 1$
times, except that when $\delta$ is set it omits the last symbol (detected by one cell of
look-ahead); then emits $\mathit{suf}$ and enters its exit state, whose action halts. It never
writes and never moves the input head.
\end{definition}

\begin{definition}[Body emitted by the emitter]
\lean{Turing.NTimeDiag.qyBody}
\label{Turing.NTimeDiag.qyBody}
\leanok
For a flag $\delta$, $r \in \mathbb{N}$ and a binary word $w$, let $w'$ be $w$ with its last
symbol removed if $\delta$ is set, and $w$ otherwise. The emitter body is the word obtained from
$w'$ by replacing each symbol $c$ with $c^{r+1}$.
\end{definition}

\begin{theorem}[Zero displacement leaves heads in place]
\lean{Turing.NTimeDiag.qy_noMove}
\label{Turing.NTimeDiag.qy_noMove}
\leanok
\difficulty{1}
For every head-position vector $P$ over $K$ work tapes, moving every head by $0$ yields $P$.
\end{theorem}

\begin{theorem}[Moving one head right]
\lean{Turing.NTimeDiag.qy_moveOne}
\label{Turing.NTimeDiag.qy_moveOne}
\leanok
\difficulty{3}
For every head-position vector $P$ over $K$ work tapes and every tape $i$, moving head $i$ by
$+1$ and every other head by $0$ yields $P$ with entry $i$ replaced by $P_i + 1$.
\end{theorem}

\begin{theorem}[Two steps of a run]
\lean{Turing.NTimeDiag.qy_run_two}
\label{Turing.NTimeDiag.qy_run_two}
\leanok
\uses{Turing.Cfg, Turing.MultiTapeTM, Turing.MultiTapeTM.runFrom, Turing.MultiTapeTM.runFrom_succ_eq_step', Turing.MultiTapeTM.step}
\difficulty{1}
For every multi-tape machine $M$ and configuration $c$, the run of $M$ from $c$ after two steps
is one step of $M$ applied to the run after one step.
\end{theorem}

\begin{theorem}[Run of the emitter]
\lean{Turing.NTimeDiag.qy_emit_run}
\label{Turing.NTimeDiag.qy_emit_run}
\leanok
\uses{Turing.FinTM.bufferTape, Turing.FinTM.bufferTape_nat, Turing.MultiTapeTM.runFrom, Turing.MultiTapeTM.runFrom_add, Turing.MultiTapeTM.runFrom_one, Turing.MultiTapeTM.runFrom_succ_eq_step', Turing.MultiTapeTM.runFrom_zero, Turing.NTimeDiag.qyBody, Turing.NTimeDiag.qyC, Turing.NTimeDiag.qyEmitTM, Turing.NTimeDiag.qy_moveOne, Turing.NTimeDiag.qy_noMove, Turing.NTimeDiag.qy_run_two, Turing.NTimeDiag.qy_step}
\difficulty{6}
Let $i$ be one of $K$ work tapes, $\mathit{pre}, \mathit{suf}, w$ binary words, $r \in \mathbb{N}$
and $\delta$ a flag. Consider the emitter that emits $\mathit{pre}$, then each symbol of the word
under head $i$ repeated $r + 1$ times (dropping the last symbol if $\delta$ is set), then
$\mathit{suf}$; let $\beta$ be the word it emits in the middle part for $w$. Write
$\langle s, p, T, P, o\rangle$ for the live configuration on input $x$ with state $s$, input head
$p$, work tapes $T$, heads $P$ and output $o$. If tape $i$ holds exactly $w$ from cell $0$ and its
head is at $0$, then there is
$t \le |\mathit{pre}| + |\mathit{suf}| + 3 + |w|(r + 3)$ such that the emitter, started in its
first prefix state, is at time $t$ in its exit state with head $i$ at $|w|$, all tapes, the input
head and all other heads unchanged, and output $o\,\mathit{pre}\,\beta\,\mathit{suf}$; and it is
not in the exit state at any earlier time.
\end{theorem}

\begin{theorem}[A released call first stops at its exit time]
\lean{Turing.NTimeDiag.qy_release_first_stop}
\label{Turing.NTimeDiag.qy_release_first_stop}
\leanok
\uses{Turing.Cfg.Stopped, Turing.Cfg.mapState, Turing.Cfg.ofWords, Turing.Cfg.ofWords_state, Turing.Cfg.stopped_iff, Turing.FinTM, Turing.MultiTapeTM.runFrom, Turing.MultiTapeTM.runFrom_state_ne_none_of_le, Turing.MultiTapeTM.runFrom_zero, Turing.seamReleaseTM, Turing.seamReleaseTM_runFrom_pos, Turing.stateWord}
\difficulty{4}
Let $C$ be a finite deterministic $k$-tape machine over bits with states $\mathit{entry}$ and
$\mathit{exit}$, and let $x, a, b$ be binary words. Write $\mathrm{seam}(q, w)$ for the
configuration on input $x$ with state $q$, input head at $1$, tape $0$ holding $w$, other tapes
blank, heads at $0$ and empty output. Suppose that for some $t > 0$ the run of $C$ from
$\mathrm{seam}(\mathit{entry}, a)$ avoids $\mathit{exit}$ at all times strictly between $0$ and
$t$ and equals $\mathrm{seam}(\mathit{exit}, b)$ at time $t$. Run the release adapter of $C$ at
$\mathit{entry}$ (a fresh start state performing $\mathit{entry}$'s action, then a tagged copy
of $C$) from $\mathrm{seam}(\mathit{entry}, a)$ with its state replaced by the fresh state, and
call a configuration stopped if it is halted or in the tagged copy of $\mathit{exit}$. Then this
run is stopped at time $t$, is not stopped at any time $u < t$, and equals
$\mathrm{seam}(\mathit{exit}, b)$ with tagged state at time $t$.
\end{theorem}

\begin{theorem}[A call is silent before its exit]
\lean{Turing.NTimeDiag.qy_call_output_nil}
\label{Turing.NTimeDiag.qy_call_output_nil}
\leanok
\uses{Turing.Cfg.ofWords, Turing.Cfg.ofWords_output, Turing.FinTM, Turing.MultiTapeTM.output_prefix, Turing.MultiTapeTM.runFrom, Turing.stateWord}
\difficulty{2}
Let $C$ be a finite deterministic $k$-tape machine over bits with states $\mathit{entry}$ and
$\mathit{exit}$, and let $x, a, b$ be binary words. Write $\mathrm{seam}(q, w)$ for the
configuration on input $x$ with state $q$, input head at $1$, tape $0$ holding $w$, other tapes
blank, heads at $0$ and empty output. If the run of $C$ from $\mathrm{seam}(\mathit{entry}, a)$
equals $\mathrm{seam}(\mathit{exit}, b)$ at time $t$, then its output is empty at every time
$v \le t$.
\end{theorem}

\begin{theorem}[The clocked released call]
\lean{Turing.NTimeDiag.qy_clocked_call}
\label{Turing.NTimeDiag.qy_clocked_call}
\leanok
\uses{Turing.Cfg, Turing.Cfg.Stopped, Turing.Cfg.mapState, Turing.Cfg.ofWords, Turing.Cfg.ofWords_output, Turing.Cfg.ofWords_state, Turing.FinTM, Turing.FinTM.bufferTape, Turing.MultiTapeTM, Turing.MultiTapeTM.runFrom, Turing.MultiTapeTM.runFrom_zero, Turing.NTimeDiag.qy_call_output_nil, Turing.NTimeDiag.qy_ones_val, Turing.NTimeDiag.qy_release_first_stop, Turing.clockCfg, Turing.exists_stepClockTMVia, Turing.seamReleaseTM, Turing.seamReleaseTM_runFrom_pos, Turing.stateWord, Turing.wordPatch}
\difficulty{5}
Let $C$ be a finite deterministic $k$-tape machine over bits with states $\mathit{entry}$ and
$\mathit{exit}$. Write $\mathrm{seam}(q, w)$ for the configuration on a given input with state
$q$, input head at $1$, tape $0$ holding $w$, other tapes blank, heads at $0$ and empty output.
There is a $(k+1)$-tape machine $H$ with finitely many states (a step clock running the release
adapter of $C$ at $\mathit{entry}$ on its first $k$ tapes, with a binary counter on its last tape),
a start-state map, and injective families of finish and timeout states, no finish state equal to
a timeout state, with the following property. Let $x, a, b$ be binary words and $t > 0$ such that
the run of $C$ on input $x$ from $\mathrm{seam}(\mathit{entry}, a)$ avoids $\mathit{exit}$
strictly between times $0$ and $t$ and equals $\mathrm{seam}(\mathit{exit}, b)$ at time $t$. Let
$m \in \mathbb{N}$ and start $H$ with $\mathrm{seam}(\mathit{entry}, a)$ (fresh state) on the body
tapes and the all-ones word $1^m$ under the counter head. If $t \le 2^m - 1$, then within
$5t + 2m + 2$ steps $H$ reaches its finish state for the tagged $\mathit{exit}$, with
$\mathrm{seam}(\mathit{exit}, b)$ on the body tapes and the counter cells overwritten by some
word of length $m$. If $t > 2^m - 1$, then within $5(2^m - 1) + 2m + 2$ steps $H$ reaches some
timeout state with empty output. In both cases no finish or timeout state occurs earlier.
\end{theorem}

\begin{theorem}[Patching over an all-ones word]
\lean{Turing.NTimeDiag.qy_wordPatch_ones}
\label{Turing.NTimeDiag.qy_wordPatch_ones}
\leanok
\uses{Turing.FinTM.bufferTape, Turing.FinTM.bufferTape_eq_none_of_not_mem, Turing.wordPatch, Turing.wordPatch_of_mem, Turing.wordPatch_of_not_mem}
\difficulty{3}
Let $m \in \mathbb{N}$ and let $w$ be a binary word of length $m$. Writing $\mathrm{buf}(v)$ for
the tape holding $v$ in cells $0, \dots, |v| - 1$ and blank elsewhere, overwriting cells
$0, \dots, m - 1$ of $\mathrm{buf}(1^m)$ with $w$ yields $\mathrm{buf}(w)$.
\end{theorem}

\begin{definition}[Specification of a clocked call]
\lean{Turing.NTimeDiag.QyClockSpec}
\label{Turing.NTimeDiag.QyClockSpec}
\leanok
\uses{Turing.Cfg, Turing.Cfg.mapState, Turing.Cfg.ofWords, Turing.FinTM, Turing.FinTM.bufferTape, Turing.MultiTapeTM, Turing.MultiTapeTM.runFrom, Turing.clockCfg, Turing.stateWord, Turing.wordPatch}
Let $C$ be a finite deterministic $k$-tape machine over bits with states $\mathit{entry}$,
$\mathit{exit}$, and let $H$ be a $(k+1)$-tape machine with a start-state map and families of
finish and timeout states. Write $\mathrm{seam}(q, w)$ for the configuration with state $q$,
input head at $1$, tape $0$ holding $w$, other tapes blank, heads at $0$ and empty output. The
clocked-call specification states: for all binary words $x, a, b$ and every $t > 0$ such that the
run of $C$ on input $x$ from $\mathrm{seam}(\mathit{entry}, a)$ avoids $\mathit{exit}$ strictly
between times $0$ and $t$ and equals $\mathrm{seam}(\mathit{exit}, b)$ at time $t$, and for every
$m$ and counter tape holding $1^m$ around its head (cells $-1$ to $m$), when $H$ is started with
the fresh-state copy of $\mathrm{seam}(\mathit{entry}, a)$ on its body tapes and that counter on
its last tape: if $t \le 2^m - 1$, then $H$ reaches within $5t + 2m + 2$ steps its finish state
for the tagged $\mathit{exit}$ with the tagged $\mathrm{seam}(\mathit{exit}, b)$ on its body tapes
and the counter overwritten from its head by a word of length $m$; if $t > 2^m - 1$, then $H$
reaches within $5(2^m - 1) + 2m + 2$ steps a timeout state with empty output; and in both cases
no finish or timeout state occurs earlier.
\end{definition}

\begin{theorem}[Start of the clocked call in normal form]
\lean{Turing.NTimeDiag.qy_clock_start}
\label{Turing.NTimeDiag.qy_clock_start}
\leanok
\uses{Turing.Cfg, Turing.Cfg.mapState, Turing.Cfg.ofWords, Turing.Cfg.ofWords_inputPos, Turing.Cfg.ofWords_output, Turing.Cfg.ofWords_workTapePos, Turing.Cfg.ofWords_workTapes, Turing.FinTM, Turing.FinTM.bufferTape, Turing.NTimeDiag.qyC, Turing.clockCfg, Turing.stateWord}
\difficulty{4}
Let $C$ be a finite deterministic $k$-tape machine over bits with a state $\mathit{entry}$, let
$q$ be a host state, and let $\sigma$ inject $k + 1$ tapes into $K$ host tapes, with
$\sigma(i) = \sigma_c(i)$ for $i < k$ and $\sigma(k) = c_t$. Let $T, P$ be host tapes and heads
such that, for each $i < k$, tape $\sigma_c(i)$ holds $a$ if $i = 0$ and is blank otherwise, with
head at $0$, and head $c_t$ is at $0$. Then the $(k+1)$-tape configuration with state $q$, input
head $1$, tapes $T \circ \sigma$, heads $P \circ \sigma$ and empty output equals the step-clock
layout with host state $q$, body configuration the seam configuration of $C$ at $\mathit{entry}$
with tape $0$ holding $a$ (its state renamed to the fresh state), and counter tape $T(c_t)$ with
head at $0$.
\end{theorem}

\begin{theorem}[Hosted clocked call that finishes]
\lean{Turing.NTimeDiag.qy_clockHost_finish}
\label{Turing.NTimeDiag.qy_clockHost_finish}
\leanok
\uses{Turing.Cfg.ofWords, Turing.Cfg.ofWords_inputPos, Turing.Cfg.ofWords_output, Turing.Cfg.ofWords_workTapePos, Turing.Cfg.ofWords_workTapes, Turing.FinTM, Turing.FinTM.bufferTape, Turing.MultiTapeTM, Turing.MultiTapeTM.runFrom, Turing.NTimeDiag.QyClockSpec, Turing.NTimeDiag.qyC, Turing.NTimeDiag.qy_clock_start, Turing.NTimeDiag.qy_embed, Turing.NTimeDiag.qy_embed_state, Turing.NTimeDiag.qy_wordPatch_ones, Turing.embedEmitTM, Turing.stateWord, Turing.stateWord_eq}
\difficulty{5}
Let $C$ be a finite deterministic machine over bits with $k > 0$ work tapes and states
$\mathit{entry}$, $\mathit{exit}$, and let $H$ be a $(k+1)$-tape step clock for $C$ satisfying the
clocked-call specification (finish within $5t + 2m + 2$ steps when the call time $t$ is at most
$2^m - 1$, timeout otherwise, no earlier finish or timeout). Let $\sigma$ inject $k + 1$ tapes
into $K$ host tapes, with body tapes $\sigma_c(i) = \sigma(i)$ for $i < k$ and counter tape
$c_t = \sigma(k)$. Let $x, a, b$ be binary words and $t > 0$ with the run of $C$ on input $x$ from
the seam configuration at $\mathit{entry}$ with argument $a$ avoiding $\mathit{exit}$ strictly
between times $0$ and $t$ and reaching the seam configuration at $\mathit{exit}$ with result $b$
at time $t$, and suppose $t \le 2^m - 1$. Start $H$ embedded along $\sigma$ in its start state
with input head $1$, output $o$, body tapes holding the argument seam (tape $\sigma_c(0)$ holds
$a$, others blank) with heads at $0$, and counter tape holding $1^m$ with head at $0$. Then for
some $T_1 \le 5t + 2m + 2$ and some word $w'$ of length $m$, at time $T_1$ the host is in the
finish state for the tagged $\mathit{exit}$, with input head $1$, tape $\sigma_c(0)$ holding $b$,
the counter tape holding $w'$, all other tapes, all heads and the output unchanged; and no finish
or timeout state occurs before $T_1$.
\end{theorem}

\begin{theorem}[Hosted clocked call that times out]
\lean{Turing.NTimeDiag.qy_clockHost_timeout}
\label{Turing.NTimeDiag.qy_clockHost_timeout}
\leanok
\uses{Turing.Cfg.mapState, Turing.Cfg.ofWords, Turing.FinTM, Turing.FinTM.bufferTape, Turing.MultiTapeTM, Turing.MultiTapeTM.runFrom, Turing.NTimeDiag.QyClockSpec, Turing.NTimeDiag.qyC, Turing.NTimeDiag.qy_clock_start, Turing.NTimeDiag.qy_embed_gen, Turing.NTimeDiag.qy_embed_state, Turing.clockCfg, Turing.embedEmitCfg, Turing.embedEmitCfg_unselected, Turing.embedEmitTM, Turing.stateWord}
\difficulty{4}
Let $C$ be a finite deterministic $k$-tape machine over bits with states $\mathit{entry}$,
$\mathit{exit}$, and let $H$ be a $(k+1)$-tape step clock for $C$ satisfying the clocked-call
specification (finish within $5t + 2m + 2$ steps when the call time $t$ is at most $2^m - 1$,
silent timeout within $5(2^m - 1) + 2m + 2$ steps otherwise, no earlier finish or timeout). Let
$\sigma$ inject $k + 1$ tapes into $K$ host tapes, with body tapes $\sigma_c(i) = \sigma(i)$ for
$i < k$ and counter tape $c_t = \sigma(k)$. Let $x, a, b$ be binary words and $t > 0$ with the run
of $C$ on input $x$ from the seam configuration at $\mathit{entry}$ with argument $a$ avoiding
$\mathit{exit}$ strictly between times $0$ and $t$ and reaching the seam configuration at
$\mathit{exit}$ with result $b$ at time $t$, and suppose $2^m - 1 < t$. Start $H$ embedded along
$\sigma$ in its start state with input head $1$, output $o$, body tapes holding the argument seam
with heads at $0$, and counter tape holding $1^m$ with head at $0$. Then for some
$T_1 \le 5(2^m - 1) + 2m + 2$, at time $T_1$ the host is in some timeout state with output still
$o$ and every tape and head outside the image of $\sigma$ unchanged; and no finish or timeout
state occurs before $T_1$.
\end{theorem}

\begin{theorem}[Start of the expansion in normal form]
\lean{Turing.NTimeDiag.qy_expand_start}
\label{Turing.NTimeDiag.qy_expand_start}
\leanok
\uses{Turing.Cfg, Turing.CounterLoopState, Turing.NTimeDiag.qyC, Turing.counterLoopCfg, Turing.counter_cfg_eq}
\difficulty{3}
Let $\sigma$ inject two tapes into $K$ host tapes with $\sigma(0) = \mathrm{dst}$ and
$\sigma(1) = \mathrm{ctr}$, let $p$ be an input-head position and $T, P$ host tapes and heads.
The two-tape configuration in the counter loop's initial decrement state, with input head $p$,
tapes $T \circ \sigma$, heads $P \circ \sigma$ and empty output, equals the counter-loop layout
in that state whose body configuration has tape $T(\mathrm{dst})$ with head $P(\mathrm{dst})$
and empty output, and whose counter tape is $T(\mathrm{ctr})$ with head $P(\mathrm{ctr})$.
\end{theorem}

\begin{theorem}[Hosted unary expansion]
\lean{Turing.NTimeDiag.qy_expandHost}
\label{Turing.NTimeDiag.qy_expandHost}
\leanok
\uses{Turing.FinTM.bufferTape, Turing.FinTM.bufferTape_eq, Turing.FinTM.bufferTape_eq_none_of_not_mem, Turing.MultiTapeTM.runFrom, Turing.NTimeDiag.qyC, Turing.NTimeDiag.qy_embed, Turing.NTimeDiag.qy_embed_state, Turing.NTimeDiag.qy_expand_start, Turing.NTimeDiag.qy_patch_end, Turing.counterOverhead, Turing.counter_cfg_eq, Turing.embedEmitTM, Turing.expandTM, Turing.expandTM_run_ofCfg}
\difficulty{5}
Consider the expansion machine with pattern $1$: a counter loop that, given a binary counter of
value $d$, writes $1$ on its body tape $d$ times from the head. Embed it along an injection
$\sigma$ of two tapes into $K$ host tapes with body tape $\sigma(0) = \mathrm{dst}$ and counter
tape $\sigma(1) = \mathrm{ctr}$. Write $\mathrm{buf}(v)$ for the tape holding $v$ in cells
$0, \dots, |v| - 1$ and blank elsewhere, and $\langle s, p, T, P, o\rangle$ for the live
configuration on input $x$ with state $s$, input head $p$, work tapes $T$, heads $P$ and output
$o$. Let $u, w$ be binary words where $w$, read as a binary numeral, has value $d$. If tape
$\mathrm{dst}$ is $\mathrm{buf}(u)$ with head at $|u|$ and tape $\mathrm{ctr}$ is
$\mathrm{buf}(w)$ with head at $0$, then from the initial decrement state, after exactly
$2d + \kappa(w, d + 1)$ steps, where $\kappa(w, r)$ is the total cost of the first $r$ decrements
of the counter $w$, the host is in the loop's exit state with tape $\mathrm{dst}$ equal to
$\mathrm{buf}(u 1^d)$ and its head at $|u| + d$, tape $\mathrm{ctr}$ equal to
$\mathrm{buf}(1^{|w|})$, and everything else unchanged; and it is not in the exit state earlier.
\end{theorem}

\begin{theorem}[Hosted input seek]
\lean{Turing.NTimeDiag.qy_seekHost}
\label{Turing.NTimeDiag.qy_seekHost}
\leanok
\uses{Turing.Cfg, Turing.FinTM.bufferTape, Turing.FinTM.bufferTape_nil, Turing.MultiTapeTM, Turing.MultiTapeTM.runFrom, Turing.NTimeDiag.qyC, Turing.NTimeDiag.qy_embed, Turing.NTimeDiag.qy_embed_state, Turing.embedEmitTM}
\difficulty{4}
Let $R$ be a two-tape machine with states $\mathit{entry}$, $\mathit{exit}$ and a constant $C$
such that, on input $x$, for every target $p \le |x| + 1$ written as a binary word $w$ under the
head of tape $0$ and every configuration $d$ in state $\mathit{entry}$ with tape $1$ blank, $R$
reaches $\mathit{exit}$ first within $C(|w| + \mathrm{pos}(d) + p + 1)$ steps, where
$\mathrm{pos}(d)$ is the input-head position of $d$, with its input head at $p$, its heads within
$C(|w| + 1)$ of their start, and its tapes, heads and output restored. Let $\sigma$ inject the two
tapes into $K$ host tapes with $\sigma(0) = s_0$, $\sigma(1) = s_1$, and write
$\langle s, q, T, P, o\rangle$ for the live configuration on input $x$ with state $s$, input head
$q$, work tapes $T$, heads $P$ and output $o$. If tape $s_0$ holds exactly the word $1$ with head
at $0$ and tape $s_1$ is blank, then for every input-head position $q$ the machine $R$ embedded
along $\sigma$ runs from $\langle \mathit{entry}, q, T, P, o\rangle$ to
$\langle \mathit{exit}, 1, T, P, o\rangle$ within $C(1 + q + 1 + 1)$ steps, without reaching
$\mathit{exit}$ earlier.
\end{theorem}

\begin{theorem}[The emitter halts at its exit]
\lean{Turing.NTimeDiag.qy_emit_halt}
\label{Turing.NTimeDiag.qy_emit_halt}
\leanok
\uses{Turing.MultiTapeTM.runFrom, Turing.MultiTapeTM.runFrom_one, Turing.MultiTapeTM.step_of_some, Turing.NTimeDiag.QyEm, Turing.NTimeDiag.qyC, Turing.NTimeDiag.qyEmitTM}
\difficulty{3}
Let $i$ be one of $K$ work tapes, $\mathit{pre}, \mathit{suf}$ binary words, $r \in \mathbb{N}$
and $\delta$ a flag, and consider the emitter that emits $\mathit{pre}$, the word under head $i$
with each symbol repeated $r + 1$ times, and $\mathit{suf}$. From any live configuration in its
exit state with output $o$, one step of the emitter halts and leaves the output equal to $o$.
\end{theorem}

\begin{definition}[Specification of an install call]
\lean{Turing.NTimeDiag.QySpec}
\label{Turing.NTimeDiag.QySpec}
\leanok
\uses{Turing.Cfg.ofWords, Turing.FinTM, Turing.MultiTapeTM.runFrom, Turing.stateWord}
Let $C$ be a finite deterministic $k$-tape machine over bits with states $\mathit{entry}$ and
$\mathit{exit}$, let $F$ map binary words to binary words, and let $B$ map binary words to
natural numbers. Write $\mathrm{seam}_x(q, w)$ for the configuration on input $x$ with state $q$,
input head at $1$, tape $0$ holding $w$, other tapes blank, heads at $0$ and empty output. The
install-call specification says that for all binary words $x$ and $a$ there is $t$ with
$0 < t \le B(a)$ such that the run of $C$ from $\mathrm{seam}_x(\mathit{entry}, a)$ is not in
$\mathit{exit}$ at any time strictly between $0$ and $t$ and equals
$\mathrm{seam}_x(\mathit{exit}, F(a))$ at time $t$.
\end{definition}

\begin{definition}[Specification of an input call]
\lean{Turing.NTimeDiag.QyInSpec}
\label{Turing.NTimeDiag.QyInSpec}
\leanok
\uses{Turing.Cfg.ofWords, Turing.FinTM, Turing.MultiTapeTM.runFrom, Turing.stateWord}
Let $C$ be a finite deterministic $k$-tape machine over bits with states $\mathit{entry}$ and
$\mathit{exit}$, let $F$ map binary words to binary words, and let
$B : \mathbb{N} \to \mathbb{N}$. Write $\mathrm{seam}_x(q, w)$ for the configuration on input $x$
with state $q$, input head at $1$, tape $0$ holding $w$, other tapes blank, heads at $0$ and empty
output. The input-call specification says that for every binary input $x$ there is $t$ with
$0 < t \le B(|x|)$ such that the run of $C$ from $\mathrm{seam}_x(\mathit{entry}, \varepsilon)$
is not in $\mathit{exit}$ at any time strictly between $0$ and $t$ and equals
$\mathrm{seam}_x(\mathit{exit}, F(x))$ at time $t$.
\end{definition}

\begin{theorem}[Install-call specification is monotone in its bound]
\lean{Turing.NTimeDiag.QySpec.mono}
\label{Turing.NTimeDiag.QySpec.mono}
\leanok
\uses{Turing.FinTM, Turing.NTimeDiag.QySpec}
\difficulty{2}
Let $C$ be a finite deterministic $k$-tape machine over bits with states $\mathit{entry}$ and
$\mathit{exit}$, let $F$ map binary words to binary words, and let $B, B'$ map binary words to
natural numbers with $B(a) \le B'(a)$ for every $a$. If for all binary words $x, a$ the run of $C$
on input $x$ from the seam configuration at $\mathit{entry}$ with argument $a$ first reaches
$\mathit{exit}$ at some positive time $t \le B(a)$, there being the seam configuration with result
$F(a)$, then the same holds with $B'$ in place of $B$.
\end{theorem}

\begin{theorem}[Install calls of function-computing machines]
\lean{Turing.NTimeDiag.qy_installSpec}
\label{Turing.NTimeDiag.qy_installSpec}
\leanok
\uses{Turing.FinTM, Turing.FinTM.ComputesFunInTime, Turing.FinTM.exists_installCallTM, Turing.NTimeDiag.QySpec}
\difficulty{2}
Let $M$ be a finite deterministic machine over bits that computes $F : \{0,1\}^* \to \{0,1\}^*$
in time $T : \mathbb{N} \to \mathbb{N}$. Write $\mathrm{seam}_x(q, w)$ for the configuration on
input $x$ with state $q$, input head at $1$, tape $0$ holding $w$, other tapes blank, heads at
$0$ and empty output. Then there are a finite deterministic machine $C$ over bits with at least one
work tape, states $\mathit{entry}, \mathit{exit}$ of $C$ and a constant $c$ such that for all
binary words $x, a$ there is $t$ with $0 < t \le c\,(T(|a|) + |a| + |F(a)| + 1)$ for which the run
of $C$ from $\mathrm{seam}_x(\mathit{entry}, a)$ avoids $\mathit{exit}$ at all times strictly
between $0$ and $t$ and equals $\mathrm{seam}_x(\mathit{exit}, F(a))$ at time $t$.
\end{theorem}

\begin{theorem}[Input calls of function-computing machines]
\lean{Turing.NTimeDiag.qy_inputSpec}
\label{Turing.NTimeDiag.qy_inputSpec}
\leanok
\uses{Turing.Cfg.ofWords, Turing.FinTM, Turing.FinTM.ComputesFunInTime, Turing.FinTM.exists_inputCallTM, Turing.MultiTapeTM.runFrom, Turing.stateWord, Turing.stateWord_eq}
\difficulty{3}
Let $M$ be a finite deterministic machine over bits that computes $F : \{0,1\}^* \to \{0,1\}^*$
in time $T : \mathbb{N} \to \mathbb{N}$. Write $\mathrm{seam}_x(q, w)$ for the configuration on
input $x$ with state $q$, input head at $1$, tape $0$ holding $w$, other tapes blank, heads at
$0$ and empty output. Then there are a finite deterministic machine $C$ over bits with at least one
work tape, states $\mathit{entry}, \mathit{exit}$ of $C$ and a constant $c$ such that for every
binary input $x$ there is $t$ with $0 < t \le c\,(T(|x|) + |x| + |F(x)| + 1)$ for which the run of
$C$ from $\mathrm{seam}_x(\mathit{entry}, \varepsilon)$ avoids $\mathit{exit}$ at all times
strictly between $0$ and $t$ and equals $\mathrm{seam}_x(\mathit{exit}, F(x))$ at time $t$.
\end{theorem}

\begin{definition}[The kit of the stage-query machine]
\lean{Turing.NTimeDiag.QyKit}
\label{Turing.NTimeDiag.QyKit}
\leanok
\uses{Turing.Cfg, Turing.FinTM, Turing.FinTM.bufferTape, Turing.MultiTapeTM, Turing.MultiTapeTM.runFrom, Turing.NTimeDiag.QyClockSpec, Turing.NTimeDiag.QyInSpec, Turing.NTimeDiag.QySpec}
For $f, g : \mathbb{N} \to \mathbb{N}$, a kit for $f$ and $g$ bundles finite deterministic machines
over bits, each with at least one work tape, entry and exit states, and constants, as follows.
Writing $\mathrm{seam}_x(q, w)$ for the configuration on input $x$ with state $q$, input head $1$,
tape $0$ holding $w$, other tapes blank, heads at $0$ and empty output, an \emph{input call} for
$F$ within $B$ reaches $\mathrm{seam}_x(\mathit{exit}, F(x))$ from
$\mathrm{seam}_x(\mathit{entry}, \varepsilon)$ for the first time at some time in $(0, B(|x|)]$,
and an \emph{install call} for $F$ within $B$ reaches $\mathrm{seam}_x(\mathit{exit}, F(a))$ from
$\mathrm{seam}_x(\mathit{entry}, a)$ for the first time at some time in $(0, B(a)]$. The kit holds:
an input call $C_b$ for $x \mapsto \mathrm{bits}(|x|)$ within $A_b(n + 1)$; an install call $C_f$
for $y \mapsto \mathrm{bits}(f(|y|))$ within
$A_f\,(c_f(f(|y|) + 1) + |y| + |\mathrm{bits}(f(|y|))| + 1)$; a step clock $H$, a $(k_f + 1)$-tape
machine with finitely many states, a start-state map and families of finish and timeout states,
no finish state equal to a timeout state, satisfying the clocked-call specification for $C_f$;
install calls $C_u$ for $y \mapsto 1^{(|y|+1)^4}$ within $A_u(|y|+1)^5$, $C_a$ for
$y \mapsto y\,1$ within $A_a(|y|+1)$, $C_p$ for $y \mapsto 1^{|\mathrm{bits}(A_f(c_f + 2))|}\,y$
within $A_p(|y|+1)$, $C_l$ for $y \mapsto \mathrm{bits}(|y|)$ within $A_l(|y|+1)$, and $C_q$ for
$y \mapsto \mathrm{bits}((|y|+1)^2)$ within $A_q(|y|+1)^3$; an input call $C_g$ for
$x \mapsto \mathrm{bits}(g(|x|))$ within $A_g(g(n) + 1)$; and an input seek, a two-tape machine
$R$ with finitely many states, states $r_E, r_X$ and a constant $r_C$ such that for every input
$x$, every binary word $w$ whose value as a binary numeral (least significant bit first) is
$p \le |x| + 1$, and every configuration $d$ in state $r_E$ whose tape $0$ holds $w$ in the cells
$-1, \dots, |w|$ relative to its head and whose tape $1$ is blank, there is
$t \le r_C(|w| + \mathrm{inputPos}(d) + p + 1)$ such that $R$ from $d$ does not visit $r_X$ before
time $t$, every work head stays within distance $r_C(|w| + 1)$ of its start up to time $t$, and at
time $t$ the machine is in $r_X$ with input head at $p$ and the tapes, work heads and output of $d$.
\end{definition}

\begin{theorem}[Powers of a successor are monotone]
\lean{Turing.NTimeDiag.qy_pow_mono}
\label{Turing.NTimeDiag.qy_pow_mono}
\leanok
\difficulty{1}
For all natural numbers $n, a, b$ with $a \le b$, one has $(n + 1)^a \le (n + 1)^b$.
\end{theorem}

\begin{theorem}[The kit exists]
\lean{Turing.NTimeDiag.qy_kit_exists}
\label{Turing.NTimeDiag.qy_kit_exists}
\leanok
\uses{Complexity.TimeConstructible, Complexity.timeConstructible_iff, Turing.FinTM.ComputesFunInTime, Turing.FinTM.computesFunInTime_appendBit, Turing.FinTM.computesFunInTime_iff, Turing.FinTM.computesFunInTime_lengthBits, Turing.FinTM.computesFunInTime_polyBits, Turing.FinTM.computesFunInTime_polyUnary, Turing.FinTM.computesFunInTime_prepend, Turing.NTimeDiag.QyInSpec, Turing.NTimeDiag.QyKit, Turing.NTimeDiag.QySpec, Turing.NTimeDiag.QySpec.mono, Turing.NTimeDiag.qy_clocked_call, Turing.NTimeDiag.qy_inputSpec, Turing.NTimeDiag.qy_installSpec, Turing.NTimeDiag.qy_pow_mono, Turing.exists_inputSeekTM, Turing.length_bits_le_self}
\difficulty{5}
Let $f, g : \mathbb{N} \to \mathbb{N}$ be time-constructible. Then there exists a kit for $f$ and
$g$, each call with its own constant $A$: input calls computing $\mathrm{bits}(|x|)$ within
$A(n + 1)$ steps and $\mathrm{bits}(g(|x|))$ within $A(g(n) + 1)$ steps on inputs of length $n$;
install calls computing $y\,1$, a fixed all-ones prefix followed by $y$, and $\mathrm{bits}(|y|)$,
each within $A(|y| + 1)$ steps, $1^{(|y|+1)^4}$ within $A(|y| + 1)^5$ steps,
$\mathrm{bits}((|y|+1)^2)$ within $A(|y| + 1)^3$ steps, and $\mathrm{bits}(f(|y|))$ within
$A(c(f(|y|) + 1) + |y| + |\mathrm{bits}(f(|y|))| + 1)$ steps for a constant $c$; together with a
step clock over the call of $f$'s witness and an input seek.
\end{theorem}

\begin{definition}[Tape count of the stage-query machine]
\lean{Turing.NTimeDiag.QyKit.K}
\label{Turing.NTimeDiag.QyKit.K}
\leanok
\uses{Turing.NTimeDiag.QyKit}
Let $\kappa$ be a kit for $f, g : \mathbb{N} \to \mathbb{N}$ whose called machines $C_b$ (bits of
the input length), $C_f$ ($f$'s witness), $C_u$ (fourth power), $C_a$ (appending a $1$), $C_p$
(prefixing), $C_l$ (bits of the length), $C_q$ (bits of the square) and $C_g$ ($g$'s witness) have
$k_b, k_f, k_u, k_a, k_p, k_l, k_q, k_g$ work tapes. The tape count of the main machine is
$K_\kappa = 16 + k_b + k_f + k_u + 3k_a + k_p + k_l + k_q + k_g$: sixteen named tapes followed by
one block per call bank.
\end{definition}

\begin{theorem}[The named tapes fit]
\lean{Turing.NTimeDiag.QyKit.le_K}
\label{Turing.NTimeDiag.QyKit.le_K}
\leanok
\uses{Turing.NTimeDiag.QyKit, Turing.NTimeDiag.QyKit.K}
\difficulty{2}
For every kit $\kappa$ for $f, g : \mathbb{N} \to \mathbb{N}$, the tape count of the main machine
satisfies $16 \le K_\kappa$.
\end{theorem}

\begin{definition}[Named tapes]
\lean{Turing.NTimeDiag.QyKit.tp}
\label{Turing.NTimeDiag.QyKit.tp}
\leanok
\uses{Turing.NTimeDiag.QyKit, Turing.NTimeDiag.QyKit.K, Turing.NTimeDiag.QyKit.le_K}
Let $\kappa$ be a kit for $f, g : \mathbb{N} \to \mathbb{N}$ with main-machine tape count
$K_\kappa \ge 16$. For $c \in \{0, \dots, 15\}$ the named tape $\tau_c$ is host tape number $c$
among the $K_\kappa$ tapes. Their roles: $\tau_0$ the unary $n + 1$, $\tau_1$ the width word
$1^{|\mathrm{bits}(n)|}$, $\tau_2$ the rung, $\tau_3$ the next rung, $\tau_4$ the word $1^{X-1}$,
$\tau_5$ its fourth power, $\tau_6$ the next rung's counter, $\tau_7$ $f$'s call tape, $\tau_8$ the
clock's counter, $\tau_9$ its saved copy, $\tau_{10}, \tau_{11}$ the pair $(j, k)$, $\tau_{12},
\tau_{13}$ the input seek's tapes, $\tau_{14}$ $g$'s call tape, and $\tau_{15}$ a tape that stays
blank.
\end{definition}

\begin{theorem}[Named tapes are injective]
\lean{Turing.NTimeDiag.QyKit.tp_inj}
\label{Turing.NTimeDiag.QyKit.tp_inj}
\leanok
\uses{Turing.NTimeDiag.QyKit, Turing.NTimeDiag.QyKit.le_K, Turing.NTimeDiag.QyKit.tp}
\difficulty{1}
Let $\kappa$ be a kit for $f, g : \mathbb{N} \to \mathbb{N}$, and for $c < 16$ let $\tau_c$ be host
tape number $c$ of the main machine. For all $a, b < 16$, $\tau_a = \tau_b$ if and only if $a = b$.
\end{theorem}

\begin{definition}[A call bank]
\lean{Turing.NTimeDiag.qyBank}
\label{Turing.NTimeDiag.qyBank}
\leanok
Let $K, \mathit{off}, k \in \mathbb{N}$ with $\mathit{off} + k \le K$, and let $t_0$ be one of $K$
tapes with $t_0 < \mathit{off}$. The bank with tape $0$ on $t_0$ and block at $\mathit{off}$ is the
injection of $\{0, \dots, k-1\}$ into $\{0, \dots, K-1\}$ sending $0$ to $t_0$ and every $i \ne 0$
to $\mathit{off} + i$.
\end{definition}

\begin{theorem}[A bank's tape zero]
\lean{Turing.NTimeDiag.qyBank_zero}
\label{Turing.NTimeDiag.qyBank_zero}
\leanok
\uses{Turing.NTimeDiag.qyBank}
\difficulty{1}
Let $K, \mathit{off}, k \in \mathbb{N}$ with $\mathit{off} + k \le K$ and $0 < k$, and let $t_0$ be
one of $K$ tapes with $t_0 < \mathit{off}$. The bank that sends $0$ to $t_0$ and each $i \ne 0$ to
$\mathit{off} + i$ maps $0$ to $t_0$.
\end{theorem}

\begin{theorem}[A bank's other tapes]
\lean{Turing.NTimeDiag.qyBank_val}
\label{Turing.NTimeDiag.qyBank_val}
\leanok
\uses{Turing.NTimeDiag.qyBank}
\difficulty{1}
Let $K, \mathit{off}, k \in \mathbb{N}$ with $\mathit{off} + k \le K$, and let $t_0$ be one of $K$
tapes with $t_0 < \mathit{off}$. Consider the bank that sends $0$ to $t_0$ and each nonzero
$i' < k$ to tape number $\mathit{off} + i'$. For every $i < k$ with $i \ne 0$, the index of the
tape it assigns to $i$ is $\mathit{off} + i$.
\end{theorem}

\begin{definition}[Two tapes as a bank]
\lean{Turing.NTimeDiag.qyEmb2}
\label{Turing.NTimeDiag.qyEmb2}
\leanok
For $K \in \mathbb{N}$ and two distinct tapes $a \ne b$ among $K$, this is the injection of
$\{0, 1\}$ into the $K$ tapes sending $0$ to $a$ and $1$ to $b$.
\end{definition}

\begin{theorem}[Index of a named tape]
\lean{Turing.NTimeDiag.QyKit.tp_val}
\label{Turing.NTimeDiag.QyKit.tp_val}
\leanok
\uses{Turing.NTimeDiag.QyKit, Turing.NTimeDiag.QyKit.tp}
\difficulty{1}
Let $\kappa$ be a kit for $f, g : \mathbb{N} \to \mathbb{N}$. For every $c < 16$, the named tape
$\tau_c$ of the main machine has index $c$.
\end{theorem}

\begin{definition}[Bank of the length-bits input call]
\lean{Turing.NTimeDiag.QyKit.bCb}
\label{Turing.NTimeDiag.QyKit.bCb}
\leanok
\uses{Turing.NTimeDiag.QyKit, Turing.NTimeDiag.QyKit.K, Turing.NTimeDiag.QyKit.tp, Turing.NTimeDiag.QyKit.tp_val, Turing.NTimeDiag.qyBank}
Let $\kappa$ be a kit for $f, g : \mathbb{N} \to \mathbb{N}$ whose input call $C_b$ computing
$\mathrm{bits}(|x|)$ has $k_b$ tapes. Its bank is the injection of its tapes into the $K_\kappa$
host tapes sending tape $0$ to the named tape $\tau_1$ and each tape $i \ne 0$ to host tape
$16 + i$.
\end{definition}

\begin{definition}[Bank of $f$'s call]
\lean{Turing.NTimeDiag.QyKit.bCf}
\label{Turing.NTimeDiag.QyKit.bCf}
\leanok
\uses{Turing.NTimeDiag.QyKit, Turing.NTimeDiag.QyKit.K, Turing.NTimeDiag.QyKit.tp, Turing.NTimeDiag.QyKit.tp_val, Turing.NTimeDiag.qyBank}
Let $\kappa$ be a kit for $f, g : \mathbb{N} \to \mathbb{N}$ whose calls $C_b$ and $C_f$ (the
latter computing $\mathrm{bits}(f(|y|))$) have $k_b$ and $k_f$ tapes. The bank of $C_f$ sends its
tape $0$ to the named tape $\tau_7$ and each tape $i \ne 0$ to host tape $16 + k_b + i$.
\end{definition}

\begin{definition}[Bank of the fourth-power call]
\lean{Turing.NTimeDiag.QyKit.bPu}
\label{Turing.NTimeDiag.QyKit.bPu}
\leanok
\uses{Turing.NTimeDiag.QyKit, Turing.NTimeDiag.QyKit.K, Turing.NTimeDiag.QyKit.tp, Turing.NTimeDiag.QyKit.tp_val, Turing.NTimeDiag.qyBank}
Let $\kappa$ be a kit for $f, g : \mathbb{N} \to \mathbb{N}$ whose calls $C_b$, $C_f$ and $C_u$ (the
last computing $y \mapsto 1^{(|y|+1)^4}$) have $k_b, k_f, k_u$ tapes. The bank of $C_u$ sends its
tape $0$ to the named tape $\tau_5$ and each tape $i \ne 0$ to host tape $16 + k_b + k_f + i$.
\end{definition}

\begin{definition}[Bank of the counter's increment]
\lean{Turing.NTimeDiag.QyKit.bAb0}
\label{Turing.NTimeDiag.QyKit.bAb0}
\leanok
\uses{Turing.NTimeDiag.QyKit, Turing.NTimeDiag.QyKit.K, Turing.NTimeDiag.QyKit.tp, Turing.NTimeDiag.QyKit.tp_val, Turing.NTimeDiag.qyBank}
Let $\kappa$ be a kit for $f, g : \mathbb{N} \to \mathbb{N}$ whose calls $C_b, C_f, C_u$ and $C_a$
(the last computing $y \mapsto y\,1$) have $k_b, k_f, k_u, k_a$ tapes. The first bank of $C_a$
sends its tape $0$ to the named tape $\tau_6$ and each tape $i \ne 0$ to host tape
$16 + k_b + k_f + k_u + i$.
\end{definition}

\begin{definition}[Bank of the increment of $j$]
\lean{Turing.NTimeDiag.QyKit.bAbJ}
\label{Turing.NTimeDiag.QyKit.bAbJ}
\leanok
\uses{Turing.NTimeDiag.QyKit, Turing.NTimeDiag.QyKit.K, Turing.NTimeDiag.QyKit.tp, Turing.NTimeDiag.QyKit.tp_val, Turing.NTimeDiag.qyBank}
Let $\kappa$ be a kit for $f, g : \mathbb{N} \to \mathbb{N}$ whose calls $C_b, C_f, C_u$ and $C_a$
(the last computing $y \mapsto y\,1$) have $k_b, k_f, k_u, k_a$ tapes. The second bank of $C_a$
sends its tape $0$ to the named tape $\tau_{10}$ and each tape $i \ne 0$ to host tape
$16 + k_b + k_f + k_u + k_a + i$.
\end{definition}

\begin{definition}[Bank of the increment of $k$]
\lean{Turing.NTimeDiag.QyKit.bAbK}
\label{Turing.NTimeDiag.QyKit.bAbK}
\leanok
\uses{Turing.NTimeDiag.QyKit, Turing.NTimeDiag.QyKit.K, Turing.NTimeDiag.QyKit.tp, Turing.NTimeDiag.QyKit.tp_val, Turing.NTimeDiag.qyBank}
Let $\kappa$ be a kit for $f, g : \mathbb{N} \to \mathbb{N}$ whose calls $C_b, C_f, C_u$ and $C_a$
(the last computing $y \mapsto y\,1$) have $k_b, k_f, k_u, k_a$ tapes. The third bank of $C_a$
sends its tape $0$ to the named tape $\tau_{11}$ and each tape $i \ne 0$ to host tape
$16 + k_b + k_f + k_u + 2k_a + i$.
\end{definition}

\begin{definition}[Bank of the prefix call]
\lean{Turing.NTimeDiag.QyKit.bPr}
\label{Turing.NTimeDiag.QyKit.bPr}
\leanok
\uses{Turing.NTimeDiag.QyKit, Turing.NTimeDiag.QyKit.K, Turing.NTimeDiag.QyKit.tp, Turing.NTimeDiag.QyKit.tp_val, Turing.NTimeDiag.qyBank}
Let $\kappa$ be a kit for $f, g : \mathbb{N} \to \mathbb{N}$ whose calls $C_b, C_f, C_u, C_a$ and
$C_p$ (the last prefixing a fixed all-ones word) have $k_b, k_f, k_u, k_a, k_p$ tapes. The bank of
$C_p$ sends its tape $0$ to the named tape $\tau_9$ and each tape $i \ne 0$ to host tape
$16 + k_b + k_f + k_u + 3k_a + i$.
\end{definition}

\begin{definition}[Bank of the length call]
\lean{Turing.NTimeDiag.QyKit.bCl}
\label{Turing.NTimeDiag.QyKit.bCl}
\leanok
\uses{Turing.NTimeDiag.QyKit, Turing.NTimeDiag.QyKit.K, Turing.NTimeDiag.QyKit.tp, Turing.NTimeDiag.QyKit.tp_val, Turing.NTimeDiag.qyBank}
Let $\kappa$ be a kit for $f, g : \mathbb{N} \to \mathbb{N}$ whose calls $C_b, C_f, C_u, C_a, C_p$
and $C_l$ (the last computing $y \mapsto \mathrm{bits}(|y|)$) have $k_b, k_f, k_u, k_a, k_p, k_l$
tapes. The bank of $C_l$ sends its tape $0$ to the named tape $\tau_{10}$ and each tape $i \ne 0$
to host tape $16 + k_b + k_f + k_u + 3k_a + k_p + i$.
\end{definition}

\begin{definition}[Bank of the square call]
\lean{Turing.NTimeDiag.QyKit.bPb}
\label{Turing.NTimeDiag.QyKit.bPb}
\leanok
\uses{Turing.NTimeDiag.QyKit, Turing.NTimeDiag.QyKit.K, Turing.NTimeDiag.QyKit.tp, Turing.NTimeDiag.QyKit.tp_val, Turing.NTimeDiag.qyBank}
Let $\kappa$ be a kit for $f, g : \mathbb{N} \to \mathbb{N}$ whose calls $C_b, C_f, C_u, C_a, C_p,
C_l$ and $C_q$ (the last computing $y \mapsto \mathrm{bits}((|y|+1)^2)$) have
$k_b, k_f, k_u, k_a, k_p, k_l, k_q$ tapes. The bank of $C_q$ sends its tape $0$ to the named tape
$\tau_4$ and each tape $i \ne 0$ to host tape $16 + k_b + k_f + k_u + 3k_a + k_p + k_l + i$.
\end{definition}

\begin{definition}[Bank of $g$'s call]
\lean{Turing.NTimeDiag.QyKit.bCg}
\label{Turing.NTimeDiag.QyKit.bCg}
\leanok
\uses{Turing.NTimeDiag.QyKit, Turing.NTimeDiag.QyKit.K, Turing.NTimeDiag.QyKit.tp, Turing.NTimeDiag.QyKit.tp_val, Turing.NTimeDiag.qyBank}
Let $\kappa$ be a kit for $f, g : \mathbb{N} \to \mathbb{N}$ whose calls $C_b, C_f, C_u, C_a, C_p,
C_l, C_q$ and $C_g$ (the last computing $\mathrm{bits}(g(|x|))$) have
$k_b, k_f, k_u, k_a, k_p, k_l, k_q, k_g$ tapes. The bank of $C_g$ sends its tape $0$ to the named
tape $\tau_{14}$ and each tape $i \ne 0$ to host tape
$16 + k_b + k_f + k_u + 3k_a + k_p + k_l + k_q + i$.
\end{definition}

\begin{definition}[The clock's bank]
\lean{Turing.NTimeDiag.QyKit.bClk}
\label{Turing.NTimeDiag.QyKit.bClk}
\leanok
\uses{Turing.NTimeDiag.QyKit, Turing.NTimeDiag.QyKit.K, Turing.NTimeDiag.QyKit.bCf, Turing.NTimeDiag.QyKit.tp, Turing.NTimeDiag.QyKit.tp_inj, Turing.NTimeDiag.QyKit.tp_val, Turing.NTimeDiag.qyBank_val, Turing.NTimeDiag.qyBank_zero}
Let $\kappa$ be a kit for $f, g : \mathbb{N} \to \mathbb{N}$ whose call $C_f$ of $f$'s witness has
$k_f$ tapes. The clock's bank is the injection of $\{0, \dots, k_f\}$ into the main machine's
tapes that sends each $i < k_f$ to the image of $i$ under the bank of $C_f$ (tape $0$ on the named
tape $\tau_7$, the others in $C_f$'s block) and sends $k_f$ to the named tape $\tau_8$, the clock's
counter.
\end{definition}

\begin{theorem}[Distinct names give distinct tapes]
\lean{Turing.NTimeDiag.QyKit.tp_ne}
\label{Turing.NTimeDiag.QyKit.tp_ne}
\leanok
\uses{Turing.NTimeDiag.QyKit, Turing.NTimeDiag.QyKit.tp, Turing.NTimeDiag.QyKit.tp_inj}
\difficulty{1}
Let $\kappa$ be a kit for $f, g : \mathbb{N} \to \mathbb{N}$, and for $c < 16$ let $\tau_c$ be host
tape number $c$ of the main machine. If $a, b < 16$ and $a \ne b$, then $\tau_a \ne \tau_b$.
\end{theorem}

\begin{definition}[Phase labels of the locator]
\lean{Turing.NTimeDiag.QyLK}
\label{Turing.NTimeDiag.QyLK}
\leanok
The finite set of phase labels of the locator's round, in round order: the anchor; the argument
copy $r_1$ and the clocked call; the length test $r_3$; the steps $r_4$--$r_7$ building and
testing the word $1^{X-1}$; the steps $r_8$, the fourth-power call and $r_{10}$ building and
testing its fourth power; the steps $r_{11}$, the increment call and $r_{13}$--$r_{16}$ building
the next rung and comparing it with $n + 1$; the cleanup steps $c_1, \dots, c_7$; the update of the
pair $(j, k)$ (two comparisons $u_1, u_2$, increments of $j$ and of $k$, clearing of $j$ and of
$k$); and three escapes: interior, top and timeout.
\end{definition}

\begin{definition}[State sets of the locator's phases]
\lean{Turing.NTimeDiag.QyLPh}
\label{Turing.NTimeDiag.QyLPh}
\leanok
\uses{Turing.CounterLoopState, Turing.NTimeDiag.QyCmp, Turing.NTimeDiag.QyKit, Turing.NTimeDiag.QyLK, Turing.RewindPhase, Turing.SweepPhase, Turing.WritePhase}
Let $\kappa$ be a kit for $f, g : \mathbb{N} \to \mathbb{N}$. Each locator phase is assigned a
state set: the clocked call gets the state set of the kit's step clock; the comparison phases
$r_3, r_7, r_{10}, r_{16}, u_1, u_2$ get the states of the length-comparison machine; the
expansions $r_4, r_{13}$ get the counter-loop states of the expansion with the one-letter pattern;
the rewinds $r_6, r_{15}$ get the rewind states; the fourth-power call gets a fresh start state
plus the states of the call $C_u$; the three increment phases get a fresh start state plus the
states of the call $C_a$; the write $r_{14}$ gets the constant-writer states for a one-letter word;
and every other phase gets the sweep states of the copy, transfer and clearing machines.
\end{definition}

\begin{definition}[The locator's phase machines]
\lean{Turing.NTimeDiag.qyLocM}
\label{Turing.NTimeDiag.qyLocM}
\leanok
\uses{Turing.MultiTapeTM, Turing.NTimeDiag.QyKit, Turing.NTimeDiag.QyKit.K, Turing.NTimeDiag.QyKit.bAb0, Turing.NTimeDiag.QyKit.bAbJ, Turing.NTimeDiag.QyKit.bAbK, Turing.NTimeDiag.QyKit.bClk, Turing.NTimeDiag.QyKit.bPu, Turing.NTimeDiag.QyKit.tp, Turing.NTimeDiag.QyKit.tp_ne, Turing.NTimeDiag.QyLK, Turing.NTimeDiag.QyLPh, Turing.NTimeDiag.qyCmpTM, Turing.NTimeDiag.qyEmb2, Turing.clearTM, Turing.copyMapTM, Turing.embedEmitTM, Turing.expandTM, Turing.rewindTM, Turing.seamReleaseTM, Turing.transferMapTM, Turing.writeConstTM}
Let $\kappa$ be a kit for $f, g : \mathbb{N} \to \mathbb{N}$, with main-machine tapes
$K_\kappa$ and named tapes $\tau_0, \dots, \tau_{15}$. Each locator phase is assigned a
$K_\kappa$-tape machine over bits; below, a \emph{release} of a call $C$ along a bank is the
release adapter of $C$ at its entry state (a fresh start state performing the entry's action, then
a copy of $C$) embedded along that bank, and an \emph{expansion} into $\tau_a$ from $\tau_b$ is the
counter loop that writes $1$ on $\tau_a$ as many times as the binary value on $\tau_b$. The
anchor and the three escapes clear $\tau_{15}$; $r_1$ copies $\tau_2$ onto $\tau_7$; the clocked
call runs the step clock embedded along the clock's bank; $r_3, r_7, r_{10}, r_{16}$ compare the
lengths on $\tau_7$, $\tau_4$, $\tau_5$, $\tau_3$ with those on $\tau_1, \tau_1, \tau_1, \tau_0$
respectively; $r_4$ expands into $\tau_4$ from $\tau_7$; $r_5$ copies $\tau_2$ onto $\tau_4$; $r_6$
rewinds $\tau_4$; $r_8$ copies $\tau_4$ onto $\tau_5$; the fourth-power phase releases $C_u$ along
its bank; $r_{11}$ copies $\tau_5$ onto $\tau_6$ mapping every letter to $0$; the increment phase
releases $C_a$ along its first bank (tape $0$ on $\tau_6$); $r_{13}$ expands into $\tau_3$ from
$\tau_6$; $r_{14}$ appends $1$ to $\tau_3$; $r_{15}$ rewinds $\tau_3$; $c_1$ clears $\tau_2$;
$c_2$ transfers $\tau_3$ onto $\tau_2$; $c_3, \dots, c_6$ clear $\tau_7, \tau_4, \tau_5, \tau_6$;
$c_7$ copies $\tau_9$ onto $\tau_8$; $u_1$ and $u_2$ compare $\tau_{10}$ with $\tau_{11}$; the
increments of $j$ and $k$ release $C_a$ along its banks with tape $0$ on $\tau_{10}$ and on
$\tau_{11}$; and the clearings of $k$ and $j$ clear $\tau_{11}$ and $\tau_{10}$.
\end{definition}

\begin{definition}[The locator's dispatch table]
\lean{Turing.NTimeDiag.qyLocRed}
\label{Turing.NTimeDiag.qyLocRed}
\leanok
\uses{Turing.NTimeDiag.QyKit, Turing.NTimeDiag.QyLK, Turing.NTimeDiag.QyLPh}
Let $\kappa$ be a kit for $f, g : \mathbb{N} \to \mathbb{N}$. The locator's dispatch table is the
partial map from (phase, state) pairs to the start of a next phase given as follows; it is
undefined everywhere else. Every state of the anchor goes to $r_1$; at their exit states,
$r_1 \to$ clocked call (at the clock's start state), $r_4 \to r_5 \to r_6 \to r_7$,
$r_8 \to$ fourth-power call, $r_{11} \to$ increment call, $r_{13} \to r_{14} \to r_{15} \to
r_{16}$, $c_1 \to \dots \to c_7 \to u_1$, and the clearing of $k$ $\to$ anchor, the clearing of
$j$ $\to$ increment of $k$; the clocked call goes to $r_3$ at a finish state and to the timeout
escape at a timeout state; the calls go on at their callee's exit state: fourth power $\to r_{10}$,
increment $\to r_{13}$, increment of $j \to u_2$, increment of $k \to$ anchor. On verdicts: $r_3$
goes to the interior escape on $>$ and to $r_4$ otherwise; $r_7$ and $r_{10}$ go to $r_8$ and
$r_{11}$ on $<$ and to the interior escape otherwise; $r_{16}$ goes to $c_1$ on $<$, to the top
escape on $=$ and to the interior escape on $>$; $u_1$ goes to the increment of $j$ on $<$, the
clearing of $j$ on $=$ and the increment of $k$ on $>$; $u_2$ goes to the clearing of $k$ on $=$
and to the anchor otherwise. The escapes are never redirected.
\end{definition}

\begin{definition}[The locator]
\lean{Turing.NTimeDiag.qyLoc}
\label{Turing.NTimeDiag.qyLoc}
\leanok
\uses{Turing.MultiTapeTM, Turing.NTimeDiag.QyKit, Turing.NTimeDiag.QyKit.K, Turing.NTimeDiag.QyLPh, Turing.NTimeDiag.qyLocM, Turing.NTimeDiag.qyLocRed, Turing.phaseHostTM}
Let $\kappa$ be a kit for $f, g : \mathbb{N} \to \mathbb{N}$. The locator is the $K_\kappa$-tape
phase host over the locator's phase machines with the locator's dispatch table, started at the
anchor phase in its initial sweep state: a redirected state takes one stationary silent step to its
target, and any other tagged state $(i, s)$ follows the phase machine of phase $i$.
\end{definition}

\begin{definition}[Phase labels of the setup]
\lean{Turing.NTimeDiag.QySK}
\label{Turing.NTimeDiag.QySK}
\leanok
The ten phase labels $s_1, \dots, s_{10}$ of the setup: computing $\mathrm{bits}(n)$; expanding it to
$1^n$; one more cell; the rewind; the first rung $1^3$; the test $n \ge 3$; the input seek's
$[1]$; the copy of the counter width; its fixed prefix; and the counter's copy.
\end{definition}

\begin{definition}[State sets of the setup's phases]
\lean{Turing.NTimeDiag.QySPh}
\label{Turing.NTimeDiag.QySPh}
\leanok
\uses{Turing.CounterLoopState, Turing.NTimeDiag.QyCmp, Turing.NTimeDiag.QyKit, Turing.NTimeDiag.QySK, Turing.RewindPhase, Turing.SweepPhase, Turing.WritePhase}
Let $\kappa$ be a kit for $f, g : \mathbb{N} \to \mathbb{N}$. The setup phases $s_1, \dots, s_{10}$
get the following state sets: $s_1$ a fresh start state plus the states of the input call $C_b$;
$s_2$ the counter-loop states of the expansion with a one-letter pattern; $s_3$ and $s_7$ the
constant-writer states for a one-letter word; $s_4$ the rewind states; $s_5$ the constant-writer
states for a three-letter word; $s_6$ the length-comparison states; $s_9$ a fresh start state plus
the states of the prefix call $C_p$; and $s_8, s_{10}$ the sweep states.
\end{definition}

\begin{definition}[The setup's phase machines]
\lean{Turing.NTimeDiag.qySetM}
\label{Turing.NTimeDiag.qySetM}
\leanok
\uses{Turing.MultiTapeTM, Turing.NTimeDiag.QyKit, Turing.NTimeDiag.QyKit.K, Turing.NTimeDiag.QyKit.bCb, Turing.NTimeDiag.QyKit.bPr, Turing.NTimeDiag.QyKit.tp, Turing.NTimeDiag.QyKit.tp_ne, Turing.NTimeDiag.QySK, Turing.NTimeDiag.QySPh, Turing.NTimeDiag.qyCmpTM, Turing.NTimeDiag.qyEmb2, Turing.copyMapTM, Turing.embedEmitTM, Turing.expandTM, Turing.rewindTM, Turing.seamReleaseTM, Turing.writeConstTM}
Let $\kappa$ be a kit for $f, g : \mathbb{N} \to \mathbb{N}$, with main-machine tapes $K_\kappa$
and named tapes $\tau_0, \dots, \tau_{15}$. The setup's phase machines are $K_\kappa$-tape machines
over bits: $s_1$ is the release adapter of the input call $C_b$ (computing $\mathrm{bits}(|x|)$)
at its entry, embedded along its bank (tape $0$ on $\tau_1$); $s_2$ is the counter loop writing $1$
on $\tau_0$ as many times as the binary value on $\tau_1$; $s_3$ appends $1$ to $\tau_0$; $s_4$
rewinds $\tau_0$; $s_5$ writes $111$ on $\tau_2$; $s_6$ compares the lengths on $\tau_0$ and
$\tau_2$; $s_7$ writes $1$ on $\tau_{12}$; $s_8$ copies $\tau_1$ onto $\tau_9$; $s_9$ is the
release adapter of the prefix call $C_p$ at its entry, embedded along its bank (tape $0$ on
$\tau_9$); and $s_{10}$ copies $\tau_9$ onto $\tau_8$.
\end{definition}

\begin{definition}[The setup's dispatch table]
\lean{Turing.NTimeDiag.qySetRed}
\label{Turing.NTimeDiag.qySetRed}
\leanok
\uses{Turing.NTimeDiag.QyKit, Turing.NTimeDiag.QySK, Turing.NTimeDiag.QySPh}
Let $\kappa$ be a kit for $f, g : \mathbb{N} \to \mathbb{N}$. The setup's dispatch table sends
$s_1$ at the exit state of the input call $C_b$ to the start of $s_2$, sends each of
$s_2, s_3, s_4, s_5, s_7, s_8$ at its exit state to the start of the next phase, sends $s_6$ to the
start of $s_7$ only at the verdict $>$, sends $s_9$ at the exit state of the prefix call $C_p$ to
the start of $s_{10}$, and is undefined on every other state (in particular on all of $s_{10}$
and on the verdicts $<$ and $=$ of $s_6$).
\end{definition}

\begin{definition}[The setup]
\lean{Turing.NTimeDiag.qySet}
\label{Turing.NTimeDiag.qySet}
\leanok
\uses{Turing.MultiTapeTM, Turing.NTimeDiag.QyKit, Turing.NTimeDiag.QyKit.K, Turing.NTimeDiag.QySPh, Turing.NTimeDiag.qySetM, Turing.NTimeDiag.qySetRed, Turing.phaseHostTM}
Let $\kappa$ be a kit for $f, g : \mathbb{N} \to \mathbb{N}$. The setup is the $K_\kappa$-tape phase
host over the setup's phase machines with the setup's dispatch table, started at phase $s_1$ in its
fresh start state.
\end{definition}

\begin{definition}[Phase labels of the top's writer]
\lean{Turing.NTimeDiag.QyTK}
\label{Turing.NTimeDiag.QyTK}
\leanok
The six phase labels $t_1, \dots, t_6$ of the top's writer: the code's argument $1^{j+1}$, its
binary length $\mathrm{bits}(j + 1)$, the budget $\mathrm{bits}(X^2)$, and three emissions.
\end{definition}

\begin{definition}[State sets of the top's phases]
\lean{Turing.NTimeDiag.QyTPh}
\label{Turing.NTimeDiag.QyTPh}
\leanok
\uses{Turing.NTimeDiag.QyEm, Turing.NTimeDiag.QyKit, Turing.NTimeDiag.QyTK}
Let $\kappa$ be a kit for $f, g : \mathbb{N} \to \mathbb{N}$. The top's phases get the state sets:
$t_1, t_2, t_3$ a fresh start state plus the states of the calls $C_a$ (appending a $1$), $C_l$
(bits of the length) and $C_q$ (bits of the square) respectively; $t_4$ the emitter states for a
prefix of length $1$, a suffix of length $4$ and repetition $3$; $t_5$ those for an empty prefix,
a suffix of length $2$ and repetition $1$; and $t_6$ those for empty prefix and suffix and
repetition $0$.
\end{definition}

\begin{definition}[The top's phase machines]
\lean{Turing.NTimeDiag.qyTopM}
\label{Turing.NTimeDiag.qyTopM}
\leanok
\uses{Turing.MultiTapeTM, Turing.NTimeDiag.QyKit, Turing.NTimeDiag.QyKit.K, Turing.NTimeDiag.QyKit.bAbJ, Turing.NTimeDiag.QyKit.bCl, Turing.NTimeDiag.QyKit.bPb, Turing.NTimeDiag.QyKit.tp, Turing.NTimeDiag.QyTK, Turing.NTimeDiag.QyTPh, Turing.NTimeDiag.qyEmitTM, Turing.embedEmitTM, Turing.seamReleaseTM}
Let $\kappa$ be a kit for $f, g : \mathbb{N} \to \mathbb{N}$, with named tapes $\tau_0, \dots,
\tau_{15}$. The top's phase machines are: $t_1$, $t_2$, $t_3$ the release adapters at their entries
of the calls $C_a$, $C_l$, $C_q$, embedded along their banks with tape $0$ on $\tau_{10}$,
$\tau_{10}$ and $\tau_4$ respectively; $t_4$ the emitter on $\tau_4$ that emits $1$, then each
symbol of $\tau_4$'s word four times, then $0011$; $t_5$ the emitter on $\tau_{10}$ that emits each
symbol of $\tau_{10}$'s word except the last twice, then $01$; and $t_6$ the emitter that emits
$\tau_2$'s word.
\end{definition}

\begin{definition}[The top's dispatch table]
\lean{Turing.NTimeDiag.qyTopRed}
\label{Turing.NTimeDiag.qyTopRed}
\leanok
\uses{Turing.NTimeDiag.QyKit, Turing.NTimeDiag.QyTK, Turing.NTimeDiag.QyTPh}
Let $\kappa$ be a kit for $f, g : \mathbb{N} \to \mathbb{N}$. The top's dispatch table sends $t_1$,
$t_2$, $t_3$ at the exit states of the calls $C_a$, $C_l$, $C_q$ to the fresh start of $t_2$, the
fresh start of $t_3$ and the first prefix state of $t_4$, sends $t_4$ and $t_5$ at their emitters'
exit states to the first prefix state of $t_5$ and $t_6$, and is undefined elsewhere; in
particular the exit of $t_6$ halts.
\end{definition}

\begin{definition}[The top's writer]
\lean{Turing.NTimeDiag.qyTop}
\label{Turing.NTimeDiag.qyTop}
\leanok
\uses{Turing.MultiTapeTM, Turing.NTimeDiag.QyKit, Turing.NTimeDiag.QyKit.K, Turing.NTimeDiag.QyTPh, Turing.NTimeDiag.qyTopM, Turing.NTimeDiag.qyTopRed, Turing.phaseHostTM}
Let $\kappa$ be a kit for $f, g : \mathbb{N} \to \mathbb{N}$. The top's writer is the
$K_\kappa$-tape phase host over the top's phase machines with the top's dispatch table, started at
phase $t_1$ in its fresh start state.
\end{definition}

\begin{definition}[Phase labels of the middle's writer]
\lean{Turing.NTimeDiag.QyMK}
\label{Turing.NTimeDiag.QyMK}
\leanok
The seven phase labels $m_0, \dots, m_6$ of the middle's writer: returning the input head home;
the code; the budget $\mathrm{bits}(g(n))$; and three emissions.
\end{definition}

\begin{definition}[State sets of the middle's phases]
\lean{Turing.NTimeDiag.QyMPh}
\label{Turing.NTimeDiag.QyMPh}
\leanok
\uses{Turing.NTimeDiag.QyEm, Turing.NTimeDiag.QyKit, Turing.NTimeDiag.QyMK}
Let $\kappa$ be a kit for $f, g : \mathbb{N} \to \mathbb{N}$. The middle's phases get the state
sets: $m_0$ the states of the kit's input seek; $m_1, m_2, m_3$ a fresh start state plus the
states of the calls $C_a$ (appending a $1$), $C_l$ (bits of the length) and $C_g$ (bits of
$g(|x|)$) respectively; $m_4$ the emitter states for a prefix of length $1$, a suffix of length $4$
and repetition $3$; $m_5$ those for an empty prefix, a suffix of length $2$ and repetition $1$; and
$m_6$ those for empty prefix and suffix and repetition $0$.
\end{definition}

\begin{definition}[The middle's phase machines]
\lean{Turing.NTimeDiag.qyMidM}
\label{Turing.NTimeDiag.qyMidM}
\leanok
\uses{Turing.MultiTapeTM, Turing.NTimeDiag.QyKit, Turing.NTimeDiag.QyKit.K, Turing.NTimeDiag.QyKit.bAbJ, Turing.NTimeDiag.QyKit.bCg, Turing.NTimeDiag.QyKit.bCl, Turing.NTimeDiag.QyKit.tp, Turing.NTimeDiag.QyKit.tp_ne, Turing.NTimeDiag.QyMK, Turing.NTimeDiag.QyMPh, Turing.NTimeDiag.qyEmb2, Turing.NTimeDiag.qyEmitTM, Turing.embedEmitTM, Turing.seamReleaseTM}
Let $\kappa$ be a kit for $f, g : \mathbb{N} \to \mathbb{N}$, with named tapes $\tau_0, \dots,
\tau_{15}$. The middle's phase machines are: $m_0$ the kit's input seek embedded on $\tau_{12}$
and $\tau_{13}$; $m_1$, $m_2$, $m_3$ the release adapters at their entries of the calls $C_a$,
$C_l$, $C_g$, embedded along their banks with tape $0$ on $\tau_{10}$, $\tau_{10}$ and $\tau_{14}$
respectively; $m_4$ the emitter on $\tau_{14}$ that emits $0$, then each symbol of $\tau_{14}$'s
word four times, then $0011$; $m_5$ the emitter on $\tau_{10}$ that emits each symbol of
$\tau_{10}$'s word except the last twice, then $01$; and $m_6$ the emitter that emits $\tau_0$'s
word.
\end{definition}

\begin{definition}[The middle's dispatch table]
\lean{Turing.NTimeDiag.qyMidRed}
\label{Turing.NTimeDiag.qyMidRed}
\leanok
\uses{Turing.NTimeDiag.QyKit, Turing.NTimeDiag.QyMK, Turing.NTimeDiag.QyMPh}
Let $\kappa$ be a kit for $f, g : \mathbb{N} \to \mathbb{N}$. The middle's dispatch table sends
$m_0$ at the seek's exit state to the fresh start of $m_1$; $m_1$, $m_2$, $m_3$ at the exit states
of the calls $C_a$, $C_l$, $C_g$ to the fresh start of $m_2$, the fresh start of $m_3$ and the
first prefix state of $m_4$; $m_4$ and $m_5$ at their emitters' exit states to the first prefix
state of $m_5$ and $m_6$; and is undefined elsewhere, so the exit of $m_6$ halts.
\end{definition}

\begin{definition}[The middle's writer]
\lean{Turing.NTimeDiag.qyMid}
\label{Turing.NTimeDiag.qyMid}
\leanok
\uses{Turing.MultiTapeTM, Turing.NTimeDiag.QyKit, Turing.NTimeDiag.QyKit.K, Turing.NTimeDiag.QyMPh, Turing.NTimeDiag.qyMidM, Turing.NTimeDiag.qyMidRed, Turing.phaseHostTM}
Let $\kappa$ be a kit for $f, g : \mathbb{N} \to \mathbb{N}$. The middle's writer is the
$K_\kappa$-tape phase host over the middle's phase machines with the middle's dispatch table,
started at phase $m_0$ in the input seek's entry state.
\end{definition}

\begin{definition}[Phase labels of the main machine]
\lean{Turing.NTimeDiag.QyMain}
\label{Turing.NTimeDiag.QyMain}
\leanok
The five phase labels of the stage-query machine: the setup, the locator, the top's writer, the
middle's writer, and the empty output.
\end{definition}

\begin{definition}[State sets of the main machine's phases]
\lean{Turing.NTimeDiag.QyMainPh}
\label{Turing.NTimeDiag.QyMainPh}
\leanok
\uses{Turing.NTimeDiag.QyEm, Turing.NTimeDiag.QyKit, Turing.NTimeDiag.QyLPh, Turing.NTimeDiag.QyMPh, Turing.NTimeDiag.QyMain, Turing.NTimeDiag.QySPh, Turing.NTimeDiag.QyTPh}
Let $\kappa$ be a kit for $f, g : \mathbb{N} \to \mathbb{N}$. The main machine's phases get as state
sets the tagged state sets of the setup, the locator, the top's writer and the middle's writer
respectively, and the empty-output phase gets the emitter states for empty prefix and suffix and
repetition $0$.
\end{definition}

\begin{definition}[The main machine's phase machines]
\lean{Turing.NTimeDiag.qyMainM}
\label{Turing.NTimeDiag.qyMainM}
\leanok
\uses{Turing.MultiTapeTM, Turing.NTimeDiag.QyKit, Turing.NTimeDiag.QyKit.K, Turing.NTimeDiag.QyKit.tp, Turing.NTimeDiag.QyMain, Turing.NTimeDiag.QyMainPh, Turing.NTimeDiag.qyEmitTM, Turing.NTimeDiag.qyLoc, Turing.NTimeDiag.qyMid, Turing.NTimeDiag.qySet, Turing.NTimeDiag.qyTop}
Let $\kappa$ be a kit for $f, g : \mathbb{N} \to \mathbb{N}$, with named tapes $\tau_0, \dots,
\tau_{15}$. The main machine's phase machines are the setup, the locator, the top's writer and the
middle's writer for their phases, and, for the empty-output phase, the emitter that emits the
word on $\tau_{15}$ (with empty prefix and suffix and no repetition).
\end{definition}

\begin{definition}[The main machine's dispatch table]
\lean{Turing.NTimeDiag.qyMainRed}
\label{Turing.NTimeDiag.qyMainRed}
\leanok
\uses{Turing.NTimeDiag.QyKit, Turing.NTimeDiag.QyMain, Turing.NTimeDiag.QyMainPh}
Let $\kappa$ be a kit for $f, g : \mathbb{N} \to \mathbb{N}$. The main machine's dispatch table
sends the setup in phase $s_{10}$ at its exit state to the locator's anchor; the setup in phase
$s_6$ at the verdict $<$ or $=$ to the first prefix state of the empty-output emitter; the locator
in any state of its interior or timeout escape to the middle's writer at $m_0$ in the input seek's
entry state; the locator in any state of its top escape to the top's writer at $t_1$ in its fresh
start state; and is undefined on all other states.
\end{definition}

\begin{definition}[The stage-query machine]
\lean{Turing.NTimeDiag.qyM₁}
\label{Turing.NTimeDiag.qyM₁}
\leanok
\uses{Turing.FinTM, Turing.NTimeDiag.QyKit, Turing.NTimeDiag.QyKit.K, Turing.NTimeDiag.QyMainPh, Turing.NTimeDiag.qyMainM, Turing.NTimeDiag.qyMainRed, Turing.phaseHostTM}
Let $\kappa$ be a kit for $f, g : \mathbb{N} \to \mathbb{N}$. The main machine $M_1$ is the finite
deterministic machine over bits with $K_\kappa$ work tapes whose states are the tagged states of
its five phases, given by the phase host over the setup, the locator, the two writers and the
empty-output emitter with the main dispatch table, started at the setup's phase $s_1$ in its fresh
start state.
\end{definition}

\begin{definition}[Width of the clock's counter]
\lean{Turing.NTimeDiag.QyKit.m}
\label{Turing.NTimeDiag.QyKit.m}
\leanok
\uses{Turing.NTimeDiag.QyKit}
Let $\kappa$ be a kit for $f, g : \mathbb{N} \to \mathbb{N}$ whose call of $f$'s witness has
constants $A_f$ and $c_f$. On inputs of length $n$ the clock's counter width is
$m(n) = |\mathrm{bits}(A_f(c_f + 2))| + |\mathrm{bits}(n)|$.
\end{definition}

\begin{definition}[The anchor layout]
\lean{Turing.NTimeDiag.qyLay}
\label{Turing.NTimeDiag.qyLay}
\leanok
\uses{Turing.NTimeDiag.QyKit, Turing.NTimeDiag.QyKit.K, Turing.NTimeDiag.QyKit.m, Turing.NTimeDiag.QyKit.tp}
Let $\kappa$ be a kit for $f, g : \mathbb{N} \to \mathbb{N}$, with named tapes $\tau_0, \dots,
\tau_{15}$ and clock counter width $m(n)$. For $n, a, j, k \in \mathbb{N}$ the anchor layout
assigns to each host tape a binary word: $1^{n+1}$ to $\tau_0$, $1^{|\mathrm{bits}(n)|}$ to
$\tau_1$, $1^a$ to $\tau_2$, $1^j$ to $\tau_{10}$, $1^k$ to $\tau_{11}$, $1$ to $\tau_{12}$,
$1^{m(n)}$ to both $\tau_8$ and $\tau_9$, and the empty word to every other tape.
\end{definition}

\begin{theorem}[Named or unnamed tapes]
\lean{Turing.NTimeDiag.QyKit.tape_cases}
\label{Turing.NTimeDiag.QyKit.tape_cases}
\leanok
\uses{Turing.NTimeDiag.QyKit, Turing.NTimeDiag.QyKit.K, Turing.NTimeDiag.QyKit.tp}
\difficulty{1}
Let $\kappa$ be a kit for $f, g : \mathbb{N} \to \mathbb{N}$, with named tapes $\tau_0, \dots,
\tau_{15}$. Every host tape $t$ of the main machine either equals $\tau_c$ for some $c < 16$ or
differs from $\tau_c$ for every $c < 16$.
\end{theorem}

\begin{theorem}[Bank tapes avoid the named tapes]
\lean{Turing.NTimeDiag.qyBank_ne_tp}
\label{Turing.NTimeDiag.qyBank_ne_tp}
\leanok
\uses{Turing.NTimeDiag.QyKit, Turing.NTimeDiag.QyKit.K, Turing.NTimeDiag.QyKit.tp, Turing.NTimeDiag.QyKit.tp_val, Turing.NTimeDiag.qyBank, Turing.NTimeDiag.qyBank_val}
\difficulty{3}
Let $\kappa$ be a kit for $f, g : \mathbb{N} \to \mathbb{N}$ with main-machine tapes $K_\kappa$ and
named tapes $\tau_0, \dots, \tau_{15}$. Let $t_0$ be a host tape, and $\mathit{off}, k$ natural
numbers with $\mathit{off} + k \le K_\kappa$, $t_0 < \mathit{off}$ and $16 \le \mathit{off}$.
For the bank sending $0$ to $t_0$ and each $i \ne 0$ to $\mathit{off} + i$, every $i < k$ with
$i \ne 0$ and every $c < 16$, the image of $i$ differs from $\tau_c$.
\end{theorem}

\begin{definition}[The setup's exits]
\lean{Turing.NTimeDiag.qySetF}
\label{Turing.NTimeDiag.qySetF}
\leanok
\uses{Turing.NTimeDiag.QyKit, Turing.NTimeDiag.QySPh}
Let $\kappa$ be a kit for $f, g : \mathbb{N} \to \mathbb{N}$. The setup's exit set consists of three
setup states: the long exit, phase $s_{10}$ in its done state, and the two short exits, phase $s_6$
(the test $n \ge 3$) with verdict $<$ and with verdict $=$.
\end{definition}

\begin{definition}[The setup's time bound]
\lean{Turing.NTimeDiag.qySetB}
\label{Turing.NTimeDiag.qySetB}
\leanok
\uses{Turing.NTimeDiag.QyKit}
Let $\kappa$ be a kit for $f, g : \mathbb{N} \to \mathbb{N}$ with input-call constant $A_b$ (for
$\mathrm{bits}(|x|)$), prefix-call constant $A_p$ and constants $A_f, c_f$ of $f$'s call. The
setup's time bound on inputs of length $n$ is
$(A_b + A_p + 2\,|\mathrm{bits}(A_f(c_f + 2))| + 60)(n + 1)$.
\end{definition}

\begin{theorem}[Value of the binary digits]
\lean{Turing.NTimeDiag.qy_bits_val}
\label{Turing.NTimeDiag.qy_bits_val}
\leanok
\uses{Turing.counterValue_bits, Turing.counterValue_eq}
\difficulty{1}
For every $n \in \mathbb{N}$, the binary digits $\mathrm{bits}(n)$ of $n$, read as a binary numeral
with the least significant bit first, have value $n$.
\end{theorem}

\begin{theorem}[A blank bank holds the empty seam]
\lean{Turing.NTimeDiag.qy_blank_bank}
\label{Turing.NTimeDiag.qy_blank_bank}
\leanok
\uses{Turing.FinTM.bufferTape, Turing.stateWord, Turing.stateWord_eq}
\difficulty{3}
Let $\sigma$ inject $k$ tapes into $K$ host tapes and let $T$ assign host tape contents. Write
$\mathrm{buf}(w)$ for the tape holding $w$ in cells $0, \dots, |w| - 1$ and blank elsewhere. If
$T(\sigma(i)) = \mathrm{buf}(\varepsilon)$ for every $i < k$, then for every $i < k$,
$T(\sigma(i)) = \mathrm{buf}(w_i)$, where $(w_i)_{i<k}$ is the seam assignment of the empty argument
(tape $0$ holds $\varepsilon$ and every other tape is blank).
\end{theorem}

\begin{theorem}[A bank holding an argument holds its seam]
\lean{Turing.NTimeDiag.qy_arg_bank}
\label{Turing.NTimeDiag.qy_arg_bank}
\leanok
\uses{Turing.FinTM.bufferTape, Turing.stateWord, Turing.stateWord_eq}
\difficulty{3}
Let $\sigma$ inject $k > 0$ tapes into $K$ host tapes, let $T$ assign host tape contents and let $a$
be a binary word. Write $\mathrm{buf}(w)$ for the tape holding $w$ in cells $0, \dots, |w| - 1$ and
blank elsewhere. If $T(\sigma(0)) = \mathrm{buf}(a)$ and $T(\sigma(i)) = \mathrm{buf}(\varepsilon)$
for every $i \ne 0$, then for every $i < k$, $T(\sigma(i)) = \mathrm{buf}(w_i)$, where $w_0 = a$
and $w_i = \varepsilon$ for $i \ne 0$ (the seam assignment of $a$).
\end{theorem}

\begin{theorem}[Prefix-call bank avoids the named tapes]
\lean{Turing.NTimeDiag.QyKit.bPr_ne}
\label{Turing.NTimeDiag.QyKit.bPr_ne}
\leanok
\uses{Turing.NTimeDiag.QyKit, Turing.NTimeDiag.QyKit.bPr, Turing.NTimeDiag.QyKit.tp, Turing.NTimeDiag.qyBank_ne_tp}
\difficulty{1}
Let $\kappa$ be a kit for $f, g : \mathbb{N} \to \mathbb{N}$, with named tapes $\tau_0, \dots,
\tau_{15}$, and consider the bank of its prefix call $C_p$ (tape $0$ on $\tau_9$, the other tapes
in $C_p$'s block). For every tape $i \ne 0$ of $C_p$ and every $c < 16$, the bank's image of $i$
differs from $\tau_c$.
\end{theorem}

\begin{definition}[The setup's tapes before its test]
\lean{Turing.NTimeDiag.qySetT5}
\label{Turing.NTimeDiag.qySetT5}
\leanok
\uses{Turing.FinTM.bufferTape, Turing.NTimeDiag.QyKit, Turing.NTimeDiag.QyKit.K, Turing.NTimeDiag.QyKit.tp}
Let $\kappa$ be a kit for $f, g : \mathbb{N} \to \mathbb{N}$, with named tapes $\tau_0, \dots,
\tau_{15}$, and let $n \in \mathbb{N}$. Write $\mathrm{buf}(w)$ for the tape holding $w$ in cells
$0, \dots, |w| - 1$ and blank elsewhere. The setup's tape contents before its test assign
$\mathrm{buf}(1^{n+1})$ to $\tau_0$, $\mathrm{buf}(1^{|\mathrm{bits}(n)|})$ to $\tau_1$,
$\mathrm{buf}(111)$ to $\tau_2$, and the blank tape $\mathrm{buf}(\varepsilon)$ to every other host
tape.
\end{definition}

\begin{theorem}[The setup's head]
\lean{Turing.NTimeDiag.qy_setup_head}
\label{Turing.NTimeDiag.qy_setup_head}
\leanok
\uses{Turing.FinTM.bufferTape, Turing.NTimeDiag.QyKit, Turing.NTimeDiag.QyKit.K, Turing.NTimeDiag.QyKit.bCb, Turing.NTimeDiag.QyKit.tp, Turing.NTimeDiag.QyKit.tp_ne, Turing.NTimeDiag.qyC, Turing.NTimeDiag.qyEmb2, Turing.NTimeDiag.qySet, Turing.NTimeDiag.qySetF, Turing.NTimeDiag.qySetM, Turing.NTimeDiag.qySetRed, Turing.NTimeDiag.qySetT5, Turing.NTimeDiag.qy_bits_val, Turing.NTimeDiag.qy_blank_bank, Turing.NTimeDiag.qy_callHost, Turing.NTimeDiag.qy_expandHost, Turing.NTimeDiag.qy_guard_of_cut, Turing.NTimeDiag.qy_hosted, Turing.NTimeDiag.qy_rewind, Turing.NTimeDiag.qy_writeConst, Turing.Segment, Turing.Segment.trans, Turing.counterOverhead_le}
\difficulty{5}
Let $\kappa$ be a kit for $f, g : \mathbb{N} \to \mathbb{N}$ with input-call constant $A_b$, and
let $x$ be a binary input of length $n$. Write $\langle s, p, T, P, o\rangle$ for the live
configuration on input $x$ with state $s$, input head $p$, work tapes $T$, heads $P$ and output
$o$. There is $t \le A_b(n + 1) + 7n + 2|\mathrm{bits}(n)| + 27$ such that the setup machine runs
in exactly $t$ steps from $\langle (s_1, \text{fresh start}), 1, \text{all blank}, 0, \varepsilon
\rangle$ to $\langle (s_6, \text{scanning}), 1, T_5, 0, \varepsilon\rangle$, where $T_5$ holds
$1^{n+1}$ on $\tau_0$, $1^{|\mathrm{bits}(n)|}$ on $\tau_1$, $111$ on $\tau_2$ and is blank
elsewhere, without visiting any of the setup's exit states (phase $s_{10}$ done, or phase $s_6$
with verdict $<$ or $=$) before time $t$.
\end{theorem}

\begin{theorem}[The setup, long path]
\lean{Turing.NTimeDiag.qy_setup_long}
\label{Turing.NTimeDiag.qy_setup_long}
\leanok
\uses{Turing.FinTM.bufferTape, Turing.NTimeDiag.QyKit, Turing.NTimeDiag.QyKit.K, Turing.NTimeDiag.QyKit.bPr, Turing.NTimeDiag.QyKit.bPr_ne, Turing.NTimeDiag.QyKit.m, Turing.NTimeDiag.QyKit.tape_cases, Turing.NTimeDiag.QyKit.tp, Turing.NTimeDiag.QyKit.tp_inj, Turing.NTimeDiag.QyKit.tp_ne, Turing.NTimeDiag.qyC, Turing.NTimeDiag.qyLay, Turing.NTimeDiag.qySet, Turing.NTimeDiag.qySetB, Turing.NTimeDiag.qySetF, Turing.NTimeDiag.qySetM, Turing.NTimeDiag.qySetRed, Turing.NTimeDiag.qySetT5, Turing.NTimeDiag.qy_arg_bank, Turing.NTimeDiag.qy_callHost, Turing.NTimeDiag.qy_cmp_run, Turing.NTimeDiag.qy_copy, Turing.NTimeDiag.qy_guard_of_cut, Turing.NTimeDiag.qy_hosted, Turing.NTimeDiag.qy_hosted', Turing.NTimeDiag.qy_hosted_last, Turing.NTimeDiag.qy_setup_head, Turing.NTimeDiag.qy_writeConst, Turing.Segment, Turing.Segment.trans, Turing.length_bits_le_self}
\difficulty{6}
Let $\kappa$ be a kit for $f, g : \mathbb{N} \to \mathbb{N}$ and let $x$ be a binary input of length
$n \ge 3$. Write $\langle s, p, T, P, o\rangle$ for the live configuration on input $x$ with state
$s$, input head $p$, work tapes $T$, heads $P$ and output $o$, and let $B(n)$ be the setup's time
bound $(A_b + A_p + 2|\mathrm{bits}(A_f(c_f + 2))| + 60)(n + 1)$. There is $t \le B(n)$ such that
the setup machine runs in exactly $t$ steps from
$\langle (s_1, \text{fresh start}), 1, \text{all blank}, 0, \varepsilon\rangle$ to
$\langle (s_{10}, \text{done}), 1, L, 0, \varepsilon\rangle$ without visiting any setup exit
(phase $s_{10}$ done, or phase $s_6$ with verdict $<$ or $=$) before time $t$, where $L$ is the
anchor layout with rung $a = 3$ and pair $(j, k) = (0, 0)$: $1^{n+1}$ on $\tau_0$,
$1^{|\mathrm{bits}(n)|}$ on $\tau_1$, $111$ on $\tau_2$, $1$ on $\tau_{12}$, $1^{m(n)}$ on
$\tau_8$ and $\tau_9$ with $m(n) = |\mathrm{bits}(A_f(c_f + 2))| + |\mathrm{bits}(n)|$, and every
other tape blank.
\end{theorem}

\begin{theorem}[The setup, short path]
\lean{Turing.NTimeDiag.qy_setup_short}
\label{Turing.NTimeDiag.qy_setup_short}
\leanok
\uses{Turing.FinTM.bufferTape, Turing.NTimeDiag.QyKit, Turing.NTimeDiag.QyKit.tp, Turing.NTimeDiag.QyKit.tp_ne, Turing.NTimeDiag.qyC, Turing.NTimeDiag.qySet, Turing.NTimeDiag.qySetB, Turing.NTimeDiag.qySetF, Turing.NTimeDiag.qySetM, Turing.NTimeDiag.qySetRed, Turing.NTimeDiag.qySetT5, Turing.NTimeDiag.qy_cmp_run, Turing.NTimeDiag.qy_hosted_last, Turing.NTimeDiag.qy_setup_head, Turing.Segment, Turing.Segment.trans, Turing.length_bits_le_self}
\difficulty{5}
Let $\kappa$ be a kit for $f, g : \mathbb{N} \to \mathbb{N}$ and let $x$ be a binary input of length
$n \le 2$. Write $\langle s, p, T, P, o\rangle$ for the live configuration on input $x$ with state
$s$, input head $p$, work tapes $T$, heads $P$ and output $o$, and let $B(n)$ be the setup's time
bound $(A_b + A_p + 2|\mathrm{bits}(A_f(c_f + 2))| + 60)(n + 1)$. There are $t \le B(n)$ and a
verdict $v \in \{<, =\}$ such that the setup machine runs in exactly $t$ steps from
$\langle (s_1, \text{fresh start}), 1, \text{all blank}, 0, \varepsilon\rangle$ to
$\langle (s_6, \text{done with } v), 1, T_5, 0, \varepsilon\rangle$, where $T_5$ holds $1^{n+1}$ on
$\tau_0$, $1^{|\mathrm{bits}(n)|}$ on $\tau_1$, $111$ on $\tau_2$ and is blank elsewhere (in
particular $\tau_{15}$ is blank and nothing has been emitted), without visiting any setup exit
before time $t$.
\end{theorem}

\begin{theorem}[Successor of a pair below the diagonal]
\lean{Turing.NTimeDiag.qy_unpair_succ_lt}
\label{Turing.NTimeDiag.qy_unpair_succ_lt}
\leanok
\difficulty{3}
Let $\mathrm{pair}(j, k) = k^2 + j$ if $j < k$ and $j^2 + j + k$ otherwise, a bijection
$\mathbb{N}^2 \to \mathbb{N}$ with inverse $\mathrm{unpair}$. If $\mathrm{unpair}(h) = (j, k)$ and
$j + 1 < k$, then $\mathrm{unpair}(h + 1) = (j + 1, k)$.
\end{theorem}

\begin{theorem}[Successor at the last pair below the diagonal]
\lean{Turing.NTimeDiag.qy_unpair_succ_last}
\label{Turing.NTimeDiag.qy_unpair_succ_last}
\leanok
\difficulty{3}
Let $\mathrm{pair}(j, k) = k^2 + j$ if $j < k$ and $j^2 + j + k$ otherwise, a bijection
$\mathbb{N}^2 \to \mathbb{N}$ with inverse $\mathrm{unpair}$. If $\mathrm{unpair}(h) = (j, k)$ and
$j + 1 = k$, then $\mathrm{unpair}(h + 1) = (k, 0)$.
\end{theorem}

\begin{theorem}[Successor on the diagonal]
\lean{Turing.NTimeDiag.qy_unpair_succ_diag}
\label{Turing.NTimeDiag.qy_unpair_succ_diag}
\leanok
\difficulty{3}
Let $\mathrm{pair}(j, k) = k^2 + j$ if $j < k$ and $j^2 + j + k$ otherwise, a bijection
$\mathbb{N}^2 \to \mathbb{N}$ with inverse $\mathrm{unpair}$. If $\mathrm{unpair}(h) = (j, j)$, then
$\mathrm{unpair}(h + 1) = (0, j + 1)$.
\end{theorem}

\begin{theorem}[Successor above the diagonal]
\lean{Turing.NTimeDiag.qy_unpair_succ_gt}
\label{Turing.NTimeDiag.qy_unpair_succ_gt}
\leanok
\difficulty{3}
Let $\mathrm{pair}(j, k) = k^2 + j$ if $j < k$ and $j^2 + j + k$ otherwise, a bijection
$\mathbb{N}^2 \to \mathbb{N}$ with inverse $\mathrm{unpair}$. If $\mathrm{unpair}(h) = (j, k)$ and
$k < j$, then $\mathrm{unpair}(h + 1) = (j, k + 1)$.
\end{theorem}

\begin{theorem}[The clock's bank on the call's tapes]
\lean{Turing.NTimeDiag.QyKit.bClk_castSucc}
\label{Turing.NTimeDiag.QyKit.bClk_castSucc}
\leanok
\uses{Turing.NTimeDiag.QyKit, Turing.NTimeDiag.QyKit.bCf, Turing.NTimeDiag.QyKit.bClk}
\difficulty{2}
Let $\kappa$ be a kit for $f, g : \mathbb{N} \to \mathbb{N}$ whose call $C_f$ of $f$'s witness has
$k_f$ tapes. For every $i < k_f$, the clock's bank (an injection of $\{0, \dots, k_f\}$ into the
main machine's tapes) maps $i$ to the same host tape as the bank of $C_f$ does.
\end{theorem}

\begin{theorem}[The clock's last tape is the counter tape]
\lean{Turing.NTimeDiag.QyKit.bClk_last}
\label{Turing.NTimeDiag.QyKit.bClk_last}
\leanok
\uses{Turing.NTimeDiag.QyKit, Turing.NTimeDiag.QyKit.bClk, Turing.NTimeDiag.QyKit.tp}
\difficulty{2}
Let $\kappa$ be a kit for $f, g : \mathbb{N} \to \mathbb{N}$ whose call $C_f$ of $f$'s witness has
$k_f$ tapes. The clock's bank, an injection of $\{0, \dots, k_f\}$ into the main machine's tapes,
maps $k_f$ to the named tape $\tau_8$ (host tape number $8$).
\end{theorem}

\begin{theorem}[$f$'s call bank avoids the named tapes]
\lean{Turing.NTimeDiag.QyKit.bCf_ne}
\label{Turing.NTimeDiag.QyKit.bCf_ne}
\leanok
\uses{Turing.NTimeDiag.QyKit, Turing.NTimeDiag.QyKit.bCf, Turing.NTimeDiag.QyKit.tp, Turing.NTimeDiag.qyBank_ne_tp}
\difficulty{1}
Let $\kappa$ be a kit for $f, g : \mathbb{N} \to \mathbb{N}$, with named tapes $\tau_0, \dots,
\tau_{15}$, and consider the bank of its call $C_f$ of $f$'s witness (tape $0$ on $\tau_7$, the
other tapes in $C_f$'s block). For every tape $i \ne 0$ of $C_f$ and every $c < 16$, the bank's
image of $i$ differs from $\tau_c$.
\end{theorem}

\begin{theorem}[Fourth-power bank avoids the named tapes]
\lean{Turing.NTimeDiag.QyKit.bPu_ne}
\label{Turing.NTimeDiag.QyKit.bPu_ne}
\leanok
\uses{Turing.NTimeDiag.QyKit, Turing.NTimeDiag.QyKit.bPu, Turing.NTimeDiag.QyKit.tp, Turing.NTimeDiag.qyBank_ne_tp}
\difficulty{1}
Let $\kappa$ be a kit for $f, g : \mathbb{N} \to \mathbb{N}$, with named tapes $\tau_0, \dots,
\tau_{15}$, and consider the bank of its fourth-power call $C_u$ (tape $0$ on $\tau_5$, the other
tapes in $C_u$'s block). For every tape $i \ne 0$ of $C_u$ and every $c < 16$, the bank's image of
$i$ differs from $\tau_c$.
\end{theorem}

\begin{theorem}[Counter-increment bank avoids the named tapes]
\lean{Turing.NTimeDiag.QyKit.bAb0_ne}
\label{Turing.NTimeDiag.QyKit.bAb0_ne}
\leanok
\uses{Turing.NTimeDiag.QyKit, Turing.NTimeDiag.QyKit.bAb0, Turing.NTimeDiag.QyKit.tp, Turing.NTimeDiag.qyBank_ne_tp}
\difficulty{1}
Let $\kappa$ be a kit for $f, g : \mathbb{N} \to \mathbb{N}$, with named tapes $\tau_0, \dots,
\tau_{15}$, and consider the first bank of its increment call $C_a$ (tape $0$ on $\tau_6$, the
other tapes in its block). For every tape $i \ne 0$ of $C_a$ and every $c < 16$, the bank's image of
$i$ differs from $\tau_c$.
\end{theorem}

\begin{theorem}[Bank of the increment of $j$ avoids the named tapes]
\lean{Turing.NTimeDiag.QyKit.bAbJ_ne}
\label{Turing.NTimeDiag.QyKit.bAbJ_ne}
\leanok
\uses{Turing.NTimeDiag.QyKit, Turing.NTimeDiag.QyKit.bAbJ, Turing.NTimeDiag.QyKit.tp, Turing.NTimeDiag.qyBank_ne_tp}
\difficulty{1}
Let $\kappa$ be a kit for $f, g : \mathbb{N} \to \mathbb{N}$, with named tapes $\tau_0, \dots,
\tau_{15}$, and consider the second bank of its increment call $C_a$ (tape $0$ on $\tau_{10}$, the
other tapes in its block). For every tape $i \ne 0$ of $C_a$ and every $c < 16$, the bank's image of
$i$ differs from $\tau_c$.
\end{theorem}

\begin{theorem}[Bank of the increment of $k$ avoids the named tapes]
\lean{Turing.NTimeDiag.QyKit.bAbK_ne}
\label{Turing.NTimeDiag.QyKit.bAbK_ne}
\leanok
\uses{Turing.NTimeDiag.QyKit, Turing.NTimeDiag.QyKit.bAbK, Turing.NTimeDiag.QyKit.tp, Turing.NTimeDiag.qyBank_ne_tp}
\difficulty{1}
Let $\kappa$ be a kit for $f, g : \mathbb{N} \to \mathbb{N}$, with named tapes $\tau_0, \dots,
\tau_{15}$, and consider the third bank of its increment call $C_a$ (tape $0$ on $\tau_{11}$, the
other tapes in its block). For every tape $i \ne 0$ of $C_a$ and every $c < 16$, the bank's image of
$i$ differs from $\tau_c$.
\end{theorem}

\begin{definition}[Tapes of the anchor layout]
\lean{Turing.NTimeDiag.qyLayT}
\label{Turing.NTimeDiag.qyLayT}
\leanok
\uses{Turing.FinTM.bufferTape, Turing.NTimeDiag.QyKit, Turing.NTimeDiag.QyKit.K, Turing.NTimeDiag.qyLay}
Let $\kappa$ be a kit for $f, g : \mathbb{N} \to \mathbb{N}$ and $n, a, j, k \in \mathbb{N}$. Write
$\mathrm{buf}(w)$ for the tape holding $w$ in cells $0, \dots, |w| - 1$ and blank elsewhere. The
anchor layout's tapes assign to each host tape $t$ the tape $\mathrm{buf}(L(t))$, where $L(t)$ is
the word of the anchor layout: $1^{n+1}$ on $\tau_0$, $1^{|\mathrm{bits}(n)|}$ on $\tau_1$, $1^a$
on $\tau_2$, $1^j$ on $\tau_{10}$, $1^k$ on $\tau_{11}$, $1$ on $\tau_{12}$, $1^{m(n)}$ on $\tau_8$
and $\tau_9$, and $\varepsilon$ elsewhere.
\end{definition}

\begin{definition}[The locator's anchor]
\lean{Turing.NTimeDiag.qyAnc}
\label{Turing.NTimeDiag.qyAnc}
\leanok
\uses{Turing.NTimeDiag.QyKit, Turing.NTimeDiag.QyLPh}
Let $\kappa$ be a kit for $f, g : \mathbb{N} \to \mathbb{N}$. The locator's anchor is its tagged
state consisting of the anchor phase in its initial sweep state.
\end{definition}

\begin{definition}[Escape phases of the locator]
\lean{Turing.NTimeDiag.QyLK.isEsc}
\label{Turing.NTimeDiag.QyLK.isEsc}
\leanok
\uses{Turing.NTimeDiag.QyLK}
The predicate on locator phase labels that holds exactly for the three escapes: the interior
escape, the top escape and the timeout escape.
\end{definition}

\begin{definition}[The locator's escape states]
\lean{Turing.NTimeDiag.qyLocE}
\label{Turing.NTimeDiag.qyLocE}
\leanok
\uses{Turing.NTimeDiag.QyKit, Turing.NTimeDiag.QyLK.isEsc, Turing.NTimeDiag.QyLPh}
Let $\kappa$ be a kit for $f, g : \mathbb{N} \to \mathbb{N}$. The locator's escape set is the set of
live locator states $(i, s)$ whose phase $i$ is one of the three escapes (interior, top or
timeout).
\end{definition}

\begin{theorem}[Non-escape locator states avoid the stop set]
\lean{Turing.NTimeDiag.qy_notStop}
\label{Turing.NTimeDiag.qy_notStop}
\leanok
\uses{Turing.NTimeDiag.QyKit, Turing.NTimeDiag.QyLK, Turing.NTimeDiag.QyLK.isEsc, Turing.NTimeDiag.QyLPh, Turing.NTimeDiag.qyAnc, Turing.NTimeDiag.qyLocE, Turing.not_mem_roundStop, Turing.roundStop}
\difficulty{2}
Let $\kappa$ be a kit for $f, g : \mathbb{N} \to \mathbb{N}$, with named tapes $\tau_0, \dots,
\tau_{15}$. Let $i$ be a phase of the locator's round other than the anchor and other than the
three escapes (interior, top and timeout), and let $s$ be any state of phase $i$. Then the locator
state $(i, s)$ does not belong to the round's stop set, which consists of the anchor state, every
escape state and the halted state.
\end{theorem}

\begin{theorem}[Escape states are round endpoints]
\lean{Turing.NTimeDiag.qy_esc_mem}
\label{Turing.NTimeDiag.qy_esc_mem}
\leanok
\uses{Turing.NTimeDiag.QyKit, Turing.NTimeDiag.QyLK, Turing.NTimeDiag.QyLK.isEsc, Turing.NTimeDiag.QyLPh, Turing.NTimeDiag.qyAnc, Turing.NTimeDiag.qyLocE}
\difficulty{1}
Let $\kappa$ be a kit for $f, g : \mathbb{N} \to \mathbb{N}$, with named tapes $\tau_0, \dots,
\tau_{15}$. For every escape phase $i$ of the locator's round (interior, top or timeout) and every
state $s$ of phase $i$, the locator state $(i, s)$ lies in the set of round endpoints, namely the
anchor state together with all escape states.
\end{theorem}

\begin{definition}[Scratch tapes of a locator round]
\lean{Turing.NTimeDiag.qyScratch}
\label{Turing.NTimeDiag.qyScratch}
\leanok
\uses{Turing.NTimeDiag.QyKit, Turing.NTimeDiag.QyKit.K, Turing.NTimeDiag.QyKit.bClk, Turing.NTimeDiag.QyKit.tp}
Let $\kappa$ be a kit for $f, g : \mathbb{N} \to \mathbb{N}$, with named tapes $\tau_0, \dots,
\tau_{15}$ and clock bank $\beta$ (the injection sending the tapes of $f$'s call $C_f$ to $\tau_7$
and $C_f$'s block, and the clock's counter to $\tau_8$). A host tape $t$ is a scratch tape of the
locator's round when $t$ lies in the range of $\beta$ or $t$ is one of $\tau_3, \tau_4, \tau_5,
\tau_6$.
\end{definition}

\begin{definition}[Escape tapes of the locator]
\lean{Turing.NTimeDiag.QyEscT}
\label{Turing.NTimeDiag.QyEscT}
\leanok
\uses{Turing.FinTM.bufferTape, Turing.NTimeDiag.QyKit, Turing.NTimeDiag.QyKit.K, Turing.NTimeDiag.qyLay, Turing.NTimeDiag.qyScratch}
Let $\kappa$ be a kit for $f, g : \mathbb{N} \to \mathbb{N}$, with named tapes $\tau_0, \dots,
\tau_{15}$, let $n, a, j, k \in \mathbb{N}$, let $T$ assign contents to the host tapes and $P$ head
positions, write $\mathrm{buf}(w)$ for the tape holding $w$ in cells $0, \dots, |w| - 1$ and blank
elsewhere, and write $\Lambda_{n,a,j,k}$ for the anchor-layout tapes, which hold $\mathrm{buf}$ of
$1^{n+1}$ on $\tau_0$, of $1^{|\mathrm{bits}(n)|}$ on $\tau_1$, of $1^a$ on $\tau_2$, of $1^j$ on
$\tau_{10}$, of $1^k$ on $\tau_{11}$, of $1$ on $\tau_{12}$, of $1^{m(n)}$ on $\tau_8$ and $\tau_9$
(with $m(n) = |\mathrm{bits}(A_f(c_f + 2))| + |\mathrm{bits}(n)|$ the clock counter width, $A_f$
and $c_f$ being the time constants of $f$'s call) and are blank on every other host tape. The pair
$(T, P)$ forms escape tapes for $(n, a, j, k)$ when every host tape $t$ outside the round's scratch
(the range of the clock's bank, which carries $f$'s call $C_f$ with tape $0$ on $\tau_7$ and its
counter on $\tau_8$, together with $\tau_3, \dots, \tau_6$) satisfies $T(t) = \Lambda_{n,a,j,k}(t)$
and $P(t) = 0$.
\end{definition}

\begin{theorem}[Call tape lies in the clock's bank]
\lean{Turing.NTimeDiag.qy_tp7_mem}
\label{Turing.NTimeDiag.qy_tp7_mem}
\leanok
\uses{Turing.NTimeDiag.QyKit, Turing.NTimeDiag.QyKit.bClk, Turing.NTimeDiag.QyKit.bClk_castSucc, Turing.NTimeDiag.QyKit.tp}
\difficulty{1}
Let $\kappa$ be a kit for $f, g : \mathbb{N} \to \mathbb{N}$, with named tapes $\tau_0, \dots,
\tau_{15}$. The named tape $\tau_7$ lies in the range of the clock's bank of $\kappa$, the
injection that sends the tapes of $f$'s call $C_f$ to $\tau_7$ (tape $0$) and $C_f$'s block, and
sends its last tape, the counter, to $\tau_8$.
\end{theorem}

\begin{theorem}[Counter tape lies in the clock's bank]
\lean{Turing.NTimeDiag.qy_tp8_mem}
\label{Turing.NTimeDiag.qy_tp8_mem}
\leanok
\uses{Turing.NTimeDiag.QyKit, Turing.NTimeDiag.QyKit.bClk, Turing.NTimeDiag.QyKit.bClk_last, Turing.NTimeDiag.QyKit.tp}
\difficulty{1}
Let $\kappa$ be a kit for $f, g : \mathbb{N} \to \mathbb{N}$, with named tapes $\tau_0, \dots,
\tau_{15}$. The named tape $\tau_8$ lies in the range of the clock's bank of $\kappa$, the
injection that sends the tapes of $f$'s call $C_f$ to $\tau_7$ (tape $0$) and $C_f$'s block, and
sends its last tape, the counter, to $\tau_8$.
\end{theorem}

\begin{theorem}[Non-scratch tapes avoid named tapes 3 to 8]
\lean{Turing.NTimeDiag.qy_not_scratch}
\label{Turing.NTimeDiag.qy_not_scratch}
\leanok
\uses{Turing.NTimeDiag.QyKit, Turing.NTimeDiag.QyKit.K, Turing.NTimeDiag.QyKit.tp, Turing.NTimeDiag.qyScratch, Turing.NTimeDiag.qy_tp7_mem, Turing.NTimeDiag.qy_tp8_mem}
\difficulty{1}
Let $\kappa$ be a kit for $f, g : \mathbb{N} \to \mathbb{N}$, with named tapes $\tau_0, \dots,
\tau_{15}$. If a host tape $t$ is not a scratch tape of the locator's round (it is outside the
range of the clock's bank and differs from $\tau_3, \dots, \tau_6$), then $t$ differs from each of
$\tau_3, \tau_4, \tau_5, \tau_6, \tau_7, \tau_8$.
\end{theorem}

\begin{theorem}[A round from a first step and a segment]
\lean{Turing.NTimeDiag.qy_round_of_seg}
\label{Turing.NTimeDiag.qy_round_of_seg}
\leanok
\uses{Turing.Cfg, Turing.MultiTapeTM, Turing.MultiTapeTM.runFrom, Turing.MultiTapeTM.runFrom_add, Turing.Round, Turing.Segment, Turing.roundStop, Turing.round_iff, Turing.segment_iff}
\difficulty{4}
Let $M$ be a multi-tape machine with state set $S$, let $a \in S$, let $E$ be a set of states
(possibly including the halted state), let $c, c_1, c'$ be configurations of $M$ on a common input
and let $t \in \mathbb{N}$. Suppose $c$ is in state $a$, one step of $M$ from $c$ yields $c_1$, $M$
runs from $c_1$ to $c'$ in exactly $t$ steps without being in $a$, in $E$ or halted strictly before
time $t$, and the state of $c'$ is $a$ or lies in $E$. Then the run of $M$ from $c$ is a round at
$a$ with escapes $E$ of length $1 + t$ ending at $c'$, with no constraint on head positions.
\end{theorem}

\begin{theorem}[The clocked call of a locator round]
\lean{Turing.NTimeDiag.qy_rnd_call}
\label{Turing.NTimeDiag.qy_rnd_call}
\leanok
\uses{Turing.FinTM.bufferTape, Turing.NTimeDiag.QyEscT, Turing.NTimeDiag.QyKit, Turing.NTimeDiag.QyKit.K, Turing.NTimeDiag.QyKit.bCf, Turing.NTimeDiag.QyKit.bCf_ne, Turing.NTimeDiag.QyKit.bClk, Turing.NTimeDiag.QyKit.bClk_castSucc, Turing.NTimeDiag.QyKit.bClk_last, Turing.NTimeDiag.QyKit.m, Turing.NTimeDiag.QyKit.tp, Turing.NTimeDiag.QyKit.tp_inj, Turing.NTimeDiag.QyKit.tp_ne, Turing.NTimeDiag.qyAnc, Turing.NTimeDiag.qyC, Turing.NTimeDiag.qyLay, Turing.NTimeDiag.qyLayT, Turing.NTimeDiag.qyLoc, Turing.NTimeDiag.qyLocE, Turing.NTimeDiag.qyLocM, Turing.NTimeDiag.qyLocRed, Turing.NTimeDiag.qy_arg_bank, Turing.NTimeDiag.qy_call_bound, Turing.NTimeDiag.qy_clockHost_finish, Turing.NTimeDiag.qy_clockHost_timeout, Turing.NTimeDiag.qy_copy, Turing.NTimeDiag.qy_guard_of_cut, Turing.NTimeDiag.qy_hosted, Turing.NTimeDiag.qy_notStop, Turing.NTimeDiag.qy_not_scratch, Turing.NTimeDiag.qy_timeout_cert, Turing.Segment, Turing.Segment.trans, Turing.roundStop, Turing.stateWord}
\difficulty{5}
Let $\kappa$ be a kit for $f, g : \mathbb{N} \to \mathbb{N}$, with named tapes $\tau_0, \dots,
\tau_{15}$, let $x$ be a binary input of length $n$, write $\langle s, p, T, P, o\rangle$ for the
live configuration on input $x$ with state $s$, input head $p$, work tapes $T$, heads $P$ and
output $o$, write $\mathrm{buf}(w)$ for the tape holding $w$ in cells $0, \dots, |w| - 1$ and blank
elsewhere, and write $\Lambda_{n,a,j,k}$ for the anchor-layout tapes, which hold $\mathrm{buf}$ of
$1^{n+1}$ on $\tau_0$, of $1^{|\mathrm{bits}(n)|}$ on $\tau_1$, of $1^a$ on $\tau_2$, of $1^j$ on
$\tau_{10}$, of $1^k$ on $\tau_{11}$, of $1$ on $\tau_{12}$, of $1^{m(n)}$ on $\tau_8$ and $\tau_9$
(with $m(n) = |\mathrm{bits}(A_f(c_f + 2))| + |\mathrm{bits}(n)|$ the clock counter width, $A_f$
and $c_f$ being the time constants of $f$'s call) and are blank on every other host tape. Let $A_f,
c_f$ be the time constants of $f$'s call, let $a, j, k \in \mathbb{N}$ with $a \le n$, and write $m
= m(n)$. Then one of the following holds. (i) There are a binary word $w'$ of length $m$ and a time
$t$ with $t \le 2a + 3 + 5(2^m - 1) + 2m + 3$ and $t \le 2a + 3 + 5A_f(c_f(f(a) + 1) + a +
|\mathrm{bits}(f(a))| + 1) + 2m + 3$ such that the locator runs in exactly $t$ steps from $\langle
(r_1, \mathrm{sweep}), 1, \Lambda_{n,a,j,k}, 0, \varepsilon\rangle$ to the length test $\langle
(r_3, \mathrm{scan}), 1, T', 0, \varepsilon \rangle$, where $T'$ is $\Lambda_{n,a,j,k}$ with
$\tau_7$ replaced by $\mathrm{buf}(\mathrm{bits}(f(a)))$ and $\tau_8$ by $\mathrm{buf}(w')$, never
in the locator's anchor, an escape state or the halted state strictly before time $t$. (ii) Or $n <
f(a)$ and there are $t \le 2a + 3 + 5(2^m - 1) + 2m + 3$, an input-head position $q$, tapes $T$ and
heads $P$ such that the locator runs in exactly $t$ steps from $\langle (r_1, \mathrm{sweep}), 1,
\Lambda_{n,a,j,k}, 0, \varepsilon\rangle$ to the timeout escape $\langle (\mathrm{timeout},
\mathrm{sweep}), q, T, P, \varepsilon\rangle$, never in the locator's anchor, an escape state or
the halted state strictly before time $t$, and every host tape outside the round's scratch (the
range of the clock's bank, which carries $f$'s call $C_f$ with tape $0$ on $\tau_7$ and its counter
on $\tau_8$, together with $\tau_3, \dots, \tau_6$) agrees with $\Lambda_{n,a,j,k}$ in $T$ and has
its head at $0$ in $P$.
\end{theorem}

\begin{definition}[Tapes after the round's call]
\lean{Turing.NTimeDiag.qyTA}
\label{Turing.NTimeDiag.qyTA}
\leanok
\uses{Turing.FinTM.bufferTape, Turing.NTimeDiag.QyKit, Turing.NTimeDiag.QyKit.K, Turing.NTimeDiag.QyKit.tp, Turing.NTimeDiag.qyLayT}
Let $\kappa$ be a kit for $f, g : \mathbb{N} \to \mathbb{N}$, with named tapes $\tau_0, \dots,
\tau_{15}$, let $n, a, j, k, F \in \mathbb{N}$ and let $w'$ be a binary word. Write
$\mathrm{buf}(w)$ for the tape holding $w$ in cells $0, \dots, |w| - 1$ and blank elsewhere, and
write $\Lambda_{n,a,j,k}$ for the anchor-layout tapes, which hold $\mathrm{buf}$ of $1^{n+1}$ on
$\tau_0$, of $1^{|\mathrm{bits}(n)|}$ on $\tau_1$, of $1^a$ on $\tau_2$, of $1^j$ on $\tau_{10}$,
of $1^k$ on $\tau_{11}$, of $1$ on $\tau_{12}$, of $1^{m(n)}$ on $\tau_8$ and $\tau_9$ (with $m(n)
= |\mathrm{bits}(A_f(c_f + 2))| + |\mathrm{bits}(n)|$ the clock counter width, $A_f$ and $c_f$
being the time constants of $f$'s call) and are blank on every other host tape. The tapes after the
round's call are $\Lambda_{n,a,j,k}$ with $\tau_7$ replaced by $\mathrm{buf}(\mathrm{bits}(F))$ and
$\tau_8$ replaced by $\mathrm{buf}(w')$.
\end{definition}

\begin{definition}[Tapes after a passed round's tests]
\lean{Turing.NTimeDiag.qyTB}
\label{Turing.NTimeDiag.qyTB}
\leanok
\uses{Turing.FinTM.bufferTape, Turing.NTimeDiag.QyKit, Turing.NTimeDiag.QyKit.K, Turing.NTimeDiag.QyKit.tp}
Let $\kappa$ be a kit for $f, g : \mathbb{N} \to \mathbb{N}$, with named tapes $\tau_0, \dots,
\tau_{15}$, let $T$ assign contents to the host tapes, let $F, a, X_4, \ell' \in \mathbb{N}$, and
write $\mathrm{buf}(w)$ for the tape holding $w$ in cells $0, \dots, |w| - 1$ and blank elsewhere.
The tapes after a passed round's tests are $T$ with $\tau_4$ replaced by $\mathrm{buf}(1^{F + a})$,
$\tau_7$ by $\mathrm{buf}(1^{|\mathrm{bits}(F)|})$, $\tau_5$ by $\mathrm{buf}(1^{X_4})$, $\tau_6$
by $\mathrm{buf}(1^{X_4 + 1})$ and $\tau_3$ by $\mathrm{buf}(1^{\ell' + 1})$.
\end{definition}

\begin{definition}[Time bound of stopping tests]
\lean{Turing.NTimeDiag.qyTestS}
\label{Turing.NTimeDiag.qyTestS}
\leanok
\uses{Turing.NTimeDiag.QyKit}
Let $\kappa$ be a kit for $f, g : \mathbb{N} \to \mathbb{N}$ with time constants $A_u$ of its
fourth-power call and $A_a$ of its increment call. For $n, a \in \mathbb{N}$, with $L =
|\mathrm{bits}(n)|$, the time bound of a round's tests when the round stops is $40(n + 1) + 4a +
16L + A_u L^5 + A_a(L + 1) + 60$.
\end{definition}

\begin{definition}[Time bound of passed tests]
\lean{Turing.NTimeDiag.qyTestP}
\label{Turing.NTimeDiag.qyTestP}
\leanok
\uses{Turing.NTimeDiag.QyKit}
Let $\kappa$ be a kit for $f, g : \mathbb{N} \to \mathbb{N}$ with time constants $A_u$ of its
fourth-power call and $A_a$ of its increment call. For $n, \ell' \in \mathbb{N}$, with $L =
|\mathrm{bits}(n)|$, the time bound of a round's tests when the round passes is $9\ell' + 25L + A_u
L^5 + A_a(L + 1) + 50$.
\end{definition}

\begin{theorem}[The tests of a locator round]
\lean{Turing.NTimeDiag.qy_rnd_tests}
\label{Turing.NTimeDiag.qy_rnd_tests}
\leanok
\uses{Complexity.LogProg.length_bits_eq_log, Complexity.LogProg.length_bits_mono, Complexity.LogProg.lt_two_pow_bits, Complexity.LogProg.two_pow_bits_le_two_mul_add_one, Turing.FinTM.bufferTape, Turing.MultiTapeTM.runFrom, Turing.NTimeDiag.QyEscT, Turing.NTimeDiag.QyKit, Turing.NTimeDiag.QyKit.K, Turing.NTimeDiag.QyKit.bAb0, Turing.NTimeDiag.QyKit.bAb0_ne, Turing.NTimeDiag.QyKit.bPu, Turing.NTimeDiag.QyKit.bPu_ne, Turing.NTimeDiag.QyKit.tp, Turing.NTimeDiag.QyKit.tp_inj, Turing.NTimeDiag.QyKit.tp_ne, Turing.NTimeDiag.qyAnc, Turing.NTimeDiag.qyC, Turing.NTimeDiag.qyEmb2, Turing.NTimeDiag.qyLay, Turing.NTimeDiag.qyLayT, Turing.NTimeDiag.qyLoc, Turing.NTimeDiag.qyLocE, Turing.NTimeDiag.qyLocM, Turing.NTimeDiag.qyLocRed, Turing.NTimeDiag.qyTA, Turing.NTimeDiag.qyTB, Turing.NTimeDiag.qyTestP, Turing.NTimeDiag.qyTestS, Turing.NTimeDiag.qy_arg_bank, Turing.NTimeDiag.qy_bits_val, Turing.NTimeDiag.qy_callHost, Turing.NTimeDiag.qy_cmp_run, Turing.NTimeDiag.qy_copy, Turing.NTimeDiag.qy_expandHost, Turing.NTimeDiag.qy_guard_of_cut, Turing.NTimeDiag.qy_hosted, Turing.NTimeDiag.qy_notStop, Turing.NTimeDiag.qy_not_scratch, Turing.NTimeDiag.qy_rewind, Turing.NTimeDiag.qy_rung_val, Turing.NTimeDiag.qy_writeConst, Turing.Segment, Turing.Segment.trans, Turing.counterOverhead, Turing.counterOverhead_le, Turing.length_bits_le_self, Turing.roundStop}
\difficulty{6}
Let $\kappa$ be a kit for $f, g : \mathbb{N} \to \mathbb{N}$, with named tapes $\tau_0, \dots,
\tau_{15}$, let $x$ be a binary input of length $n$, write $\langle s, p, T, P, o\rangle$ for the
live configuration on input $x$ with state $s$, input head $p$, work tapes $T$, heads $P$ and
output $o$, write $\mathrm{buf}(w)$ for the tape holding $w$ in cells $0, \dots, |w| - 1$ and blank
elsewhere, and write $\Lambda_{n,a,j,k}$ for the anchor-layout tapes, which hold $\mathrm{buf}$ of
$1^{n+1}$ on $\tau_0$, of $1^{|\mathrm{bits}(n)|}$ on $\tau_1$, of $1^a$ on $\tau_2$, of $1^j$ on
$\tau_{10}$, of $1^k$ on $\tau_{11}$, of $1$ on $\tau_{12}$, of $1^{m(n)}$ on $\tau_8$ and $\tau_9$
(with $m(n) = |\mathrm{bits}(A_f(c_f + 2))| + |\mathrm{bits}(n)|$ the clock counter width, $A_f$
and $c_f$ being the time constants of $f$'s call) and are blank on every other host tape. Let $A_u$
and $A_a$ be the time constants of the fourth-power call and the increment call, let $a, j, k \in
\mathbb{N}$, let $w'$ be a binary word, put $L = |\mathrm{bits}(n)|$, $X = f(a) + a + 1$ and $\ell'
= 2^{X^4}$, let $T_A$ be $\Lambda_{n,a,j,k}$ with $\tau_7$ replaced by
$\mathrm{buf}(\mathrm{bits}(f(a)))$ and $\tau_8$ by $\mathrm{buf}(w')$, let $c_0 = \langle (r_3,
\mathrm{scan}), 1, T_A, 0, \varepsilon \rangle$, and let $B_S = 40(n + 1) + 4a + 16L + A_u L^5 +
A_a(L + 1) + 60$. In each case below the run lasts exactly $t$ steps and is never in the locator's
anchor, an escape state or the halted state strictly before time $t$. One of the following holds.
(i) $n < \ell'$, and for some $t \le B_S$ and $(T, P)$ the locator runs from $c_0$ to $\langle
(\mathrm{interior}, \mathrm{sweep}), 1, T, P, \varepsilon\rangle$ with every tape outside the
round's scratch (the range of the clock's bank, which carries $f$'s call $C_f$ with tape $0$ on
$\tau_7$ and its counter on $\tau_8$, together with $\tau_3, \dots, \tau_6$) agreeing with
$\Lambda_{n,a,j,k}$ and heads there at $0$. (ii) $\ell' = n$, and for some $t \le B_S$ and $(T, P)$
of that kind the locator runs from $c_0$ to $\langle (\mathrm{top}, \mathrm{sweep}), 1, T, P,
\varepsilon\rangle$ with additionally $T(\tau_4) = \mathrm{buf}(1^{f(a) + a})$ and $P(\tau_4) = 0$.
(iii) $\ell' < n$, $X \le L$, $X^4 < L$, and for some $t \le 9\ell' + 25L + A_u L^5 + A_a(L + 1) +
50$ the locator runs from $c_0$ to $\langle (c_1, \mathrm{sweep}), 1, T_B, 0, \varepsilon\rangle$,
where $T_B$ is $T_A$ with $\tau_4, \tau_7, \tau_5, \tau_6, \tau_3$ replaced by $\mathrm{buf}$ of
$1^{f(a) + a}$, $1^{|\mathrm{bits}(f(a))|}$, $1^{X^4}$, $1^{X^4 + 1}$, $1^{\ell' + 1}$
respectively.
\end{theorem}

\begin{theorem}[Cleanup after a passed round]
\lean{Turing.NTimeDiag.qy_rnd_clean}
\label{Turing.NTimeDiag.qy_rnd_clean}
\leanok
\uses{Turing.FinTM.bufferTape, Turing.NTimeDiag.QyKit, Turing.NTimeDiag.QyKit.K, Turing.NTimeDiag.QyKit.m, Turing.NTimeDiag.QyKit.tape_cases, Turing.NTimeDiag.QyKit.tp, Turing.NTimeDiag.QyKit.tp_inj, Turing.NTimeDiag.QyKit.tp_ne, Turing.NTimeDiag.qyAnc, Turing.NTimeDiag.qyC, Turing.NTimeDiag.qyLay, Turing.NTimeDiag.qyLayT, Turing.NTimeDiag.qyLoc, Turing.NTimeDiag.qyLocE, Turing.NTimeDiag.qyLocM, Turing.NTimeDiag.qyLocRed, Turing.NTimeDiag.qyTA, Turing.NTimeDiag.qyTB, Turing.NTimeDiag.qy_clear, Turing.NTimeDiag.qy_copy_over, Turing.NTimeDiag.qy_guard_of_cut, Turing.NTimeDiag.qy_hosted, Turing.NTimeDiag.qy_notStop, Turing.NTimeDiag.qy_transfer, Turing.Segment, Turing.Segment.trans, Turing.roundStop}
\difficulty{4}
Let $\kappa$ be a kit for $f, g : \mathbb{N} \to \mathbb{N}$, with named tapes $\tau_0, \dots,
\tau_{15}$, let $x$ be a binary input of length $n$, write $\langle s, p, T, P, o\rangle$ for the
live configuration on input $x$ with state $s$, input head $p$, work tapes $T$, heads $P$ and
output $o$, write $\mathrm{buf}(w)$ for the tape holding $w$ in cells $0, \dots, |w| - 1$ and blank
elsewhere, and write $\Lambda_{n,a,j,k}$ for the anchor-layout tapes, which hold $\mathrm{buf}$ of
$1^{n+1}$ on $\tau_0$, of $1^{|\mathrm{bits}(n)|}$ on $\tau_1$, of $1^a$ on $\tau_2$, of $1^j$ on
$\tau_{10}$, of $1^k$ on $\tau_{11}$, of $1$ on $\tau_{12}$, of $1^{m(n)}$ on $\tau_8$ and $\tau_9$
(with $m(n) = |\mathrm{bits}(A_f(c_f + 2))| + |\mathrm{bits}(n)|$ the clock counter width, $A_f$
and $c_f$ being the time constants of $f$'s call) and are blank on every other host tape. Let $a,
j, k, F, X_4, \ell' \in \mathbb{N}$ and let $w'$ be a binary word of length $m(n)$. Let $T_B$ be
$\Lambda_{n,a,j,k}$ with $\tau_7$ replaced by $\mathrm{buf}(1^{|\mathrm{bits}(F)|})$, $\tau_8$ by
$\mathrm{buf}(w')$, $\tau_4$ by $\mathrm{buf}(1^{F + a})$, $\tau_5$ by $\mathrm{buf}(1^{X_4})$,
$\tau_6$ by $\mathrm{buf}(1^{X_4 + 1})$ and $\tau_3$ by $\mathrm{buf}(1^{\ell' + 1})$. Then for
some $t \le 2a + 2(\ell' + 1) + 2|\mathrm{bits}(F)| + 2(F + a) + 2X_4 + 2(X_4 + 1) + 2m(n) + 21$
the locator runs in exactly $t$ steps from $\langle (c_1, \mathrm{sweep}), 1, T_B, 0, \varepsilon
\rangle$ to $\langle (u_1, \mathrm{scan}), 1, \Lambda_{n,\ell' + 1,j,k}, 0, \varepsilon\rangle$,
never in the locator's anchor, an escape state or the halted state strictly before time $t$.
\end{theorem}

\begin{theorem}[Changing the pair in the anchor layout]
\lean{Turing.NTimeDiag.qy_layT_pair}
\label{Turing.NTimeDiag.qy_layT_pair}
\leanok
\uses{Turing.FinTM.bufferTape, Turing.NTimeDiag.QyKit, Turing.NTimeDiag.QyKit.K, Turing.NTimeDiag.QyKit.tp, Turing.NTimeDiag.QyKit.tp_inj, Turing.NTimeDiag.qyLay, Turing.NTimeDiag.qyLayT}
\difficulty{3}
Let $\kappa$ be a kit for $f, g : \mathbb{N} \to \mathbb{N}$, with named tapes $\tau_0, \dots,
\tau_{15}$, let $n, a, j, k, j', k' \in \mathbb{N}$, write $\mathrm{buf}(w)$ for the tape holding
$w$ in cells $0, \dots, |w| - 1$ and blank elsewhere, and write $\Lambda_{n,a,j,k}$ for the
anchor-layout tapes, which hold $\mathrm{buf}$ of $1^{n+1}$ on $\tau_0$, of
$1^{|\mathrm{bits}(n)|}$ on $\tau_1$, of $1^a$ on $\tau_2$, of $1^j$ on $\tau_{10}$, of $1^k$ on
$\tau_{11}$, of $1$ on $\tau_{12}$, of $1^{m(n)}$ on $\tau_8$ and $\tau_9$ (with $m(n) =
|\mathrm{bits}(A_f(c_f + 2))| + |\mathrm{bits}(n)|$ the clock counter width, $A_f$ and $c_f$ being
the time constants of $f$'s call) and are blank on every other host tape. If host tape contents $T$
satisfy $T(\tau_{10}) = \mathrm{buf}(1^{j'})$, $T(\tau_{11}) = \mathrm{buf}(1^{k'})$ and $T(t) =
\Lambda_{n,a,j,k}(t)$ for every other host tape $t$, then $T = \Lambda_{n,a,j',k'}$.
\end{theorem}

\begin{theorem}[Pair update of a passed round]
\lean{Turing.NTimeDiag.qy_rnd_pair}
\label{Turing.NTimeDiag.qy_rnd_pair}
\leanok
\uses{Turing.FinTM.bufferTape, Turing.MultiTapeTM.runFrom, Turing.NTimeDiag.QyKit, Turing.NTimeDiag.QyKit.K, Turing.NTimeDiag.QyKit.bAbJ, Turing.NTimeDiag.QyKit.bAbJ_ne, Turing.NTimeDiag.QyKit.bAbK, Turing.NTimeDiag.QyKit.bAbK_ne, Turing.NTimeDiag.QyKit.tp, Turing.NTimeDiag.QyKit.tp_inj, Turing.NTimeDiag.QyKit.tp_ne, Turing.NTimeDiag.qyAnc, Turing.NTimeDiag.qyC, Turing.NTimeDiag.qyLay, Turing.NTimeDiag.qyLayT, Turing.NTimeDiag.qyLoc, Turing.NTimeDiag.qyLocE, Turing.NTimeDiag.qyLocM, Turing.NTimeDiag.qyLocRed, Turing.NTimeDiag.qy_arg_bank, Turing.NTimeDiag.qy_callHost, Turing.NTimeDiag.qy_clear, Turing.NTimeDiag.qy_cmp_run, Turing.NTimeDiag.qy_guard_of_cut, Turing.NTimeDiag.qy_hosted, Turing.NTimeDiag.qy_layT_pair, Turing.NTimeDiag.qy_notStop, Turing.NTimeDiag.qy_unpair_succ_diag, Turing.NTimeDiag.qy_unpair_succ_gt, Turing.NTimeDiag.qy_unpair_succ_last, Turing.NTimeDiag.qy_unpair_succ_lt, Turing.Segment, Turing.Segment.trans, Turing.roundStop}
\difficulty{5}
Let $\kappa$ be a kit for $f, g : \mathbb{N} \to \mathbb{N}$, with named tapes $\tau_0, \dots,
\tau_{15}$, let $x$ be a binary input of length $n$, write $\langle s, p, T, P, o\rangle$ for the
live configuration on input $x$ with state $s$, input head $p$, work tapes $T$, heads $P$ and
output $o$, write $\mathrm{buf}(w)$ for the tape holding $w$ in cells $0, \dots, |w| - 1$ and blank
elsewhere, and write $\Lambda_{n,a,j,k}$ for the anchor-layout tapes, which hold $\mathrm{buf}$ of
$1^{n+1}$ on $\tau_0$, of $1^{|\mathrm{bits}(n)|}$ on $\tau_1$, of $1^a$ on $\tau_2$, of $1^j$ on
$\tau_{10}$, of $1^k$ on $\tau_{11}$, of $1$ on $\tau_{12}$, of $1^{m(n)}$ on $\tau_8$ and $\tau_9$
(with $m(n) = |\mathrm{bits}(A_f(c_f + 2))| + |\mathrm{bits}(n)|$ the clock counter width, $A_f$
and $c_f$ being the time constants of $f$'s call) and are blank on every other host tape. Let $A_a$
be the time constant of the increment call, and let $a, h, j, k \in \mathbb{N}$ with
$\mathrm{unpair}(h) = (j, k)$ for the standard unpairing of $\mathbb{N}$, and let $(j', k') =
\mathrm{unpair}(h + 1)$. Then for some $t \le A_a(j + k + 2) + 6(j + k) + 30$ the locator runs in
exactly $t$ steps from $\langle (u_1, \mathrm{scan}), 1, \Lambda_{n,a,j,k}, 0, \varepsilon\rangle$
to its anchor $\langle \mathrm{anchor}, 1, \Lambda_{n,a,j',k'}, 0, \varepsilon\rangle$, never in
the locator's anchor, an escape state or the halted state strictly before time $t$.
\end{theorem}

\begin{definition}[Time bound of a passed round]
\lean{Turing.NTimeDiag.qyRndP}
\label{Turing.NTimeDiag.qyRndP}
\leanok
\uses{Turing.NTimeDiag.QyKit, Turing.NTimeDiag.QyKit.m}
Let $\kappa$ be a kit for $f, g : \mathbb{N} \to \mathbb{N}$ with time constants $A_f, c_f$ of
$f$'s call, $A_u$ of the fourth-power call and $A_a$ of the increment call, and let $m(n) =
|\mathrm{bits}(A_f(c_f + 2))| + |\mathrm{bits}(n)|$. For $n, \ell', j, k \in \mathbb{N}$, with $L =
|\mathrm{bits}(n)|$, the time bound of a passed round is \[ 11\ell' + 50L + 4m(n) + 5A_f(c_f + 3)(L
+ 1) + A_u L^5 + A_a(L + 1) + A_a(j + k + 2) + 6(j + k) + 150. \]
\end{definition}

\begin{definition}[Time bound of a stopping round]
\lean{Turing.NTimeDiag.qyRndS}
\label{Turing.NTimeDiag.qyRndS}
\leanok
\uses{Turing.NTimeDiag.QyKit, Turing.NTimeDiag.QyKit.m}
Let $\kappa$ be a kit for $f, g : \mathbb{N} \to \mathbb{N}$ with time constants $A_f, c_f$ of
$f$'s call, $A_u$ of the fourth-power call and $A_a$ of the increment call, and let $m(n) =
|\mathrm{bits}(A_f(c_f + 2))| + |\mathrm{bits}(n)|$. For $n, a \in \mathbb{N}$, with $L =
|\mathrm{bits}(n)|$, the time bound of a stopping round is \[ 6a + 5 \cdot 2^{m(n)} + 2m(n) + 40(n
+ 1) + 16L + A_u L^5 + A_a(L + 1) + 70. \]
\end{definition}

\begin{theorem}[One round of the locator]
\lean{Turing.NTimeDiag.qy_round}
\label{Turing.NTimeDiag.qy_round}
\leanok
\uses{Turing.Cfg, Turing.FinTM.bufferTape, Turing.MultiTapeTM.runFrom, Turing.MultiTapeTM.runFrom_zero, Turing.NTimeDiag.QyEscT, Turing.NTimeDiag.QyKit, Turing.NTimeDiag.QyKit.K, Turing.NTimeDiag.QyKit.tp, Turing.NTimeDiag.QyLPh, Turing.NTimeDiag.qyAnc, Turing.NTimeDiag.qyC, Turing.NTimeDiag.qyLayT, Turing.NTimeDiag.qyLoc, Turing.NTimeDiag.qyLocE, Turing.NTimeDiag.qyLocM, Turing.NTimeDiag.qyLocRed, Turing.NTimeDiag.qyRndP, Turing.NTimeDiag.qyRndS, Turing.NTimeDiag.qyTestP, Turing.NTimeDiag.qyTestS, Turing.NTimeDiag.qy_esc_mem, Turing.NTimeDiag.qy_hosted_run, Turing.NTimeDiag.qy_rnd_call, Turing.NTimeDiag.qy_rnd_clean, Turing.NTimeDiag.qy_rnd_pair, Turing.NTimeDiag.qy_rnd_tests, Turing.NTimeDiag.qy_round_of_seg, Turing.Round, Turing.Segment.trans, Turing.length_bits_le_self}
\difficulty{5}
Let $\kappa$ be a kit for $f, g : \mathbb{N} \to \mathbb{N}$, with named tapes $\tau_0, \dots,
\tau_{15}$, let $x$ be a binary input of length $n$, write $\langle s, p, T, P, o\rangle$ for the
live configuration on input $x$ with state $s$, input head $p$, work tapes $T$, heads $P$ and
output $o$, write $\mathrm{buf}(w)$ for the tape holding $w$ in cells $0, \dots, |w| - 1$ and blank
elsewhere, and write $\Lambda_{n,a,j,k}$ for the anchor-layout tapes, which hold $\mathrm{buf}$ of
$1^{n+1}$ on $\tau_0$, of $1^{|\mathrm{bits}(n)|}$ on $\tau_1$, of $1^a$ on $\tau_2$, of $1^j$ on
$\tau_{10}$, of $1^k$ on $\tau_{11}$, of $1$ on $\tau_{12}$, of $1^{m(n)}$ on $\tau_8$ and $\tau_9$
(with $m(n) = |\mathrm{bits}(A_f(c_f + 2))| + |\mathrm{bits}(n)|$ the clock counter width, $A_f$
and $c_f$ being the time constants of $f$'s call) and are blank on every other host tape. Let $A_f,
c_f$ be the time constants of $f$'s call and $A_u, A_a$ those of the fourth-power and increment
calls; put $L = |\mathrm{bits}(n)|$, $R_P(\ell', j, k) = 11\ell' + 50L + 4m(n) + 5A_f(c_f + 3)(L +
1) + A_u L^5 + A_a(L + 1) + A_a(j + k + 2) + 6(j + k) + 150$ and $R_S(a) = 6a + 5 \cdot 2^{m(n)} +
2m(n) + 40(n + 1) + 16L + A_u L^5 + A_a(L + 1) + 70$. Let $a, h, j, k \in \mathbb{N}$ with $a \le
n$ and $\mathrm{unpair}(h) = (j, k)$, and put $X = f(a) + a + 1$, $\ell' = 2^{X^4}$. Then there are
a configuration $c'$ and $t \in \mathbb{N}$ such that the locator's run from $\langle
\mathrm{anchor}, 1, \Lambda_{n,a,j,k}, 0, \varepsilon\rangle$ is a round of length $t$ at the
anchor with the three escapes, ending at $c'$, and one of the following holds: (i) $\ell' < n$, $t
\le R_P(\ell', j, k)$ and $c' = \langle \mathrm{anchor}, 1, \Lambda_{n,\ell' + 1,j',k'}, 0,
\varepsilon\rangle$ with $(j', k') = \mathrm{unpair}(h + 1)$; (ii) $n < \ell'$, $t \le R_S(a)$ and
$c'$ is the interior or the timeout escape in its sweep state, with some input head, tapes $T$ and
heads $P$ such that every host tape outside the round's scratch (the range of the clock's bank,
which carries $f$'s call $C_f$ with tape $0$ on $\tau_7$ and its counter on $\tau_8$, together with
$\tau_3, \dots, \tau_6$) agrees with $\Lambda_{n,a,j,k}$ and has its head at $0$; (iii) $\ell' =
n$, $t \le R_S(a)$ and $c' = \langle (\mathrm{top}, \mathrm{sweep}), 1, T, P, \varepsilon\rangle$
with $(T, P)$ of that kind and moreover $T(\tau_4) = \mathrm{buf}(1^{f(a) + a})$ and $P(\tau_4) =
0$.
\end{theorem}

\begin{theorem}[Sixth power below an exponential]
\lean{Turing.NTimeDiag.qy_pow6_le}
\label{Turing.NTimeDiag.qy_pow6_le}
\leanok
\difficulty{4}
For every natural number $x$, \[ x^6 \le 6^6 \cdot 2^x. \]
\end{theorem}

\begin{theorem}[Sum of an affine family]
\lean{Turing.NTimeDiag.qy_sum_affine}
\label{Turing.NTimeDiag.qy_sum_affine}
\leanok
\difficulty{3}
For every function $u : \mathbb{N} \to \mathbb{N}$ and all $c, s \in \mathbb{N}$, \[ \sum_{r < s}
\bigl(11\,u(r) + c\bigr) = 11 \sum_{r < s} u(r) + s\,c. \]
\end{theorem}

\begin{definition}[Polylogarithmic constant of a kit]
\lean{Turing.NTimeDiag.qyCk}
\label{Turing.NTimeDiag.qyCk}
\leanok
\uses{Turing.NTimeDiag.QyKit}
Let $\kappa$ be a kit for $f, g : \mathbb{N} \to \mathbb{N}$ with time constants $A_f, c_f$ of
$f$'s call, $A_u$ of the fourth-power call and $A_a$ of the increment call. Its polylogarithmic
constant is $C_\kappa = 210 + 4|\mathrm{bits}(A_f(c_f + 2))| + 5A_f(c_f + 3) + A_u + 3A_a$.
\end{definition}

\begin{definition}[Linear constant of a kit]
\lean{Turing.NTimeDiag.qyDk}
\label{Turing.NTimeDiag.qyDk}
\leanok
\uses{Turing.NTimeDiag.QyKit}
Let $\kappa$ be a kit for $f, g : \mathbb{N} \to \mathbb{N}$ with time constants $A_f, c_f$ of
$f$'s call, and let $D_0 = |\mathrm{bits}(A_f(c_f + 2))|$. Its linear constant is $D_\kappa = 134 +
10 \cdot 2^{D_0} + 2D_0$.
\end{definition}

\begin{definition}[The locator's constant]
\lean{Turing.NTimeDiag.qyLoopA}
\label{Turing.NTimeDiag.qyLoopA}
\leanok
\uses{Turing.NTimeDiag.QyKit, Turing.NTimeDiag.qyCk, Turing.NTimeDiag.qyDk}
Let $\kappa$ be a kit for $f, g : \mathbb{N} \to \mathbb{N}$ with time constants $A_f, c_f$ of
$f$'s call, $A_u$ of the fourth-power call and $A_a$ of the increment call; let $D_0 =
|\mathrm{bits}(A_f(c_f + 2))|$, $C_\kappa = 210 + 4D_0 + 5A_f(c_f + 3) + A_u + 3A_a$ and $D_\kappa
= 134 + 10 \cdot 2^{D_0} + 2D_0$. The locator's constant is $22 + D_\kappa + 8 \cdot 6^6 \cdot
C_\kappa$.
\end{definition}

\begin{theorem}[Simplified bound of a passed round]
\lean{Turing.NTimeDiag.qy_rndP_le}
\label{Turing.NTimeDiag.qy_rndP_le}
\leanok
\uses{Turing.NTimeDiag.QyKit, Turing.NTimeDiag.QyKit.m, Turing.NTimeDiag.qyCk, Turing.NTimeDiag.qyRndP}
\difficulty{4}
Let $\kappa$ be a kit for $f, g : \mathbb{N} \to \mathbb{N}$ with time constants $A_f, c_f$ of
$f$'s call, $A_u$ of the fourth-power call and $A_a$ of the increment call; let $D_0 =
|\mathrm{bits}(A_f(c_f + 2))|$ and $C_\kappa = 210 + 4D_0 + 5A_f(c_f + 3) + A_u + 3A_a$. Let $n,
\ell', j, k \in \mathbb{N}$ with $j + k \le L$, where $L = |\mathrm{bits}(n)|$ and $m(n) = D_0 +
L$. Then \[ 11\ell' + 50L + 4m(n) + 5A_f(c_f + 3)(L + 1) + A_u L^5 + A_a(L + 1) + A_a(j + k + 2) +
6(j + k) + 150 \le 11\ell' + C_\kappa (L + 1)^5. \]
\end{theorem}

\begin{theorem}[Simplified bound of a stopping round]
\lean{Turing.NTimeDiag.qy_rndS_le}
\label{Turing.NTimeDiag.qy_rndS_le}
\leanok
\uses{Complexity.LogProg.two_pow_bits_le_two_mul_add_one, Turing.NTimeDiag.QyKit, Turing.NTimeDiag.QyKit.m, Turing.NTimeDiag.qyCk, Turing.NTimeDiag.qyDk, Turing.NTimeDiag.qyRndS, Turing.length_bits_le_self}
\difficulty{4}
Let $\kappa$ be a kit for $f, g : \mathbb{N} \to \mathbb{N}$ with time constants $A_f, c_f$ of
$f$'s call, $A_u$ of the fourth-power call and $A_a$ of the increment call; let $D_0 =
|\mathrm{bits}(A_f(c_f + 2))|$, $C_\kappa = 210 + 4D_0 + 5A_f(c_f + 3) + A_u + 3A_a$ and $D_\kappa
= 134 + 10 \cdot 2^{D_0} + 2D_0$. Let $n, a \in \mathbb{N}$ with $a \le n$, $L =
|\mathrm{bits}(n)|$ and $m(n) = D_0 + L$. Then \[ 6a + 5 \cdot 2^{m(n)} + 2m(n) + 40(n + 1) + 16L +
A_u L^5 + A_a(L + 1) + 70 \le D_\kappa (n + 1) + C_\kappa (L + 1)^5. \]
\end{theorem}

\begin{definition}[Locator anchor configuration at a stage]
\lean{Turing.NTimeDiag.qyAt}
\label{Turing.NTimeDiag.qyAt}
\leanok
\uses{Turing.Cfg, Turing.NTimeDiag.QyKit, Turing.NTimeDiag.QyKit.K, Turing.NTimeDiag.QyLPh, Turing.NTimeDiag.ladder, Turing.NTimeDiag.qyAnc, Turing.NTimeDiag.qyC, Turing.NTimeDiag.qyLayT}
Let $\kappa$ be a kit for $f, g : \mathbb{N} \to \mathbb{N}$, with named tapes $\tau_0, \dots,
\tau_{15}$, let $f$ have ladder $(\ell_i)$, let $x$ be a binary input of length $n$, and write
$\Lambda_{n,a,j,k}$ for the anchor-layout tapes, which hold $\mathrm{buf}$ of $1^{n+1}$ on
$\tau_0$, of $1^{|\mathrm{bits}(n)|}$ on $\tau_1$, of $1^a$ on $\tau_2$, of $1^j$ on $\tau_{10}$,
of $1^k$ on $\tau_{11}$, of $1$ on $\tau_{12}$, of $1^{m(n)}$ on $\tau_8$ and $\tau_9$ (with $m(n)
= |\mathrm{bits}(A_f(c_f + 2))| + |\mathrm{bits}(n)|$ the clock counter width, $A_f$ and $c_f$
being the time constants of $f$'s call) and are blank on every other host tape. For $i \in
\mathbb{N}$ with $\mathrm{unpair}(i) = (j, k)$, the locator's anchor configuration at stage $i$ is
the live configuration on $x$ in the anchor state with input head $1$, tapes $\Lambda_{n,\ell_i +
1,j,k}$, every work head at $0$ and empty output.
\end{definition}

\begin{theorem}[The locator's loop]
\lean{Turing.NTimeDiag.qy_loop}
\label{Turing.NTimeDiag.qy_loop}
\leanok
\uses{Complexity.LogProg.length_bits_eq_log, Complexity.LogProg.two_pow_bits_le_two_mul_add_one, Turing.Cfg, Turing.FinTM.bufferTape, Turing.NTimeDiag.QyEscT, Turing.NTimeDiag.QyKit, Turing.NTimeDiag.QyKit.K, Turing.NTimeDiag.QyKit.tp, Turing.NTimeDiag.QyLPh, Turing.NTimeDiag.ladder, Turing.NTimeDiag.ladder_strictMono, Turing.NTimeDiag.qyAnc, Turing.NTimeDiag.qyAt, Turing.NTimeDiag.qyC, Turing.NTimeDiag.qyCk, Turing.NTimeDiag.qyDk, Turing.NTimeDiag.qyLoc, Turing.NTimeDiag.qyLocE, Turing.NTimeDiag.qyLoopA, Turing.NTimeDiag.qy_ladder_succ_eq, Turing.NTimeDiag.qy_ledger, Turing.NTimeDiag.qy_pow6_le, Turing.NTimeDiag.qy_rndP_le, Turing.NTimeDiag.qy_rndS_le, Turing.NTimeDiag.qy_round, Turing.NTimeDiag.qy_stage_lt_log, Turing.NTimeDiag.qy_sum_affine, Turing.NTimeDiag.stage, Turing.NTimeDiag.stage_spec, Turing.Round, Turing.Round.iterate_seq, Turing.Round.segment, Turing.Segment, Turing.Segment.trans}
\difficulty{6}
Let $\kappa$ be a kit for $f, g : \mathbb{N} \to \mathbb{N}$, with named tapes $\tau_0, \dots,
\tau_{15}$, let $x$ be a binary input of length $n$, write $\langle s, p, T, P, o\rangle$ for the
live configuration on input $x$ with state $s$, input head $p$, work tapes $T$, heads $P$ and
output $o$, write $\mathrm{buf}(w)$ for the tape holding $w$ in cells $0, \dots, |w| - 1$ and blank
elsewhere, and write $\Lambda_{n,a,j,k}$ for the anchor-layout tapes, which hold $\mathrm{buf}$ of
$1^{n+1}$ on $\tau_0$, of $1^{|\mathrm{bits}(n)|}$ on $\tau_1$, of $1^a$ on $\tau_2$, of $1^j$ on
$\tau_{10}$, of $1^k$ on $\tau_{11}$, of $1$ on $\tau_{12}$, of $1^{m(n)}$ on $\tau_8$ and $\tau_9$
(with $m(n) = |\mathrm{bits}(A_f(c_f + 2))| + |\mathrm{bits}(n)|$ the clock counter width, $A_f$
and $c_f$ being the time constants of $f$'s call) and are blank on every other host tape. Let
$(\ell_i)$ be the ladder of $f$, suppose $n \ge 3$, and let $s$ be the stage of $n$ (the least $i$
with $n \le \ell_{i+1}$), $(j, k) = \mathrm{unpair}(s)$ and $a = \ell_s + 1$. Let $A_\kappa$ be the
locator's constant (a number determined by $\kappa$ alone). Then there are $t \le A_\kappa (n + 1)$
and a configuration $c'$ such that the locator runs in exactly $t$ steps from $\langle
\mathrm{anchor}, 1, \Lambda_{n,\ell_0 + 1,0,0}, 0, \varepsilon\rangle$ to $c'$ without being in an
escape state strictly before time $t$, and either (i) $n < \ell_{s+1}$ and $c'$ is the interior or
the timeout escape in its sweep state with some input head, tapes $T$ and heads $P$ such that every
host tape outside the round's scratch (the range of the clock's bank, which carries $f$'s call
$C_f$ with tape $0$ on $\tau_7$ and its counter on $\tau_8$, together with $\tau_3, \dots, \tau_6$)
agrees with $\Lambda_{n,a,j,k}$ and has its head at $0$; or (ii) $\ell_{s+1} = n$ and $c' = \langle
(\mathrm{top}, \mathrm{sweep}), 1, T, P, \varepsilon\rangle$ with $(T, P)$ of that kind, $T(\tau_4)
= \mathrm{buf}(1^{f(a) + a})$ and $P(\tau_4) = 0$.
\end{theorem}

\begin{theorem}[High host tapes are not scratch]
\lean{Turing.NTimeDiag.qy_not_scratch_ge}
\label{Turing.NTimeDiag.qy_not_scratch_ge}
\leanok
\uses{Turing.NTimeDiag.QyKit, Turing.NTimeDiag.QyKit.K, Turing.NTimeDiag.QyKit.bCf, Turing.NTimeDiag.QyKit.bClk, Turing.NTimeDiag.QyKit.tp_val, Turing.NTimeDiag.qyBank_val, Turing.NTimeDiag.qyBank_zero, Turing.NTimeDiag.qyScratch}
\difficulty{4}
Let $\kappa$ be a kit for $f, g : \mathbb{N} \to \mathbb{N}$, with named tapes $\tau_0, \dots,
\tau_{15}$, whose input call $C_b$ and $f$'s call $C_f$ have $k_b$ and $k_f$ tapes. Every host tape
$t$ with index $t \ge 16 + k_b + k_f$ lies outside the round's scratch (the range of the clock's
bank, which carries $f$'s call $C_f$ with tape $0$ on $\tau_7$ and its counter on $\tau_8$,
together with $\tau_3, \dots, \tau_6$).
\end{theorem}

\begin{theorem}[Named tapes off 3 to 8 are not scratch]
\lean{Turing.NTimeDiag.qy_not_scratch_tp}
\label{Turing.NTimeDiag.qy_not_scratch_tp}
\leanok
\uses{Turing.NTimeDiag.QyKit, Turing.NTimeDiag.QyKit.bCf_ne, Turing.NTimeDiag.QyKit.bClk, Turing.NTimeDiag.QyKit.tp, Turing.NTimeDiag.QyKit.tp_inj, Turing.NTimeDiag.qyScratch}
\difficulty{4}
Let $\kappa$ be a kit for $f, g : \mathbb{N} \to \mathbb{N}$, with named tapes $\tau_0, \dots,
\tau_{15}$. For every $c < 16$ with $c < 3$ or $c > 8$, the named tape $\tau_c$ lies outside the
round's scratch (the range of the clock's bank, which carries $f$'s call $C_f$ with tape $0$ on
$\tau_7$ and its counter on $\tau_8$, together with $\tau_3, \dots, \tau_6$).
\end{theorem}

\begin{theorem}[Late bank tapes are not scratch]
\lean{Turing.NTimeDiag.qy_bank_not_scratch}
\label{Turing.NTimeDiag.qy_bank_not_scratch}
\leanok
\uses{Turing.NTimeDiag.QyKit, Turing.NTimeDiag.QyKit.K, Turing.NTimeDiag.qyBank, Turing.NTimeDiag.qyBank_val, Turing.NTimeDiag.qyScratch, Turing.NTimeDiag.qy_not_scratch_ge}
\difficulty{1}
Let $\kappa$ be a kit for $f, g : \mathbb{N} \to \mathbb{N}$, with named tapes $\tau_0, \dots,
\tau_{15}$, whose input call $C_b$ and $f$'s call $C_f$ have $k_b$ and $k_f$ tapes, and let
$K_\kappa$ be its number of host tapes. Let $\mathit{off}, k \in \mathbb{N}$ with $\mathit{off} + k
\le K_\kappa$ and $\mathit{off} \ge 16 + k_b + k_f$, and let $t_0$ be a host tape with index below
$\mathit{off}$. For the bank sending tape $0$ to $t_0$ and each tape $i \ne 0$ (of $k$ tapes) to
host tape $\mathit{off} + i$, every image of a tape $i \ne 0$ lies outside the round's scratch (the
range of the clock's bank, which carries $f$'s call $C_f$ with tape $0$ on $\tau_7$ and its counter
on $\tau_8$, together with $\tau_3, \dots, \tau_6$).
\end{theorem}

\begin{theorem}[Escape tapes on named tapes]
\lean{Turing.NTimeDiag.QyEscT.tp}
\label{Turing.NTimeDiag.QyEscT.tp}
\leanok
\uses{Turing.FinTM.bufferTape, Turing.NTimeDiag.QyEscT, Turing.NTimeDiag.QyKit, Turing.NTimeDiag.QyKit.K, Turing.NTimeDiag.QyKit.tp, Turing.NTimeDiag.qyLay, Turing.NTimeDiag.qy_not_scratch_tp}
\difficulty{1}
Let $\kappa$ be a kit for $f, g : \mathbb{N} \to \mathbb{N}$, with named tapes $\tau_0, \dots,
\tau_{15}$, let $n, a, j, k \in \mathbb{N}$, write $\mathrm{buf}(w)$ for the tape holding $w$ in
cells $0, \dots, |w| - 1$ and blank elsewhere, write $\Lambda_{n,a,j,k}$ for the anchor-layout
tapes, which hold $\mathrm{buf}$ of $1^{n+1}$ on $\tau_0$, of $1^{|\mathrm{bits}(n)|}$ on $\tau_1$,
of $1^a$ on $\tau_2$, of $1^j$ on $\tau_{10}$, of $1^k$ on $\tau_{11}$, of $1$ on $\tau_{12}$, of
$1^{m(n)}$ on $\tau_8$ and $\tau_9$ (with $m(n) = |\mathrm{bits}(A_f(c_f + 2))| +
|\mathrm{bits}(n)|$ the clock counter width, $A_f$ and $c_f$ being the time constants of $f$'s
call) and are blank on every other host tape, and say that $(T, P)$ are escape tapes for $(n, a, j,
k)$ when every host tape outside the round's scratch (the range of the clock's bank, which carries
$f$'s call $C_f$ with tape $0$ on $\tau_7$ and its counter on $\tau_8$, together with $\tau_3,
\dots, \tau_6$) agrees with $\Lambda_{n,a,j,k}$ and has its head at $0$. If $(T, P)$ are escape
tapes for $(n, a, j, k)$, then for every $c < 16$ with $c < 3$ or $c > 8$ we have $T(\tau_c) =
\Lambda_{n,a,j,k}(\tau_c)$ and $P(\tau_c) = 0$.
\end{theorem}

\begin{theorem}[Escape tapes on late bank tapes]
\lean{Turing.NTimeDiag.QyEscT.bank}
\label{Turing.NTimeDiag.QyEscT.bank}
\leanok
\uses{Turing.FinTM.bufferTape, Turing.NTimeDiag.QyEscT, Turing.NTimeDiag.QyKit, Turing.NTimeDiag.QyKit.K, Turing.NTimeDiag.QyKit.tp, Turing.NTimeDiag.qyBank, Turing.NTimeDiag.qyBank_ne_tp, Turing.NTimeDiag.qyLay, Turing.NTimeDiag.qy_bank_not_scratch}
\difficulty{3}
Let $\kappa$ be a kit for $f, g : \mathbb{N} \to \mathbb{N}$, with named tapes $\tau_0, \dots,
\tau_{15}$, whose input call $C_b$ and $f$'s call $C_f$ have $k_b$ and $k_f$ tapes, let $K_\kappa$
be its number of host tapes, let $n, a, j, k \in \mathbb{N}$, write $\mathrm{buf}(w)$ for the tape
holding $w$ in cells $0, \dots, |w| - 1$ and blank elsewhere, write $\Lambda_{n,a,j,k}$ for the
anchor-layout tapes, which hold $\mathrm{buf}$ of $1^{n+1}$ on $\tau_0$, of
$1^{|\mathrm{bits}(n)|}$ on $\tau_1$, of $1^a$ on $\tau_2$, of $1^j$ on $\tau_{10}$, of $1^k$ on
$\tau_{11}$, of $1$ on $\tau_{12}$, of $1^{m(n)}$ on $\tau_8$ and $\tau_9$ (with $m(n) =
|\mathrm{bits}(A_f(c_f + 2))| + |\mathrm{bits}(n)|$ the clock counter width, $A_f$ and $c_f$ being
the time constants of $f$'s call) and are blank on every other host tape, and say that $(T, P)$ are
escape tapes for $(n, a, j, k)$ when every host tape outside the round's scratch (the range of the
clock's bank, which carries $f$'s call $C_f$ with tape $0$ on $\tau_7$ and its counter on $\tau_8$,
together with $\tau_3, \dots, \tau_6$) agrees with $\Lambda_{n,a,j,k}$ and has its head at $0$. Let
$\mathit{off}, k' \in \mathbb{N}$ with $\mathit{off} + k' \le K_\kappa$ and $\mathit{off} \ge 16 +
k_b + k_f$, let $t_0$ be a host tape with index below $\mathit{off}$, and let $\beta$ be the bank
sending tape $0$ to $t_0$ and each tape $i \ne 0$ of $k'$ tapes to $\mathit{off} + i$. If $(T, P)$
are escape tapes for $(n, a, j, k)$, then for every $i \ne 0$ the tape $\beta(i)$ is blank in $T$
and has its head at $0$ in $P$.
\end{theorem}

\begin{theorem}[Bank of the length call avoids named tapes]
\lean{Turing.NTimeDiag.QyKit.bCl_ne}
\label{Turing.NTimeDiag.QyKit.bCl_ne}
\leanok
\uses{Turing.NTimeDiag.QyKit, Turing.NTimeDiag.QyKit.bCl, Turing.NTimeDiag.QyKit.tp, Turing.NTimeDiag.qyBank_ne_tp}
\difficulty{1}
Let $\kappa$ be a kit for $f, g : \mathbb{N} \to \mathbb{N}$, with named tapes $\tau_0, \dots,
\tau_{15}$, and consider the bank of its call $C_l$ (computing $y \mapsto \mathrm{bits}(|y|)$),
which sends tape $0$ to $\tau_{10}$ and the other tapes into its block. For every tape $i \ne 0$ of
$C_l$ and every $c < 16$, the bank's image of $i$ differs from $\tau_c$.
\end{theorem}

\begin{theorem}[Stage-code calls after an escape]
\lean{Turing.NTimeDiag.qy_code_calls}
\label{Turing.NTimeDiag.qy_code_calls}
\leanok
\uses{Turing.FinTM.bufferTape, Turing.MultiTapeTM.runFrom, Turing.NTimeDiag.QyEscT, Turing.NTimeDiag.QyEscT.bank, Turing.NTimeDiag.QyEscT.tp, Turing.NTimeDiag.QyKit, Turing.NTimeDiag.QyKit.K, Turing.NTimeDiag.QyKit.bAbJ, Turing.NTimeDiag.QyKit.bCl, Turing.NTimeDiag.QyKit.bCl_ne, Turing.NTimeDiag.QyKit.tp, Turing.NTimeDiag.QyKit.tp_inj, Turing.NTimeDiag.qyC, Turing.NTimeDiag.qyLay, Turing.NTimeDiag.qy_arg_bank, Turing.NTimeDiag.qy_callHost, Turing.embedEmitTM, Turing.seamReleaseTM}
\difficulty{4}
Let $\kappa$ be a kit for $f, g : \mathbb{N} \to \mathbb{N}$, with named tapes $\tau_0, \dots,
\tau_{15}$, let $x$ be a binary input of length $n$, write $\langle s, p, T, P, o\rangle$ for the
live configuration on input $x$ with state $s$, input head $p$, work tapes $T$, heads $P$ and
output $o$, write $\mathrm{buf}(w)$ for the tape holding $w$ in cells $0, \dots, |w| - 1$ and blank
elsewhere, and write $\Lambda_{n,a,j,k}$ for the anchor-layout tapes, which hold $\mathrm{buf}$ of
$1^{n+1}$ on $\tau_0$, of $1^{|\mathrm{bits}(n)|}$ on $\tau_1$, of $1^a$ on $\tau_2$, of $1^j$ on
$\tau_{10}$, of $1^k$ on $\tau_{11}$, of $1$ on $\tau_{12}$, of $1^{m(n)}$ on $\tau_8$ and $\tau_9$
(with $m(n) = |\mathrm{bits}(A_f(c_f + 2))| + |\mathrm{bits}(n)|$ the clock counter width, $A_f$
and $c_f$ being the time constants of $f$'s call) and are blank on every other host tape. Let $a,
j, k \in \mathbb{N}$, and say that $(T, P)$ are escape tapes for $(n, a, j, k)$ when every host
tape outside the round's scratch (the range of the clock's bank, which carries $f$'s call $C_f$
with tape $0$ on $\tau_7$ and its counter on $\tau_8$, together with $\tau_3, \dots, \tau_6$)
agrees with $\Lambda_{n,a,j,k}$ and has its head at $0$. Let $A_a$ and $A_l$ be the time constants
of the increment call $C_a$ ($y \mapsto y\,1$) and the length call $C_l$ ($y \mapsto
\mathrm{bits}(|y|)$), and suppose $(T, P)$ are escape tapes for $(n, a, j, k)$. Consider the
release adapter of $C_a$ embedded along its bank with tape $0$ on $\tau_{10}$ and the release
adapter of $C_l$ embedded along its bank (tape $0$ also on $\tau_{10}$), each started in its fresh
start state. Then (i) for some $t_1 \le A_a(j + 1)$ the first runs from $\langle \mathrm{start}, 1,
T, P, \varepsilon\rangle$ to the tagged exit state of $C_a$ with input head $1$, tapes $T$ with
$\tau_{10}$ replaced by $\mathrm{buf}(1^{j+1})$, heads $P$ and empty output at time $t_1$, and is
not in that exit state earlier; and (ii) for some $t_2 \le A_l(j + 2)$ the second runs from the
start state on those tapes to the tagged exit state of $C_l$ with input head $1$, tapes $T$ with
$\tau_{10}$ replaced by $\mathrm{buf}(\mathrm{bits}(j + 1))$, heads $P$ and empty output at time
$t_2$, and is not in that exit state earlier.
\end{theorem}

\begin{theorem}[Input seek after an escape]
\lean{Turing.NTimeDiag.qy_seek_run}
\label{Turing.NTimeDiag.qy_seek_run}
\leanok
\uses{Turing.MultiTapeTM.runFrom, Turing.NTimeDiag.QyEscT, Turing.NTimeDiag.QyEscT.tp, Turing.NTimeDiag.QyKit, Turing.NTimeDiag.QyKit.K, Turing.NTimeDiag.QyKit.tp, Turing.NTimeDiag.QyKit.tp_inj, Turing.NTimeDiag.QyKit.tp_ne, Turing.NTimeDiag.qyC, Turing.NTimeDiag.qyEmb2, Turing.NTimeDiag.qyLay, Turing.NTimeDiag.qy_seekHost, Turing.embedEmitTM}
\difficulty{3}
Let $\kappa$ be a kit for $f, g : \mathbb{N} \to \mathbb{N}$, with named tapes $\tau_0, \dots,
\tau_{15}$, let $x$ be a binary input of length $n$, write $\langle s, p, T, P, o\rangle$ for the
live configuration on input $x$ with state $s$, input head $p$, work tapes $T$, heads $P$ and
output $o$, write $\mathrm{buf}(w)$ for the tape holding $w$ in cells $0, \dots, |w| - 1$ and blank
elsewhere, and write $\Lambda_{n,a,j,k}$ for the anchor-layout tapes, which hold $\mathrm{buf}$ of
$1^{n+1}$ on $\tau_0$, of $1^{|\mathrm{bits}(n)|}$ on $\tau_1$, of $1^a$ on $\tau_2$, of $1^j$ on
$\tau_{10}$, of $1^k$ on $\tau_{11}$, of $1$ on $\tau_{12}$, of $1^{m(n)}$ on $\tau_8$ and $\tau_9$
(with $m(n) = |\mathrm{bits}(A_f(c_f + 2))| + |\mathrm{bits}(n)|$ the clock counter width, $A_f$
and $c_f$ being the time constants of $f$'s call) and are blank on every other host tape. Let $a,
j, k \in \mathbb{N}$, and say that $(T, P)$ are escape tapes for $(n, a, j, k)$ when every host
tape outside the round's scratch (the range of the clock's bank, which carries $f$'s call $C_f$
with tape $0$ on $\tau_7$ and its counter on $\tau_8$, together with $\tau_3, \dots, \tau_6$)
agrees with $\Lambda_{n,a,j,k}$ and has its head at $0$. Let $R$ be the kit's input seek, a
two-tape machine with entry state $r_E$, exit state $r_X$ and constant $r_C$, embedded on the host
tapes $\tau_{12}$ and $\tau_{13}$. If $(T, P)$ are escape tapes for $(n, a, j, k)$, then for every
input-head position $q \le n + 1$ there is $t \le r_C(n + 4)$ such that the embedded seek runs from
$\langle r_E, q, T, P, \varepsilon \rangle$ to $\langle r_X, 1, T, P, \varepsilon\rangle$ in
exactly $t$ steps and is not in $r_X$ at any earlier time.
\end{theorem}

\begin{theorem}[Head of the middle writer]
\lean{Turing.NTimeDiag.qy_mid_code}
\label{Turing.NTimeDiag.qy_mid_code}
\leanok
\uses{Turing.FinTM.bufferTape, Turing.NTimeDiag.QyEscT, Turing.NTimeDiag.QyKit, Turing.NTimeDiag.QyKit.K, Turing.NTimeDiag.QyKit.tp, Turing.NTimeDiag.qyC, Turing.NTimeDiag.qyMid, Turing.NTimeDiag.qyMidM, Turing.NTimeDiag.qyMidRed, Turing.NTimeDiag.qy_code_calls, Turing.NTimeDiag.qy_guard_of_cut, Turing.NTimeDiag.qy_hosted, Turing.NTimeDiag.qy_seek_run, Turing.Segment, Turing.Segment.trans}
\difficulty{4}
Let $\kappa$ be a kit for $f, g : \mathbb{N} \to \mathbb{N}$, with named tapes $\tau_0, \dots,
\tau_{15}$, let $x$ be a binary input of length $n$, write $\langle s, p, T, P, o\rangle$ for the
live configuration on input $x$ with state $s$, input head $p$, work tapes $T$, heads $P$ and
output $o$, write $\mathrm{buf}(w)$ for the tape holding $w$ in cells $0, \dots, |w| - 1$ and blank
elsewhere, and write $\Lambda_{n,a,j,k}$ for the anchor-layout tapes, which hold $\mathrm{buf}$ of
$1^{n+1}$ on $\tau_0$, of $1^{|\mathrm{bits}(n)|}$ on $\tau_1$, of $1^a$ on $\tau_2$, of $1^j$ on
$\tau_{10}$, of $1^k$ on $\tau_{11}$, of $1$ on $\tau_{12}$, of $1^{m(n)}$ on $\tau_8$ and $\tau_9$
(with $m(n) = |\mathrm{bits}(A_f(c_f + 2))| + |\mathrm{bits}(n)|$ the clock counter width, $A_f$
and $c_f$ being the time constants of $f$'s call) and are blank on every other host tape. Let $a,
j, k \in \mathbb{N}$, and say that $(T, P)$ are escape tapes for $(n, a, j, k)$ when every host
tape outside the round's scratch (the range of the clock's bank, which carries $f$'s call $C_f$
with tape $0$ on $\tau_7$ and its counter on $\tau_8$, together with $\tau_3, \dots, \tau_6$)
agrees with $\Lambda_{n,a,j,k}$ and has its head at $0$. Let $r_E$ and $r_C$ be the entry state and
constant of the kit's input seek, and $A_a$, $A_l$ the time constants of its increment and length
calls. Let $(T, P)$ be escape tapes for $(n, a, j, k)$ and $q \le n + 1$ an input-head position.
Then for some $t \le r_C(n + 4) + A_a(j + 1) + A_l(j + 2) + 3$ the middle's writer runs in exactly
$t$ steps from $\langle (m_0, r_E), q, T, P, \varepsilon\rangle$ to $\langle (m_3, \mathrm{start}),
1, T', P, \varepsilon\rangle$, the fresh start of the call of $g$'s witness, where $T'$ is $T$ with
$\tau_{10}$ replaced by $\mathrm{buf}(\mathrm{bits}(j + 1))$.
\end{theorem}

\begin{theorem}[Head of the top writer]
\lean{Turing.NTimeDiag.qy_top_code}
\label{Turing.NTimeDiag.qy_top_code}
\leanok
\uses{Turing.FinTM.bufferTape, Turing.NTimeDiag.QyEscT, Turing.NTimeDiag.QyKit, Turing.NTimeDiag.QyKit.K, Turing.NTimeDiag.QyKit.tp, Turing.NTimeDiag.qyC, Turing.NTimeDiag.qyTop, Turing.NTimeDiag.qyTopM, Turing.NTimeDiag.qyTopRed, Turing.NTimeDiag.qy_code_calls, Turing.NTimeDiag.qy_guard_of_cut, Turing.NTimeDiag.qy_hosted, Turing.Segment, Turing.Segment.trans}
\difficulty{3}
Let $\kappa$ be a kit for $f, g : \mathbb{N} \to \mathbb{N}$, with named tapes $\tau_0, \dots,
\tau_{15}$, let $x$ be a binary input of length $n$, write $\langle s, p, T, P, o\rangle$ for the
live configuration on input $x$ with state $s$, input head $p$, work tapes $T$, heads $P$ and
output $o$, write $\mathrm{buf}(w)$ for the tape holding $w$ in cells $0, \dots, |w| - 1$ and blank
elsewhere, and write $\Lambda_{n,a,j,k}$ for the anchor-layout tapes, which hold $\mathrm{buf}$ of
$1^{n+1}$ on $\tau_0$, of $1^{|\mathrm{bits}(n)|}$ on $\tau_1$, of $1^a$ on $\tau_2$, of $1^j$ on
$\tau_{10}$, of $1^k$ on $\tau_{11}$, of $1$ on $\tau_{12}$, of $1^{m(n)}$ on $\tau_8$ and $\tau_9$
(with $m(n) = |\mathrm{bits}(A_f(c_f + 2))| + |\mathrm{bits}(n)|$ the clock counter width, $A_f$
and $c_f$ being the time constants of $f$'s call) and are blank on every other host tape. Let $a,
j, k \in \mathbb{N}$, and say that $(T, P)$ are escape tapes for $(n, a, j, k)$ when every host
tape outside the round's scratch (the range of the clock's bank, which carries $f$'s call $C_f$
with tape $0$ on $\tau_7$ and its counter on $\tau_8$, together with $\tau_3, \dots, \tau_6$)
agrees with $\Lambda_{n,a,j,k}$ and has its head at $0$. Let $A_a$ and $A_l$ be the time constants
of the kit's increment and length calls and let $(T, P)$ be escape tapes for $(n, a, j, k)$. Then
for some $t \le A_a(j + 1) + A_l(j + 2) + 2$ the top's writer runs in exactly $t$ steps from
$\langle (t_1, \mathrm{start}), 1, T, P, \varepsilon\rangle$ to $\langle (t_3, \mathrm{start}), 1,
T', P, \varepsilon\rangle$, where $T'$ is $T$ with $\tau_{10}$ replaced by
$\mathrm{buf}(\mathrm{bits}(j + 1))$.
\end{theorem}

\begin{theorem}[Doubling twice is quadrupling]
\lean{Turing.NTimeDiag.qy_flatMap_dbl_dbl}
\label{Turing.NTimeDiag.qy_flatMap_dbl_dbl}
\leanok
\difficulty{3}
For every binary word $B = b_1 \cdots b_r$, doubling every symbol of $B$ and then doubling every
symbol of the result gives the word $b_1^4 b_2^4 \cdots b_r^4$ in which each symbol of $B$ is
repeated four times.
\end{theorem}

\begin{theorem}[The writer's output word]
\lean{Turing.NTimeDiag.qy_writer_word}
\label{Turing.NTimeDiag.qy_writer_word}
\leanok
\uses{Turing.NTimeDiag.qyBody, Turing.NTimeDiag.qy_flatMap_dbl_dbl, Turing.pairEncode, Turing.pairEncode_eq}
\difficulty{4}
For a binary word $w$, a bit $\delta$ and $r \in \mathbb{N}$, let $\mathrm{body}_{\delta,r}(w)$ be
the word obtained from $w$ (with its last symbol removed if $\delta = 1$) by repeating every symbol
$r + 1$ times, and write $\langle u, v\rangle$ for the pairing encoding of binary words. For every
bit $b$ and binary words $B, W, N$, \[
b\;\mathrm{body}_{0,3}(B)\;0011\;\mathrm{body}_{1,1}(W)\;01\;\mathrm{body}_{0,0}(N) =
b\,\langle\langle B, W'\rangle, N\rangle, \] where $W'$ is $W$ with its last symbol removed.
\end{theorem}

\begin{theorem}[A hosted phase that halts]
\lean{Turing.NTimeDiag.qy_hosted_halt}
\label{Turing.NTimeDiag.qy_hosted_halt}
\leanok
\uses{Turing.Cfg, Turing.MultiTapeTM, Turing.MultiTapeTM.runFrom, Turing.NTimeDiag.qyC, Turing.NTimeDiag.qyC_mapState, Turing.phaseHostTM, Turing.phaseHostTM_run_le}
\difficulty{3}
Let $(M_i)_{i \in I}$ be a family of $K$-tape machines over bits with state sets $\Phi_i$, let $H$
be the phase host with start state $q_0$ and redirect table $r$ (a redirected state takes one
stationary silent step to its target; any other tagged state $(i, s)$ follows $M_i$), and write
$\langle s, p, T, P, o\rangle$ for the live configuration on input $x$ with state $s$, input head
$p$, work tapes $T$, heads $P$ and output $o$. Let $k \in I$, $e \in \Phi_k$, and suppose that
$M_k$, run from $\langle e, p, T, P, o\rangle$ for $t$ steps, reaches a halted configuration $c'$,
while $r_k$ is undefined on every state that this run occupies at times before $t$. Then $H$, run
from $\langle (k, e), p, T, P, o\rangle$ for $t$ steps, is halted and its output equals the output
of $c'$.
\end{theorem}

\begin{theorem}[An emitter runs to its halt]
\lean{Turing.NTimeDiag.qy_emit_stop}
\label{Turing.NTimeDiag.qy_emit_stop}
\leanok
\uses{Turing.FinTM.bufferTape, Turing.MultiTapeTM.runFrom, Turing.MultiTapeTM.runFrom_add, Turing.NTimeDiag.qyBody, Turing.NTimeDiag.qyC, Turing.NTimeDiag.qyEmitTM, Turing.NTimeDiag.qy_emit_halt, Turing.NTimeDiag.qy_emit_run}
\difficulty{3}
Let $i$ be one of $K$ work tapes, let $\mathit{pre}, \mathit{suf}, w$ be binary words, $r \in
\mathbb{N}$ and $\delta$ a bit, and let $\mathrm{body}_{\delta,r}(w)$ be $w$ (with its last symbol
removed if $\delta = 1$) with every symbol repeated $r + 1$ times. Consider the emitter that emits
$\mathit{pre}$, then the word under head $i$ processed in this way, then $\mathit{suf}$, and write
$\langle s, p, T, P, o\rangle$ for the live configuration on input $x$ with state $s$, input head
$p$, work tapes $T$, heads $P$ and output $o$. If tape $i$ of $T$ holds exactly $w$ from cell $0$
and $P(i) = 0$, then there is $t \le |\mathit{pre}| + |\mathit{suf}| + 4 + |w|(r + 3)$ such that
the emitter, run for $t$ steps from $\langle \text{first prefix state}, p, T, P, o\rangle$, is
halted with output $o\,\mathit{pre}\,\mathrm{body}_{\delta,r}(w)\,\mathit{suf}$.
\end{theorem}

\begin{theorem}[Bank of the call of g avoids named tapes]
\lean{Turing.NTimeDiag.QyKit.bCg_ne}
\label{Turing.NTimeDiag.QyKit.bCg_ne}
\leanok
\uses{Turing.NTimeDiag.QyKit, Turing.NTimeDiag.QyKit.bCg, Turing.NTimeDiag.QyKit.tp, Turing.NTimeDiag.qyBank_ne_tp}
\difficulty{1}
Let $\kappa$ be a kit for $f, g : \mathbb{N} \to \mathbb{N}$, with named tapes $\tau_0, \dots,
\tau_{15}$, and consider the bank of its input call $C_g$ (computing $\mathrm{bits}(g(|x|))$),
which sends tape $0$ to $\tau_{14}$ and the other tapes into its block. For every tape $i \ne 0$ of
$C_g$ and every $c < 16$, the bank's image of $i$ differs from $\tau_c$.
\end{theorem}

\begin{theorem}[Bank of the square call avoids named tapes]
\lean{Turing.NTimeDiag.QyKit.bPb_ne}
\label{Turing.NTimeDiag.QyKit.bPb_ne}
\leanok
\uses{Turing.NTimeDiag.QyKit, Turing.NTimeDiag.QyKit.bPb, Turing.NTimeDiag.QyKit.tp, Turing.NTimeDiag.qyBank_ne_tp}
\difficulty{1}
Let $\kappa$ be a kit for $f, g : \mathbb{N} \to \mathbb{N}$, with named tapes $\tau_0, \dots,
\tau_{15}$, and consider the bank of its call $C_q$ (computing $y \mapsto \mathrm{bits}((|y| +
1)^2)$), which sends tape $0$ to $\tau_4$ and the other tapes into its block. For every tape $i \ne
0$ of $C_q$ and every $c < 16$, the bank's image of $i$ differs from $\tau_c$.
\end{theorem}

\begin{definition}[Time bound of the middle writer]
\lean{Turing.NTimeDiag.qyMidB}
\label{Turing.NTimeDiag.qyMidB}
\leanok
\uses{Turing.NTimeDiag.QyKit}
Let $\kappa$ be a kit for $f, g : \mathbb{N} \to \mathbb{N}$ with input-seek constant $r_C$ and
time constants $A_a$, $A_l$, $A_g$ of its increment call, length call and call of $g$'s witness.
For $n, j \in \mathbb{N}$ the middle writer's time bound is \[ r_C(n + 4) + A_a(j + 1) + A_l(j + 2)
+ A_g(g(n) + 1) + 6g(n) + 4(j + 1) + 3(n + 1) + 30. \]
\end{definition}

\begin{definition}[Time bound of the top writer]
\lean{Turing.NTimeDiag.qyTopB}
\label{Turing.NTimeDiag.qyTopB}
\leanok
\uses{Turing.NTimeDiag.QyKit}
Let $\kappa$ be a kit for $f, g : \mathbb{N} \to \mathbb{N}$ with time constants $A_a$, $A_l$,
$A_q$ of its increment call, length call and squaring call ($y \mapsto \mathrm{bits}((|y| +
1)^2)$). For $X, a, j \in \mathbb{N}$ the top writer's time bound is \[ A_a(j + 1) + A_l(j + 2) +
A_q X^3 + 6X^2 + 4(j + 1) + 3a + 30. \]
\end{definition}

\begin{theorem}[The middle writer's run]
\lean{Turing.NTimeDiag.qy_mid_run}
\label{Turing.NTimeDiag.qy_mid_run}
\leanok
\uses{Turing.FinTM.bufferTape, Turing.MultiTapeTM.runFrom, Turing.MultiTapeTM.runFrom_add, Turing.NTimeDiag.QyEscT, Turing.NTimeDiag.QyEscT.bank, Turing.NTimeDiag.QyEscT.tp, Turing.NTimeDiag.QyKit, Turing.NTimeDiag.QyKit.K, Turing.NTimeDiag.QyKit.bCg, Turing.NTimeDiag.QyKit.bCg_ne, Turing.NTimeDiag.QyKit.tp, Turing.NTimeDiag.QyKit.tp_inj, Turing.NTimeDiag.QyKit.tp_ne, Turing.NTimeDiag.qyC, Turing.NTimeDiag.qyLay, Turing.NTimeDiag.qyMid, Turing.NTimeDiag.qyMidB, Turing.NTimeDiag.qyMidM, Turing.NTimeDiag.qyMidRed, Turing.NTimeDiag.qy_blank_bank, Turing.NTimeDiag.qy_callHost, Turing.NTimeDiag.qy_emit_run, Turing.NTimeDiag.qy_emit_stop, Turing.NTimeDiag.qy_guard_of_cut, Turing.NTimeDiag.qy_hosted, Turing.NTimeDiag.qy_hosted_halt, Turing.NTimeDiag.qy_mid_code, Turing.NTimeDiag.qy_writer_word, Turing.Segment.trans, Turing.length_bits_le_self, Turing.pairEncode, Turing.segment_iff}
\difficulty{5}
Let $\kappa$ be a kit for $f, g : \mathbb{N} \to \mathbb{N}$, with named tapes $\tau_0, \dots,
\tau_{15}$, let $x$ be a binary input of length $n$, write $\langle s, p, T, P, o\rangle$ for the
live configuration on input $x$ with state $s$, input head $p$, work tapes $T$, heads $P$ and
output $o$, write $\mathrm{buf}(w)$ for the tape holding $w$ in cells $0, \dots, |w| - 1$ and blank
elsewhere, and write $\Lambda_{n,a,j,k}$ for the anchor-layout tapes, which hold $\mathrm{buf}$ of
$1^{n+1}$ on $\tau_0$, of $1^{|\mathrm{bits}(n)|}$ on $\tau_1$, of $1^a$ on $\tau_2$, of $1^j$ on
$\tau_{10}$, of $1^k$ on $\tau_{11}$, of $1$ on $\tau_{12}$, of $1^{m(n)}$ on $\tau_8$ and $\tau_9$
(with $m(n) = |\mathrm{bits}(A_f(c_f + 2))| + |\mathrm{bits}(n)|$ the clock counter width, $A_f$
and $c_f$ being the time constants of $f$'s call) and are blank on every other host tape. Let $a,
j, k \in \mathbb{N}$, and say that $(T, P)$ are escape tapes for $(n, a, j, k)$ when every host
tape outside the round's scratch (the range of the clock's bank, which carries $f$'s call $C_f$
with tape $0$ on $\tau_7$ and its counter on $\tau_8$, together with $\tau_3, \dots, \tau_6$)
agrees with $\Lambda_{n,a,j,k}$ and has its head at $0$. Let $r_E, r_C$ be the entry state and
constant of the kit's input seek, $A_a, A_l, A_g$ the time constants of its increment call, length
call and call of $g$'s witness, and $B = r_C(n + 4) + A_a(j + 1) + A_l(j + 2) + A_g(g(n) + 1) +
6g(n) + 4(j + 1) + 3(n + 1) + 30$. Write $\langle u, v\rangle$ for the pairing encoding of binary
words. If $(T, P)$ are escape tapes for $(n, a, j, k)$ and $q \le n + 1$, then for some $t \le B$
the middle's writer, run for $t$ steps from $\langle (m_0, r_E), q, T, P, \varepsilon\rangle$, is
halted with output $0\,\langle\langle \mathrm{bits}(g(n)), \alpha\rangle, 1^{n+1}\rangle$, where
$\alpha$ is $\mathrm{bits}(j + 1)$ with its last symbol removed.
\end{theorem}

\begin{theorem}[The top writer's run]
\lean{Turing.NTimeDiag.qy_top_run}
\label{Turing.NTimeDiag.qy_top_run}
\leanok
\uses{Turing.FinTM.bufferTape, Turing.MultiTapeTM.runFrom, Turing.MultiTapeTM.runFrom_add, Turing.NTimeDiag.QyEscT, Turing.NTimeDiag.QyEscT.bank, Turing.NTimeDiag.QyEscT.tp, Turing.NTimeDiag.QyKit, Turing.NTimeDiag.QyKit.K, Turing.NTimeDiag.QyKit.bPb, Turing.NTimeDiag.QyKit.bPb_ne, Turing.NTimeDiag.QyKit.tp, Turing.NTimeDiag.QyKit.tp_inj, Turing.NTimeDiag.QyKit.tp_ne, Turing.NTimeDiag.qyC, Turing.NTimeDiag.qyLay, Turing.NTimeDiag.qyTop, Turing.NTimeDiag.qyTopB, Turing.NTimeDiag.qyTopM, Turing.NTimeDiag.qyTopRed, Turing.NTimeDiag.qy_arg_bank, Turing.NTimeDiag.qy_callHost, Turing.NTimeDiag.qy_emit_run, Turing.NTimeDiag.qy_emit_stop, Turing.NTimeDiag.qy_guard_of_cut, Turing.NTimeDiag.qy_hosted, Turing.NTimeDiag.qy_hosted_halt, Turing.NTimeDiag.qy_top_code, Turing.NTimeDiag.qy_writer_word, Turing.Segment.trans, Turing.length_bits_le_self, Turing.pairEncode, Turing.segment_iff}
\difficulty{5}
Let $\kappa$ be a kit for $f, g : \mathbb{N} \to \mathbb{N}$, with named tapes $\tau_0, \dots,
\tau_{15}$, let $x$ be a binary input of length $n$, write $\langle s, p, T, P, o\rangle$ for the
live configuration on input $x$ with state $s$, input head $p$, work tapes $T$, heads $P$ and
output $o$, write $\mathrm{buf}(w)$ for the tape holding $w$ in cells $0, \dots, |w| - 1$ and blank
elsewhere, and write $\Lambda_{n,a,j,k}$ for the anchor-layout tapes, which hold $\mathrm{buf}$ of
$1^{n+1}$ on $\tau_0$, of $1^{|\mathrm{bits}(n)|}$ on $\tau_1$, of $1^a$ on $\tau_2$, of $1^j$ on
$\tau_{10}$, of $1^k$ on $\tau_{11}$, of $1$ on $\tau_{12}$, of $1^{m(n)}$ on $\tau_8$ and $\tau_9$
(with $m(n) = |\mathrm{bits}(A_f(c_f + 2))| + |\mathrm{bits}(n)|$ the clock counter width, $A_f$
and $c_f$ being the time constants of $f$'s call) and are blank on every other host tape. Let $a,
j, k \in \mathbb{N}$, and say that $(T, P)$ are escape tapes for $(n, a, j, k)$ when every host
tape outside the round's scratch (the range of the clock's bank, which carries $f$'s call $C_f$
with tape $0$ on $\tau_7$ and its counter on $\tau_8$, together with $\tau_3, \dots, \tau_6$)
agrees with $\Lambda_{n,a,j,k}$ and has its head at $0$. Let $A_a, A_l, A_q$ be the time constants
of the kit's increment, length and squaring calls, put $X = f(a) + a + 1$ and $B = A_a(j + 1) +
A_l(j + 2) + A_q X^3 + 6X^2 + 4(j + 1) + 3a + 30$, and write $\langle u, v\rangle$ for the pairing
encoding of binary words. If $(T, P)$ are escape tapes for $(n, a, j, k)$, $T(\tau_4) =
\mathrm{buf}(1^{f(a) + a})$ and $P(\tau_4) = 0$, then for some $t \le B$ the top's writer, run for
$t$ steps from $\langle (t_1, \mathrm{start}), 1, T, P, \varepsilon\rangle$, is halted with output
$1\,\langle\langle \mathrm{bits}(X^2), \alpha\rangle, 1^a\rangle$, where $\alpha$ is
$\mathrm{bits}(j + 1)$ with its last symbol removed.
\end{theorem}

\begin{theorem}[Initial configuration of the main machine]
\lean{Turing.NTimeDiag.qy_initCfg}
\label{Turing.NTimeDiag.qy_initCfg}
\leanok
\uses{Turing.Cfg.ofWords, Turing.Cfg.ofWords_inputPos, Turing.Cfg.ofWords_output, Turing.Cfg.ofWords_state, Turing.Cfg.ofWords_workTapePos, Turing.Cfg.ofWords_workTapes, Turing.FinTM.bufferTape, Turing.MultiTapeTM.initCfg, Turing.NTimeDiag.QyKit, Turing.NTimeDiag.qyC, Turing.NTimeDiag.qyMainM, Turing.NTimeDiag.qyMainRed, Turing.NTimeDiag.qyM₁, Turing.initCfg_ofWords, Turing.phaseHostTM, Turing.phaseHostTM_q₀}
\difficulty{3}
Let $\kappa$ be a kit for $f, g : \mathbb{N} \to \mathbb{N}$, let $M_1$ be its main machine (the
phase host of the setup, the locator, the two writers and the empty-output emitter), let $x$ be a
binary input, and write $\langle s, p, T, P, o\rangle$ for the live configuration on input $x$ with
state $s$, input head $p$, work tapes $T$, heads $P$ and output $o$. The initial configuration of
$M_1$ on $x$ is $\langle (\mathrm{setup}, (s_1, \mathrm{start})), 1, \text{all tapes blank}, 0,
\varepsilon \rangle$.
\end{theorem}

\begin{theorem}[Main dispatch from the setup]
\lean{Turing.NTimeDiag.qy_mainRed_setup}
\label{Turing.NTimeDiag.qy_mainRed_setup}
\leanok
\uses{Turing.NTimeDiag.QyKit, Turing.NTimeDiag.QySPh, Turing.NTimeDiag.qyMainRed, Turing.NTimeDiag.qySetF}
\difficulty{4}
Let $\kappa$ be a kit for $f, g : \mathbb{N} \to \mathbb{N}$. For every state $s$ of the setup that
is not a setup exit (the done state of phase $s_{10}$, or phase $s_6$ with verdict $<$ or $=$), the
main machine's dispatch table is undefined on the setup state $s$.
\end{theorem}

\begin{theorem}[Main dispatch from the locator]
\lean{Turing.NTimeDiag.qy_mainRed_loc}
\label{Turing.NTimeDiag.qy_mainRed_loc}
\leanok
\uses{Turing.NTimeDiag.QyKit, Turing.NTimeDiag.QyLK.isEsc, Turing.NTimeDiag.QyLPh, Turing.NTimeDiag.qyLocE, Turing.NTimeDiag.qyMainRed}
\difficulty{3}
Let $\kappa$ be a kit for $f, g : \mathbb{N} \to \mathbb{N}$. For every locator state $s$ that is
not an escape state (its phase is not the interior, top or timeout escape), the main machine's
dispatch table is undefined on the locator state $s$.
\end{theorem}

\begin{definition}[The writers' constant]
\lean{Turing.NTimeDiag.qyWk}
\label{Turing.NTimeDiag.qyWk}
\leanok
\uses{Turing.NTimeDiag.QyKit}
Let $\kappa$ be a kit for $f, g : \mathbb{N} \to \mathbb{N}$ with input-seek constant $r_C$ and
time constants $A_a, A_l, A_g, A_q$ of its increment call, length call, call of $g$'s witness and
squaring call. The writers' constant is $W_\kappa = 4r_C + A_a + 2A_l + A_g + A_q + 60$.
\end{definition}

\begin{definition}[The main machine's constant]
\lean{Turing.NTimeDiag.qyMainA}
\label{Turing.NTimeDiag.qyMainA}
\leanok
\uses{Turing.NTimeDiag.QyKit, Turing.NTimeDiag.qyLoopA, Turing.NTimeDiag.qyWk}
Let $\kappa$ be a kit for $f, g : \mathbb{N} \to \mathbb{N}$ with time constants $A_b$ of its input
call for $\mathrm{bits}(|x|)$, $A_p$ of its prefix call and $A_f, c_f$ of $f$'s call. Its main
constant is $(A_b + A_p + 2|\mathrm{bits}(A_f(c_f + 2))| + 60) + A^{\mathrm{loop}}_\kappa +
W_\kappa + 8$, where $A^{\mathrm{loop}}_\kappa$ is the locator's constant and $W_\kappa$ the
writers' constant.
\end{definition}

\begin{theorem}[Middle writer's bound is linear]
\lean{Turing.NTimeDiag.qy_midB_le}
\label{Turing.NTimeDiag.qy_midB_le}
\leanok
\uses{Turing.NTimeDiag.QyKit, Turing.NTimeDiag.qyMidB, Turing.NTimeDiag.qyWk}
\difficulty{4}
Let $\kappa$ be a kit for $f, g : \mathbb{N} \to \mathbb{N}$ with input-seek constant $r_C$ and
time constants $A_a, A_l, A_g, A_q$ of its increment, length, $g$ and squaring calls, and let
$W_\kappa = 4r_C + A_a + 2A_l + A_g + A_q + 60$. For all $n, j \in \mathbb{N}$ with $j \le n \le
g(n)$, \[ r_C(n + 4) + A_a(j + 1) + A_l(j + 2) + A_g(g(n) + 1) + 6g(n) + 4(j + 1) + 3(n + 1) + 30
\le W_\kappa (g(n) + 1). \]
\end{theorem}

\begin{theorem}[Top writer's bound is linear]
\lean{Turing.NTimeDiag.qy_topB_le}
\label{Turing.NTimeDiag.qy_topB_le}
\leanok
\uses{Turing.NTimeDiag.QyKit, Turing.NTimeDiag.qyTopB, Turing.NTimeDiag.qyWk}
\difficulty{4}
Let $\kappa$ be a kit for $f, g : \mathbb{N} \to \mathbb{N}$ with input-seek constant $r_C$ and
time constants $A_a, A_l, A_g, A_q$ of its increment, length, $g$ and squaring calls, and let
$W_\kappa = 4r_C + A_a + 2A_l + A_g + A_q + 60$. For all $n, X, a, j \in \mathbb{N}$ with $X^3 \le
n$, $X^2 \le n$, $a \le n$ and $j \le n$, \[ A_a(j + 1) + A_l(j + 2) + A_q X^3 + 6X^2 + 4(j + 1) +
3a + 30 \le W_\kappa (n + 1). \]
\end{theorem}

\begin{theorem}[Main machine on inputs of length at least 3]
\lean{Turing.NTimeDiag.qy_main_long}
\label{Turing.NTimeDiag.qy_main_long}
\leanok
\uses{Turing.FinTM.bufferTape, Turing.MultiTapeTM.initCfg, Turing.MultiTapeTM.runFrom, Turing.MultiTapeTM.runFrom_add, Turing.NTimeDiag.QyKit, Turing.NTimeDiag.budget_eq, Turing.NTimeDiag.ladder, Turing.NTimeDiag.ladder_zero, Turing.NTimeDiag.query, Turing.NTimeDiag.query_eq, Turing.NTimeDiag.qyAnc, Turing.NTimeDiag.qyAt, Turing.NTimeDiag.qyC, Turing.NTimeDiag.qyLay, Turing.NTimeDiag.qyLayT, Turing.NTimeDiag.qyLoopA, Turing.NTimeDiag.qyMainA, Turing.NTimeDiag.qyMainM, Turing.NTimeDiag.qyMainRed, Turing.NTimeDiag.qyM₁, Turing.NTimeDiag.qySetB, Turing.NTimeDiag.qyWk, Turing.NTimeDiag.qy_hosted, Turing.NTimeDiag.qy_hosted_halt, Turing.NTimeDiag.qy_initCfg, Turing.NTimeDiag.qy_ladder_succ_eq, Turing.NTimeDiag.qy_loop, Turing.NTimeDiag.qy_mainRed_loc, Turing.NTimeDiag.qy_mainRed_setup, Turing.NTimeDiag.qy_midB_le, Turing.NTimeDiag.qy_mid_run, Turing.NTimeDiag.qy_setup_long, Turing.NTimeDiag.qy_stage_lt_log, Turing.NTimeDiag.qy_topB_le, Turing.NTimeDiag.qy_top_run, Turing.NTimeDiag.stage, Turing.NTimeDiag.stageCode_eq, Turing.NTimeDiag.stage_spec, Turing.Segment, Turing.Segment.trans, Turing.phaseHostTM, Turing.segment_iff}
\difficulty{6}
Let $\kappa$ be a kit for $f, g : \mathbb{N} \to \mathbb{N}$ with $n \le g(n)$ for every $n$, let
$M_1$ be its main machine and $A_\kappa$ its main constant (a number determined by $\kappa$ alone),
and let $q_{f,g}$ be the stage query of the lazy diagonalizer. For every binary input $x$ of length
$n \ge 3$ there is $t \le A_\kappa (g(n) + 1)$ such that $M_1$, run on $x$ from its initial
configuration for $t$ steps, is halted with output $q_{f,g}(1^n)$.
\end{theorem}

\begin{theorem}[Main machine on inputs of length at most 2]
\lean{Turing.NTimeDiag.qy_main_short}
\label{Turing.NTimeDiag.qy_main_short}
\leanok
\uses{Turing.MultiTapeTM.initCfg, Turing.MultiTapeTM.runFrom, Turing.MultiTapeTM.runFrom_add, Turing.NTimeDiag.QyKit, Turing.NTimeDiag.QyKit.tp, Turing.NTimeDiag.QyKit.tp_inj, Turing.NTimeDiag.query, Turing.NTimeDiag.query_of_not_guard, Turing.NTimeDiag.qyBody, Turing.NTimeDiag.qyLoopA, Turing.NTimeDiag.qyMainA, Turing.NTimeDiag.qyMainM, Turing.NTimeDiag.qyMainRed, Turing.NTimeDiag.qyM₁, Turing.NTimeDiag.qySetB, Turing.NTimeDiag.qySetT5, Turing.NTimeDiag.qyWk, Turing.NTimeDiag.qy_emit_stop, Turing.NTimeDiag.qy_hosted, Turing.NTimeDiag.qy_hosted_halt, Turing.NTimeDiag.qy_initCfg, Turing.NTimeDiag.qy_mainRed_setup, Turing.NTimeDiag.qy_setup_short, Turing.phaseHostTM, Turing.segment_iff}
\difficulty{5}
Let $\kappa$ be a kit for $f, g : \mathbb{N} \to \mathbb{N}$ with $n \le g(n)$ for every $n$, let
$M_1$ be its main machine and $A_\kappa$ its main constant (a number determined by $\kappa$ alone),
and let $q_{f,g}$ be the stage query of the lazy diagonalizer. For every binary input $x$ of length
$n \le 2$ there is $t \le A_\kappa (g(n) + 1)$ such that $M_1$, run on $x$ from its initial
configuration for $t$ steps, is halted with output $q_{f,g}(1^n)$ (the empty word).
\end{theorem}

\begin{theorem}[A machine computing the query on unary words]
\lean{Turing.NTimeDiag.qy_main}
\label{Turing.NTimeDiag.qy_main}
\leanok
\uses{Complexity.TimeConstructible, Complexity.timeConstructible_iff, Turing.FinTM, Turing.FinTM.ComputesFunInTime, Turing.FinTM.ComputesInTime.mono, Turing.FinTM.computesFunInTime_iff, Turing.FinTM.computesInTime_iff, Turing.NTimeDiag.query, Turing.NTimeDiag.qyMainA, Turing.NTimeDiag.qyM₁, Turing.NTimeDiag.qy_kit_exists, Turing.NTimeDiag.qy_main_long, Turing.NTimeDiag.qy_main_short}
\difficulty{3}
Let $f, g : \mathbb{N} \to \mathbb{N}$ be time-constructible and let $q_{f,g}$ be the stage query
of the lazy diagonalizer. Then there are a finite deterministic Turing machine $M$ over bits and a
constant $A$ such that $M$ computes the map $x \mapsto q_{f,g}(1^{|x|})$ within $A(g(n) + 1)$ steps
on inputs of length $n$.
\end{theorem}

\begin{theorem}[The stage query within a multiple of $g(n) + 1$]
\lean{Turing.NTimeDiag.exists_queryTM}
\label{Turing.NTimeDiag.exists_queryTM}
\leanok
\uses{Complexity.TimeConstructible, Turing.FinTM, Turing.FinTM.ComputesFunInTime, Turing.NTimeDiag.query, Turing.NTimeDiag.qy_main, Turing.NTimeDiag.qy_query_of_main}
\difficulty{2}
Let $f, g : \mathbb{N} \to \mathbb{N}$ be time-constructible and let $q_{f,g}$ be the stage query
of the lazy diagonalizer. Then there are a finite deterministic Turing machine $G$ over bits and a
constant $C$ such that $G$ computes $q_{f,g}$ within $C(g(n) + 1)$ steps on inputs of length $n$.
\end{theorem}

\begin{definition}[The one-cell take]
\lean{Turing.NTimeDiag.dzTakeTM}
\label{Turing.NTimeDiag.dzTakeTM}
\leanok
\uses{Turing.MultiTapeTM}
For $K \in \mathbb{N}$ and a work tape $i < K$, the one-cell take is the $K$-tape machine over bits
whose states are the optional values of an optional bit, with start state $\mathrm{none}$. From
$\mathrm{none}$, in one step it reads the cell under head $i$ (a bit or blank, call it $o$), writes
a blank there, moves head $i$ one cell to the right and enters state $\mathrm{some}(o)$; it leaves
the input head and all other tapes and heads in place and emits nothing. In a state
$\mathrm{some}(o)$ it stays in that state without writing, moving or emitting.
\end{definition}

\begin{theorem}[One step of the take]
\lean{Turing.NTimeDiag.dz_take_run}
\label{Turing.NTimeDiag.dz_take_run}
\leanok
\uses{Turing.Action.apply_workTapes, Turing.Cfg, Turing.Cfg.inputSymbol, Turing.Cfg.workTapeSymbols, Turing.Cfg.workTapeSymbols_eq, Turing.MultiTapeTM.runFrom, Turing.MultiTapeTM.runFrom_one, Turing.MultiTapeTM.step_of_some, Turing.NTimeDiag.dzTakeTM, Turing.NTimeDiag.qyC, Turing.moveInputPos_zero}
\difficulty{4}
Let $i$ be one of $K$ work tapes, let $x$ be a binary input, and write $\langle s, p, T, P,
o\rangle$ for the live configuration on input $x$ with state $s$, input head $p$, work tapes $T$,
heads $P$ and output $o$. For every input-head position $p$, tapes $T$, heads $P$ and output $o$,
one step of the one-cell take on tape $i$ from $\langle \mathrm{none}, p, T, P, o \rangle$ gives
$\langle \mathrm{some}(T_i(P_i)), p, T', P', o\rangle$, where $T'$ is $T$ with cell $P_i$ of tape
$i$ made blank and $P'$ is $P$ with $P'_i = P_i + 1$.
\end{theorem}

\begin{theorem}[Recognising a silent transport]
\lean{Turing.NTimeDiag.dz_silentCfg_eq}
\label{Turing.NTimeDiag.dz_silentCfg_eq}
\leanok
\uses{Turing.Cfg, Turing.FinTM.bufferTape, Turing.embedSilentCfg, Turing.embedSilentCfg_cap, Turing.embedSilentCfg_selected_pos, Turing.embedSilentCfg_selected_tape, Turing.embedSilentCfg_unselected}
\difficulty{4}
Let $\iota$ inject $m$ tapes into $K$ tapes, let $\mathit{cap}$ be a host tape outside the range of
$\iota$, let $T$ and $P$ be host tape contents and heads, let $\mathit{pre}, o$ be binary words,
and write $\mathrm{buf}(w)$ for the tape holding $w$ in cells $0, \dots, |w| - 1$ and blank
elsewhere. Let $c$ be an $m$-tape configuration and $d$ a $K$-tape configuration on a common input
with the same state and input head, such that $d$ has output $o$; each tape $i$ of $c$ and its head
equal tape $\iota(i)$ of $d$ and its head; tape $\mathit{cap}$ of $d$ is
$\mathrm{buf}(\mathit{pre}\,o_c)$, with $o_c$ the output of $c$, and its head is at
$|\mathit{pre}\,o_c|$; and every other host tape $j$ of $d$ equals $T_j$ with head $P_j$. Then $d$
is the silent transport of $c$ with frame $(T, P)$, recorded prefix $\mathit{pre}$ and output $o$.
\end{theorem}

\begin{theorem}[A captured run in normal form]
\lean{Turing.NTimeDiag.dz_silent}
\label{Turing.NTimeDiag.dz_silent}
\leanok
\uses{Turing.FinTM.bufferTape, Turing.MultiTapeTM, Turing.MultiTapeTM.runFrom, Turing.NTimeDiag.dz_silentCfg_eq, Turing.NTimeDiag.qyC, Turing.embedSilentCfg, Turing.embedSilentTM, Turing.embedSilentTM_runFrom}
\difficulty{3}
Let $\iota$ inject $m$ tapes into $K$ tapes, let $\mathit{cap}$ be a host tape outside the range of
$\iota$, let $M$ be an $m$-tape machine over bits with states $s, s'$, let $x$ be a binary input,
write $\langle s, p, T, P, o\rangle$ for the live configuration on input $x$ with state $s$, input
head $p$, work tapes $T$, heads $P$ and output $o$, and write $\mathrm{buf}(w)$ for the tape
holding $w$ in cells $0, \dots, |w| - 1$ and blank elsewhere. Let $T, T'$ be host tapes, $P, P'$
host heads, $o, \mathit{pre}, o'$ binary words, $T_0, P_0$ tapes and heads for $M$, $q$ an
input-head position and $t \in \mathbb{N}$. Suppose $T(\mathit{cap}) = \mathrm{buf}(\mathit{pre})$
and $P(\mathit{cap}) = |\mathit{pre}|$; that $M$ runs in $t$ steps from $\langle s, q, T \circ
\iota, P \circ \iota, \varepsilon\rangle$ to $\langle s', q, T_0, P_0, o'\rangle$; that
$T'(\iota(i)) = T_0(i)$ and $P'(\iota(i)) = P_0(i)$ for every $i$; that $T'(\mathit{cap}) =
\mathrm{buf}(\mathit{pre}\,o')$ with $P'(\mathit{cap}) = |\mathit{pre}\,o'|$; and that $T'$, $P'$
agree with $T$, $P$ on every other tape. Then the capture wrapper of $M$ along $\iota$ with capture
tape $\mathit{cap}$ runs in $t$ steps from $\langle s, q, T, P, o\rangle$ to $\langle s', q, T',
P', o\rangle$, and at every time $u$ its state equals the state of $M$ at time $u$ from $\langle s,
q, T \circ \iota, P \circ \iota, \varepsilon\rangle$.
\end{theorem}

\begin{definition}[Tape map of the hosted clock]
\lean{Turing.NTimeDiag.dzClkEmb}
\label{Turing.NTimeDiag.dzClkEmb}
\leanok
Let $K, m \in \mathbb{N}$, let $\mathit{buf}, \mathit{cap}, \mathit{ctr}$ be three pairwise
distinct host tapes among $K$, and let $\beta$ inject $m$ tapes into the $K$ host tapes with none
of $\mathit{buf}, \mathit{cap}, \mathit{ctr}$ in its range. The hosted clock's tape map is the
injection of the $(1 + m) + 1 + 1$ tapes of a step clock over a virtually hosted $m$-tape machine
into the host tapes that sends the virtual-input buffer (the first tape) to $\mathit{buf}$, the
hosted tape $i$ to $\beta(i)$, the capture tape to $\mathit{cap}$ and the counter (the last tape)
to $\mathit{ctr}$.
\end{definition}

\begin{theorem}[The hosted clock's tape map, tape by tape]
\lean{Turing.NTimeDiag.dzClkEmb_eqs}
\label{Turing.NTimeDiag.dzClkEmb_eqs}
\leanok
\uses{Turing.NTimeDiag.dzClkEmb}
\difficulty{1}
Let $K, m \in \mathbb{N}$, let $\mathit{buf}, \mathit{cap}, \mathit{ctr}$ be pairwise distinct host
tapes among $K$, and let $\beta$ inject $m$ tapes into the host tapes avoiding all three. Consider
the hosted clock's tape map for these data, an injection of the $(1 + m) + 1 + 1$ tapes (indexed
$0, \dots, m + 2$) of a step clock over a virtually hosted $m$-tape machine into the host tapes. It
sends the last tape (index $m + 2$) to $\mathit{ctr}$, the tape of index $m + 1$ to $\mathit{cap}$,
the tape of index $0$ to $\mathit{buf}$, and the tape of index $1 + i$ to $\beta(i)$ for every $i <
m$.
\end{theorem}

\begin{definition}[The hosted clock's clause]
\lean{Turing.NTimeDiag.DzClockSpec}
\label{Turing.NTimeDiag.DzClockSpec}
\leanok
\uses{Turing.Cfg, Turing.Cfg.Stopped, Turing.FinTM.bufferTape, Turing.MultiTapeTM, Turing.MultiTapeTM.runFrom, Turing.NTimeDiag.dzClkEmb, Turing.NTimeDiag.qyC, Turing.embedEmitTM}
Let $M$ be an $m$-tape machine over bits with state set $S$, let $\mathrm{stop} : S \to \{0, 1\}$,
and call a configuration of $M$ stopped when it is halted or its state satisfies $\mathrm{stop}$.
Let $H$ be a machine with $(1 + m) + 1 + 1$ tapes and state set $Q$, with maps $\mathrm{inj} : S
\times \{0, 1\} \to Q$ and $\mathrm{fin}, \mathrm{tmo}$ from optional elements of $S \times \{0,
1\}$ to $Q$. Write $\langle s, p, T, P, o\rangle$ for the live configuration on input $x$ with
state $s$, input head $p$, work tapes $T$, heads $P$ and output $o$, and write $\mathrm{buf}(w)$
for the tape holding $w$ in cells $0, \dots, |w| - 1$ and blank elsewhere. The hosted clock's
clause holds when the following is true for every number $K$ of host tapes, binary words $x, y$,
pairwise distinct host tapes $\mathit{buf}, \mathit{cap}, \mathit{ctr}$, injection $\beta$ of $M$'s
tapes into the host tapes avoiding all three, host tapes $T$, heads $P$ and output $o$, live
configuration $c$ of $M$ on input $y$ in a state $q$ with input head $1$ and empty output, and
binary word $w$ whose value as a binary numeral (least significant bit first) is $d$. Let $\sigma$
send $H$'s first tape to $\mathit{buf}$, its hosted tapes along $\beta$, its capture tape to
$\mathit{cap}$ and its last tape to $\mathit{ctr}$, and run $H$ embedded along $\sigma$ on input
$x$ from $\langle \mathrm{inj}(q, 1), 1, T, P, o\rangle$, assuming that $\mathit{buf}$ holds
$\mathrm{buf}(y)$, each $\beta(i)$ holds tape $i$ of $c$ with $c$'s head, $\mathit{cap}$ is blank,
$\mathit{ctr}$ holds $\mathrm{buf}(w)$, and the heads of $\mathit{buf}$, $\mathit{cap}$,
$\mathit{ctr}$ are at $0$. (a) If $M$ from $c$ is first stopped at time $\tau \le d$, then at some
time $t \le 5\tau + 2|w| + 2$ the run is at $\langle \mathrm{fin}(r), 1, T', P', o\rangle$, where
$r$ is $M$'s state at time $\tau$ paired with some bit (or none if $M$ has halted), with
$\mathit{cap}$ holding $\mathrm{buf}$ of $M$'s output at time $\tau$ and its head just after it.
(b) If $M$ from $c$ is stopped at no time $u \le d$, then at some time $t \le 5d + 2|w| + 2$ the
run is at $\langle \mathrm{tmo}(r), 1, T', P', o\rangle$ for some $r$, $T'$, $P'$. In both cases no
state $\mathrm{fin}(r')$ or $\mathrm{tmo}(r')$ occurs before time $t$.
\end{definition}

\begin{theorem}[Existence of a hosted clock]
\lean{Turing.NTimeDiag.dz_clock_exists}
\label{Turing.NTimeDiag.dz_clock_exists}
\leanok
\uses{Turing.Cfg, Turing.Cfg.Stopped, Turing.Cfg.stopped_iff, Turing.FinTM.VirtualTag, Turing.FinTM.VirtualTag_iff, Turing.FinTM.bufferTape, Turing.MultiTapeTM, Turing.MultiTapeTM.runFrom, Turing.NTimeDiag.DzClockSpec, Turing.NTimeDiag.dzClkEmb, Turing.NTimeDiag.dzClkEmb_eqs, Turing.NTimeDiag.qyC, Turing.NTimeDiag.qy_embed_gen, Turing.NTimeDiag.qy_embed_state, Turing.clockCfg, Turing.embedEmitCfg, Turing.embedEmitCfg_selected_pos, Turing.embedEmitCfg_selected_tape, Turing.embedEmitTM, Turing.embedSilentCfg_cap, Turing.embedSilentCfg_selected_pos, Turing.embedSilentCfg_selected_tape, Turing.exists_stepClockTMVia, Turing.vhostCap, Turing.vhostCap_eq, Turing.vhostCfg, Turing.vhostSilentCfg, Turing.vhostSilentCfg_eq, Turing.vhostSilentTM, Turing.vhostSilentTM_runFrom, Turing.wordPatch}
\difficulty{6}
Let $M$ be an $m$-tape machine over bits with a finite state set $S$ (with decidable equality) and
let $\mathrm{stop} : S \to \{0, 1\}$. Then there are a finite state set $Q$ with decidable
equality, a machine $H$ over bits with $(1 + m) + 1 + 1$ tapes and state set $Q$, and maps
$\mathrm{inj} : S \times \{0, 1\} \to Q$ and $\mathrm{fin}, \mathrm{tmo}$ from optional elements of
$S \times \{0, 1\}$ to $Q$, such that $\mathrm{fin}$ is injective, $\mathrm{fin}(a) \ne
\mathrm{tmo}(b)$ for all $a, b$, and $(H, \mathrm{inj}, \mathrm{fin}, \mathrm{tmo})$ satisfies the
hosted clock's clause for $M$ and $\mathrm{stop}$: run with $M$'s configuration on its hosted tapes
and a counter of value $d$, it reaches $\mathrm{fin}$ of $M$'s stopped state within $5\tau + 2|w| +
2$ steps when $M$ first stops at a time $\tau \le d$ (capturing $M$'s output), and reaches a
timeout state within $5d + 2|w| + 2$ steps otherwise, with $w$ the counter's word.
\end{theorem}

\begin{definition}[Inner second component of a pair]
\lean{Turing.NTimeDiag.dzAlpha}
\label{Turing.NTimeDiag.dzAlpha}
\leanok
\uses{Turing.pairDecode}
Write $\langle u, v\rangle$ for the pairing encoding of binary words. For a binary word $z$, decode
$z$ as a pair $\langle u, v\rangle$ (taking $u$ to be empty if decoding fails) and decode $u$ as a
pair $\langle u_1, u_2\rangle$ (taking $u_2$ to be empty if decoding fails); the inner second
component of $z$ is $u_2$.
\end{definition}

\begin{theorem}[Inner second component of a nested pair]
\lean{Turing.NTimeDiag.dz_alpha_eq}
\label{Turing.NTimeDiag.dz_alpha_eq}
\leanok
\uses{Turing.NTimeDiag.dzAlpha, Turing.pairDecode_pairEncode, Turing.pairEncode}
\difficulty{1}
Write $\langle u, v\rangle$ for the pairing encoding of binary words. For all binary words $B,
\alpha, N$, the inner second component of $\langle\langle B, \alpha\rangle, N\rangle$ (the second
component of the first component of its decoding) equals $\alpha$.
\end{theorem}

\begin{theorem}[Output length bounded by running time]
\lean{Turing.NTimeDiag.dz_out_le}
\label{Turing.NTimeDiag.dz_out_le}
\leanok
\uses{Turing.FinTM, Turing.FinTM.ComputesFunInTime, Turing.FinTM.computesFunInTime_iff, Turing.FinTM.computesInTime_iff, Turing.MultiTapeTM.output_length_le}
\difficulty{2}
Let $M$ be a finite deterministic Turing machine over bits, let $F$ map binary words to binary
words, and let $T : \mathbb{N} \to \mathbb{N}$. If $M$ computes $F$ within time $T$, then for every
binary word $y$ we have $|F(y)| \le T(|y|)$.
\end{theorem}

\begin{definition}[The diagonalizer's kit]
\lean{Turing.NTimeDiag.DzKit}
\label{Turing.NTimeDiag.DzKit}
\leanok
\uses{Turing.CodeNDTM.toFinNDTM, Turing.EffectiveNDMachineCode, Turing.FinTM, Turing.MultiTapeTM, Turing.MultiTapeTM.runFrom, Turing.NDTM.initCfg, Turing.NDTM.runWith, Turing.NTimeDiag.DzClockSpec, Turing.NTimeDiag.QyInSpec, Turing.NTimeDiag.QySpec, Turing.NTimeDiag.dzAlpha, Turing.NTimeDiag.query, Turing.pairEncode, Turing.streamStartCfg}
Let $f, g : \mathbb{N} \to \mathbb{N}$, let $c$ be an effective coding scheme for nondeterministic
machines, let $B_F$ be a finite deterministic machine over bits, let $q_{f,g}$ be the stage query,
write $\langle u, v\rangle$ for the pairing encoding of binary words, and write $\mathrm{seam}_x(e,
w)$ for the configuration on input $x$ in state $e$ with input head $1$, tape $0$ holding $w$,
other tapes blank, heads at $0$ and empty output. A kit for $f, g, c, B_F$ consists of: the fact $n
\le g(n)$ for all $n$; input calls (each a finite deterministic machine with at least one tape and
entry and exit states that reaches $\mathrm{seam}_x(\mathit{exit}, F(x))$ from
$\mathrm{seam}_x(\mathit{entry}, \varepsilon)$ for the first time within the given bound on $|x|$)
for $q_{f,g}$, for $x \mapsto \mathrm{bits}(g(|x|))$ and for $x \mapsto 1x$, each within $A(g(n) +
1)$ for its own constant $A$, the first with $|q_{f,g}(x)| \le A(g(|x|) + 1)$; install calls
(reaching $\mathrm{seam}_x(\mathit{exit}, F(y))$ from $\mathrm{seam}_x(\mathit{entry}, y)$ for the
first time within the bound on $y$) for $y \mapsto \mathrm{bits}(|y|)$ and for the inner second
component $y \mapsto \alpha$ of $y = \langle\langle B, \alpha\rangle, N\rangle$, each within $A(|y|
+ 1)$, the latter with output length at most $A(|y| + 1)$; an interpreter core, a $k_R$-tape
deterministic machine $R$ with finitely many states, a stream tape $j_R$ and states
$\mathit{entry}$, $\mathit{acc} \ne \mathit{rej}$ such that for each code $\alpha$ there is $C$
with: for all $x$, $t$ and choice words $s$ with $t \le |s|$, $R$ started in $\mathit{entry}$ on
input $\langle\langle \mathrm{bits}(t), \alpha\rangle, x\rangle$ with $s$ on tape $j_R$ and other
tapes blank is, at some time $T \le C(t + 1)$, in $\mathit{acc}$ if the machine decoded from
$\alpha$, run on $x$ with the first $t$ choices of $s$, has halted with output $1$, and in
$\mathit{rej}$ otherwise, and in neither state earlier; a hosted clock over $R$ with stop predicate
``the state is $\mathit{acc}$ or $\mathit{rej}$'' whose finish map is injective and never equals
its timeout map, satisfying the hosted clock's clause; and a hosted clock over $B_F$ with the stop
predicate that never holds, whose finish and timeout maps never agree, satisfying the hosted
clock's clause.
\end{definition}

\begin{theorem}[Existence of the diagonalizer's kit]
\lean{Turing.NTimeDiag.dz_kit_exists}
\label{Turing.NTimeDiag.dz_kit_exists}
\leanok
\uses{Complexity.TimeConstructible, Complexity.timeConstructible_iff, Turing.Cfg, Turing.CodeNDTM.toFinNDTM, Turing.EffectiveNDMachineCode, Turing.FinTM, Turing.FinTM.ComputesFunInTime, Turing.FinTM.bufferTape, Turing.FinTM.computesFunInTime_comp, Turing.FinTM.computesFunInTime_iff, Turing.FinTM.computesFunInTime_lengthBits, Turing.FinTM.computesFunInTime_prepend, Turing.FinTM.pairExtractTM_fst, Turing.FinTM.pairExtractTM_snd, Turing.MultiTapeTM.runFrom, Turing.NDTM.initCfg, Turing.NDTM.runWith, Turing.NTimeDiag.DzKit, Turing.NTimeDiag.QyInSpec, Turing.NTimeDiag.QySpec, Turing.NTimeDiag.QySpec.mono, Turing.NTimeDiag.dzAlpha, Turing.NTimeDiag.dz_clock_exists, Turing.NTimeDiag.dz_out_le, Turing.NTimeDiag.exists_queryTM, Turing.NTimeDiag.query, Turing.NTimeDiag.qy_inputSpec, Turing.NTimeDiag.qy_installSpec, Turing.exists_ndInterpreterCore, Turing.length_bits_le_self, Turing.pairDecode, Turing.pairEncode, Turing.streamStartCfg}
\difficulty{5}
Let $f, g : \mathbb{N} \to \mathbb{N}$ be time-constructible, let $c$ be an effective coding scheme
for nondeterministic machines and let $B_F$ be a finite deterministic machine over bits. Then a
diagonalizer's kit for $f, g, c, B_F$ exists: input calls for the stage query, for
$\mathrm{bits}(g(|x|))$ and for prepending $1$, each running within $A(g(n) + 1)$ steps on inputs of
length $n$ for its own constant $A$ (and the stage query's output has length at most
$A(g(|x|) + 1)$); install calls for $\mathrm{bits}(|y|)$ and for the inner second component of a
nested pair, each within $A(|y| + 1)$ steps (and the latter's output has length at most
$A(|y| + 1)$); an interpreter core for the codes of $c$; and hosted clocks over the core and over
$B_F$.
\end{theorem}

\begin{definition}[Tape count of the diagonalizer]
\lean{Turing.NTimeDiag.DzKit.K}
\label{Turing.NTimeDiag.DzKit.K}
\leanok
\uses{Turing.EffectiveNDMachineCode, Turing.FinTM, Turing.NTimeDiag.DzKit}
Let $\kappa$ be a diagonalizer's kit for $f, g : \mathbb{N} \to \mathbb{N}$, a coding scheme and a
decider $B_F$, whose calls for the stage query, $g$'s witness, the prefix, the length and the inner
second component have $k_q, k_g, k_p, k_l, k_a$ tapes and whose interpreter core has $k_R$ tapes.
The diagonalizer's tape count is $K_\kappa = 14 + k_q + k_g + k_p + k_l + k_a + k_R + k_{B_F}$,
where $k_{B_F}$ is the number of tapes of $B_F$.
\end{definition}

\begin{theorem}[At least fourteen tapes]
\lean{Turing.NTimeDiag.DzKit.le_K}
\label{Turing.NTimeDiag.DzKit.le_K}
\leanok
\uses{Turing.EffectiveNDMachineCode, Turing.FinTM, Turing.NTimeDiag.DzKit, Turing.NTimeDiag.DzKit.K}
\difficulty{2}
For every diagonalizer's kit $\kappa$ for $f, g : \mathbb{N} \to \mathbb{N}$, a coding scheme and a
decider, its tape count $K_\kappa$ satisfies $14 \le K_\kappa$.
\end{theorem}

\begin{definition}[Named tapes of the diagonalizer]
\lean{Turing.NTimeDiag.DzKit.tp}
\label{Turing.NTimeDiag.DzKit.tp}
\leanok
\uses{Turing.EffectiveNDMachineCode, Turing.FinTM, Turing.NTimeDiag.DzKit, Turing.NTimeDiag.DzKit.K, Turing.NTimeDiag.DzKit.le_K}
Let $\kappa$ be a diagonalizer's kit for $f, g : \mathbb{N} \to \mathbb{N}$, a coding scheme and a
decider, with tape count $K_\kappa \ge 14$. For $i < 14$ the named tape $\tau_i$ is the host tape
of index $i$; the roles are: $\tau_0$ the query, $\tau_1$ the bits of $g(n)$, $\tau_2$ the word
$1^{n+1}$, $\tau_3$ the clock's counter, $\tau_4, \tau_5$ the block's expansion counters, $\tau_6$
the block, $\tau_7$ the ruler, $\tau_8$ the query's tail, $\tau_9$ the budget word, $\tau_{10}$ its
expansion and bits, $\tau_{11}$ the core's input, $\tau_{12}$ the hosted capture and $\tau_{13}$
the decider's first tape.
\end{definition}

\begin{theorem}[Index of a named tape]
\lean{Turing.NTimeDiag.DzKit.tp_val}
\label{Turing.NTimeDiag.DzKit.tp_val}
\leanok
\uses{Turing.EffectiveNDMachineCode, Turing.FinTM, Turing.NTimeDiag.DzKit, Turing.NTimeDiag.DzKit.tp}
\difficulty{1}
Let $\kappa$ be a diagonalizer's kit for $f, g : \mathbb{N} \to \mathbb{N}$, a coding scheme and a
decider. For every $i < 14$ the named tape $\tau_i$ of $\kappa$ has index $i$.
\end{theorem}

\begin{theorem}[Named tapes are distinct]
\lean{Turing.NTimeDiag.DzKit.tp_ne}
\label{Turing.NTimeDiag.DzKit.tp_ne}
\leanok
\uses{Turing.EffectiveNDMachineCode, Turing.FinTM, Turing.NTimeDiag.DzKit, Turing.NTimeDiag.DzKit.le_K, Turing.NTimeDiag.DzKit.tp}
\difficulty{1}
Let $\kappa$ be a diagonalizer's kit for $f, g : \mathbb{N} \to \mathbb{N}$, a coding scheme and a
decider. For all $a, b < 14$ with $a \ne b$, the named tapes satisfy $\tau_a \ne \tau_b$.
\end{theorem}

\begin{definition}[A tape layout]
\lean{Turing.NTimeDiag.dzLay}
\label{Turing.NTimeDiag.dzLay}
\leanok
Let $K \in \mathbb{N}$, let $L$ assign a value of some type to each of the indices $0, \dots, 13$,
and let $d$ be a value of that type. The layout of $L$ with default $d$ assigns to each host tape
$t < K$ the value $L(t)$ if $t < 14$ and $d$ otherwise.
\end{definition}

\begin{theorem}[A layout on a named tape]
\lean{Turing.NTimeDiag.dzLay_tp}
\label{Turing.NTimeDiag.dzLay_tp}
\leanok
\uses{Turing.EffectiveNDMachineCode, Turing.FinTM, Turing.NTimeDiag.DzKit, Turing.NTimeDiag.DzKit.tp, Turing.NTimeDiag.DzKit.tp_val, Turing.NTimeDiag.dzLay}
\difficulty{2}
Let $\kappa$ be a diagonalizer's kit for $f, g : \mathbb{N} \to \mathbb{N}$, a coding scheme and a
decider, let $L$ assign a value to each index $i < 14$ and let $d$ be a default value. For every $i
< 14$, the layout of $L$ with default $d$ takes the value $L(i)$ at the named tape $\tau_i$.
\end{theorem}

\begin{theorem}[A layout off the named tapes]
\lean{Turing.NTimeDiag.dzLay_off}
\label{Turing.NTimeDiag.dzLay_off}
\leanok
\uses{Turing.NTimeDiag.dzLay}
\difficulty{2}
Let $K \in \mathbb{N}$, let $L$ assign a value of some type to each index $i < 14$ and let $d$ be a value of that type. For every host tape $t < K$ with $14 \le t$, the layout of $L$ with default $d$ (which gives tape $t$ the value $L(t)$ if $t < 14$ and $d$ otherwise) takes the value $d$ at $t$.
\end{theorem}

\begin{theorem}[Updating a layout on a named tape]
\lean{Turing.NTimeDiag.dzLay_update}
\label{Turing.NTimeDiag.dzLay_update}
\leanok
\uses{Turing.EffectiveNDMachineCode, Turing.FinTM, Turing.NTimeDiag.DzKit, Turing.NTimeDiag.DzKit.tp, Turing.NTimeDiag.DzKit.tp_ne, Turing.NTimeDiag.DzKit.tp_val, Turing.NTimeDiag.dzLay, Turing.NTimeDiag.dzLay_off, Turing.NTimeDiag.dzLay_tp}
\difficulty{4}
Let $\kappa$ be a diagonalizer's kit for $f, g : \mathbb{N} \to \mathbb{N}$, an effective coding scheme $c$ and a finite deterministic decider $B_F$ over bits, with $K_\kappa$ host tapes and named tapes $\tau_0, \dots, \tau_{13}$ (the host tapes of indices $0, \dots, 13$). Let $L$ assign a value to each index $i < 14$, let $d$ be a default value, and write $\mathrm{lay}(L, d)$ for the layout giving each host tape $\tau_i$ the value $L(i)$ and every other host tape the value $d$. For every $i < 14$ and every value $v$, updating $\mathrm{lay}(L, d)$ at $\tau_i$ to $v$ gives $\mathrm{lay}(L', d)$, where $L'$ is $L$ updated at $i$ to $v$.
\end{theorem}

\begin{theorem}[Constant tapes form a layout]
\lean{Turing.NTimeDiag.dzLay_const}
\label{Turing.NTimeDiag.dzLay_const}
\leanok
\uses{Turing.NTimeDiag.dzLay}
\difficulty{2}
Let $K \in \mathbb{N}$ and let $d$ be a value. The constant assignment of $d$ to all $K$ host tapes equals the layout of the constant assignment $i \mapsto d$ on the indices $i < 14$ with default $d$.
\end{theorem}

\begin{theorem}[Bank tapes lie off the named tapes]
\lean{Turing.NTimeDiag.dz_bank_off}
\label{Turing.NTimeDiag.dz_bank_off}
\leanok
\uses{Turing.NTimeDiag.qyBank, Turing.NTimeDiag.qyBank_val}
\difficulty{2}
Let $K, \mathit{off}, k \in \mathbb{N}$ with $\mathit{off} + k \le K$ and $14 \le \mathit{off}$, and let $t_0$ be one of $K$ tapes with $t_0 < \mathit{off}$. Consider the bank sending $0$ to $t_0$ and each $i \ne 0$ with $i < k$ to tape $\mathit{off} + i$. For every $i < k$ with $i \ne 0$, the tape it assigns to $i$ has index at least $14$.
\end{theorem}

\begin{definition}[Bank of the stage-query call]
\lean{Turing.NTimeDiag.DzKit.bQ}
\label{Turing.NTimeDiag.DzKit.bQ}
\leanok
\uses{Turing.EffectiveNDMachineCode, Turing.FinTM, Turing.NTimeDiag.DzKit, Turing.NTimeDiag.DzKit.K, Turing.NTimeDiag.DzKit.tp, Turing.NTimeDiag.DzKit.tp_val, Turing.NTimeDiag.qyBank}
Let $\kappa$ be a diagonalizer's kit for $f, g : \mathbb{N} \to \mathbb{N}$, an effective coding scheme $c$ and a finite deterministic decider $B_F$ over bits, with $K_\kappa$ host tapes and named tapes $\tau_0, \dots, \tau_{13}$ (the host tapes of indices $0, \dots, 13$). The bank of the stage-query call is the injection of the $k_q$ tapes of the kit's input call for the stage query into the $K_\kappa$ host tapes that sends tape $0$ to $\tau_0$ and each tape $i \ne 0$ to host tape $14 + i$.
\end{definition}

\begin{definition}[Bank of the call for g]
\lean{Turing.NTimeDiag.DzKit.bG}
\label{Turing.NTimeDiag.DzKit.bG}
\leanok
\uses{Turing.EffectiveNDMachineCode, Turing.FinTM, Turing.NTimeDiag.DzKit, Turing.NTimeDiag.DzKit.K, Turing.NTimeDiag.DzKit.tp, Turing.NTimeDiag.DzKit.tp_val, Turing.NTimeDiag.qyBank}
Let $\kappa$ be a diagonalizer's kit for $f, g : \mathbb{N} \to \mathbb{N}$, an effective coding scheme $c$ and a finite deterministic decider $B_F$ over bits, with $K_\kappa$ host tapes and named tapes $\tau_0, \dots, \tau_{13}$ (the host tapes of indices $0, \dots, 13$). With $k_q, k_g, k_p, k_l, k_a$ the numbers of tapes of the kit's calls for the stage query, for $\mathrm{bits}(g(|x|))$, for prepending $1$, for $y \mapsto \mathrm{bits}(|y|)$ and for the inner second component, and $k_R$ the number of tapes of its interpreter core, the bank of $g$'s call is the injection of the $k_g$ tapes of the call computing $\mathrm{bits}(g(|x|))$ into the host tapes that sends tape $0$ to $\tau_1$ and each tape $i \ne 0$ to host tape $14 + k_q + i$.
\end{definition}

\begin{definition}[Bank of the prepend call]
\lean{Turing.NTimeDiag.DzKit.bP}
\label{Turing.NTimeDiag.DzKit.bP}
\leanok
\uses{Turing.EffectiveNDMachineCode, Turing.FinTM, Turing.NTimeDiag.DzKit, Turing.NTimeDiag.DzKit.K, Turing.NTimeDiag.DzKit.tp, Turing.NTimeDiag.DzKit.tp_val, Turing.NTimeDiag.qyBank}
Let $\kappa$ be a diagonalizer's kit for $f, g : \mathbb{N} \to \mathbb{N}$, an effective coding scheme $c$ and a finite deterministic decider $B_F$ over bits, with $K_\kappa$ host tapes and named tapes $\tau_0, \dots, \tau_{13}$ (the host tapes of indices $0, \dots, 13$). With $k_q, k_g, k_p, k_l, k_a$ the numbers of tapes of the kit's calls for the stage query, for $\mathrm{bits}(g(|x|))$, for prepending $1$, for $y \mapsto \mathrm{bits}(|y|)$ and for the inner second component, and $k_R$ the number of tapes of its interpreter core, the bank of the prepend call is the injection of the $k_p$ tapes of the call computing $x \mapsto 1x$ into the host tapes that sends tape $0$ to $\tau_2$ and each tape $i \ne 0$ to host tape $14 + k_q + k_g + i$.
\end{definition}

\begin{definition}[Bank of the length call]
\lean{Turing.NTimeDiag.DzKit.bL}
\label{Turing.NTimeDiag.DzKit.bL}
\leanok
\uses{Turing.EffectiveNDMachineCode, Turing.FinTM, Turing.NTimeDiag.DzKit, Turing.NTimeDiag.DzKit.K, Turing.NTimeDiag.DzKit.tp, Turing.NTimeDiag.DzKit.tp_val, Turing.NTimeDiag.qyBank}
Let $\kappa$ be a diagonalizer's kit for $f, g : \mathbb{N} \to \mathbb{N}$, an effective coding scheme $c$ and a finite deterministic decider $B_F$ over bits, with $K_\kappa$ host tapes and named tapes $\tau_0, \dots, \tau_{13}$ (the host tapes of indices $0, \dots, 13$). With $k_q, k_g, k_p, k_l, k_a$ the numbers of tapes of the kit's calls for the stage query, for $\mathrm{bits}(g(|x|))$, for prepending $1$, for $y \mapsto \mathrm{bits}(|y|)$ and for the inner second component, and $k_R$ the number of tapes of its interpreter core, the bank of the length call is the injection of the $k_l$ tapes of the call computing $y \mapsto \mathrm{bits}(|y|)$ into the host tapes that sends tape $0$ to $\tau_{10}$ and each tape $i \ne 0$ to host tape $14 + k_q + k_g + k_p + i$.
\end{definition}

\begin{definition}[Bank of the extractor call]
\lean{Turing.NTimeDiag.DzKit.bA}
\label{Turing.NTimeDiag.DzKit.bA}
\leanok
\uses{Turing.EffectiveNDMachineCode, Turing.FinTM, Turing.NTimeDiag.DzKit, Turing.NTimeDiag.DzKit.K, Turing.NTimeDiag.DzKit.tp, Turing.NTimeDiag.DzKit.tp_val, Turing.NTimeDiag.qyBank}
Let $\kappa$ be a diagonalizer's kit for $f, g : \mathbb{N} \to \mathbb{N}$, an effective coding scheme $c$ and a finite deterministic decider $B_F$ over bits, with $K_\kappa$ host tapes and named tapes $\tau_0, \dots, \tau_{13}$ (the host tapes of indices $0, \dots, 13$). With $k_q, k_g, k_p, k_l, k_a$ the numbers of tapes of the kit's calls for the stage query, for $\mathrm{bits}(g(|x|))$, for prepending $1$, for $y \mapsto \mathrm{bits}(|y|)$ and for the inner second component, and $k_R$ the number of tapes of its interpreter core, the bank of the extractor call is the injection of the $k_a$ tapes of the call computing the inner second component into the host tapes that sends tape $0$ to $\tau_8$ and each tape $i \ne 0$ to host tape $14 + k_q + k_g + k_p + k_l + i$.
\end{definition}

\begin{definition}[Tape map of the interpreter core]
\lean{Turing.NTimeDiag.DzKit.bR}
\label{Turing.NTimeDiag.DzKit.bR}
\leanok
\uses{Turing.EffectiveNDMachineCode, Turing.FinTM, Turing.NTimeDiag.DzKit, Turing.NTimeDiag.DzKit.K, Turing.NTimeDiag.DzKit.tp, Turing.NTimeDiag.DzKit.tp_val, Turing.NTimeDiag.qyBank}
Let $\kappa$ be a diagonalizer's kit for $f, g : \mathbb{N} \to \mathbb{N}$, an effective coding scheme $c$ and a finite deterministic decider $B_F$ over bits, with $K_\kappa$ host tapes and named tapes $\tau_0, \dots, \tau_{13}$ (the host tapes of indices $0, \dots, 13$). With $k_q, k_g, k_p, k_l, k_a$ the numbers of tapes of the kit's calls for the stage query, for $\mathrm{bits}(g(|x|))$, for prepending $1$, for $y \mapsto \mathrm{bits}(|y|)$ and for the inner second component, and $k_R$ the number of tapes of its interpreter core and $j_R$ its stream tape, let $\mathit{off} = 14 + k_q + k_g + k_p + k_l + k_a$. The core's tape map is the injection of the $k_R$ core tapes into the host tapes obtained by first exchanging tapes $0$ and $j_R$ and then applying the bank that sends $0$ to $\tau_6$ and each $i \ne 0$ to host tape $\mathit{off} + i$; so the stream tape $j_R$ lands on the block tape $\tau_6$.
\end{definition}

\begin{definition}[Tape map of the decider]
\lean{Turing.NTimeDiag.DzKit.bB}
\label{Turing.NTimeDiag.DzKit.bB}
\leanok
\uses{Turing.EffectiveNDMachineCode, Turing.FinTM, Turing.NTimeDiag.DzKit, Turing.NTimeDiag.DzKit.K, Turing.NTimeDiag.DzKit.tp, Turing.NTimeDiag.DzKit.tp_val, Turing.NTimeDiag.qyBank}
Let $\kappa$ be a diagonalizer's kit for $f, g : \mathbb{N} \to \mathbb{N}$, an effective coding scheme $c$ and a finite deterministic decider $B_F$ over bits, with $K_\kappa$ host tapes and named tapes $\tau_0, \dots, \tau_{13}$ (the host tapes of indices $0, \dots, 13$). With $k_q, k_g, k_p, k_l, k_a$ the numbers of tapes of the kit's calls for the stage query, for $\mathrm{bits}(g(|x|))$, for prepending $1$, for $y \mapsto \mathrm{bits}(|y|)$ and for the inner second component, and $k_R$ the number of tapes of its interpreter core and $k_{B_F}$ the number of tapes of $B_F$, let $\mathit{off} = 14 + k_q + k_g + k_p + k_l + k_a + k_R$. The decider's tape map is the injection of the $k_{B_F}$ tapes of $B_F$ into the host tapes that sends tape $0$ to $\tau_{13}$ and each tape $i \ne 0$ to host tape $\mathit{off} + i$.
\end{definition}

\begin{theorem}[Core's stream tape is the block tape]
\lean{Turing.NTimeDiag.DzKit.bR_j}
\label{Turing.NTimeDiag.DzKit.bR_j}
\leanok
\uses{Turing.EffectiveNDMachineCode, Turing.FinTM, Turing.NTimeDiag.DzKit, Turing.NTimeDiag.DzKit.bR, Turing.NTimeDiag.DzKit.tp}
\difficulty{2}
Let $\kappa$ be a diagonalizer's kit for $f, g : \mathbb{N} \to \mathbb{N}$, an effective coding scheme $c$ and a finite deterministic decider $B_F$ over bits, with $K_\kappa$ host tapes and named tapes $\tau_0, \dots, \tau_{13}$ (the host tapes of indices $0, \dots, 13$), and let $j_R$ be the stream tape of the kit's interpreter core. The core's tape map (exchange tapes $0$ and $j_R$, then send $0$ to $\tau_6$ and the other core tapes into their own block of host tapes) sends $j_R$ to the block tape $\tau_6$.
\end{theorem}

\begin{theorem}[Core's other tapes avoid the named tapes]
\lean{Turing.NTimeDiag.DzKit.bR_off}
\label{Turing.NTimeDiag.DzKit.bR_off}
\leanok
\uses{Turing.EffectiveNDMachineCode, Turing.FinTM, Turing.NTimeDiag.DzKit, Turing.NTimeDiag.DzKit.bR, Turing.NTimeDiag.dz_bank_off}
\difficulty{3}
Let $\kappa$ be a diagonalizer's kit for $f, g : \mathbb{N} \to \mathbb{N}$, an effective coding scheme $c$ and a finite deterministic decider $B_F$ over bits, with $K_\kappa$ host tapes and named tapes $\tau_0, \dots, \tau_{13}$ (the host tapes of indices $0, \dots, 13$), and let $j_R$ be the stream tape of the kit's interpreter core. Under the core's tape map (exchange tapes $0$ and $j_R$, then send $0$ to $\tau_6$ and each $i \ne 0$ to host tape $14 + k_q + k_g + k_p + k_l + k_a + i$, the $k$'s being the tape counts of the kit's calls), every core tape $i \ne j_R$ is sent to a host tape of index at least $14$.
\end{theorem}

\begin{definition}[Phases of the setup]
\lean{Turing.NTimeDiag.DzPK}
\label{Turing.NTimeDiag.DzPK}
\leanok
The phases of the diagonalizer's setup are twelve labels: the three input calls (the stage query, the bits of $g(n)$ and the word $1^{n+1}$); two copies of $\mathrm{bits}(g(n))$; the expansion writing the block's stream cells; the write of the separator's counter and the separator's expansion; the copy of the budget word into the block and the block's rewind; and the ruler's expansion and the ruler's rewind.
\end{definition}

\begin{definition}[State sets of the setup's phases]
\lean{Turing.NTimeDiag.DzPPh}
\label{Turing.NTimeDiag.DzPPh}
\leanok
\uses{Turing.CounterLoopState, Turing.EffectiveNDMachineCode, Turing.FinTM, Turing.NTimeDiag.DzKit, Turing.NTimeDiag.DzPK, Turing.RewindPhase, Turing.SweepPhase, Turing.WritePhase}
Let $\kappa$ be a diagonalizer's kit for $f, g : \mathbb{N} \to \mathbb{N}$, an effective coding scheme $c$ and a finite deterministic decider $B_F$ over bits. The state sets of the setup's phases are: for each of the three input-call phases, a fresh start state together with the states of the corresponding call of $\kappa$ (for the stage query, for $\mathrm{bits}(g(|x|))$, for prepending $1$); for the two copy phases and the budget-word copy, the states of a sweep; for the three expansion phases, the counter-loop states of an expansion of the one-symbol word $1$; for the separator's counter, the states of writing a one-symbol constant; and for the two rewind phases, the rewind states.
\end{definition}

\begin{definition}[Phase machines of the setup]
\lean{Turing.NTimeDiag.dzSetM}
\label{Turing.NTimeDiag.dzSetM}
\leanok
\uses{Turing.EffectiveNDMachineCode, Turing.FinTM, Turing.MultiTapeTM, Turing.NTimeDiag.DzKit, Turing.NTimeDiag.DzKit.K, Turing.NTimeDiag.DzKit.bG, Turing.NTimeDiag.DzKit.bP, Turing.NTimeDiag.DzKit.bQ, Turing.NTimeDiag.DzKit.tp, Turing.NTimeDiag.DzKit.tp_ne, Turing.NTimeDiag.DzPK, Turing.NTimeDiag.DzPPh, Turing.NTimeDiag.qyEmb2, Turing.copyMapTM, Turing.embedEmitTM, Turing.expandTM, Turing.rewindTM, Turing.seamReleaseTM, Turing.writeConstTM}
Let $\kappa$ be a diagonalizer's kit for $f, g : \mathbb{N} \to \mathbb{N}$, an effective coding scheme $c$ and a finite deterministic decider $B_F$ over bits, with $K_\kappa$ host tapes and named tapes $\tau_0, \dots, \tau_{13}$ (the host tapes of indices $0, \dots, 13$). The setup's phase machines, all over the $K_\kappa$ host tapes, are: the release adapters (at their entry states) of the kit's input calls for the stage query, for $\mathrm{bits}(g(|x|))$ and for prepending $1$, embedded along the banks with tape $0$ on $\tau_0$, $\tau_1$, $\tau_2$ respectively; the identity copies from $\tau_1$ to $\tau_3$ and from $\tau_1$ to $\tau_4$; the expansion of the word $1$ by the binary counter on $\tau_4$ onto $\tau_6$; the write of the constant word $1$ on $\tau_5$; the expansion by the counter on $\tau_5$ onto $\tau_6$; the identity copy from $\tau_1$ to $\tau_6$; the rewind of $\tau_6$; the expansion by the counter on $\tau_1$ onto $\tau_7$; and the rewind of $\tau_7$.
\end{definition}

\begin{definition}[Dispatch table of the setup]
\lean{Turing.NTimeDiag.dzSetRed}
\label{Turing.NTimeDiag.dzSetRed}
\leanok
\uses{Turing.EffectiveNDMachineCode, Turing.FinTM, Turing.NTimeDiag.DzKit, Turing.NTimeDiag.DzPK, Turing.NTimeDiag.DzPPh}
Let $\kappa$ be a diagonalizer's kit for $f, g : \mathbb{N} \to \mathbb{N}$, an effective coding scheme $c$ and a finite deterministic decider $B_F$ over bits. The setup's dispatch table takes the phases in the order: stage-query call, $g$'s call, prepend call, first copy, second copy, stream expansion, separator write, separator expansion, budget copy, block rewind, ruler expansion, ruler rewind. It redirects the exit state of each phase except the last to the entry state of the next (a call phase exits at its call's exit state, every other phase at its done state; a call phase is entered at its fresh start, a sweep at its sweeping state, an expansion at its first decrement, the write at its first cell and a rewind at its start), and leaves every other state, including every state of the ruler's rewind, without redirect.
\end{definition}

\begin{definition}[The diagonalizer's setup machine]
\lean{Turing.NTimeDiag.dzSet}
\label{Turing.NTimeDiag.dzSet}
\leanok
\uses{Turing.EffectiveNDMachineCode, Turing.FinTM, Turing.MultiTapeTM, Turing.NTimeDiag.DzKit, Turing.NTimeDiag.DzKit.K, Turing.NTimeDiag.DzPPh, Turing.NTimeDiag.dzSetM, Turing.NTimeDiag.dzSetRed, Turing.phaseHostTM}
Let $\kappa$ be a diagonalizer's kit for $f, g : \mathbb{N} \to \mathbb{N}$, an effective coding scheme $c$ and a finite deterministic decider $B_F$ over bits. The setup machine is the phase host over the setup's twelve phases and their phase machines (three released input calls, copies, expansions, a constant write and rewinds), started at the fresh start state of the stage-query call phase and switching phases according to the setup's dispatch table.
\end{definition}

\begin{definition}[Exit state of the setup]
\lean{Turing.NTimeDiag.dzPexit}
\label{Turing.NTimeDiag.dzPexit}
\leanok
\uses{Turing.EffectiveNDMachineCode, Turing.FinTM, Turing.NTimeDiag.DzKit, Turing.NTimeDiag.DzPPh}
Let $\kappa$ be a diagonalizer's kit for $f, g : \mathbb{N} \to \mathbb{N}$, an effective coding scheme $c$ and a finite deterministic decider $B_F$ over bits. The setup's exit state is the done state of its last phase, the ruler's rewind.
\end{definition}

\begin{definition}[The block]
\lean{Turing.NTimeDiag.dzBlk}
\label{Turing.NTimeDiag.dzBlk}
\leanok
For $G \in \mathbb{N}$, the block is the binary word $1^G\,1\,\mathrm{bits}(G)$: $G$ stream cells, a separator and the budget word, where $\mathrm{bits}(G)$ is the binary expansion of $G$.
\end{definition}

\begin{theorem}[Width of the block]
\lean{Turing.NTimeDiag.dzBlk_length}
\label{Turing.NTimeDiag.dzBlk_length}
\leanok
\uses{Turing.NTimeDiag.dzBlk}
\difficulty{2}
For every $G \in \mathbb{N}$, the block $1^G\,1\,\mathrm{bits}(G)$ has length $G + 1 + |\mathrm{bits}(G)|$.
\end{theorem}

\begin{definition}[Named tapes at the setup's exit]
\lean{Turing.NTimeDiag.dzSetL}
\label{Turing.NTimeDiag.dzSetL}
\leanok
\uses{Turing.FinTM.bufferTape, Turing.NTimeDiag.dzBlk, Turing.NTimeDiag.query}
Let $f, g : \mathbb{N} \to \mathbb{N}$, let $q_{f,g}$ be the stage query, let $z$ be a binary word, write $\mathrm{buf}(w)$ for the tape holding $w$ in cells $0, \dots, |w| - 1$ and blank elsewhere, and let $m = |\mathrm{bits}(g(|z|))|$. The named-tape contents at the setup's exit assign $\mathrm{buf}(q_{f,g}(z))$ to index $0$, $\mathrm{buf}(1^m)$ to $1$, $\mathrm{buf}(1z)$ to $2$, $\mathrm{buf}(\mathrm{bits}(g(|z|)))$ to $3$, $\mathrm{buf}(1^m)$ to $4$, $\mathrm{buf}(1)$ to $5$, the block $\mathrm{buf}(1^{g(|z|)}\,1\,\mathrm{bits}(g(|z|)))$ to $6$, $\mathrm{buf}(1^{g(|z|)})$ to $7$, and the blank tape to the indices $8, \dots, 13$.
\end{definition}

\begin{theorem}[A bank holding an argument]
\lean{Turing.NTimeDiag.dz_arg_bank}
\label{Turing.NTimeDiag.dz_arg_bank}
\leanok
\uses{Turing.EffectiveNDMachineCode, Turing.FinTM, Turing.FinTM.bufferTape, Turing.NTimeDiag.DzKit, Turing.NTimeDiag.DzKit.K, Turing.NTimeDiag.DzKit.tp, Turing.NTimeDiag.dzLay, Turing.NTimeDiag.dzLay_off, Turing.NTimeDiag.dzLay_tp, Turing.NTimeDiag.dz_bank_off, Turing.NTimeDiag.qyBank, Turing.NTimeDiag.qyBank_zero, Turing.NTimeDiag.qy_arg_bank, Turing.stateWord}
\difficulty{2}
Let $\kappa$ be a diagonalizer's kit for $f, g : \mathbb{N} \to \mathbb{N}$, an effective coding scheme $c$ and a finite deterministic decider $B_F$ over bits, with $K_\kappa$ host tapes and named tapes $\tau_0, \dots, \tau_{13}$ (the host tapes of indices $0, \dots, 13$). Let $L$ assign a tape to each index $i < 14$, let $i_0 < 14$, let $\mathit{off}, k \in \mathbb{N}$ with $\mathit{off} + k \le K_\kappa$, $\tau_{i_0} < \mathit{off}$, $14 \le \mathit{off}$ and $0 < k$, and let $a$ be a binary word with $L(i_0) = \mathrm{buf}(a)$, where $\mathrm{buf}(w)$ for the tape holding the binary word $w$ in cells $0, \dots, |w| - 1$ and blank elsewhere. Let $\hat L$ give $\tau_i$ the contents $L(i)$ and every other host tape the blank tape, and let $\sigma$ be the bank sending $0$ to $\tau_{i_0}$ and each $i \ne 0$ to host tape $\mathit{off} + i$. Then for every $i < k$, $\hat L(\sigma(i)) = \mathrm{buf}(w_i)$, where $w_0 = a$ and $w_i = \varepsilon$ for $i \ne 0$.
\end{theorem}

\begin{theorem}[A bank with all heads at zero]
\lean{Turing.NTimeDiag.dz_zero_bank}
\label{Turing.NTimeDiag.dz_zero_bank}
\leanok
\uses{Turing.EffectiveNDMachineCode, Turing.FinTM, Turing.NTimeDiag.DzKit, Turing.NTimeDiag.DzKit.K, Turing.NTimeDiag.DzKit.tp, Turing.NTimeDiag.dzLay, Turing.NTimeDiag.dzLay_off, Turing.NTimeDiag.dzLay_tp, Turing.NTimeDiag.dz_bank_off, Turing.NTimeDiag.qyBank, Turing.NTimeDiag.qyBank_zero}
\difficulty{3}
Let $\kappa$ be a diagonalizer's kit for $f, g : \mathbb{N} \to \mathbb{N}$, an effective coding scheme $c$ and a finite deterministic decider $B_F$ over bits, with $K_\kappa$ host tapes and named tapes $\tau_0, \dots, \tau_{13}$ (the host tapes of indices $0, \dots, 13$). Let $P$ assign a head position to each index $i < 14$, let $i_0 < 14$ with $P(i_0) = 0$, and let $\mathit{off}, k \in \mathbb{N}$ with $\mathit{off} + k \le K_\kappa$, $\tau_{i_0} < \mathit{off}$ and $14 \le \mathit{off}$. Let $\hat P$ give $\tau_i$ the position $P(i)$ and every other head the position $0$, and let $\sigma$ be the bank sending $0$ to $\tau_{i_0}$ and each $i \ne 0$ to host tape $\mathit{off} + i$. Then $\hat P(\sigma(i)) = 0$ for every $i < k$.
\end{theorem}

\begin{definition}[Time bound of the setup]
\lean{Turing.NTimeDiag.dzSetB}
\label{Turing.NTimeDiag.dzSetB}
\leanok
\uses{Turing.EffectiveNDMachineCode, Turing.FinTM, Turing.NTimeDiag.DzKit}
Let $\kappa$ be a diagonalizer's kit for $f, g : \mathbb{N} \to \mathbb{N}$, an effective coding scheme $c$ and a finite deterministic decider $B_F$ over bits, and let $A_q, A_g, A_p$ be the constants of the kit's input calls for the stage query, for $\mathrm{bits}(g(|x|))$ and for prepending $1$. For $G \in \mathbb{N}$ (standing for $g(n)$), the setup's time bound is $(A_q + A_g + A_p + 40)(G + 1)$.
\end{definition}

\begin{theorem}[Run of the setup]
\lean{Turing.NTimeDiag.dz_setup}
\label{Turing.NTimeDiag.dz_setup}
\leanok
\uses{Turing.EffectiveNDMachineCode, Turing.FinTM, Turing.FinTM.bufferTape, Turing.NTimeDiag.DzKit, Turing.NTimeDiag.DzKit.K, Turing.NTimeDiag.DzKit.bG, Turing.NTimeDiag.DzKit.bP, Turing.NTimeDiag.DzKit.bQ, Turing.NTimeDiag.DzKit.tp, Turing.NTimeDiag.DzKit.tp_ne, Turing.NTimeDiag.DzPK, Turing.NTimeDiag.DzPPh, Turing.NTimeDiag.dzLay, Turing.NTimeDiag.dzLay_const, Turing.NTimeDiag.dzLay_tp, Turing.NTimeDiag.dzLay_update, Turing.NTimeDiag.dzPexit, Turing.NTimeDiag.dzSet, Turing.NTimeDiag.dzSetB, Turing.NTimeDiag.dzSetL, Turing.NTimeDiag.dzSetM, Turing.NTimeDiag.dzSetRed, Turing.NTimeDiag.dz_arg_bank, Turing.NTimeDiag.dz_zero_bank, Turing.NTimeDiag.query, Turing.NTimeDiag.qyC, Turing.NTimeDiag.qyEmb2, Turing.NTimeDiag.qy_bits_val, Turing.NTimeDiag.qy_callHost, Turing.NTimeDiag.qy_copy, Turing.NTimeDiag.qy_expandHost, Turing.NTimeDiag.qy_guard_of_cut, Turing.NTimeDiag.qy_hosted, Turing.NTimeDiag.qy_hosted_last, Turing.NTimeDiag.qy_rewind, Turing.NTimeDiag.qy_writeConst, Turing.Segment, Turing.Segment.trans, Turing.counterOverhead_le, Turing.length_bits_le_self}
\difficulty{5}
Let $\kappa$ be a diagonalizer's kit for $f, g : \mathbb{N} \to \mathbb{N}$, an effective coding scheme $c$ and a finite deterministic decider $B_F$ over bits, with $K_\kappa$ host tapes and named tapes $\tau_0, \dots, \tau_{13}$ (the host tapes of indices $0, \dots, 13$), let $A_q, A_g, A_p$ be the constants of its input calls for the stage query, for $\mathrm{bits}(g(|x|))$ and for prepending $1$, and let $z$ be a binary input. Write $\mathrm{buf}(w)$ for the tape holding the binary word $w$ in cells $0, \dots, |w| - 1$ and blank elsewhere, $\langle s, p, T, P, o\rangle$ for the live configuration on input $z$ with state $s$, input head $p$, work tapes $T$, heads $P$ and output $o$, and $\Lambda_z$ for the named-tape contents at the setup's exit: $\mathrm{buf}(q_{f,g}(z))$ on $\tau_0$, where $q_{f,g}$ is the stage query; $\mathrm{buf}(1^m)$ on $\tau_1$ and $\tau_4$, where $m = |\mathrm{bits}(g(|z|))|$; $\mathrm{buf}(1z)$ on $\tau_2$; $\mathrm{buf}(\mathrm{bits}(g(|z|)))$ on $\tau_3$; $\mathrm{buf}(1)$ on $\tau_5$; the block $\mathrm{buf}(1^{g(|z|)}\,1\,\mathrm{bits}(g(|z|)))$ on $\tau_6$; $\mathrm{buf}(1^{g(|z|)})$ on $\tau_7$; and the blank tape on $\tau_8, \dots, \tau_{13}$, extended blank to the other host tapes. Then there is $t \le (A_q + A_g + A_p + 40)(g(|z|) + 1)$ such that the setup machine runs in exactly $t$ steps from $\langle (\text{query call}, \text{start}), 1, \text{all blank}, 0, \varepsilon\rangle$ to $\langle \mathit{exit}, 1, \Lambda_z, 0, \varepsilon\rangle$, where $\mathit{exit}$ is the done state of the ruler's rewind, and is not in $\mathit{exit}$ at any earlier time.
\end{theorem}

\begin{definition}[A single tape as a bank]
\lean{Turing.NTimeDiag.dzOne}
\label{Turing.NTimeDiag.dzOne}
\leanok
For $K \in \mathbb{N}$ and a tape $t$ among $K$ tapes, the one-tape bank is the injection of the single index $0$ into the $K$ tapes sending $0$ to $t$.
\end{definition}

\begin{theorem}[Core's tapes avoid other named tapes]
\lean{Turing.NTimeDiag.DzKit.bR_ne}
\label{Turing.NTimeDiag.DzKit.bR_ne}
\leanok
\uses{Turing.EffectiveNDMachineCode, Turing.FinTM, Turing.NTimeDiag.DzKit, Turing.NTimeDiag.DzKit.bR, Turing.NTimeDiag.DzKit.bR_j, Turing.NTimeDiag.DzKit.bR_off, Turing.NTimeDiag.DzKit.le_K, Turing.NTimeDiag.DzKit.tp, Turing.NTimeDiag.DzKit.tp_val}
\difficulty{4}
Let $\kappa$ be a diagonalizer's kit for $f, g : \mathbb{N} \to \mathbb{N}$, an effective coding scheme $c$ and a finite deterministic decider $B_F$ over bits, with $K_\kappa$ host tapes and named tapes $\tau_0, \dots, \tau_{13}$ (the host tapes of indices $0, \dots, 13$). For every tape $i$ of the kit's interpreter core and every $j < 14$ with $j \ne 6$, the core's tape map (which sends the stream tape to $\tau_6$ and every other core tape to a host tape of index at least $14$) does not send $i$ to $\tau_j$.
\end{theorem}

\begin{theorem}[Decider's tapes avoid other named tapes]
\lean{Turing.NTimeDiag.DzKit.bB_ne}
\label{Turing.NTimeDiag.DzKit.bB_ne}
\leanok
\uses{Turing.EffectiveNDMachineCode, Turing.FinTM, Turing.NTimeDiag.DzKit, Turing.NTimeDiag.DzKit.bB, Turing.NTimeDiag.DzKit.le_K, Turing.NTimeDiag.DzKit.tp, Turing.NTimeDiag.DzKit.tp_val, Turing.NTimeDiag.dz_bank_off}
\difficulty{4}
Let $\kappa$ be a diagonalizer's kit for $f, g : \mathbb{N} \to \mathbb{N}$, an effective coding scheme $c$ and a finite deterministic decider $B_F$ over bits, with $K_\kappa$ host tapes and named tapes $\tau_0, \dots, \tau_{13}$ (the host tapes of indices $0, \dots, 13$). For every tape $i$ of $B_F$ and every $j < 14$ with $j \ne 13$, the decider's tape map (which sends tape $0$ to $\tau_{13}$ and every other tape of $B_F$ to a host tape of index at least $14$) does not send $i$ to $\tau_j$.
\end{theorem}

\begin{definition}[Tape map of the hosted core's clock]
\lean{Turing.NTimeDiag.DzKit.sR}
\label{Turing.NTimeDiag.DzKit.sR}
\leanok
\uses{Turing.EffectiveNDMachineCode, Turing.FinTM, Turing.NTimeDiag.DzKit, Turing.NTimeDiag.DzKit.K, Turing.NTimeDiag.DzKit.bR, Turing.NTimeDiag.DzKit.bR_ne, Turing.NTimeDiag.DzKit.tp, Turing.NTimeDiag.DzKit.tp_ne, Turing.NTimeDiag.dzClkEmb}
Let $\kappa$ be a diagonalizer's kit for $f, g : \mathbb{N} \to \mathbb{N}$, an effective coding scheme $c$ and a finite deterministic decider $B_F$ over bits, with $K_\kappa$ host tapes and named tapes $\tau_0, \dots, \tau_{13}$ (the host tapes of indices $0, \dots, 13$), and let $k_R$ be the number of tapes of its interpreter core. The hosted core's tape map is the hosted clock's tape map for the buffer tape $\tau_{11}$, the capture tape $\tau_{12}$, the counter tape $\tau_3$ and the core's tape map: the injection of the $(1 + k_R) + 1 + 1$ tapes of the clock over the hosted core into the host tapes sending the virtual-input buffer to $\tau_{11}$, each core tape $i$ to its image under the core's tape map (stream tape on $\tau_6$, the others at index at least $14$), the capture tape to $\tau_{12}$ and the counter to $\tau_3$.
\end{definition}

\begin{definition}[Tape map of the hosted decider's clock]
\lean{Turing.NTimeDiag.DzKit.sB}
\label{Turing.NTimeDiag.DzKit.sB}
\leanok
\uses{Turing.EffectiveNDMachineCode, Turing.FinTM, Turing.NTimeDiag.DzKit, Turing.NTimeDiag.DzKit.K, Turing.NTimeDiag.DzKit.bB, Turing.NTimeDiag.DzKit.bB_ne, Turing.NTimeDiag.DzKit.tp, Turing.NTimeDiag.DzKit.tp_ne, Turing.NTimeDiag.dzClkEmb}
Let $\kappa$ be a diagonalizer's kit for $f, g : \mathbb{N} \to \mathbb{N}$, an effective coding scheme $c$ and a finite deterministic decider $B_F$ over bits, with $K_\kappa$ host tapes and named tapes $\tau_0, \dots, \tau_{13}$ (the host tapes of indices $0, \dots, 13$), and let $k_{B_F}$ be the number of tapes of $B_F$. The hosted decider's tape map is the hosted clock's tape map for the buffer tape $\tau_8$, the capture tape $\tau_{12}$, the counter tape $\tau_3$ and the decider's tape map: the injection of the $(1 + k_{B_F}) + 1 + 1$ tapes of the clock over the hosted decider into the host tapes sending the virtual-input buffer to $\tau_8$, each tape $i$ of $B_F$ to its image under the decider's tape map (tape $0$ on $\tau_{13}$, the others at index at least $14$), the capture tape to $\tau_{12}$ and the counter to $\tau_3$.
\end{definition}

\begin{definition}[Phases of the run]
\lean{Turing.NTimeDiag.DzQK}
\label{Turing.NTimeDiag.DzQK}
\leanok
The phases of the diagonalizer's run are twenty-three labels: the flag; in the middle, the tail's copy, the split (a skip back, a lockstep walk, the take of the separator, the move of the budget word, a second skip back), the budget's test (an expansion, a rewind, a comparison), the normalization, the extraction of the code $\alpha$, three emissions, the rewind of the capture and the hosted interpreter core; at the top, the tail's copy, the hosted decider, the rewind of the capture and its read; and the two verdicts, accept and reject.
\end{definition}

\begin{definition}[State sets of the run's phases]
\lean{Turing.NTimeDiag.DzQPh}
\label{Turing.NTimeDiag.DzQPh}
\leanok
\uses{Turing.CounterLoopState, Turing.EffectiveNDMachineCode, Turing.FinTM, Turing.NTimeDiag.DzKit, Turing.NTimeDiag.DzQK, Turing.NTimeDiag.QyCmp, Turing.NTimeDiag.QyEm, Turing.RewindPhase, Turing.SeekPhase, Turing.SweepPhase}
Let $\kappa$ be a diagonalizer's kit for $f, g : \mathbb{N} \to \mathbb{N}$, an effective coding scheme $c$ and a finite deterministic decider $B_F$ over bits. The state sets of the run's phases are: for the flag, the separator's take and the read, the states of a one-cell take (optional values of an optional bit); for the two tail copies, the budget word's move and the two verdicts, sweep states; for the two skip backs and the walk, seek states; for the expansion, counter-loop states of an expansion of the word $1$; for the three rewinds, rewind states; for the comparison, the comparison's states; for the normalization and the extraction, a fresh start state together with the states of the kit's call for $y \mapsto \mathrm{bits}(|y|)$, respectively for the inner second component; for the three emissions, the emitter states with empty prefix, suffix lengths $4$, $2$, $0$ and repetition parameters $3$, $1$, $0$; for the hosted core, the states of the kit's clock over the core; and for the hosted decider, the states of the kit's clock over $B_F$.
\end{definition}

\begin{definition}[Phase machines of the run]
\lean{Turing.NTimeDiag.dzQM}
\label{Turing.NTimeDiag.dzQM}
\leanok
\uses{Turing.EffectiveNDMachineCode, Turing.FinTM, Turing.MultiTapeTM, Turing.NTimeDiag.DzKit, Turing.NTimeDiag.DzKit.K, Turing.NTimeDiag.DzKit.bA, Turing.NTimeDiag.DzKit.bL, Turing.NTimeDiag.DzKit.sB, Turing.NTimeDiag.DzKit.sR, Turing.NTimeDiag.DzKit.tp, Turing.NTimeDiag.DzKit.tp_ne, Turing.NTimeDiag.DzQK, Turing.NTimeDiag.DzQPh, Turing.NTimeDiag.dzOne, Turing.NTimeDiag.dzTakeTM, Turing.NTimeDiag.qyCmpTM, Turing.NTimeDiag.qyEmb2, Turing.NTimeDiag.qyEmitTM, Turing.clearTM, Turing.copyMapTM, Turing.embedEmitTM, Turing.embedSilentTM, Turing.expandTM, Turing.rewindTM, Turing.seamReleaseTM, Turing.skipBackTM, Turing.transferMapTM, Turing.walkTM}
Let $\kappa$ be a diagonalizer's kit for $f, g : \mathbb{N} \to \mathbb{N}$, an effective coding scheme $c$ and a finite deterministic decider $B_F$ over bits, with $K_\kappa$ host tapes and named tapes $\tau_0, \dots, \tau_{13}$ (the host tapes of indices $0, \dots, 13$). The run's phase machines, all over the $K_\kappa$ host tapes, are: for the flag, the one-cell take on $\tau_0$; for the middle's tail copy, the identity copy from $\tau_0$ to $\tau_8$; the skip back on $\tau_6$; the walk led by $\tau_7$ moving the heads of $\tau_7$ and $\tau_6$ rightwards in lockstep; the one-cell take on $\tau_6$; the identity move from $\tau_6$ to $\tau_9$, blanking the source; the skip back on $\tau_6$; the expansion of the word $1$ onto $\tau_{10}$ by the counter on $\tau_9$; the rewind of $\tau_{10}$; the length comparison of $\tau_{10}$ with $\tau_6$; the released calls for $y \mapsto \mathrm{bits}(|y|)$ and for the inner second component, embedded along the banks with tape $0$ on $\tau_{10}$ and on $\tau_8$; three emitters under the capture wrapper with capture tape $\tau_{11}$, reading the single tape $\tau_{10}$ (each symbol $4$ times, then $0011$), $\tau_8$ (each symbol twice, then $01$) and $\tau_2$ (each symbol once); the rewind of $\tau_{11}$; the kit's clock over the core embedded along the hosted core's tape map; for the top's tail copy, the identity copy from $\tau_0$ to $\tau_8$; the kit's clock over $B_F$ embedded along the hosted decider's tape map; the rewind of $\tau_{12}$; the one-cell take on $\tau_{12}$; and, for both verdicts, the clearing of $\tau_{13}$.
\end{definition}

\begin{definition}[Accepting verdict of the run]
\lean{Turing.NTimeDiag.dzAcc}
\label{Turing.NTimeDiag.dzAcc}
\leanok
\uses{Turing.EffectiveNDMachineCode, Turing.FinTM, Turing.NTimeDiag.DzKit, Turing.NTimeDiag.DzQPh}
Let $\kappa$ be a diagonalizer's kit for $f, g : \mathbb{N} \to \mathbb{N}$, an effective coding scheme $c$ and a finite deterministic decider $B_F$ over bits. The run's accepting verdict is the sweeping state of the accept phase.
\end{definition}

\begin{definition}[Rejecting verdict of the run]
\lean{Turing.NTimeDiag.dzRej}
\label{Turing.NTimeDiag.dzRej}
\leanok
\uses{Turing.EffectiveNDMachineCode, Turing.FinTM, Turing.NTimeDiag.DzKit, Turing.NTimeDiag.DzQPh}
Let $\kappa$ be a diagonalizer's kit for $f, g : \mathbb{N} \to \mathbb{N}$, an effective coding scheme $c$ and a finite deterministic decider $B_F$ over bits. The run's rejecting verdict is the sweeping state of the reject phase.
\end{definition}

\begin{definition}[Dispatch table of the run]
\lean{Turing.NTimeDiag.dzQRed}
\label{Turing.NTimeDiag.dzQRed}
\leanok
\uses{Turing.EffectiveNDMachineCode, Turing.FinTM, Turing.NTimeDiag.DzKit, Turing.NTimeDiag.DzQK, Turing.NTimeDiag.DzQPh, Turing.NTimeDiag.dzAcc, Turing.NTimeDiag.dzRej}
Let $\kappa$ be a diagonalizer's kit for $f, g : \mathbb{N} \to \mathbb{N}$, an effective coding scheme $c$ and a finite deterministic decider $B_F$ over bits, with interpreter core $R$ (entry state $r_E$, accepting state $r_A$), clock over $R$ with injection, finish and timeout maps $\iota_R, \phi_R, \theta_R$, and clock over $B_F$ with maps $\iota_B, \phi_B, \theta_B$. The run's dispatch table redirects: the flag, after reading a blank, to reject, after reading $1$ to the top's tail copy and after reading $0$ to the middle's tail copy; in the middle, the done state of each of the tail copy, skip back, walk, move, second skip back, expansion and rewind to the entry of the next phase, the separator's take after its step to the move, the comparison's verdict $>$ to reject and any other verdict to the normalization, the exit of the length call to the extraction, the exit of the extraction call to the first emission, each emission's done state to the next emission and the last to the capture's rewind, whose done state goes to the hosted core at $\iota_R(r_E, 1)$; the hosted core's state $\phi_R(r_A, b)$ to accept, any other finish state $\phi_R(r)$ or timeout state $\theta_R(r)$ to reject; at the top, the tail copy's done state to the hosted decider at $\iota_B(q_0, 1)$, $q_0$ the start state of $B_F$, any finish state $\phi_B(r)$ to the capture's rewind and any timeout state $\theta_B(r)$ to reject, the rewind's done state to the read, and the read, after reading $0$, to accept and, after reading $1$ or a blank, to reject. All other states, including the verdicts, have no redirect.
\end{definition}

\begin{definition}[The diagonalizer's run machine]
\lean{Turing.NTimeDiag.dzQ}
\label{Turing.NTimeDiag.dzQ}
\leanok
\uses{Turing.EffectiveNDMachineCode, Turing.FinTM, Turing.MultiTapeTM, Turing.NTimeDiag.DzKit, Turing.NTimeDiag.DzKit.K, Turing.NTimeDiag.DzQPh, Turing.NTimeDiag.dzQM, Turing.NTimeDiag.dzQRed, Turing.phaseHostTM}
Let $\kappa$ be a diagonalizer's kit for $f, g : \mathbb{N} \to \mathbb{N}$, an effective coding scheme $c$ and a finite deterministic decider $B_F$ over bits. The run machine is the phase host over the run's twenty-three phases and their phase machines, started at the flag's start state and switching phases according to the run's dispatch table.
\end{definition}

\begin{definition}[Verdicts of the run]
\lean{Turing.NTimeDiag.dzQF}
\label{Turing.NTimeDiag.dzQF}
\leanok
\uses{Turing.EffectiveNDMachineCode, Turing.FinTM, Turing.NTimeDiag.DzKit, Turing.NTimeDiag.DzQPh, Turing.NTimeDiag.dzAcc, Turing.NTimeDiag.dzRej}
Let $\kappa$ be a diagonalizer's kit for $f, g : \mathbb{N} \to \mathbb{N}$, an effective coding scheme $c$ and a finite deterministic decider $B_F$ over bits. The run's verdict set is the set of host states consisting of the accepting verdict and the rejecting verdict.
\end{definition}

\begin{theorem}[Only the verdict phases are verdicts]
\lean{Turing.NTimeDiag.dz_notF}
\label{Turing.NTimeDiag.dz_notF}
\leanok
\uses{Turing.EffectiveNDMachineCode, Turing.FinTM, Turing.NTimeDiag.DzKit, Turing.NTimeDiag.DzQK, Turing.NTimeDiag.DzQPh, Turing.NTimeDiag.dzQF}
\difficulty{3}
Let $\kappa$ be a diagonalizer's kit for $f, g : \mathbb{N} \to \mathbb{N}$, an effective coding scheme $c$ and a finite deterministic decider $B_F$ over bits. For every phase $i$ of the run other than the accept and reject phases, and every state $s$ of phase $i$, the tagged state $(i, s)$ is not in the run's verdict set $\{\mathit{acc}, \mathit{rej}\}$.
\end{theorem}

\begin{theorem}[Accepting and rejecting verdicts differ]
\lean{Turing.NTimeDiag.dz_acc_ne_rej}
\label{Turing.NTimeDiag.dz_acc_ne_rej}
\leanok
\uses{Turing.EffectiveNDMachineCode, Turing.FinTM, Turing.NTimeDiag.DzKit, Turing.NTimeDiag.dzAcc, Turing.NTimeDiag.dzRej}
\difficulty{2}
Let $\kappa$ be a diagonalizer's kit for $f, g : \mathbb{N} \to \mathbb{N}$, an effective coding scheme $c$ and a finite deterministic decider $B_F$ over bits. The run's accepting verdict and its rejecting verdict are different states.
\end{theorem}

\begin{theorem}[Handoff from the setup to the run]
\lean{Turing.NTimeDiag.dz_handoff}
\label{Turing.NTimeDiag.dz_handoff}
\leanok
\uses{Turing.EffectiveNDMachineCode, Turing.FinTM, Turing.FinTM.bufferTape, Turing.FinTM.bufferTape_eq_none_of_not_mem, Turing.NTimeDiag.DzKit, Turing.NTimeDiag.DzKit.tp, Turing.NTimeDiag.DzQPh, Turing.NTimeDiag.dzBlk, Turing.NTimeDiag.dzLay, Turing.NTimeDiag.dzLay_tp, Turing.NTimeDiag.dzLay_update, Turing.NTimeDiag.dzPexit, Turing.NTimeDiag.dzSetL, Turing.NTimeDiag.qyC, Turing.guessRunCfg}
\difficulty{4}
Let $\kappa$ be a diagonalizer's kit for $f, g : \mathbb{N} \to \mathbb{N}$, an effective coding scheme $c$ and a finite deterministic decider $B_F$ over bits, with $K_\kappa$ host tapes and named tapes $\tau_0, \dots, \tau_{13}$ (the host tapes of indices $0, \dots, 13$), let $z$ be a binary input and let $s$ be a binary word of length $|1^{g(|z|)}\,1\,\mathrm{bits}(g(|z|))|$. Write $\mathrm{buf}(w)$ for the tape holding the binary word $w$ in cells $0, \dots, |w| - 1$ and blank elsewhere, $\langle s, p, T, P, o\rangle$ for the live configuration on input $z$ with state $s$, input head $p$, work tapes $T$, heads $P$ and output $o$, and $\Lambda_z$ for the named-tape contents at the setup's exit: $\mathrm{buf}(q_{f,g}(z))$ on $\tau_0$, where $q_{f,g}$ is the stage query; $\mathrm{buf}(1^m)$ on $\tau_1$ and $\tau_4$, where $m = |\mathrm{bits}(g(|z|))|$; $\mathrm{buf}(1z)$ on $\tau_2$; $\mathrm{buf}(\mathrm{bits}(g(|z|)))$ on $\tau_3$; $\mathrm{buf}(1)$ on $\tau_5$; the block $\mathrm{buf}(1^{g(|z|)}\,1\,\mathrm{bits}(g(|z|)))$ on $\tau_6$; $\mathrm{buf}(1^{g(|z|)})$ on $\tau_7$; and the blank tape on $\tau_8, \dots, \tau_{13}$, extended blank to the other host tapes. Writing $s$ on $\tau_6$ from its head onwards in the setup's exit configuration $\langle \mathit{exit}, 1, \Lambda_z, 0, \varepsilon\rangle$ and switching to the flag's start state gives the configuration $\langle (\text{flag}, \text{start}), 1, \Lambda_z', 0, \varepsilon\rangle$, where $\Lambda_z'$ agrees with $\Lambda_z$ except that $\tau_6$ holds $\mathrm{buf}(s)$.
\end{theorem}

\begin{theorem}[The run reads the query's flag]
\lean{Turing.NTimeDiag.dz_flag}
\label{Turing.NTimeDiag.dz_flag}
\leanok
\uses{Turing.EffectiveNDMachineCode, Turing.FinTM, Turing.FinTM.bufferTape, Turing.MultiTapeTM.runFrom_zero, Turing.NTimeDiag.DzKit, Turing.NTimeDiag.DzKit.tp, Turing.NTimeDiag.DzQPh, Turing.NTimeDiag.dzLay, Turing.NTimeDiag.dzLay_tp, Turing.NTimeDiag.dzLay_update, Turing.NTimeDiag.dzQ, Turing.NTimeDiag.dzQF, Turing.NTimeDiag.dzQM, Turing.NTimeDiag.dzQRed, Turing.NTimeDiag.dz_notF, Turing.NTimeDiag.dz_take_run, Turing.NTimeDiag.qyC, Turing.NTimeDiag.qy_hosted, Turing.Segment}
\difficulty{3}
Let $\kappa$ be a diagonalizer's kit for $f, g : \mathbb{N} \to \mathbb{N}$, an effective coding scheme $c$ and a finite deterministic decider $B_F$ over bits, with $K_\kappa$ host tapes and named tapes $\tau_0, \dots, \tau_{13}$ (the host tapes of indices $0, \dots, 13$), let $z$ be a binary input, let $L$ assign a tape to each index $i < 14$ and $P$ a head position to each index. Write $\mathrm{buf}(w)$ for the tape holding the binary word $w$ in cells $0, \dots, |w| - 1$ and blank elsewhere, $L[i \mapsto v]$ for $L$ updated at $i$ to $v$, $\hat L$ for the tape layout giving $\tau_i$ the contents $L(i)$ and every other host tape the blank tape, and $\hat P$ for the head layout giving $\tau_i$ the position $P(i)$ and every other head the position $0$, $\langle s, p, T, P, o\rangle$ for the live configuration on input $z$ with state $s$, input head $p$, work tapes $T$, heads $P$ and output $o$, and $Q_\kappa$ for the diagonalizer's run machine (the phase host over the run's phases), with accepting verdict $\mathit{acc}$ and rejecting verdict $\mathit{rej}$; a verdict-avoiding segment of length $t$ from $a$ to $b$ means that $Q_\kappa$ reaches $b$ from $a$ in exactly $t$ steps and is in neither verdict at any earlier time. Let $q$ be a binary word with $L(0) = \mathrm{buf}(q)$ and $P(0) = 0$, and let $q'$ be the state to which the run's dispatch table sends the flag after reading the cell $\mathrm{buf}(q)(0)$. Then there is a verdict-avoiding segment of length $2$ from $\langle (\text{flag}, \text{start}), 1, \hat L, \hat P, \varepsilon\rangle$ to $\langle q', 1, \hat{L'}, \hat{P'}, \varepsilon\rangle$, where $L' = L[0 \mapsto T_q]$ with $T_q$ the tape $\mathrm{buf}(q)$ with cell $0$ blanked, and $P' = P[0 \mapsto 1]$.
\end{theorem}

\begin{theorem}[Copy of the query's tail]
\lean{Turing.NTimeDiag.dz_tail_run}
\label{Turing.NTimeDiag.dz_tail_run}
\leanok
\uses{Turing.FinTM.bufferTape, Turing.FinTM.bufferTape_append_shift, Turing.FinTM.bufferTape_eq_none_of_not_mem, Turing.FinTM.bufferTape_nil, Turing.FinTM.bufferTape_of_neg, Turing.MultiTapeTM.runFrom, Turing.NTimeDiag.qyC, Turing.copyMapTM, Turing.copyMapTM_run_ofCfg}
\difficulty{4}
Let $K \in \mathbb{N}$, let $\mathit{src} \ne \mathit{dst}$ be two of $K$ tapes, let $x$ be a binary input, and write $\langle s, p, T, P, o\rangle$ for the live configuration on input $x$ with state $s$, input head $p$, work tapes $T$, heads $P$ and output $o$, and $\mathrm{buf}(w)$ for the tape holding the binary word $w$ in cells $0, \dots, |w| - 1$ and blank elsewhere. Let $p$ be an input-head position, $T$ tapes, $P$ heads, $o$ an output, $b$ a bit and $w$ a binary word, such that $T(\mathit{src})$ is $\mathrm{buf}(b\,w)$ with cell $0$ blanked, $P(\mathit{src}) = 1$, $T(\mathit{dst}) = \mathrm{buf}(\varepsilon)$ and $P(\mathit{dst}) = 0$. Then the identity copy machine from $\mathit{src}$ to $\mathit{dst}$ runs from $\langle \text{sweep}, p, T, P, o\rangle$ in $2|w| + 2$ steps to $\langle \text{done}, p, T[\mathit{dst} \mapsto \mathrm{buf}(w)], P, o\rangle$ and is not in its done state at any earlier time.
\end{theorem}

\begin{theorem}[The block after the split]
\lean{Turing.NTimeDiag.dz_split_tape}
\label{Turing.NTimeDiag.dz_split_tape}
\leanok
\uses{Turing.FinTM.bufferTape, Turing.FinTM.bufferTape_eq_none_of_not_mem, Turing.FinTM.bufferTape_take}
\difficulty{4}
Let $s_1, w$ be binary words, let $b$ be a bit, and write $\mathrm{buf}(w)$ for the tape holding the binary word $w$ in cells $0, \dots, |w| - 1$ and blank elsewhere. Take the tape $\mathrm{buf}(s_1\,b\,w)$, blank its cell $|s_1|$ and blank its cells $|s_1| + 1, \dots, |s_1| + |w|$. The resulting tape equals $\mathrm{buf}(s_1)$.
\end{theorem}

\begin{theorem}[The split of the block]
\lean{Turing.NTimeDiag.dz_split}
\label{Turing.NTimeDiag.dz_split}
\leanok
\uses{Turing.EffectiveNDMachineCode, Turing.FinTM, Turing.FinTM.bufferTape, Turing.FinTM.bufferTape_append_mid, Turing.FinTM.bufferTape_append_shift, Turing.FinTM.bufferTape_eq_none_of_not_mem, Turing.FinTM.bufferTape_nil, Turing.FinTM.bufferTape_of_neg, Turing.MultiTapeTM.runFrom, Turing.MultiTapeTM.runFrom_zero, Turing.NTimeDiag.DzKit, Turing.NTimeDiag.DzKit.K, Turing.NTimeDiag.DzKit.tp, Turing.NTimeDiag.DzKit.tp_ne, Turing.NTimeDiag.dzLay, Turing.NTimeDiag.dzLay_tp, Turing.NTimeDiag.dzLay_update, Turing.NTimeDiag.dzQ, Turing.NTimeDiag.dzQF, Turing.NTimeDiag.dzQM, Turing.NTimeDiag.dzQRed, Turing.NTimeDiag.dz_notF, Turing.NTimeDiag.dz_split_tape, Turing.NTimeDiag.dz_take_run, Turing.NTimeDiag.qyC, Turing.NTimeDiag.qy_guard_of_cut, Turing.NTimeDiag.qy_hosted, Turing.Segment, Turing.Segment.trans, Turing.skipBackTM, Turing.skipBackTM_run_ofCfg, Turing.transferMapTM, Turing.transferMapTM_run_ofCfg, Turing.walkTM, Turing.walkTM_run_ofCfg}
\difficulty{6}
Let $\kappa$ be a diagonalizer's kit for $f, g : \mathbb{N} \to \mathbb{N}$, an effective coding scheme $c$ and a finite deterministic decider $B_F$ over bits, with $K_\kappa$ host tapes and named tapes $\tau_0, \dots, \tau_{13}$ (the host tapes of indices $0, \dots, 13$), let $z$ be a binary input, let $L$ assign a tape to each index $i < 14$ and $P$ a head position to each index. Write $\mathrm{buf}(w)$ for the tape holding the binary word $w$ in cells $0, \dots, |w| - 1$ and blank elsewhere, $L[i \mapsto v]$ for $L$ updated at $i$ to $v$, $\hat L$ for the tape layout giving $\tau_i$ the contents $L(i)$ and every other host tape the blank tape, and $\hat P$ for the head layout giving $\tau_i$ the position $P(i)$ and every other head the position $0$, $\langle s, p, T, P, o\rangle$ for the live configuration on input $z$ with state $s$, input head $p$, work tapes $T$, heads $P$ and output $o$, and $Q_\kappa$ for the diagonalizer's run machine (the phase host over the run's phases), with accepting verdict $\mathit{acc}$ and rejecting verdict $\mathit{rej}$; a verdict-avoiding segment of length $t$ from $a$ to $b$ means that $Q_\kappa$ reaches $b$ from $a$ in exactly $t$ steps and is in neither verdict at any earlier time. Let $s_1, w$ be binary words and $b$ a bit with $L(6) = \mathrm{buf}(s_1\,b\,w)$, $L(7) = \mathrm{buf}(1^{|s_1|})$, $L(9) = \mathrm{buf}(\varepsilon)$ and $P(6) = P(7) = P(9) = 0$. Then there is $t \le 2|s_1| + 2|w| + 15$ and a verdict-avoiding segment of length $t$ from the entry $\langle (\text{first skip back}, \text{run}), 1, \hat L, \hat P, \varepsilon\rangle$ to $\langle (\text{expansion}, \text{start}), 1, \hat{L'}, \hat{P'}, \varepsilon\rangle$, where $L' = L[6 \mapsto \mathrm{buf}(s_1)][9 \mapsto \mathrm{buf}(w)]$ and $P' = P[7 \mapsto |s_1| + 1]$.
\end{theorem}

\begin{theorem}[Test of the budget]
\lean{Turing.NTimeDiag.dz_test}
\label{Turing.NTimeDiag.dz_test}
\leanok
\uses{Turing.EffectiveNDMachineCode, Turing.FinTM, Turing.FinTM.bufferTape, Turing.NTimeDiag.DzKit, Turing.NTimeDiag.DzKit.tp, Turing.NTimeDiag.DzKit.tp_ne, Turing.NTimeDiag.dzLay, Turing.NTimeDiag.dzLay_tp, Turing.NTimeDiag.dzLay_update, Turing.NTimeDiag.dzQ, Turing.NTimeDiag.dzQF, Turing.NTimeDiag.dzQM, Turing.NTimeDiag.dzQRed, Turing.NTimeDiag.dzRej, Turing.NTimeDiag.dz_notF, Turing.NTimeDiag.qyC, Turing.NTimeDiag.qyEmb2, Turing.NTimeDiag.qy_cmp_run, Turing.NTimeDiag.qy_expandHost, Turing.NTimeDiag.qy_guard_of_cut, Turing.NTimeDiag.qy_hosted, Turing.NTimeDiag.qy_rewind, Turing.Segment, Turing.Segment.trans, Turing.counterOverhead_le}
\difficulty{5}
Let $\kappa$ be a diagonalizer's kit for $f, g : \mathbb{N} \to \mathbb{N}$, an effective coding scheme $c$ and a finite deterministic decider $B_F$ over bits, with $K_\kappa$ host tapes and named tapes $\tau_0, \dots, \tau_{13}$ (the host tapes of indices $0, \dots, 13$), let $z$ be a binary input, let $L$ assign a tape to each index $i < 14$ and $P$ a head position to each index. Write $\mathrm{buf}(w)$ for the tape holding the binary word $w$ in cells $0, \dots, |w| - 1$ and blank elsewhere, $L[i \mapsto v]$ for $L$ updated at $i$ to $v$, $\hat L$ for the tape layout giving $\tau_i$ the contents $L(i)$ and every other host tape the blank tape, and $\hat P$ for the head layout giving $\tau_i$ the position $P(i)$ and every other head the position $0$, $\langle s, p, T, P, o\rangle$ for the live configuration on input $z$ with state $s$, input head $p$, work tapes $T$, heads $P$ and output $o$, and $Q_\kappa$ for the diagonalizer's run machine (the phase host over the run's phases), with accepting verdict $\mathit{acc}$ and rejecting verdict $\mathit{rej}$; a verdict-avoiding segment of length $t$ from $a$ to $b$ means that $Q_\kappa$ reaches $b$ from $a$ in exactly $t$ steps and is in neither verdict at any earlier time. Let $s_1, w$ be binary words with $L(6) = \mathrm{buf}(s_1)$, $L(9) = \mathrm{buf}(w)$, $L(10) = \mathrm{buf}(\varepsilon)$ and $P(6) = P(9) = P(10) = 0$, and let $b = \sum_i w_i 2^i$ be the value of $w$ read least significant bit first. Then there is $t \le 9b + 2|w| + 12$ and a verdict-avoiding segment of length $t$ from $\langle (\text{expansion}, \text{start}), 1, \hat L, \hat P, \varepsilon\rangle$ to $\langle e, 1, \hat{L'}, \hat P, \varepsilon\rangle$, where $L' = L[10 \mapsto \mathrm{buf}(1^b)][9 \mapsto \mathrm{buf}(1^{|w|})]$, and $e$ is the normalization's start state if $b \le |s_1|$ and $\mathit{rej}$ otherwise.
\end{theorem}

\begin{theorem}[Normalization of the budget]
\lean{Turing.NTimeDiag.dz_norm}
\label{Turing.NTimeDiag.dz_norm}
\leanok
\uses{Turing.EffectiveNDMachineCode, Turing.FinTM, Turing.FinTM.bufferTape, Turing.NTimeDiag.DzKit, Turing.NTimeDiag.DzKit.bL, Turing.NTimeDiag.DzKit.tp, Turing.NTimeDiag.dzLay, Turing.NTimeDiag.dzLay_update, Turing.NTimeDiag.dzQ, Turing.NTimeDiag.dzQF, Turing.NTimeDiag.dzQM, Turing.NTimeDiag.dzQRed, Turing.NTimeDiag.dz_arg_bank, Turing.NTimeDiag.dz_notF, Turing.NTimeDiag.dz_zero_bank, Turing.NTimeDiag.qyC, Turing.NTimeDiag.qy_callHost, Turing.NTimeDiag.qy_guard_of_cut, Turing.NTimeDiag.qy_hosted, Turing.Segment}
\difficulty{4}
Let $\kappa$ be a diagonalizer's kit for $f, g : \mathbb{N} \to \mathbb{N}$, an effective coding scheme $c$ and a finite deterministic decider $B_F$ over bits, with $K_\kappa$ host tapes and named tapes $\tau_0, \dots, \tau_{13}$ (the host tapes of indices $0, \dots, 13$), let $A_l$ be the constant of its install call for $y \mapsto \mathrm{bits}(|y|)$, let $z$ be a binary input, let $L$ assign a tape to each index $i < 14$ and $P$ a head position to each index. Write $\mathrm{buf}(w)$ for the tape holding the binary word $w$ in cells $0, \dots, |w| - 1$ and blank elsewhere, $L[i \mapsto v]$ for $L$ updated at $i$ to $v$, $\hat L$ for the tape layout giving $\tau_i$ the contents $L(i)$ and every other host tape the blank tape, and $\hat P$ for the head layout giving $\tau_i$ the position $P(i)$ and every other head the position $0$, $\langle s, p, T, P, o\rangle$ for the live configuration on input $z$ with state $s$, input head $p$, work tapes $T$, heads $P$ and output $o$, and $Q_\kappa$ for the diagonalizer's run machine (the phase host over the run's phases), with accepting verdict $\mathit{acc}$ and rejecting verdict $\mathit{rej}$; a verdict-avoiding segment of length $t$ from $a$ to $b$ means that $Q_\kappa$ reaches $b$ from $a$ in exactly $t$ steps and is in neither verdict at any earlier time. Let $b \in \mathbb{N}$ with $L(10) = \mathrm{buf}(1^b)$ and $P(10) = 0$. Then there is $t \le A_l(b + 1) + 1$ and a verdict-avoiding segment of length $t$ from $\langle (\text{normalization}, \text{start}), 1, \hat L, \hat P, \varepsilon\rangle$ to $\langle (\text{extraction}, \text{start}), 1, \widehat{L[10 \mapsto \mathrm{buf}(\mathrm{bits}(b))]}, \hat P, \varepsilon\rangle$.
\end{theorem}

\begin{theorem}[The core's input word from three emissions]
\lean{Turing.NTimeDiag.dz_core_word}
\label{Turing.NTimeDiag.dz_core_word}
\leanok
\uses{Turing.NTimeDiag.qyBody, Turing.NTimeDiag.qy_writer_word, Turing.pairEncode}
\difficulty{4}
For a binary word $w$ and $r \in \mathbb{N}$, let $\mathrm{rep}_r(w)$ be the word obtained from $w$ by repeating every symbol $r + 1$ times, and write $\langle u, v\rangle$ for the pairing encoding of binary words. For all binary words $B, \alpha, N$, \[ \mathrm{rep}_3(B)\;0011\;\mathrm{rep}_1(\alpha)\;01\;\mathrm{rep}_0(N) = \langle\langle B, \alpha\rangle, N\rangle. \]
\end{theorem}

\begin{theorem}[An emission under the capture wrapper]
\lean{Turing.NTimeDiag.dz_emit_run}
\label{Turing.NTimeDiag.dz_emit_run}
\leanok
\uses{Turing.FinTM.bufferTape, Turing.MultiTapeTM.runFrom, Turing.NTimeDiag.dzOne, Turing.NTimeDiag.dz_silent, Turing.NTimeDiag.qyBody, Turing.NTimeDiag.qyC, Turing.NTimeDiag.qyEmitTM, Turing.NTimeDiag.qy_emit_run, Turing.embedSilentTM}
\difficulty{3}
Let $K \in \mathbb{N}$, let $\mathit{src} \ne \mathit{cap}$ be two of $K$ tapes, let $x$ be a binary input, and write $\langle s, p, T, P, o\rangle$ for the live configuration on input $x$ with state $s$, input head $p$, work tapes $T$, heads $P$ and output $o$, and $\mathrm{buf}(w)$ for the tape holding the binary word $w$ in cells $0, \dots, |w| - 1$ and blank elsewhere. For a word $v$ and $r \in \mathbb{N}$ let $\mathrm{rep}_r(v)$ repeat every symbol of $v$ $r + 1$ times. Let $\mathit{suf}$, $\mathit{pre}$, $w$, $o$ be binary words, $r \in \mathbb{N}$, and $T, P$ tapes and heads with $T(\mathit{src}) = \mathrm{buf}(w)$, $P(\mathit{src}) = 0$, $T(\mathit{cap}) = \mathrm{buf}(\mathit{pre})$ and $P(\mathit{cap}) = |\mathit{pre}|$. Let $M$ be the emitter that reads the single tape $\mathit{src}$, emits $\mathrm{rep}_r$ of its word followed by $\mathit{suf}$ (with empty prefix), run under the capture wrapper with capture tape $\mathit{cap}$. Then there is $t \le |\mathit{suf}| + 3 + |w|(r + 3)$ such that $M$ runs from $\langle \text{prefix}_0, 1, T, P, o\rangle$ in $t$ steps to $\langle \text{done}, 1, T[\mathit{cap} \mapsto \mathrm{buf}(u)], P[\mathit{src} \mapsto |w|][\mathit{cap} \mapsto |u|], o\rangle$ with $u = \mathit{pre}\,\mathrm{rep}_r(w)\,\mathit{suf}$, and is not in its done state at any earlier time.
\end{theorem}

\begin{theorem}[Rewriting the core's input]
\lean{Turing.NTimeDiag.dz_rewrite}
\label{Turing.NTimeDiag.dz_rewrite}
\leanok
\uses{Turing.EffectiveNDMachineCode, Turing.FinTM, Turing.FinTM.bufferTape, Turing.NTimeDiag.DzKit, Turing.NTimeDiag.DzKit.bA, Turing.NTimeDiag.DzKit.tp, Turing.NTimeDiag.DzKit.tp_ne, Turing.NTimeDiag.dzAlpha, Turing.NTimeDiag.dzLay, Turing.NTimeDiag.dzLay_tp, Turing.NTimeDiag.dzLay_update, Turing.NTimeDiag.dzQ, Turing.NTimeDiag.dzQF, Turing.NTimeDiag.dzQM, Turing.NTimeDiag.dzQRed, Turing.NTimeDiag.dz_arg_bank, Turing.NTimeDiag.dz_core_word, Turing.NTimeDiag.dz_emit_run, Turing.NTimeDiag.dz_notF, Turing.NTimeDiag.dz_zero_bank, Turing.NTimeDiag.qyBody, Turing.NTimeDiag.qyC, Turing.NTimeDiag.qy_callHost, Turing.NTimeDiag.qy_guard_of_cut, Turing.NTimeDiag.qy_hosted, Turing.NTimeDiag.qy_rewind, Turing.Segment, Turing.Segment.trans, Turing.length_pairEncode, Turing.pairEncode}
\difficulty{5}
Let $\kappa$ be a diagonalizer's kit for $f, g : \mathbb{N} \to \mathbb{N}$, an effective coding scheme $c$ and a finite deterministic decider $B_F$ over bits, with $K_\kappa$ host tapes and named tapes $\tau_0, \dots, \tau_{13}$ (the host tapes of indices $0, \dots, 13$), let $A_a$ be the constant of its install call for the inner second component, let $z$ be a binary input, let $L$ assign a tape to each index $i < 14$ and $P$ a head position to each index. Write $\mathrm{buf}(w)$ for the tape holding the binary word $w$ in cells $0, \dots, |w| - 1$ and blank elsewhere, $L[i \mapsto v]$ for $L$ updated at $i$ to $v$, $\hat L$ for the tape layout giving $\tau_i$ the contents $L(i)$ and every other host tape the blank tape, and $\hat P$ for the head layout giving $\tau_i$ the position $P(i)$ and every other head the position $0$, $\langle s, p, T, P, o\rangle$ for the live configuration on input $z$ with state $s$, input head $p$, work tapes $T$, heads $P$ and output $o$, and $Q_\kappa$ for the diagonalizer's run machine (the phase host over the run's phases), with accepting verdict $\mathit{acc}$ and rejecting verdict $\mathit{rej}$; a verdict-avoiding segment of length $t$ from $a$ to $b$ means that $Q_\kappa$ reaches $b$ from $a$ in exactly $t$ steps and is in neither verdict at any earlier time. Write $\langle u, v\rangle$ for the pairing encoding and $\alpha$ for the inner second component of a binary word $y$ (the second component of the first component of its decoding). Let $y, B$ be binary words with $L(2) = \mathrm{buf}(1z)$, $L(8) = \mathrm{buf}(y)$, $L(10) = \mathrm{buf}(B)$, $L(11) = \mathrm{buf}(\varepsilon)$ and $P(2) = P(8) = P(10) = P(11) = 0$. Then there is $t \le A_a(|y| + 1) + 10|B| + 6|\alpha| + 4|z| + 40$ and a verdict-avoiding segment of length $t$ from $\langle (\text{extraction}, \text{start}), 1, \hat L, \hat P, \varepsilon\rangle$ to the hosted core's entry $\langle h, 1, \hat{L'}, \hat{P'}, \varepsilon\rangle$, where $h$ is the clock's injection of the core's entry state with flag $1$, $L' = L[8 \mapsto \mathrm{buf}(\alpha)][11 \mapsto \mathrm{buf}(\langle\langle B, \alpha\rangle, 1z\rangle)]$ and $P' = P[10 \mapsto |B|][8 \mapsto |\alpha|][2 \mapsto |z| + 1]$.
\end{theorem}

\begin{theorem}[The middle's hosted interpreter run]
\lean{Turing.NTimeDiag.dz_runR}
\label{Turing.NTimeDiag.dz_runR}
\leanok
\uses{Turing.Cfg.Stopped, Turing.Cfg.stopped_iff, Turing.EffectiveNDMachineCode, Turing.FinTM, Turing.FinTM.bufferTape, Turing.FinTM.bufferTape_nil, Turing.MultiTapeTM.runFrom, Turing.MultiTapeTM.runFrom_state_ne_none_of_le, Turing.NTimeDiag.DzKit, Turing.NTimeDiag.DzKit.K, Turing.NTimeDiag.DzKit.bR, Turing.NTimeDiag.DzKit.bR_j, Turing.NTimeDiag.DzKit.bR_ne, Turing.NTimeDiag.DzKit.bR_off, Turing.NTimeDiag.DzKit.sR, Turing.NTimeDiag.DzKit.tp, Turing.NTimeDiag.DzKit.tp_ne, Turing.NTimeDiag.dzAcc, Turing.NTimeDiag.dzLay, Turing.NTimeDiag.dzLay_off, Turing.NTimeDiag.dzLay_tp, Turing.NTimeDiag.dzQ, Turing.NTimeDiag.dzQF, Turing.NTimeDiag.dzQM, Turing.NTimeDiag.dzQRed, Turing.NTimeDiag.dzRej, Turing.NTimeDiag.dz_notF, Turing.NTimeDiag.qyC, Turing.NTimeDiag.qy_bits_val, Turing.NTimeDiag.qy_hosted, Turing.Segment, Turing.embedEmitTM, Turing.streamStartCfg}
\difficulty{6}
Let $\kappa$ be a diagonalizer's kit for $f, g : \mathbb{N} \to \mathbb{N}$, an effective coding scheme $c$ and a finite deterministic decider $B_F$ over bits, with $K_\kappa$ host tapes and named tapes $\tau_0, \dots, \tau_{13}$ (the host tapes of indices $0, \dots, 13$), let $z$ be a binary input, let $L$ assign a tape to each index $i < 14$ and $P$ a head position to each index. Write $\mathrm{buf}(w)$ for the tape holding the binary word $w$ in cells $0, \dots, |w| - 1$ and blank elsewhere, $L[i \mapsto v]$ for $L$ updated at $i$ to $v$, $\hat L$ for the tape layout giving $\tau_i$ the contents $L(i)$ and every other host tape the blank tape, and $\hat P$ for the head layout giving $\tau_i$ the position $P(i)$ and every other head the position $0$, $\langle s, p, T, P, o\rangle$ for the live configuration on input $z$ with state $s$, input head $p$, work tapes $T$, heads $P$ and output $o$, and $Q_\kappa$ for the diagonalizer's run machine (the phase host over the run's phases), with accepting verdict $\mathit{acc}$ and rejecting verdict $\mathit{rej}$; a verdict-avoiding segment of length $t$ from $a$ to $b$ means that $Q_\kappa$ reaches $b$ from $a$ in exactly $t$ steps and is in neither verdict at any earlier time. Let $z', s_1$ be binary words and $G \in \mathbb{N}$ with $L(3) = \mathrm{buf}(\mathrm{bits}(G))$, $L(6) = \mathrm{buf}(s_1)$, $L(11) = \mathrm{buf}(z')$, $L(12) = \mathrm{buf}(\varepsilon)$ and $P(3) = P(6) = P(11) = P(12) = 0$. Let $R$ be the kit's interpreter core with entry state $r_E$, accepting state $r_A$ and rejecting state $r_J$, started in $r_E$ on input $z'$ with $s_1$ on its stream tape, other tapes blank, heads at $0$ and input head $1$; suppose that at time $T_R$ it is in a state $v \in \{r_A, r_J\}$ and at every earlier time in neither. Then there are $t \le 5G + 2|\mathrm{bits}(G)| + 3$, tapes $T'$ and heads $P'$ with a verdict-avoiding segment of length $t$ from $\langle h, 1, \hat L, \hat P, \varepsilon\rangle$, $h$ the clock's injection of $r_E$ with flag $1$ in the hosted core phase, to $\langle e, 1, T', P', \varepsilon\rangle$, where $e = \mathit{acc}$ if $T_R \le G$ and $v = r_A$, and $e = \mathit{rej}$ otherwise.
\end{theorem}

\begin{theorem}[The top's hosted decider run and read]
\lean{Turing.NTimeDiag.dz_top_run}
\label{Turing.NTimeDiag.dz_top_run}
\leanok
\uses{Turing.Cfg.Stopped, Turing.Cfg.stopped_iff, Turing.EffectiveNDMachineCode, Turing.FinTM, Turing.FinTM.bufferTape, Turing.FinTM.bufferTape_nat, Turing.FinTM.bufferTape_nil, Turing.MultiTapeTM.exists_first_halt, Turing.MultiTapeTM.initCfg, Turing.MultiTapeTM.initCfg_eq, Turing.MultiTapeTM.output_length_le, Turing.MultiTapeTM.runFrom, Turing.MultiTapeTM.runFrom_output_eq_of_halt, Turing.MultiTapeTM.runFrom_zero, Turing.NTimeDiag.DzKit, Turing.NTimeDiag.DzKit.K, Turing.NTimeDiag.DzKit.bB, Turing.NTimeDiag.DzKit.bB_ne, Turing.NTimeDiag.DzKit.sB, Turing.NTimeDiag.DzKit.tp, Turing.NTimeDiag.DzKit.tp_ne, Turing.NTimeDiag.DzQPh, Turing.NTimeDiag.dzAcc, Turing.NTimeDiag.dzLay, Turing.NTimeDiag.dzLay_off, Turing.NTimeDiag.dzLay_tp, Turing.NTimeDiag.dzQ, Turing.NTimeDiag.dzQF, Turing.NTimeDiag.dzQM, Turing.NTimeDiag.dzQRed, Turing.NTimeDiag.dzRej, Turing.NTimeDiag.dz_acc_ne_rej, Turing.NTimeDiag.dz_bank_off, Turing.NTimeDiag.dz_notF, Turing.NTimeDiag.dz_take_run, Turing.NTimeDiag.qyC, Turing.NTimeDiag.qy_bits_val, Turing.NTimeDiag.qy_guard_of_cut, Turing.NTimeDiag.qy_hosted, Turing.NTimeDiag.qy_rewind, Turing.Segment, Turing.Segment.trans, Turing.embedEmitTM}
\difficulty{6}
Let $\kappa$ be a diagonalizer's kit for $f, g : \mathbb{N} \to \mathbb{N}$, an effective coding scheme $c$ and a finite deterministic decider $B_F$ over bits, with $K_\kappa$ host tapes and named tapes $\tau_0, \dots, \tau_{13}$ (the host tapes of indices $0, \dots, 13$), let $z$ be a binary input, let $L$ assign a tape to each index $i < 14$ and $P$ a head position to each index. Write $\mathrm{buf}(w)$ for the tape holding the binary word $w$ in cells $0, \dots, |w| - 1$ and blank elsewhere, $L[i \mapsto v]$ for $L$ updated at $i$ to $v$, $\hat L$ for the tape layout giving $\tau_i$ the contents $L(i)$ and every other host tape the blank tape, and $\hat P$ for the head layout giving $\tau_i$ the position $P(i)$ and every other head the position $0$, $\langle s, p, T, P, o\rangle$ for the live configuration on input $z$ with state $s$, input head $p$, work tapes $T$, heads $P$ and output $o$, and $Q_\kappa$ for the diagonalizer's run machine (the phase host over the run's phases), with accepting verdict $\mathit{acc}$ and rejecting verdict $\mathit{rej}$; a verdict-avoiding segment of length $t$ from $a$ to $b$ means that $Q_\kappa$ reaches $b$ from $a$ in exactly $t$ steps and is in neither verdict at any earlier time. Let $y$ be a binary word and $G \in \mathbb{N}$ with $L(3) = \mathrm{buf}(\mathrm{bits}(G))$, $L(8) = \mathrm{buf}(y)$, $L(12) = L(13) = \mathrm{buf}(\varepsilon)$ and $P(3) = P(8) = P(12) = P(13) = 0$. Then there are $t \le 6G + 2|\mathrm{bits}(G)| + 10$, a verdict $e \in \{\mathit{acc}, \mathit{rej}\}$, tapes $T'$ and heads $P'$ with a verdict-avoiding segment of length $t$ from $\langle h, 1, \hat L, \hat P, \varepsilon\rangle$, $h$ the clock's injection of the start state of $B_F$ with flag $1$ in the hosted decider phase, to $\langle e, 1, T', P', \varepsilon\rangle$, such that for every $\tau \le G$ at which $B_F$, run from its initial configuration on $y$, has halted, $e = \mathit{acc}$ holds if and only if the output of $B_F$ at time $\tau$ begins with $0$.
\end{theorem}

\begin{theorem}[Bound on the value of a budget word]
\lean{Turing.NTimeDiag.dz_budget_le}
\label{Turing.NTimeDiag.dz_budget_le}
\leanok
\uses{Complexity.LogProg.ofDigits_lt_two_pow, Complexity.LogProg.two_pow_bits_le_two_mul_add_one}
\difficulty{3}
Let $G \in \mathbb{N}$ and let $w$ be a binary word with $|w| = |\mathrm{bits}(G)|$, where $\mathrm{bits}(G)$ is the binary expansion of $G$. Then the value $\sum_i w_i 2^i$ of $w$, read least significant bit first, is at most $2G$.
\end{theorem}

\begin{theorem}[Value of zero-padded binary digits]
\lean{Turing.NTimeDiag.dz_padded_value}
\label{Turing.NTimeDiag.dz_padded_value}
\leanok
\uses{Turing.NTimeDiag.qy_bits_val}
\difficulty{2}
For all $t, r \in \mathbb{N}$, the binary word $\mathrm{bits}(t)\,0^r$ (the binary expansion of $t$, least significant bit first, followed by $r$ zeros), read as a binary numeral least significant bit first, has value $t$.
\end{theorem}

\begin{definition}[Time bound of the run]
\lean{Turing.NTimeDiag.dzQB}
\label{Turing.NTimeDiag.dzQB}
\leanok
\uses{Turing.EffectiveNDMachineCode, Turing.FinTM, Turing.NTimeDiag.DzKit}
Let $\kappa$ be a diagonalizer's kit for $f, g : \mathbb{N} \to \mathbb{N}$, an effective coding scheme $c$ and a finite deterministic decider $B_F$ over bits, and let $A_q, A_a, A_l$ be the constants of the kit's calls for the stage query, for the inner second component and for $y \mapsto \mathrm{bits}(|y|)$. For $G \in \mathbb{N}$ (standing for $g(n)$), the run's time bound is $(A_q(7A_a + 8) + 2A_l + 140)(G + 1)$.
\end{definition}

\begin{theorem}[Named tapes at the run's entry]
\lean{Turing.NTimeDiag.dz_entry_vals}
\label{Turing.NTimeDiag.dz_entry_vals}
\leanok
\uses{Turing.FinTM.bufferTape, Turing.NTimeDiag.dzSetL, Turing.NTimeDiag.query}
\difficulty{2}
Let $f, g : \mathbb{N} \to \mathbb{N}$, let $q_{f,g}$ be the stage query, let $z, s$ be binary words, write $\mathrm{buf}(w)$ for the tape holding the binary word $w$ in cells $0, \dots, |w| - 1$ and blank elsewhere, and let $L$ be the named-tape contents at the setup's exit for $z$ with index $6$ replaced by $\mathrm{buf}(s)$. Then $L(0) = \mathrm{buf}(q_{f,g}(z))$, $L(2) = \mathrm{buf}(1z)$, $L(3) = \mathrm{buf}(\mathrm{bits}(g(|z|)))$, $L(6) = \mathrm{buf}(s)$, $L(7) = \mathrm{buf}(1^{g(|z|)})$, and $L(i) = \mathrm{buf}(\varepsilon)$ for $i = 8, \dots, 13$.
\end{theorem}

\begin{theorem}[The run off the stages]
\lean{Turing.NTimeDiag.dz_qoff}
\label{Turing.NTimeDiag.dz_qoff}
\leanok
\uses{Turing.EffectiveNDMachineCode, Turing.FinTM, Turing.FinTM.bufferTape, Turing.FinTM.bufferTape_nil, Turing.NTimeDiag.DzKit, Turing.NTimeDiag.DzKit.K, Turing.NTimeDiag.dzLay, Turing.NTimeDiag.dzLay_const, Turing.NTimeDiag.dzQ, Turing.NTimeDiag.dzQB, Turing.NTimeDiag.dzQF, Turing.NTimeDiag.dzRej, Turing.NTimeDiag.dzSetL, Turing.NTimeDiag.dz_entry_vals, Turing.NTimeDiag.dz_flag, Turing.NTimeDiag.query, Turing.NTimeDiag.qyC, Turing.Segment}
\difficulty{3}
Let $\kappa$ be a diagonalizer's kit for $f, g : \mathbb{N} \to \mathbb{N}$, an effective coding scheme $c$ and a finite deterministic decider $B_F$ over bits, with $K_\kappa$ host tapes and named tapes $\tau_0, \dots, \tau_{13}$ (the host tapes of indices $0, \dots, 13$), and let $z, s$ be binary words. Write $\mathrm{buf}(w)$ for the tape holding the binary word $w$ in cells $0, \dots, |w| - 1$ and blank elsewhere, $L[i \mapsto v]$ for $L$ updated at $i$ to $v$, $\hat L$ for the tape layout giving $\tau_i$ the contents $L(i)$ and every other host tape the blank tape, and $\hat P$ for the head layout giving $\tau_i$ the position $P(i)$ and every other head the position $0$, $\langle s, p, T, P, o\rangle$ for the live configuration on input $z$ with state $s$, input head $p$, work tapes $T$, heads $P$ and output $o$, and $Q_\kappa$ for the diagonalizer's run machine (the phase host over the run's phases), with accepting verdict $\mathit{acc}$ and rejecting verdict $\mathit{rej}$; a verdict-avoiding segment of length $t$ from $a$ to $b$ means that $Q_\kappa$ reaches $b$ from $a$ in exactly $t$ steps and is in neither verdict at any earlier time. Write $\Lambda_z$ for the named-tape contents at the setup's exit: $\mathrm{buf}(q_{f,g}(z))$ on $\tau_0$, where $q_{f,g}$ is the stage query; $\mathrm{buf}(1^m)$ on $\tau_1$ and $\tau_4$, where $m = |\mathrm{bits}(g(|z|))|$; $\mathrm{buf}(1z)$ on $\tau_2$; $\mathrm{buf}(\mathrm{bits}(g(|z|)))$ on $\tau_3$; $\mathrm{buf}(1)$ on $\tau_5$; the block $\mathrm{buf}(1^{g(|z|)}\,1\,\mathrm{bits}(g(|z|)))$ on $\tau_6$; $\mathrm{buf}(1^{g(|z|)})$ on $\tau_7$; and the blank tape on $\tau_8, \dots, \tau_{13}$, and let $B_Q(G) = (A_q(7A_a + 8) + 2A_l + 140)(G + 1)$, where $A_q$, $A_a$, $A_l$ are the constants of the kit's calls for the stage query, for the inner second component and for $y \mapsto \mathrm{bits}(|y|)$. If $q_{f,g}(z) = \varepsilon$, then there are $t \le B_Q(g(|z|))$, tapes $T'$ and heads $P'$ with a verdict-avoiding segment of length $t$ from the run's entry configuration $E_s = \langle (\text{flag}, \text{start}), 1, \widehat{\Lambda_z[6 \mapsto \mathrm{buf}(s)]}, 0, \varepsilon\rangle$ to $\langle \mathit{rej}, 1, T', P', \varepsilon\rangle$.
\end{theorem}

\begin{theorem}[First cell of a buffered word]
\lean{Turing.NTimeDiag.dz_cell0}
\label{Turing.NTimeDiag.dz_cell0}
\leanok
\uses{Turing.FinTM.bufferTape, Turing.FinTM.bufferTape_nat}
\difficulty{1}
For every bit $b$ and binary word $w$, the tape holding $b\,w$ in cells $0, \dots, |w|$ and blank elsewhere holds $b$ at cell $0$.
\end{theorem}

\begin{theorem}[The run at a stage's top]
\lean{Turing.NTimeDiag.dz_qtop}
\label{Turing.NTimeDiag.dz_qtop}
\leanok
\uses{Turing.EffectiveNDMachineCode, Turing.FinTM, Turing.FinTM.bufferTape, Turing.MultiTapeTM.initCfg, Turing.MultiTapeTM.runFrom, Turing.NTimeDiag.DzKit, Turing.NTimeDiag.DzKit.K, Turing.NTimeDiag.DzKit.tp, Turing.NTimeDiag.DzKit.tp_ne, Turing.NTimeDiag.DzQPh, Turing.NTimeDiag.dzAcc, Turing.NTimeDiag.dzLay, Turing.NTimeDiag.dzLay_const, Turing.NTimeDiag.dzLay_tp, Turing.NTimeDiag.dzLay_update, Turing.NTimeDiag.dzQ, Turing.NTimeDiag.dzQB, Turing.NTimeDiag.dzQF, Turing.NTimeDiag.dzQM, Turing.NTimeDiag.dzQRed, Turing.NTimeDiag.dzRej, Turing.NTimeDiag.dzSetL, Turing.NTimeDiag.dz_cell0, Turing.NTimeDiag.dz_entry_vals, Turing.NTimeDiag.dz_flag, Turing.NTimeDiag.dz_notF, Turing.NTimeDiag.dz_tail_run, Turing.NTimeDiag.dz_top_run, Turing.NTimeDiag.query, Turing.NTimeDiag.qyC, Turing.NTimeDiag.qy_guard_of_cut, Turing.NTimeDiag.qy_hosted, Turing.Segment, Turing.Segment.trans, Turing.length_bits_le_self}
\difficulty{5}
Let $\kappa$ be a diagonalizer's kit for $f, g : \mathbb{N} \to \mathbb{N}$, an effective coding scheme $c$ and a finite deterministic decider $B_F$ over bits, with $K_\kappa$ host tapes and named tapes $\tau_0, \dots, \tau_{13}$ (the host tapes of indices $0, \dots, 13$), and let $z, y, s$ be binary words. Write $\mathrm{buf}(w)$ for the tape holding the binary word $w$ in cells $0, \dots, |w| - 1$ and blank elsewhere, $L[i \mapsto v]$ for $L$ updated at $i$ to $v$, $\hat L$ for the tape layout giving $\tau_i$ the contents $L(i)$ and every other host tape the blank tape, and $\hat P$ for the head layout giving $\tau_i$ the position $P(i)$ and every other head the position $0$, $\langle s, p, T, P, o\rangle$ for the live configuration on input $z$ with state $s$, input head $p$, work tapes $T$, heads $P$ and output $o$, and $Q_\kappa$ for the diagonalizer's run machine (the phase host over the run's phases), with accepting verdict $\mathit{acc}$ and rejecting verdict $\mathit{rej}$; a verdict-avoiding segment of length $t$ from $a$ to $b$ means that $Q_\kappa$ reaches $b$ from $a$ in exactly $t$ steps and is in neither verdict at any earlier time. Write $\Lambda_z$ for the named-tape contents at the setup's exit: $\mathrm{buf}(q_{f,g}(z))$ on $\tau_0$, where $q_{f,g}$ is the stage query; $\mathrm{buf}(1^m)$ on $\tau_1$ and $\tau_4$, where $m = |\mathrm{bits}(g(|z|))|$; $\mathrm{buf}(1z)$ on $\tau_2$; $\mathrm{buf}(\mathrm{bits}(g(|z|)))$ on $\tau_3$; $\mathrm{buf}(1)$ on $\tau_5$; the block $\mathrm{buf}(1^{g(|z|)}\,1\,\mathrm{bits}(g(|z|)))$ on $\tau_6$; $\mathrm{buf}(1^{g(|z|)})$ on $\tau_7$; and the blank tape on $\tau_8, \dots, \tau_{13}$, and let $B_Q(G) = (A_q(7A_a + 8) + 2A_l + 140)(G + 1)$, where $A_q$, $A_a$, $A_l$ are the constants of the kit's calls for the stage query, for the inner second component and for $y \mapsto \mathrm{bits}(|y|)$. If $q_{f,g}(z) = 1\,y$, then there are $t \le B_Q(g(|z|))$, a verdict $e \in \{\mathit{acc}, \mathit{rej}\}$, tapes $T'$ and heads $P'$ with a verdict-avoiding segment of length $t$ from the run's entry configuration $E_s = \langle (\text{flag}, \text{start}), 1, \widehat{\Lambda_z[6 \mapsto \mathrm{buf}(s)]}, 0, \varepsilon\rangle$ to $\langle e, 1, T', P', \varepsilon\rangle$, such that for every $\tau \le g(|z|)$ at which $B_F$, run from its initial configuration on $y$, has halted, $e = \mathit{acc}$ holds if and only if the output of $B_F$ at time $\tau$ begins with $0$.
\end{theorem}

\begin{theorem}[The run in a stage's middle]
\lean{Turing.NTimeDiag.dz_qmid}
\label{Turing.NTimeDiag.dz_qmid}
\leanok
\uses{Turing.CodeNDTM.toFinNDTM, Turing.EffectiveNDMachineCode, Turing.FinTM, Turing.FinTM.bufferTape, Turing.MultiTapeTM.runFrom, Turing.NDTM.initCfg, Turing.NDTM.runWith, Turing.NTimeDiag.DzKit, Turing.NTimeDiag.DzKit.K, Turing.NTimeDiag.DzKit.tp, Turing.NTimeDiag.DzKit.tp_ne, Turing.NTimeDiag.DzQPh, Turing.NTimeDiag.dzAcc, Turing.NTimeDiag.dzAlpha, Turing.NTimeDiag.dzLay, Turing.NTimeDiag.dzLay_const, Turing.NTimeDiag.dzLay_tp, Turing.NTimeDiag.dzLay_update, Turing.NTimeDiag.dzQ, Turing.NTimeDiag.dzQB, Turing.NTimeDiag.dzQF, Turing.NTimeDiag.dzQM, Turing.NTimeDiag.dzQRed, Turing.NTimeDiag.dzRej, Turing.NTimeDiag.dzSetL, Turing.NTimeDiag.dz_acc_ne_rej, Turing.NTimeDiag.dz_budget_le, Turing.NTimeDiag.dz_cell0, Turing.NTimeDiag.dz_entry_vals, Turing.NTimeDiag.dz_flag, Turing.NTimeDiag.dz_norm, Turing.NTimeDiag.dz_notF, Turing.NTimeDiag.dz_rewrite, Turing.NTimeDiag.dz_runR, Turing.NTimeDiag.dz_split, Turing.NTimeDiag.dz_tail_run, Turing.NTimeDiag.dz_test, Turing.NTimeDiag.query, Turing.NTimeDiag.qyC, Turing.NTimeDiag.qy_guard_of_cut, Turing.NTimeDiag.qy_hosted, Turing.Segment, Turing.Segment.trans, Turing.length_bits_le_self, Turing.pairEncode, Turing.streamStartCfg}
\difficulty{6}
Let $\kappa$ be a diagonalizer's kit for $f, g : \mathbb{N} \to \mathbb{N}$, an effective coding scheme $c$ and a finite deterministic decider $B_F$ over bits, with $K_\kappa$ host tapes and named tapes $\tau_0, \dots, \tau_{13}$ (the host tapes of indices $0, \dots, 13$), and let $z, y$ be binary words with $q_{f,g}(z) = 0\,y$. Write $\mathrm{buf}(w)$ for the tape holding the binary word $w$ in cells $0, \dots, |w| - 1$ and blank elsewhere, $L[i \mapsto v]$ for $L$ updated at $i$ to $v$, $\hat L$ for the tape layout giving $\tau_i$ the contents $L(i)$ and every other host tape the blank tape, and $\hat P$ for the head layout giving $\tau_i$ the position $P(i)$ and every other head the position $0$, $\langle s, p, T, P, o\rangle$ for the live configuration on input $z$ with state $s$, input head $p$, work tapes $T$, heads $P$ and output $o$, and $Q_\kappa$ for the diagonalizer's run machine (the phase host over the run's phases), with accepting verdict $\mathit{acc}$ and rejecting verdict $\mathit{rej}$; a verdict-avoiding segment of length $t$ from $a$ to $b$ means that $Q_\kappa$ reaches $b$ from $a$ in exactly $t$ steps and is in neither verdict at any earlier time. Write $\Lambda_z$ for the named-tape contents at the setup's exit: $\mathrm{buf}(q_{f,g}(z))$ on $\tau_0$, where $q_{f,g}$ is the stage query; $\mathrm{buf}(1^m)$ on $\tau_1$ and $\tau_4$, where $m = |\mathrm{bits}(g(|z|))|$; $\mathrm{buf}(1z)$ on $\tau_2$; $\mathrm{buf}(\mathrm{bits}(g(|z|)))$ on $\tau_3$; $\mathrm{buf}(1)$ on $\tau_5$; the block $\mathrm{buf}(1^{g(|z|)}\,1\,\mathrm{bits}(g(|z|)))$ on $\tau_6$; $\mathrm{buf}(1^{g(|z|)})$ on $\tau_7$; and the blank tape on $\tau_8, \dots, \tau_{13}$, let $B_Q(G) = (A_q(7A_a + 8) + 2A_l + 140)(G + 1)$, where $A_q$, $A_a$, $A_l$ are the constants of the kit's calls for the stage query, for the inner second component and for $y \mapsto \mathrm{bits}(|y|)$, write $\langle u, v\rangle$ for the pairing encoding, let $\alpha$ be the inner second component of $y$ (the second component of the first component of its decoding), and let $M_\alpha$ be the nondeterministic machine decoded from $\alpha$ by $c$. Let $C \in \mathbb{N}$ be such that for all binary words $x, s$ and every $t \le |s|$ there is $T \le C(t + 1)$ at which the kit's interpreter core $R$, started in its entry state on input $\langle\langle \mathrm{bits}(t), \alpha\rangle, x\rangle$ with $s$ on its stream tape, other tapes blank and heads at $0$, is in its accepting state $r_A$ if $M_\alpha$ run on $x$ along the first $t$ symbols of $s$ has halted with output $1$, and in its rejecting state $r_J$ otherwise, and is in neither of these states before $T$. Let $s_1, w$ be binary words and $b'$ a bit with $|s_1| = g(|z|)$ and $|w| = |\mathrm{bits}(g(|z|))|$, let $b = \sum_i w_i 2^i$, and let $s = s_1\,b'\,w$. Then there are $t \le B_Q(g(|z|))$, a verdict $e \in \{\mathit{acc}, \mathit{rej}\}$, tapes $T'$ and heads $P'$ with a verdict-avoiding segment of length $t$ from the run's entry configuration $E_s = \langle (\text{flag}, \text{start}), 1, \widehat{\Lambda_z[6 \mapsto \mathrm{buf}(s)]}, 0, \varepsilon\rangle$ to $\langle e, 1, T', P', \varepsilon\rangle$, such that: if $e = \mathit{acc}$, then $b \le g(|z|)$ and $M_\alpha$, run on $1z$ along the choice word consisting of the first $b$ symbols of $s_1$, has halted with output $1$; and if $b \le g(|z|)$, $C(b + 1) \le g(|z|)$ and $M_\alpha$ run on $1z$ along the first $b$ symbols of $s_1$ has halted with output $1$, then $e = \mathit{acc}$.
\end{theorem}

\begin{definition}[The diagonalizer]
\lean{Turing.NTimeDiag.dzD}
\label{Turing.NTimeDiag.dzD}
\leanok
\uses{Turing.EffectiveNDMachineCode, Turing.FinNDTM, Turing.FinTM, Turing.GuessRunPhase.State, Turing.NTimeDiag.DzKit, Turing.NTimeDiag.DzKit.K, Turing.NTimeDiag.DzKit.tp, Turing.NTimeDiag.DzPPh, Turing.NTimeDiag.DzQPh, Turing.NTimeDiag.dzAcc, Turing.NTimeDiag.dzPexit, Turing.NTimeDiag.dzQ, Turing.NTimeDiag.dzRej, Turing.NTimeDiag.dzSet, Turing.guessRunTM}
Let $\kappa$ be a diagonalizer's kit for $f, g : \mathbb{N} \to \mathbb{N}$, an effective coding scheme $c$ and a finite deterministic decider $B_F$ over bits. The diagonalizer $D_\kappa$ is the finite nondeterministic machine over bits with the kit's $K_\kappa$ host tapes whose states are the phase-tagged states of the guess-then-run host, and whose transition is the guess-then-run host over the setup machine (with the done state of the ruler's rewind as its exit), the block tape $\tau_6$ as the guess tape, and the run machine entered at the flag's start state with the run's accepting and rejecting verdicts.
\end{definition}

\begin{definition}[The diagonalizer's uniform constant]
\lean{Turing.NTimeDiag.dzK}
\label{Turing.NTimeDiag.dzK}
\leanok
\uses{Turing.EffectiveNDMachineCode, Turing.FinTM, Turing.NTimeDiag.DzKit}
Let $\kappa$ be a diagonalizer's kit for $f, g : \mathbb{N} \to \mathbb{N}$, an effective coding scheme $c$ and a finite deterministic decider $B_F$ over bits, and let $A_q, A_g, A_p, A_l, A_a$ be the constants of the kit's calls for the stage query, for $\mathrm{bits}(g(|x|))$, for prepending $1$, for $y \mapsto \mathrm{bits}(|y|)$ and for the inner second component. The diagonalizer's uniform constant is $K_D = (A_q + A_g + A_p + 40) + (A_q(7A_a + 8) + 2A_l + 140) + 15$.
\end{definition}

\begin{theorem}[Initial configuration of the setup]
\lean{Turing.NTimeDiag.dz_initCfg}
\label{Turing.NTimeDiag.dz_initCfg}
\leanok
\uses{Turing.Cfg.ofWords_inputPos, Turing.Cfg.ofWords_output, Turing.Cfg.ofWords_state, Turing.Cfg.ofWords_workTapePos, Turing.Cfg.ofWords_workTapes, Turing.EffectiveNDMachineCode, Turing.FinTM, Turing.FinTM.bufferTape, Turing.MultiTapeTM.initCfg, Turing.NTimeDiag.DzKit, Turing.NTimeDiag.dzSet, Turing.NTimeDiag.qyC, Turing.initCfg_ofWords, Turing.phaseHostTM_q₀}
\difficulty{3}
Let $\kappa$ be a diagonalizer's kit for $f, g : \mathbb{N} \to \mathbb{N}$, an effective coding scheme $c$ and a finite deterministic decider $B_F$ over bits, and let $z$ be a binary input. The initial configuration of the diagonalizer's setup machine on $z$ is the live configuration on input $z$ in the fresh start state of the stage-query call phase, with input head $1$, every work tape blank, every work head at $0$ and empty output.
\end{theorem}

\begin{theorem}[Splitting a word at a position]
\lean{Turing.NTimeDiag.dz_split_stream}
\label{Turing.NTimeDiag.dz_split_stream}
\leanok
\difficulty{2}
Let $s$ be a binary word and $G \in \mathbb{N}$ with $G < |s|$. Then $s$ equals the concatenation of its first $G$ symbols, its symbol at position $G$, and its symbols from position $G + 1$ on.
\end{theorem}

\begin{theorem}[Reading a verdict off a segment]
\lean{Turing.NTimeDiag.dz_seg_read}
\label{Turing.NTimeDiag.dz_seg_read}
\leanok
\uses{Turing.Cfg, Turing.EffectiveNDMachineCode, Turing.FinTM, Turing.MultiTapeTM.runFrom, Turing.NTimeDiag.DzKit, Turing.NTimeDiag.DzKit.K, Turing.NTimeDiag.DzQPh, Turing.NTimeDiag.dzAcc, Turing.NTimeDiag.dzQ, Turing.NTimeDiag.dzQF, Turing.NTimeDiag.dzRej, Turing.NTimeDiag.qyC, Turing.Segment, Turing.segment_iff}
\difficulty{3}
Let $\kappa$ be a diagonalizer's kit for $f, g : \mathbb{N} \to \mathbb{N}$, an effective coding scheme $c$ and a finite deterministic decider $B_F$ over bits, with $K_\kappa$ host tapes and named tapes $\tau_0, \dots, \tau_{13}$ (the host tapes of indices $0, \dots, 13$), and let $z$ be a binary input. Write $\mathrm{buf}(w)$ for the tape holding the binary word $w$ in cells $0, \dots, |w| - 1$ and blank elsewhere, $L[i \mapsto v]$ for $L$ updated at $i$ to $v$, $\hat L$ for the tape layout giving $\tau_i$ the contents $L(i)$ and every other host tape the blank tape, and $\hat P$ for the head layout giving $\tau_i$ the position $P(i)$ and every other head the position $0$, $\langle s, p, T, P, o\rangle$ for the live configuration on input $z$ with state $s$, input head $p$, work tapes $T$, heads $P$ and output $o$, and $Q_\kappa$ for the diagonalizer's run machine (the phase host over the run's phases), with accepting verdict $\mathit{acc}$ and rejecting verdict $\mathit{rej}$; a verdict-avoiding segment of length $t$ from $a$ to $b$ means that $Q_\kappa$ reaches $b$ from $a$ in exactly $t$ steps and is in neither verdict at any earlier time. Let $a$ be a configuration of $Q_\kappa$ on $z$, $t \in \mathbb{N}$, $v$ a state of $Q_\kappa$, and $T', P'$ tapes and heads. If there is a verdict-avoiding segment of length $t$ from $a$ to $\langle v, 1, T', P', \varepsilon\rangle$, then $Q_\kappa$ run from $a$ is in state $v$ at time $t$, is in neither $\mathit{acc}$ nor $\mathit{rej}$ at any time before $t$, and has empty output at time $t$.
\end{theorem}

\begin{theorem}[Every guessed stream's run reaches a verdict]
\lean{Turing.NTimeDiag.dz_hQ}
\label{Turing.NTimeDiag.dz_hQ}
\leanok
\uses{Turing.Cfg, Turing.EffectiveNDMachineCode, Turing.FinTM, Turing.FinTM.bufferTape, Turing.MultiTapeTM.runFrom, Turing.NTimeDiag.DzKit, Turing.NTimeDiag.DzKit.K, Turing.NTimeDiag.DzKit.tp, Turing.NTimeDiag.DzQPh, Turing.NTimeDiag.dzAcc, Turing.NTimeDiag.dzAlpha, Turing.NTimeDiag.dzBlk, Turing.NTimeDiag.dzBlk_length, Turing.NTimeDiag.dzLay, Turing.NTimeDiag.dzPexit, Turing.NTimeDiag.dzQ, Turing.NTimeDiag.dzQB, Turing.NTimeDiag.dzQF, Turing.NTimeDiag.dzRej, Turing.NTimeDiag.dzSetL, Turing.NTimeDiag.dz_handoff, Turing.NTimeDiag.dz_qmid, Turing.NTimeDiag.dz_qoff, Turing.NTimeDiag.dz_qtop, Turing.NTimeDiag.dz_seg_read, Turing.NTimeDiag.dz_split_stream, Turing.NTimeDiag.query, Turing.NTimeDiag.qyC, Turing.Segment, Turing.guessRunCfg}
\difficulty{5}
Let $\kappa$ be a diagonalizer's kit for $f, g : \mathbb{N} \to \mathbb{N}$, an effective coding scheme $c$ and a finite deterministic decider $B_F$ over bits, with $K_\kappa$ host tapes and named tapes $\tau_0, \dots, \tau_{13}$ (the host tapes of indices $0, \dots, 13$), and let $z, s$ be binary words with $|s| = |1^{g(|z|)}\,1\,\mathrm{bits}(g(|z|))|$. Write $\mathrm{buf}(w)$ for the tape holding the binary word $w$ in cells $0, \dots, |w| - 1$ and blank elsewhere, $L[i \mapsto v]$ for $L$ updated at $i$ to $v$, $\hat L$ for the tape layout giving $\tau_i$ the contents $L(i)$ and every other host tape the blank tape, and $\hat P$ for the head layout giving $\tau_i$ the position $P(i)$ and every other head the position $0$, $\langle s, p, T, P, o\rangle$ for the live configuration on input $z$ with state $s$, input head $p$, work tapes $T$, heads $P$ and output $o$, and $Q_\kappa$ for the diagonalizer's run machine (the phase host over the run's phases), with accepting verdict $\mathit{acc}$ and rejecting verdict $\mathit{rej}$; a verdict-avoiding segment of length $t$ from $a$ to $b$ means that $Q_\kappa$ reaches $b$ from $a$ in exactly $t$ steps and is in neither verdict at any earlier time. Write $\Lambda_z$ for the named-tape contents at the setup's exit: $\mathrm{buf}(q_{f,g}(z))$ on $\tau_0$, where $q_{f,g}$ is the stage query; $\mathrm{buf}(1^m)$ on $\tau_1$ and $\tau_4$, where $m = |\mathrm{bits}(g(|z|))|$; $\mathrm{buf}(1z)$ on $\tau_2$; $\mathrm{buf}(\mathrm{bits}(g(|z|)))$ on $\tau_3$; $\mathrm{buf}(1)$ on $\tau_5$; the block $\mathrm{buf}(1^{g(|z|)}\,1\,\mathrm{bits}(g(|z|)))$ on $\tau_6$; $\mathrm{buf}(1^{g(|z|)})$ on $\tau_7$; and the blank tape on $\tau_8, \dots, \tau_{13}$, extended blank to the other host tapes, and let $B_Q(G) = (A_q(7A_a + 8) + 2A_l + 140)(G + 1)$, where $A_q$, $A_a$, $A_l$ are the constants of the kit's calls for the stage query, for the inner second component and for $y \mapsto \mathrm{bits}(|y|)$. Let $c_s$ be the configuration obtained from the setup's exit configuration $\langle \mathit{exit}, 1, \Lambda_z, 0, \varepsilon\rangle$ by writing $s$ on $\tau_6$ from its head onwards and switching to the flag's start state. Then there is $T \le B_Q(g(|z|))$ such that $Q_\kappa$ run from $c_s$ is in $\mathit{acc}$ or in $\mathit{rej}$ at time $T$ and in neither at any earlier time.
\end{theorem}

\begin{theorem}[Facts about the setup]
\lean{Turing.NTimeDiag.dz_setup_facts}
\label{Turing.NTimeDiag.dz_setup_facts}
\leanok
\uses{Turing.EffectiveNDMachineCode, Turing.FinTM, Turing.FinTM.bufferTape, Turing.FinTM.bufferTape_eq_none_of_not_mem, Turing.FinTM.bufferTape_nat, Turing.MultiTapeTM.initCfg, Turing.MultiTapeTM.runFrom, Turing.NTimeDiag.DzKit, Turing.NTimeDiag.DzKit.tp, Turing.NTimeDiag.dzBlk, Turing.NTimeDiag.dzLay, Turing.NTimeDiag.dzLay_tp, Turing.NTimeDiag.dzPexit, Turing.NTimeDiag.dzSet, Turing.NTimeDiag.dzSetB, Turing.NTimeDiag.dzSetL, Turing.NTimeDiag.dz_initCfg, Turing.NTimeDiag.dz_setup, Turing.NTimeDiag.qyC, Turing.segment_iff}
\difficulty{4}
Let $\kappa$ be a diagonalizer's kit for $f, g : \mathbb{N} \to \mathbb{N}$, an effective coding scheme $c$ and a finite deterministic decider $B_F$ over bits, with $K_\kappa$ host tapes and named tapes $\tau_0, \dots, \tau_{13}$ (the host tapes of indices $0, \dots, 13$), let $A_q, A_g, A_p$ be the constants of its input calls for the stage query, for $\mathrm{bits}(g(|x|))$ and for prepending $1$, and let $z$ be a binary input. Write $\mathrm{buf}(w)$ for the tape holding the binary word $w$ in cells $0, \dots, |w| - 1$ and blank elsewhere, $\langle s, p, T, P, o\rangle$ for the live configuration on input $z$ with state $s$, input head $p$, work tapes $T$, heads $P$ and output $o$, $\Lambda_z$ for the named-tape contents at the setup's exit: $\mathrm{buf}(q_{f,g}(z))$ on $\tau_0$, where $q_{f,g}$ is the stage query; $\mathrm{buf}(1^m)$ on $\tau_1$ and $\tau_4$, where $m = |\mathrm{bits}(g(|z|))|$; $\mathrm{buf}(1z)$ on $\tau_2$; $\mathrm{buf}(\mathrm{bits}(g(|z|)))$ on $\tau_3$; $\mathrm{buf}(1)$ on $\tau_5$; the block $\mathrm{buf}(1^{g(|z|)}\,1\,\mathrm{bits}(g(|z|)))$ on $\tau_6$; $\mathrm{buf}(1^{g(|z|)})$ on $\tau_7$; and the blank tape on $\tau_8, \dots, \tau_{13}$, extended blank to the other host tapes, and let $\beta = 1^{g(|z|)}\,1\,\mathrm{bits}(g(|z|))$. Then there is $t \le (A_q + A_g + A_p + 40)(g(|z|) + 1)$ such that the setup machine, run from its initial configuration on $z$, equals $\langle \mathit{exit}, 1, \Lambda_z, 0, \varepsilon\rangle$ at time $t$ (with $\mathit{exit}$ the done state of the ruler's rewind) and is not in $\mathit{exit}$ at any earlier time; and in this configuration the cells of $\tau_6$ at positions $h + p$ are nonblank for every $p < |\beta|$, where $h$ is the head position of $\tau_6$, while the cells at $h + |\beta|$ and $h - 1$ are blank.
\end{theorem}

\begin{theorem}[The budget covers every branch]
\lean{Turing.NTimeDiag.dz_K_bound}
\label{Turing.NTimeDiag.dz_K_bound}
\leanok
\uses{Turing.EffectiveNDMachineCode, Turing.FinTM, Turing.NTimeDiag.DzKit, Turing.NTimeDiag.dzBlk, Turing.NTimeDiag.dzBlk_length, Turing.NTimeDiag.dzK, Turing.NTimeDiag.dzQB, Turing.NTimeDiag.dzSetB, Turing.length_bits_le_self}
\difficulty{3}
Let $\kappa$ be a diagonalizer's kit for $f, g : \mathbb{N} \to \mathbb{N}$, an effective coding scheme $c$ and a finite deterministic decider $B_F$ over bits, with $K_\kappa$ host tapes and named tapes $\tau_0, \dots, \tau_{13}$ (the host tapes of indices $0, \dots, 13$), let $K_D = (A_q + A_g + A_p + 40) + (A_q(7A_a + 8) + 2A_l + 140) + 15$, where $A_q, A_g, A_p, A_l, A_a$ are the constants of the kit's calls for the stage query, for $\mathrm{bits}(g(|x|))$, for prepending $1$, for $y \mapsto \mathrm{bits}(|y|)$ and for the inner second component, and write $B_P(G) = (A_q + A_g + A_p + 40)(G + 1)$ for the setup's time bound and $B_Q(G) = (A_q(7A_a + 8) + 2A_l + 140)(G + 1)$ for the run's. For all $G, t \in \mathbb{N}$ with $t \le B_P(G)$, \[ t + 2|1^G\,1\,\mathrm{bits}(G)| + B_Q(G) + 11 \le K_D (G + 1). \]
\end{theorem}

\begin{theorem}[All branches of the diagonalizer halt]
\lean{Turing.NTimeDiag.dz_halts}
\label{Turing.NTimeDiag.dz_halts}
\leanok
\uses{Turing.EffectiveNDMachineCode, Turing.FinTM, Turing.NDTM.HaltsWithin, Turing.NDTM.HaltsWithin.mono, Turing.NTimeDiag.DzKit, Turing.NTimeDiag.DzKit.tp, Turing.NTimeDiag.dzAcc, Turing.NTimeDiag.dzD, Turing.NTimeDiag.dzK, Turing.NTimeDiag.dzPexit, Turing.NTimeDiag.dzQ, Turing.NTimeDiag.dzQB, Turing.NTimeDiag.dzRej, Turing.NTimeDiag.dzSet, Turing.NTimeDiag.dz_K_bound, Turing.NTimeDiag.dz_hQ, Turing.NTimeDiag.dz_setup_facts, Turing.guessRunTM_haltsWithin}
\difficulty{3}
Let $\kappa$ be a diagonalizer's kit for $f, g : \mathbb{N} \to \mathbb{N}$, an effective coding scheme $c$ and a finite deterministic decider $B_F$ over bits, with $K_\kappa$ host tapes and named tapes $\tau_0, \dots, \tau_{13}$ (the host tapes of indices $0, \dots, 13$), let $D_\kappa$ be the diagonalizer built from $\kappa$, and let $K_D = (A_q + A_g + A_p + 40) + (A_q(7A_a + 8) + 2A_l + 140) + 15$, where $A_q, A_g, A_p, A_l, A_a$ are the constants of the kit's calls for the stage query, for $\mathrm{bits}(g(|x|))$, for prepending $1$, for $y \mapsto \mathrm{bits}(|y|)$ and for the inner second component. For every binary input $z$, every branch of $D_\kappa$ on $z$ halts within $K_D(g(|z|) + 1)$ steps.
\end{theorem}

\begin{theorem}[Acceptance of the diagonalizer via streams]
\lean{Turing.NTimeDiag.dz_accepts_iff}
\label{Turing.NTimeDiag.dz_accepts_iff}
\leanok
\uses{Turing.EffectiveNDMachineCode, Turing.FinNDTM.AcceptsWithin, Turing.FinNDTM.acceptsWithin_iff, Turing.FinTM, Turing.FinTM.bufferTape, Turing.MultiTapeTM.runFrom, Turing.NTimeDiag.DzKit, Turing.NTimeDiag.DzKit.tp, Turing.NTimeDiag.dzAcc, Turing.NTimeDiag.dzBlk, Turing.NTimeDiag.dzD, Turing.NTimeDiag.dzK, Turing.NTimeDiag.dzLay, Turing.NTimeDiag.dzPexit, Turing.NTimeDiag.dzQ, Turing.NTimeDiag.dzQB, Turing.NTimeDiag.dzRej, Turing.NTimeDiag.dzSet, Turing.NTimeDiag.dzSetL, Turing.NTimeDiag.dz_K_bound, Turing.NTimeDiag.dz_hQ, Turing.NTimeDiag.dz_handoff, Turing.NTimeDiag.dz_setup_facts, Turing.NTimeDiag.qyC, Turing.guessRunTM_accepts_iff}
\difficulty{4}
Let $\kappa$ be a diagonalizer's kit for $f, g : \mathbb{N} \to \mathbb{N}$, an effective coding scheme $c$ and a finite deterministic decider $B_F$ over bits, with $K_\kappa$ host tapes and named tapes $\tau_0, \dots, \tau_{13}$ (the host tapes of indices $0, \dots, 13$), let $D_\kappa$ be the diagonalizer built from $\kappa$, let $K_D = (A_q + A_g + A_p + 40) + (A_q(7A_a + 8) + 2A_l + 140) + 15$, where $A_q, A_g, A_p, A_l, A_a$ are the constants of the kit's calls for the stage query, for $\mathrm{bits}(g(|x|))$, for prepending $1$, for $y \mapsto \mathrm{bits}(|y|)$ and for the inner second component, and let $z$ be a binary input. Write $\mathrm{buf}(w)$ for the tape holding the binary word $w$ in cells $0, \dots, |w| - 1$ and blank elsewhere, $L[i \mapsto v]$ for $L$ updated at $i$ to $v$, $\hat L$ for the tape layout giving $\tau_i$ the contents $L(i)$ and every other host tape the blank tape, and $\hat P$ for the head layout giving $\tau_i$ the position $P(i)$ and every other head the position $0$, $\langle s, p, T, P, o\rangle$ for the live configuration on input $z$ with state $s$, input head $p$, work tapes $T$, heads $P$ and output $o$, and $Q_\kappa$ for the diagonalizer's run machine (the phase host over the run's phases), with accepting verdict $\mathit{acc}$ and rejecting verdict $\mathit{rej}$; a verdict-avoiding segment of length $t$ from $a$ to $b$ means that $Q_\kappa$ reaches $b$ from $a$ in exactly $t$ steps and is in neither verdict at any earlier time. Write $\Lambda_z$ for the named-tape contents at the setup's exit: $\mathrm{buf}(q_{f,g}(z))$ on $\tau_0$, where $q_{f,g}$ is the stage query; $\mathrm{buf}(1^m)$ on $\tau_1$ and $\tau_4$, where $m = |\mathrm{bits}(g(|z|))|$; $\mathrm{buf}(1z)$ on $\tau_2$; $\mathrm{buf}(\mathrm{bits}(g(|z|)))$ on $\tau_3$; $\mathrm{buf}(1)$ on $\tau_5$; the block $\mathrm{buf}(1^{g(|z|)}\,1\,\mathrm{bits}(g(|z|)))$ on $\tau_6$; $\mathrm{buf}(1^{g(|z|)})$ on $\tau_7$; and the blank tape on $\tau_8, \dots, \tau_{13}$. Then $D_\kappa$ accepts $z$ within $K_D(g(|z|) + 1)$ steps if and only if there are a binary word $s$ with $|s| = |1^{g(|z|)}\,1\,\mathrm{bits}(g(|z|))|$ and $T \in \mathbb{N}$ such that $Q_\kappa$, run from the run's entry configuration $E_s = \langle (\text{flag}, \text{start}), 1, \widehat{\Lambda_z[6 \mapsto \mathrm{buf}(s)]}, 0, \varepsilon\rangle$, is in $\mathit{acc}$ at time $T$, is in neither $\mathit{acc}$ nor $\mathit{rej}$ at any earlier time, and has empty output at time $T$.
\end{theorem}

\begin{theorem}[The first verdict is unique]
\lean{Turing.NTimeDiag.dz_first}
\label{Turing.NTimeDiag.dz_first}
\leanok
\uses{Turing.Cfg, Turing.EffectiveNDMachineCode, Turing.FinTM, Turing.MultiTapeTM.runFrom, Turing.NTimeDiag.DzKit, Turing.NTimeDiag.DzKit.K, Turing.NTimeDiag.DzQPh, Turing.NTimeDiag.dzAcc, Turing.NTimeDiag.dzQ, Turing.NTimeDiag.dzQF, Turing.NTimeDiag.dzRej, Turing.NTimeDiag.qyC, Turing.Segment, Turing.Segment.time_eq, Turing.segment_iff}
\difficulty{4}
Let $\kappa$ be a diagonalizer's kit for $f, g : \mathbb{N} \to \mathbb{N}$, an effective coding scheme $c$ and a finite deterministic decider $B_F$ over bits, with $K_\kappa$ host tapes and named tapes $\tau_0, \dots, \tau_{13}$ (the host tapes of indices $0, \dots, 13$), and let $z$ be a binary input. Write $\mathrm{buf}(w)$ for the tape holding the binary word $w$ in cells $0, \dots, |w| - 1$ and blank elsewhere, $L[i \mapsto v]$ for $L$ updated at $i$ to $v$, $\hat L$ for the tape layout giving $\tau_i$ the contents $L(i)$ and every other host tape the blank tape, and $\hat P$ for the head layout giving $\tau_i$ the position $P(i)$ and every other head the position $0$, $\langle s, p, T, P, o\rangle$ for the live configuration on input $z$ with state $s$, input head $p$, work tapes $T$, heads $P$ and output $o$, and $Q_\kappa$ for the diagonalizer's run machine (the phase host over the run's phases), with accepting verdict $\mathit{acc}$ and rejecting verdict $\mathit{rej}$; a verdict-avoiding segment of length $t$ from $a$ to $b$ means that $Q_\kappa$ reaches $b$ from $a$ in exactly $t$ steps and is in neither verdict at any earlier time. Let $a$ be a configuration of $Q_\kappa$ on $z$, let $T, t \in \mathbb{N}$, let $v \in \{\mathit{acc}, \mathit{rej}\}$ and let $T', P'$ be tapes and heads. If $Q_\kappa$ run from $a$ is in $\mathit{acc}$ at time $T$ and in neither verdict at any earlier time, and there is a verdict-avoiding segment of length $t$ from $a$ to $\langle v, 1, T', P', \varepsilon\rangle$, then $v = \mathit{acc}$.
\end{theorem}

\begin{theorem}[The stage query in a middle]
\lean{Turing.NTimeDiag.dz_query_mid}
\label{Turing.NTimeDiag.dz_query_mid}
\leanok
\uses{Turing.NTimeDiag.ladder, Turing.NTimeDiag.query, Turing.NTimeDiag.query_eq, Turing.NTimeDiag.stage, Turing.NTimeDiag.stageCode, Turing.pairEncode}
\difficulty{1}
Let $f, g : \mathbb{N} \to \mathbb{N}$, let $q_{f,g}$ be the stage query, write $\langle u, v\rangle$ for the pairing encoding of binary words, $\alpha_i$ for the code of stage $i$, $s_f(n)$ for the stage of a length $n$, $(\ell_i)$ for the ladder of $f$, and $M_\alpha$ for the nondeterministic machine decoded from a code $\alpha$ by $c$. For every $n \ge 3$ with $n \ne \ell_{s_f(n) + 1}$, \[ q_{f,g}(1^n) = 0\,\langle\langle \mathrm{bits}(g(n)), \alpha_{s_f(n)}\rangle, 1^{n+1}\rangle. \]
\end{theorem}

\begin{theorem}[Prepending a one to a block of ones]
\lean{Turing.NTimeDiag.dz_ones_succ}
\label{Turing.NTimeDiag.dz_ones_succ}
\leanok
\difficulty{1}
For every $n \in \mathbb{N}$, the binary word $1$ followed by $1^n$ equals $1^{n+1}$.
\end{theorem}

\begin{theorem}[Soundness of the diagonalizer in a middle]
\lean{Turing.NTimeDiag.dz_mid_sound}
\label{Turing.NTimeDiag.dz_mid_sound}
\leanok
\uses{Turing.CodeNDTM.toFinNDTM, Turing.EffectiveNDMachineCode, Turing.FinNDTM.AcceptsWithin, Turing.FinNDTM.AcceptsWithin.mono, Turing.FinNDTM.acceptsWithin_iff, Turing.FinTM, Turing.NTimeDiag.DzKit, Turing.NTimeDiag.dzAlpha, Turing.NTimeDiag.dzBlk_length, Turing.NTimeDiag.dzD, Turing.NTimeDiag.dzK, Turing.NTimeDiag.dz_accepts_iff, Turing.NTimeDiag.dz_alpha_eq, Turing.NTimeDiag.dz_first, Turing.NTimeDiag.dz_ones_succ, Turing.NTimeDiag.dz_qmid, Turing.NTimeDiag.dz_query_mid, Turing.NTimeDiag.dz_split_stream, Turing.NTimeDiag.ladder, Turing.NTimeDiag.stage, Turing.NTimeDiag.stageCode, Turing.pairEncode}
\difficulty{5}
Let $\kappa$ be a diagonalizer's kit for $f, g : \mathbb{N} \to \mathbb{N}$, an effective coding scheme $c$ and a finite deterministic decider $B_F$ over bits, with $K_\kappa$ host tapes and named tapes $\tau_0, \dots, \tau_{13}$ (the host tapes of indices $0, \dots, 13$), let $D_\kappa$ be the diagonalizer built from $\kappa$, let $K_D = (A_q + A_g + A_p + 40) + (A_q(7A_a + 8) + 2A_l + 140) + 15$, where $A_q, A_g, A_p, A_l, A_a$ are the constants of the kit's calls for the stage query, for $\mathrm{bits}(g(|x|))$, for prepending $1$, for $y \mapsto \mathrm{bits}(|y|)$ and for the inner second component, and write $\langle u, v\rangle$ for the pairing encoding of binary words, $\alpha_i$ for the code of stage $i$, $s_f(n)$ for the stage of a length $n$, $(\ell_i)$ for the ladder of $f$, and $M_\alpha$ for the nondeterministic machine decoded from a code $\alpha$ by $c$. Let $n \ge 3$ with $n \ne \ell_{s_f(n) + 1}$. If $D_\kappa$ accepts $1^n$ within $K_D(g(n) + 1)$ steps, then $M_{\alpha_{s_f(n)}}$ accepts $1^{n+1}$ within $g(n)$ steps.
\end{theorem}

\begin{theorem}[Completeness of the diagonalizer in a middle]
\lean{Turing.NTimeDiag.dz_mid_complete}
\label{Turing.NTimeDiag.dz_mid_complete}
\leanok
\uses{Complexity.LogProg.length_bits_mono, Turing.CodeNDTM.toFinNDTM, Turing.EffectiveNDMachineCode, Turing.FinNDTM.AcceptsWithin, Turing.FinNDTM.acceptsWithin_iff, Turing.FinTM, Turing.MultiTapeTM.runFrom, Turing.NDTM.initCfg, Turing.NDTM.runWith, Turing.NTimeDiag.DzKit, Turing.NTimeDiag.dzAlpha, Turing.NTimeDiag.dzBlk_length, Turing.NTimeDiag.dzD, Turing.NTimeDiag.dzK, Turing.NTimeDiag.dz_accepts_iff, Turing.NTimeDiag.dz_alpha_eq, Turing.NTimeDiag.dz_ones_succ, Turing.NTimeDiag.dz_padded_value, Turing.NTimeDiag.dz_qmid, Turing.NTimeDiag.dz_query_mid, Turing.NTimeDiag.dz_seg_read, Turing.NTimeDiag.ladder, Turing.NTimeDiag.stage, Turing.NTimeDiag.stageCode, Turing.pairEncode, Turing.streamStartCfg}
\difficulty{5}
Let $\kappa$ be a diagonalizer's kit for $f, g : \mathbb{N} \to \mathbb{N}$, an effective coding scheme $c$ and a finite deterministic decider $B_F$ over bits, with $K_\kappa$ host tapes and named tapes $\tau_0, \dots, \tau_{13}$ (the host tapes of indices $0, \dots, 13$), let $D_\kappa$ be the diagonalizer built from $\kappa$, let $K_D = (A_q + A_g + A_p + 40) + (A_q(7A_a + 8) + 2A_l + 140) + 15$, where $A_q, A_g, A_p, A_l, A_a$ are the constants of the kit's calls for the stage query, for $\mathrm{bits}(g(|x|))$, for prepending $1$, for $y \mapsto \mathrm{bits}(|y|)$ and for the inner second component, and write $\langle u, v\rangle$ for the pairing encoding of binary words, $\alpha_i$ for the code of stage $i$, $s_f(n)$ for the stage of a length $n$, $(\ell_i)$ for the ladder of $f$, and $M_\alpha$ for the nondeterministic machine decoded from a code $\alpha$ by $c$. Let $n \ge 3$ with $n \ne \ell_{s_f(n) + 1}$, let $\alpha = \alpha_{s_f(n)}$, and let $C \in \mathbb{N}$ be such that for all binary words $x, s$ and every $t \le |s|$ there is $T \le C(t + 1)$ at which the kit's interpreter core $R$, started in its entry state on input $\langle\langle \mathrm{bits}(t), \alpha\rangle, x\rangle$ with $s$ on its stream tape, other tapes blank and heads at $0$, is in its accepting state $r_A$ if $M_\alpha$ run on $x$ along the first $t$ symbols of $s$ has halted with output $1$, and in its rejecting state $r_J$ otherwise, and is in neither of these states before $T$. If $M_\alpha$ accepts $1^{n+1}$ within $t'$ steps, $C(t' + 1) \le g(n)$ and $t' \le g(n)$, then $D_\kappa$ accepts $1^n$ within $K_D(g(n) + 1)$ steps.
\end{theorem}

\begin{theorem}[The diagonalizer at a stage's top]
\lean{Turing.NTimeDiag.dz_top_iff}
\label{Turing.NTimeDiag.dz_top_iff}
\leanok
\uses{Turing.CodeNDTM.toFinNDTM, Turing.EffectiveNDMachineCode, Turing.FinNDTM.AcceptsWithin, Turing.FinTM, Turing.FinTM.ComputesInTime, Turing.FinTM.computesInTime_iff, Turing.MultiTapeTM.initCfg, Turing.MultiTapeTM.runFrom, Turing.NTimeDiag.DzKit, Turing.NTimeDiag.DzQPh, Turing.NTimeDiag.budget, Turing.NTimeDiag.dzAcc, Turing.NTimeDiag.dzBlk, Turing.NTimeDiag.dzD, Turing.NTimeDiag.dzK, Turing.NTimeDiag.dz_accepts_iff, Turing.NTimeDiag.dz_first, Turing.NTimeDiag.dz_qtop, Turing.NTimeDiag.dz_seg_read, Turing.NTimeDiag.ladder, Turing.NTimeDiag.query, Turing.NTimeDiag.query_eq, Turing.NTimeDiag.stage, Turing.NTimeDiag.stageCode, Turing.pairEncode}
\difficulty{5}
Let $\kappa$ be a diagonalizer's kit for $f, g : \mathbb{N} \to \mathbb{N}$, an effective coding scheme $c$ and a finite deterministic decider $B_F$ over bits, with $K_\kappa$ host tapes and named tapes $\tau_0, \dots, \tau_{13}$ (the host tapes of indices $0, \dots, 13$), let $D_\kappa$ be the diagonalizer built from $\kappa$, let $K_D = (A_q + A_g + A_p + 40) + (A_q(7A_a + 8) + 2A_l + 140) + 15$, where $A_q, A_g, A_p, A_l, A_a$ are the constants of the kit's calls for the stage query, for $\mathrm{bits}(g(|x|))$, for prepending $1$, for $y \mapsto \mathrm{bits}(|y|)$ and for the inner second component, and write $\langle u, v\rangle$ for the pairing encoding of binary words, $\alpha_i$ for the code of stage $i$, $s_f(n)$ for the stage of a length $n$, $(\ell_i)$ for the ladder of $f$, and $M_\alpha$ for the nondeterministic machine decoded from a code $\alpha$ by $c$. Let $n \ge 3$ with $n = \ell_{i+1}$ where $i = s_f(n)$; let $\ell = \ell_i$, $\alpha = \alpha_i$, $T^* = (f(\ell + 1) + \ell + 2)^2$ the stage budget and $y = \langle\langle \mathrm{bits}(T^*), \alpha\rangle, 1^{\ell+1}\rangle$. Let $C_B \in \mathbb{N}$ be such that $B_F$ computes the output $1$ from $y$ within $C_B 2^{C_B(T^* + 1)}$ steps if $M_\alpha$ accepts $1^{\ell+1}$ within $T^*$ steps, and computes the output $0$ from $y$ within $C_B 2^{C_B(T^* + 1)}$ steps otherwise. If $C_B 2^{C_B(T^* + 1)} \le g(n)$, then $D_\kappa$ accepts $1^n$ within $K_D(g(n) + 1)$ steps if and only if $M_\alpha$ does not accept $1^{\ell+1}$ within $T^*$ steps.
\end{theorem}

\begin{theorem}[Existence of the lazy diagonalizer]
\lean{Turing.NTimeDiag.exists_diagonalizer}
\label{Turing.NTimeDiag.exists_diagonalizer}
\leanok
\uses{Complexity.TimeConstructible, Turing.CodeNDTM.toFinNDTM, Turing.EffectiveNDMachineCode, Turing.FinNDTM, Turing.FinNDTM.AcceptsWithin, Turing.FinTM, Turing.FinTM.ComputesInTime, Turing.NDTM.HaltsWithin, Turing.NTimeDiag.budget, Turing.NTimeDiag.dzD, Turing.NTimeDiag.dzK, Turing.NTimeDiag.dz_halts, Turing.NTimeDiag.dz_kit_exists, Turing.NTimeDiag.dz_mid_complete, Turing.NTimeDiag.dz_mid_sound, Turing.NTimeDiag.dz_top_iff, Turing.NTimeDiag.ladder, Turing.NTimeDiag.stage, Turing.NTimeDiag.stageCode, Turing.pairEncode}
\difficulty{4}
Let $f, g : \mathbb{N} \to \mathbb{N}$ be time-constructible and let $c$ be an effective coding scheme for nondeterministic machines; write $\langle u, v\rangle$ for the pairing encoding of binary words, $\alpha_i$ for the code of stage $i$, $s_f(n)$ for the stage of a length $n$, $(\ell_i)$ for the ladder of $f$, and $M_\alpha$ for the nondeterministic machine decoded from a code $\alpha$ by $c$, and $T^*_i = (f(\ell_i + 1) + \ell_i + 2)^2$ for the budget of stage $i$. Suppose there is a finite deterministic machine $B_F$ over bits such that for every code $\alpha$ there is $C$ with: for all binary words $x$ and all $t \in \mathbb{N}$, $B_F$ computes the output $1$ from $\langle\langle \mathrm{bits}(t), \alpha\rangle, x\rangle$ within $C 2^{C(t+1)}$ steps if $M_\alpha$ accepts $x$ within $t$ steps, and the output $0$ within the same time otherwise. Then there are a finite nondeterministic machine $D$ over bits and $K \in \mathbb{N}$ such that every branch of $D$ on every input $x$ halts within $K(g(|x|) + 1)$ steps, and for every code $\alpha$ there is $C$ such that for every $n \ge 3$ with $\alpha_{s_f(n)} = \alpha$: (i) if $n \ne \ell_{s_f(n)+1}$, then $D$ accepting $1^n$ within $K(g(n) + 1)$ steps implies that $M_\alpha$ accepts $1^{n+1}$ within $g(n)$ steps, and for every $t'$, if $M_\alpha$ accepts $1^{n+1}$ within $t'$ steps and $C(t' + 1) \le g(n)$, then $D$ accepts $1^n$ within $K(g(n) + 1)$ steps; (ii) if $n = \ell_{s_f(n)+1}$ and $C 2^{C(T^*_{s_f(n)} + 1)} \le n$, then $D$ accepts $1^n$ within $K(g(n) + 1)$ steps if and only if $M_\alpha$ does not accept $1^{\ell_{s_f(n)}+1}$ within $T^*_{s_f(n)}$ steps.
\end{theorem}
